% =====================================================================
%  GMN3 《带框 BPS 态》中译本  主文件
%  原文：arXiv:1006.0146v2（2010），124 页
%  当前版本 v1.0：全篇 126 页（含附录 A–G），2026-11-21 合稿
% =====================================================================
\documentclass[UTF8,a4paper,12pt,fontset=fandol,scheme=chinese]{ctexart}
% !TEX program = xelatex
% GMN3 中译 —— 主控导言区（自 GMN2 改写，保留原文字母宏以便逐式核对）
\usepackage[a4paper,margin=2.6cm,top=2.8cm,bottom=2.8cm]{geometry}
\usepackage{amsmath,amssymb,amsfonts,amsthm,mathtools}
\usepackage{mathrsfs}
\usepackage{bm}
\usepackage{graphicx}
\usepackage{booktabs}
\usepackage{array}
\usepackage{enumitem}
\usepackage{xcolor}
\usepackage{tcolorbox}
\usepackage{hyperref}
\usepackage{xspace}
\usepackage{tikz}
\usepackage{epigraph}

\graphicspath{{figures/}}

% ---------- 字体与行距 ----------
\linespread{1.18}
\setlength{\parskip}{0.25em}
\emergencystretch=1.2em
\sloppy

% 公式按节编号（与原刊排版一致：(1.1)、(2.18)…）
\numberwithin{equation}{section}

% ---------- 超链接 ----------
\hypersetup{
  colorlinks=true,
  linkcolor=blue!45!black,
  citecolor=green!40!black,
  urlcolor=magenta!60!black,
  pdftitle={带框 BPS 态（Framed BPS States）},
  pdfauthor={Gaiotto--Moore--Neitzke},
}

% ---------- 定理类环境（与数译体例一致：各类独立连续编号） ----------
\theoremstyle{plain}
\newtheorem{thm}{定理}
\newtheorem{lem}[thm]{引理}
\newtheorem{prop}[thm]{命题}
\newtheorem{cor}[thm]{推论}
\newtheorem{conj}[thm]{猜想}
\theoremstyle{definition}
\newtheorem{defn}{定义}
\newtheorem{exmp}{例}
\newtheorem{prob}{问题}
\theoremstyle{remark}
\newtheorem{rem}{注}
\newtheorem*{conv}{约定}
\newtheorem*{notn}{记号}

% ---------- 译注 ----------
\newcounter{transnote}
\newenvironment{yizhu}
  {\par\medskip\refstepcounter{transnote}%
   \begin{tcolorbox}[colback=orange!4,colframe=orange!55!black,boxrule=0.5pt,
                     left=4pt,right=4pt,top=3pt,bottom=3pt,arc=1pt]
   \noindent\textbf{译注 \thetransnote：}}
  {\end{tcolorbox}\par\medskip}
% 行内短译注：\zy{...} 生成脚注
\newcommand{\zy}[1]{\footnote{\textcolor{orange!70!black}{【译注】}#1}}

% ---------- 原文字母宏（与 wcf-paper.tex 对齐，便于逐式核对） ----------
\newcommand{\IC}{\mathbb{C}}
\newcommand{\IN}{\mathbb{N}}
\newcommand{\IZ}{\mathbb{Z}}
\newcommand{\IR}{\mathbb{R}}
\newcommand{\IP}{\mathbb{P}}
\newcommand{\IQ}{\mathbb{Q}}

\newcommand{\CA}{\mathcal{A}}
\newcommand{\CB}{\mathcal{B}}
\newcommand{\CC}{\mathcal{C}}
\newcommand{\CD}{\mathcal{D}}
\newcommand{\CE}{\mathcal{E}}
\newcommand{\CF}{\mathcal{F}}
\newcommand{\CG}{\mathcal{G}}
\newcommand{\CH}{\mathcal{H}}
\newcommand{\CI}{\mathcal{I}}
\newcommand{\CJ}{\mathcal{J}}
\newcommand{\CK}{\mathcal{K}}
\newcommand{\CL}{\mathcal{L}}
\newcommand{\CM}{\mathcal{M}}
\newcommand{\CN}{\mathcal{N}}
\newcommand{\CO}{\mathcal{O}}
\newcommand{\CP}{\mathcal{P}}
\newcommand{\CQ}{\mathcal{Q}}
\newcommand{\CR}{\mathcal{R}}
\newcommand{\CS}{\mathcal{S}}
\newcommand{\CT}{\mathcal{T}}
\newcommand{\CU}{\mathcal{U}}
\newcommand{\CV}{\mathcal{V}}
\newcommand{\CW}{\mathcal{W}}
\newcommand{\CX}{\mathcal{X}}
\newcommand{\CY}{\mathcal{Y}}
\newcommand{\CZ}{\mathcal{Z}}

\newcommand{\N}{\mathcal{N}}
\newcommand{\la}{\langle}
\newcommand{\ra}{\rangle}
\newcommand{\half}{\tfrac{1}{2}}
\newcommand{\abs}[1]{\lvert #1 \rvert}
\newcommand{\norm}[1]{\lVert #1 \rVert}
\newcommand{\inprod}[1]{\langle #1 \rangle}
\newcommand{\pz}{(\zeta)}
\newcommand{\ti}[1]{\textit{#1}}
\newcommand{\re}{\operatorname{Re}}
\newcommand{\im}{\operatorname{Im}}
\newcommand{\tr}{\operatorname{tr}}
\newcommand{\sgn}{\operatorname{sgn}}
\newcommand{\vol}{\operatorname{vol}}
\newcommand{\Res}{\operatorname{Res}}
\newcommand{\Tr}{\operatorname{Tr}}
\newcommand{\sfm}{\mathrm{sf}}
\newcommand{\hk}{超凯勒}
\newcommand{\qk}{四元数凯勒}
\newcommand{\kahler}{凯勒}
\newcommand{\elec}{\mathrm{e}}
\newcommand{\magn}{\mathrm{m}}
\newcommand{\degen}{\Omega}
\newcommand{\hsymp}{\varpi}
\newcommand{\clat}{\Gamma}
\newcommand{\vm}{\CB}
\newcommand{\p}{\partial}
\newcommand{\pb}{\bar\partial}
\newcommand{\inst}{\mathrm{inst}}
\newcommand{\elecm}{\mathrm{e}}
\newcommand{\cwarrow}{\text{\Large$\curvearrowright$}}
\newcommand{\ccwarrow}{\text{\Large$\curvearrowleft$}}

% —— 原文自定义宏（与原文同名，便于逐式核对；勿在各 chunk 里重复定义）——
\renewcommand\sf{{\mathrm{sf}}}
\def\tp{\tilde\phi}
\def\n{\vec{n}}
\def\nd{\vec{n}^{\dagger}}
\newcommand{\plog}[1]{#1\partial_{#1}}
\newcommand{\eps}{\epsilon}

% —— GMN3（1006.0146）追加的原文宏 ——
% 注意：原文另有 \def\be{\begin{equation}} / \def\ee{\end{equation}} 与大量单字母宏
% （\a \b \c \p \la \ra \half ...），一律不抄进来，避免覆盖译者常用命令。
\newcommand{\RH}{{\rm RH}}
\newcommand{\WKB}{{\rm WKB}}
\newcommand{\hyper}{{\rm hyper}}
\newcommand{\vect}{{\rm vector}}
\newcommand{\core}{{\rm core}}
\newcommand{\halo}{{\rm halo}}
\newcommand{\BPS}{{\rm BPS}}
\newcommand{\fro}{\overline{\underline{\Omega}}}
\newcommand{\reg}{{\rm reg}}
\newcommand{\flav}{{\rm flavor}}
\newcommand{\juggle}{juggle\xspace}
\newcommand{\juggles}{juggles\xspace}
\newcommand{\juggleing}{juggling\xspace}
\newcommand{\nsing}{{l}}
\newcommand{\ts}{\tilde{s}}
\newcommand{\id}{\mathbf{1}}
\newcommand{\hkq}{/\!/\!/}
\newcommand{\bS}{\mathbf S}
\newcommand{\IS}{{\mathbb S}}
\newcommand{\tCM}{\widetilde{\cal M}}
\newcommand{\tCY}{\widetilde{\cal Y}}
\newcommand{\lieg}{\mathfrak{g}}
\newcommand{\liet}{\mathfrak{t}}
% 原文 \sect{...}：重设公式计数器后再起节（原文用它实现按节编号）
\newcommand{\sect}[1]{\setcounter{equation}{0}\section{#1}}

\DeclareMathOperator{\rank}{rank}
\DeclareMathOperator{\Area}{Area}
\DeclareMathOperator{\End}{End}
\DeclareMathOperator{\Pexp}{Pexp}
\DeclareMathOperator{\sTr}{sTr}
\DeclareMathOperator{\Hom}{Hom}
\DeclareMathOperator{\Hol}{Hol}
\DeclareMathOperator{\sdim}{sdim}
\DeclareMathOperator{\Edges}{Edges}
\DeclareMathOperator{\fund}{fund}

% 插图：原文 \insfig{名}{题注} 对应 figures/名-crop.pdf（scale=0.275）。
% 本译本改用 \insfig[宽]{名}{题注}，默认 0.8\textwidth，便于统一版面。
\newcommand{\insfig}[3][0.8\textwidth]{\begin{figure}[htbp]\centering\includegraphics[width=#1]{figures/#2-crop.pdf}\caption{#3}\label{fig:#2}\end{figure}}

% 原文 \instwofigs{宽1}{名1}{注1}{宽2}{名2}{注2}
\newcommand{\instwofigs}[6]{
\begin{figure}[htbp]
\begin{minipage}[b]{0.5\linewidth}
\centering \includegraphics[width=#1]{figures/#2-crop.pdf} \caption{#3} \label{fig:#2}
\end{minipage}
\hspace{0.4cm}
\begin{minipage}[b]{0.45\linewidth}
\centering \includegraphics[width=#4]{figures/#5-crop.pdf} \caption{#6} \label{fig:#5}
\end{minipage}
\end{figure}
}

% 原文 \insthreefigs{高1}{名1}{注1}{高2}{名2}{注2}{高3}{名3}{注3}
\newcommand{\insthreefigs}[9]{
\begin{figure}[htbp]
\begin{minipage}[b]{0.31\linewidth}
\centering \includegraphics[height=#1]{figures/#2-crop.pdf} \caption{#3} \label{fig:#2}
\end{minipage}
\hspace{0.25cm}
\begin{minipage}[b]{0.31\linewidth}
\centering \includegraphics[height=#4]{figures/#5-crop.pdf} \caption{#6} \label{fig:#5}
\end{minipage}
\hspace{0.25cm}
\begin{minipage}[b]{0.31\linewidth}
\centering \includegraphics[height=#7]{figures/#8-crop.pdf} \caption{#9} \label{fig:#8}
\end{minipage}
\end{figure}
}

% 译者常用的几家译名（统一在此，便于全系列一致）
\newcommand{\KS}{Kontsevich--Soibelman\xspace}
\newcommand{\SW}{Seiberg--Witten\xspace}
\newcommand{\CFT}{CFT\xspace}

% 原文自定义宏（ch04/ch05 用到，原稿未定义）
\newcommand{\ch}{{\rm ch}}
\providecommand{\hW}{\widehat W}
\newcommand{\qbinom}[2]{\genfrac{[}{]}{0pt}{}{#1}{#2}}


\hypersetup{pdftitle={带框 BPS 态（中译注本）},
  pdfauthor={Davide Gaiotto, Gregory W.~Moore, Andrew Neitzke；译：数译（DSH agent）},
  pdfsubject={arXiv:1006.0146v2 中译注本 v1.0},
  pdfkeywords={BPS states, wall-crossing, cluster algebras, spectral networks}}

\begin{document}

% ------------------------------ 标题页 ------------------------------
\pagestyle{plain}
\thispagestyle{empty}

\begin{center}

{\LARGE\bfseries 带框 BPS 态}\\[0.35em]
{\large Framed BPS States}\\[1.0em]

{\large Davide Gaiotto$^{1}$,\quad Gregory W.~Moore$^{2}$,\quad Andrew Neitzke$^{3}$}\\[0.5em]
{\small $^{1}$ 普林斯顿高等研究院自然科学学院\quad $^{2}$ 罗格斯大学 NHETC 与物理与天文系\quad $^{3}$ 得克萨斯大学奥斯汀分校数学系}\\[0.6em]

{\small 原文：arXiv:1006.0146v2\quad 124 页\quad 2010 年}\\[0.3em]
{\small 译者：数译（DSH agent）\quad 版本：v1.0\quad 2026-11-21}

\end{center}

\vspace{1.2em}
\begin{center}\rule{0.85\textwidth}{0.4pt}\end{center}
\vspace{0.4em}

% ------------------------------ 摘要 ------------------------------
\section*{摘要}
\addcontentsline{toc}{section}{摘要}

我们考虑 $d=4$、$\CN=2$ 超对称场论中保持四个超对称不被破缺的一类线算子。这类线算子支持一类新的 BPS 态，我们称之为「带框 BPS 态」（framed BPS states）。它们包含类似 $d=4$、$\CN=2$ 超引力中的 halo 束缚态——其中（普通的）BPS 粒子松散地束缚在线算子上。利用这一构造，我们把 Kontsevich–Soibelman 穿墙公式归约到半本原穿墙公式，从而给出普通 BPS 粒子穿墙公式的一个新证明。在 $S^1$ 上约化之后，线算子 vev 的红外展开为理论模空间 $\CM$ 上的「达布坐标」提供了新的物理解释。此外，我们引入「保护自旋特征标」，用以记录带框 BPS 态的自旋自由度。我们证明，保护自旋特征标的生成函数之间有一种乘法，它定义了 $\CM$ 上全纯函数代数的一个形变。作为这些想法的示例，我们考虑 $A_1$ 型的六维 $(2,0)$ 场论在黎曼面 $C$ 上的紧化。这里我们证明（推广了先前的结果）：线算子由 $C$ 的一个适当装饰版本上的某些叶层（lamination）分类；并在若干显式例子中计算了带框 BPS 态的谱。最后我们指出与簇代数理论的一些有趣联系。

\vspace{0.6em}
\begin{center}\rule{0.85\textwidth}{0.4pt}\end{center}

% ------------------------------ 目录 ------------------------------
\newpage
\tableofcontents
\newpage

% ------------------------------ 正文 ------------------------------
% =====================================================================
%  GMN3 §1  引言与摘要（Introduction and Summary）
%  原文：Lamination-Paper.tex 360--558 行；PDF p.4--6（印刷 p.3--5）
%  编号公式 (1.1)--(1.5) 共 5 条，全部在本块内，故 \eqref 可用；
%  跨块引用（节号、式号）暂写原文编号并加 % TODO，合稿时统一改 \ref/\eqref
% =====================================================================

\section{引言与摘要}
\label{sec:IntroductionAndSummary}

本文延续一项关于模空间 $\CM$ 之性质的研究\cite{Gaiotto:2008cd,Gaiotto:2009hg}；这些模空间自然联系于四维 $\CN=2$ 超对称量子场论\cite{Seiberg:1994rs,Seiberg:1996nz}。这些模空间正是这类理论在 $\IR^3\times S^1$ 上的真空模空间。它们可以被赋予完全可积系统的结构，并且一直是大量密集研究的对象。文献 \cite{Gaiotto:2008cd,Gaiotto:2009hg} 的关注点，是 $\CM$ 上的一族函数，即所谓的「达布坐标」；这族坐标既实用又有趣。说它实用，是因为它提供了计算 $\CM$ 上\hk 度量的简洁途径，提供了可以证明 Kontsevich--Soibelman 穿墙公式\cite{ks1} 的框架，还可以用来构造强耦合下 $\CN=4$ 理论的散射振幅\cite{Alday:2009yn,Alday:2009dv}。说它有趣，是因为它具有许多漂亮的渐近性质与解析性质，把它与若干其他主题联系起来，其中包括可积系统理论的热力学 Bethe 拟设，以及（高阶）Teichm\"uller 理论的 Fock--Goncharov 坐标。

在本引言的余下部分，我们对本文的四个主要结果作一简短的说明性描述。

鉴于 $\CM$ 上达布坐标的核心重要性，为它们给出更直接的物理解释是可取的。本文的主要目标之一正是提供这样一个解释：它们是 $\CN=2$ 理论中线算子的 vev。这是我们的第一个主要结果。这一解释浓缩在第 6 节% TODO \ref{sec:linehol}
的式 (6.3)% TODO \eqref{eq:trh-large-r}
之中：
\begin{equation} \label{eq:intro-eq}
\inprod{L_\zeta} = \sum_\gamma \fro(L_\zeta, \gamma) \CY_\gamma.
\end{equation}
此式左端是一个超对称线算子 $L_\zeta$ 的真空期望值。该算子位于空间中的单个点，沿时间方向伸展，并保持由角度 $\zeta$ 指定的四个庞加莱超荷——与相位为 $\zeta$ 的 BPS 粒子所保持的正是同样的四个超对称。（第 2 节% TODO \ref{sec:Susy-BPS-line}
详细讨论这一定义，并概述超对称线算子理论的若干基本方面。）超对称 't Hooft--Wilson 算子给出了可能的 $L_\zeta$ 的例子，但正如我们将看到的，还有别的例子。式 \eqref{eq:intro-eq} 的右端是达布坐标 $\CY_\gamma$ 的展开。形象地说，$\CY_\gamma$ 可以看成低能红外理论中线算子的 vev，它由插入一个无穷重的双荷子得到；该双荷子带电磁荷 $\gamma$、中心荷相位为 $\zeta$。在这个意义上，式 \eqref{eq:intro-eq} 表达了紫外理论中的线算子 $L_\zeta$ 如何分解为红外理论中的线算子。

展开式 \eqref{eq:intro-eq} 的系数 $\fro(L_\zeta,\gamma)$ 是整数，并且实际上定义了一种新的 BPS 指标，它计数一类新的 BPS 态，我们称之为\emph{带框 BPS 态}（framed BPS state）。这些态在第 3 节% TODO \ref{sec:Hilbert-Halo-WC}
中引入。简言之，在线算子 $L$ 存在时，理论的希尔伯特空间 $\CH_L$ 是 $\N=2$ 超对称之未破缺部分的一个表示。它仍按电磁荷 $\gamma$ 分次；在标号为 $\gamma$ 的扇区中，能量满足一个修正过的 BPS 界
\begin{equation} \label{eq:bps-bound}
E\geq - {\rm Re} (Z_\gamma/\zeta).
\end{equation}
带框 BPS 态就是 $\CH_L$ 中使这个界取到等号的态。为减少混淆，我们把原 $\CN=2$ 理论中（没有线算子时）的 BPS 态称为「原味 BPS 态」（vanilla BPS state）。

我们可以精细化 BPS 指标，以计入带框 BPS 态的自旋信息。这个非常有趣的指标定义于第 3.2 节% TODO \ref{subsec:BPS-Deg-PSC}
的式 (3.12) 与式 (3.14)% TODO \eqref{eq:Def-Framed-PSC}、\eqref{eq:PSC-one}
，称为\emph{保护自旋特征标}（protected spin character，PSC）。带框 BPS 态与「原味」BPS 态都可以定义 PSC。PSC 是变量 $y$ 的函数，记作 $\fro(L_\zeta,\gamma;y)$。取 $y=-1$ 即得 BPS 指标。

本文的第二个与第三个主要结果，都围绕带框保护自旋特征标的生成函数（例如式 (3.38)% TODO \eqref{eq:form-gen}
）的性质展开：
\begin{equation}\label{eq:intro-eq-2}
F( L_\zeta ) := \sum_{\gamma} \fro( L_\zeta, \gamma;y) X_\gamma.
\end{equation}
这里生成函数中的形式变量 $X_{\gamma}$ 应视为某个非交换海森堡代数的元素，因而满足
\begin{equation} \label{eq:heis}
X_{\gamma} X_{\gamma'} = y^{\langle \gamma, \gamma' \rangle} X_{\gamma+\gamma'}
\end{equation}
其中 $\langle \gamma, \gamma' \rangle$ 是通常的 Dirac--Schwinger--Zwanziger 荷的反对称乘积。本文的第二个主要结果是：生成函数 \eqref{eq:intro-eq-2} 所成的代数，给出 $\CM$ 上函数代数的一个非交换形变，其中 $y$ 是形变参数。

这条乘法法则在第 3.6 节% TODO \ref{subsec:Deformed-Ring}
中在物理上得到论证。所得的非交换代数在第 4 节% TODO \ref{sec:Formal-line-operators}
的许多具体例子中得到展示。在第 8 节% TODO \ref{sec:Examples-of-M}
中，我们以一种便于看清「第 4 节的代数确实是第 8 节代数之非交换形变」的方式，描述达布坐标的交换环。在一类重要的例子（$A_1$ 理论）中，这个代数与量子 Teichm\"uller 理论中量子测地长度算子的代数相关，见第 11 节% TODO \ref{sec:Quantum-Holonomy}
。

PSC 与近年来被大量研究的「refined BPS 指标」\cite{Gukov:2004hz,Diaconescu:2007bf,Dimofte:2009bv,Dimofte:2009tm} 之间的主要区别在于：从\emph{先验}的理由看，PSC 是一个指标。然而在我们检查过的例子中，这两个量相符。这种相符并非平凡，它蕴涵着：无论是带框的还是「原味」的 BPS 态，都不存在超对称的某些表示。这一出人意料的相符在第 3.2 节% TODO \ref{subsec:BPS-Deg-PSC}
中被表述为一条「强正性猜想」。该猜想若成立，将产生深远影响。我们在第 4 节% TODO \ref{sec:Formal-line-operators}
中证明：强正性蕴涵线算子之环及其 vev 本质上可以用形式代数手段算出。考虑到 PSC 的几何解释，这是一个高度非平凡的结果。此外，线算子的强正性让人联想到一些深刻的数学，包括 Lusztig 与 Kashiwara 量子群的典范基、Fomin 与 Zelevinsky 簇代数的洛朗现象，以及 Fock 与 Goncharov 的泛洛朗多项式。我们并未试图把这些联系做得很精确，而是把它们留待将来的工作。

带框 BPS 态经历与「原味」BPS 态类似的穿墙。边缘稳定墙的类比物是式 (3.9)% TODO \eqref{eq:Def-BPS-Walls}
定义的 BPS 墙。在这些墙附近，一部分带框 BPS 态有一个非常简单的物理描述，它密切类比于 Frederik Denef 发现的多中心超引力解中的「halo」构型。穿墙机制与文献 \cite{Denef:2007vg} 所述本质上相同：当参数越过一道 BPS 墙时，整个 halo 构型的 Fock 空间被产生或被湮灭。本文的第三个主要结果是描述生成函数 \eqref{eq:intro-eq-2} 如何穿墙的公式。它本质上由量子双对数的共轭给出（精确公式陈述于式 (3.47) 与式 (3.48)% TODO \eqref{eq:Spin-KS-tmn-S}、\eqref{eq:Spin-KS-ii}
）。

这条带框穿墙公式的一个推论，是对 Kontsevich 与 Soibelman 关于\emph{原味} PSC 的所谓「motivic 穿墙公式」的一个简单物理推导，见第 3.5 节% TODO \ref{subsec:Motivic-WCF}
。粗略地说，想法是：通过「原味」BPS 态对带框 BPS 态的影响，来间接研究「原味」BPS 态的谱。（这与我们早先对穿墙公式的推导\cite{Gaiotto:2008cd} 有关；那个推导同样是通过「原味」BPS 态对别的东西的影响来间接研究它们——在那种情形，受影响的是维数约化理论的有效拉氏量。）这样我们就证明了：motivic KSWCF 直接来自 halo Fock 空间的简单物理。

最后，在本文余下部分，我们展示第四个主要结果：在一类有趣的 $\CN=2$ 理论例子中，如何计算带框 BPS 指标与带框 PSC；这类理论即 $\CS$ 类，由 $d=6$ 非阿贝尔 $(2,0)$ 理论在亏格 $g$、$n$ 个穿刺点的黎曼面 $C$ 上作部分扭转紧化而得到\cite{Witten:1997sc,Gaiotto:2009we,Gaiotto:2009hg}。在第 7 节% TODO \ref{sec:Six-Dim-View}
简要回顾 $d=6$ 的 $(2,0)$ 理论之后，我们利用 $d=6$ 理论的表面算子，得到使计算成为可能的关键结果 (7.13)% TODO \eqref{eq:L-5d-Holon}
。该结果断言：可以把 $\CM$ 等同于 $C$ 上平坦联络的模空间，并把线算子的量子期望值表达为平坦联络的经典和乐。在第 10 节% TODO \ref{sec:Examples}
中，我们详细展开取自「$A_1$ 理论」这一类的许多例子；这类理论是 $\CS$ 中基于 $\lieg=su(2)$ 的 $d=6$ 理论的那些。在此过程中，第 9 节% TODO \ref{sec:A1}
把共形 $A_1$ 理论中线算子的分类（由 \cite{Drukker:2009tz} 提出）推广到渐近自由理论。结论是：$C$ 上闭曲线的同痕类被所谓「叶层」（lamination）的对象取代（定义见第 9.4 节% TODO \ref{subsec:Lams-Def}
）。我们由此与 Fock 和 Goncharov 的一些非常有趣的工作\cite{MR2349682,MR2567745} 建立联系。一旦建立了这一联系，带框 BPS 指标的计算在原则上就是直接的，可以化归为一个本质上众所周知的算法——我们称之为\mbox{\emph{traffic rule}} 算法——详见附录 F% TODO \ref{app:traffic}
。

相当直接的 \mbox{traffic-rule} 计算给出了一些物理上出人意料、数学上非平凡的结果。例如，$SU(2)$、$N_f=0$ 理论中基础表示下的 Wilson 线振幅，由形如
\begin{equation} \label{eq:nf0-wilson}
\CY_{\gamma_e} + \CY_{-\gamma_e} + \CY_{\gamma_m},
\end{equation}
的三项表达式 (10.33)% TODO \eqref{eq:wilson-vev-nf0}
给出，其中 $\gamma_e$ 是 $W$ 玻色子电荷的一半，$\gamma_m$ 是一个与之互不局域的电荷。前两项是半经典推理所预期的，但第三项出人意料。它在弱耦合区域是指数小的。沿着这条思路还可以给出若干例子，见第 10 节% TODO \ref{sec:Examples}
。这里我们只再提一个例子：在强耦合的 $SU(2)$、$N_f=4$ 理论中，Wilson 线支持 $7$ 个带框 BPS 超多重态，这一点由简单的计算 (10.44)% TODO \eqref{eq:Nf4-Wilson}
揭示。要通过研究单极子模空间的几何来重现这个结果，大概会相当不平凡。

在第 5 节与第 12 节% TODO \ref{sec:Cluster-Algebras}、\ref{sec:Tropic}
中，我们简要探讨与一些有趣数学的联系，并说明这些联系有实用的物理应用。在第 5 节% TODO \ref{sec:Cluster-Algebras}
中我们说明，形式线算子的代数是一个量子簇代数的例子。这一观点在构造 $\CN=2^*$ 理论的形式线算子时很有用，并对簇代数文献中一些正性猜想提出了精细化。我们引入「簇 $\CN=2$ 理论」的概念。在第 12 节% TODO \ref{sec:Tropic}
中，我们类似地引入「热带 $\CN=2$ 理论」的概念，并为线算子引入「热带标签」。应当把它们看成是对（单）线算子之枚举的一种紫外到红外的映射。在这两节中，我们只是触及了也许非常深刻的联系的表层。

一些技术性论证与一些约定被放进附录。

尽管篇幅（令人遗憾地）很长，本文仍留下许多漂亮的石头没有翻动。我们在第 13 节% TODO \ref{sec:Open-Problems}
指出若干开放问题与未来的研究方向。


% 以下块按 BRIEF.md 分块表逐块补齐后解除注释：
% =====================================================================
%  GMN3 §2  BPS 线算子（BPS line operators）
%  原文：Lamination-Paper.tex 561--1028 行；PDF p.7--13（印刷 p.6--12）
%  编号公式 (2.1)--(2.15) 共 15 条，全部在本块内，故 \eqref 可用；
%  跨块引用（节号、式号、附录）暂写原文编号并加 % TODO，合稿时统一改 \ref/\eqref
% =====================================================================

\section{BPS 线算子}
\label{sec:Susy-BPS-line}

\subsection{$\zeta$ 型线算子的定义}

本文沿用 \cite{Kapustin:2005py,Kapustin:2006pk} 所主张的线算子定义。我们考虑的是由紫外不动点理论经相关微扰而得到的量子场论；所谓线算子，就是该紫外理论在 $AdS_2\times S^2$ 上的一个共形不变边界条件。

本文聚焦于具有 $d=4$、$\CN=2$ 超对称的理论。因此其紫外不动点具有 $su(2,2\vert 2)$ 超共形对称性。（我们的约定见附录 A% TODO 合稿改 \ref{app:Susy-Conventions}
。）我们通常考虑 $\IR^{1,3}$ 中位于空间原点 $x^1=x^2=x^3=0$、沿时间方向伸展的线算子。此时，共形代数 $so(2,4)$ 的未破缺子代数由 $D,P_0,K_0$（它们生成 $sl(2,\IR)$ 的一个拷贝）与空间转动 $M_{ij}$（它们生成 $so(3)$ 的一个拷贝）生成。于是，定义我们线算子的那个紫外边界条件保持 $sl(2,\IR)\oplus so(3)\cong so(4^*)$，即 $AdS_2\times S^2$ 的等距变换李代数。此外，我们只研究保持超共形代数之 $su(2)_R$ R 对称性的算子。我们还要求保持一半超对称，因此要研究保持超代数 $osp(4^*\vert 2)$ 的线算子。事实上，选取该子代数的奇生成元有许多种方式，而这个族的存在在后续讨论中颇为重要。超代数上有一个对合 $I_\zeta$，它由超共形群的某个 $U(1)_R$ 变换共轭空间反射 $x^i\to -x^i$ 而诱导。在庞加莱超对称上，该对合的作用为
\begin{equation}
\begin{split}
Q_\alpha^A & \rightarrow \zeta^{-1}  \epsilon^{AB}
\sigma^0_{\alpha\dot \beta}\bar
Q^{\dot \beta}_{~ B} \\
\bar Q^{\dot \alpha}_{~ A} & \rightarrow  \zeta  \epsilon_{AB}
\bar\sigma^{0 \dot \alpha\beta} Q_\beta^{~B}
\end{split}
\end{equation}
而对特殊共形超对称 $S_\alpha^{~A},\bar S^{\dot \alpha A}$ 有类似的作用。这里 $\zeta$ 是由 $U(1)_R$ 转动得到的相位。该自同构的不动子代数记作 $osp(4^*\vert 2)_\zeta$；对一切 $|\zeta|=1$，它同构于 $osp(4^*\vert 2)$。事实上，所有这样的子代数都由 $U(1)_R$ 变换互相转过去。我们把保持上述超代数的线算子记作 $L(\zeta;\cdots)$ 或 $L_\zeta(\cdots)$，其中 $\cdots$ 指刻画某个具体线算子的其他参数，如下文第 \ref{subsec:LabelsForLines} 节所讨论的那些。正如第 4 节% TODO 合稿改 \ref{sec:Formal-line-operators}
与第 10 节% TODO 合稿改 \ref{sec:Examples}
的例子将显示的，线算子一般只在 $\zeta$ 圆周的万有覆盖上才单值。

我们要研究的场论，由各自的超共形不动点出发、通过沿库仑分支移动或通过开启质量而微扰得到。在理论的红外描述中有 $d=4$、$\CN=2$ 的庞加莱超对称。对合 $I_\zeta$ 也作用在这个超代数上，并附加一条规则 $I_\zeta: Z\to\zeta^2\bar Z$，其中 $Z$ 是中心荷算子。为方便起见，我们取 $\zeta^{-1}$ 的一个平方根，记作 $\xi$，并用下式定义在该对合下不变的超荷：
\begin{equation}\label{eq:Q-zeta-susy}
\CR_{\alpha}^{~A} = \xi^{-1} Q_{\alpha}^{~A} + \xi \sigma^0_{\alpha
\dot \beta } \bar Q^{\dot \beta A}
\end{equation}

沿时间方向伸展、且带线算子 $L$ 的理论，其希尔伯特空间记作 $\CH_L$。$\CR_\alpha^{~A}$ 就是作用在这个希尔伯特空间上的算子。

\subsection{线算子的半环}\label{subsubsec:RingStructure}

一个有趣的观察 \cite{Kapustin:2006hi}（在 \cite{Kapustin:2007wm} 中有更详细的展开）是：$\CN=2$ 规范理论中的超对称线算子构成一个半环（semiring），其结构相当精致，并细致地依赖于理论的物质内容。

为定义加法，我们说：两个线算子之和 $L+L'$ 的关联函数，就是 $L$ 与 $L'$ 各自关联函数之和。存在两个算子之和时，理论的希尔伯特空间就是两个超选择扇区之和，即分别与这两个算子相伴的两个希尔伯特空间：
\begin{equation}\label{eq:sum-line}
\CH_{L+L'}:= \CH_L \oplus \CH_{L'}.
\end{equation}
我们把不是另外两个线算子之和的线算子称为\emph{单线算子}。

此外，还可以通过在路径积分中插入两个线算子来定义算子的乘积 $LL'$。我们所考虑的线算子有一条特殊性质：它们的关联函数与两个算子的间距无关。为看清这一点，注意到那些在对合 $I_\zeta$ 下为奇的超对称 $T_{\alpha}^A = \xi^{-1} Q_\alpha^A -  \xi \sigma^0_{\alpha \dot \beta } \bar Q^{\dot \beta A}$ 满足对易关系
\begin{equation}
[\CR_\alpha^A, T_\beta^B ] = 2 \epsilon_{\alpha\beta} \epsilon^{AB}
(\bar   Z - \zeta^{-2} Z) + 8 \zeta^{-1} \epsilon^{AB}
\sigma^{0m}_{\alpha\beta} P_m.
\end{equation}
取 $\alpha\beta$ 的对称部分，我们就能把空间平移算子 $P_i$ 写成某个 $\CR_\alpha^A$ 的对易子；因此，所有沿 $x^0$ 轴对齐的型 $I_\zeta$ 线算子，其关联函数都与空间坐标无关。我们可以让这些算子互相靠近，而由局域性，其乘积应当等价于某个别的线缺陷。于是，算子乘积展开应当定义一种环乘法。下文我们将对若干有趣的例子把它确定下来。把两个线算子之乘积对应的希尔伯特空间与这两个线算子各自的希尔伯特空间联系起来，是一件有趣的事。这可以做，但出人意料地微妙。它将引导我们在第 3.6 节% TODO 合稿改 \ref{subsec:Deformed-Ring}
定义线算子乘积的一种有趣的非交换形变。

\subsection{拉氏规范理论中线算子的标签}\label{subsec:LabelsForLines}

许多 $d=4$、$\CN=2$ 理论都容许用矢量多重态与超多重态写出的拉氏量描述。对这类理论，按照 \cite{Kapustin:2005py} 与 \cite{Kapustin:2006hi}，我们可以对我们所关心的线算子的定义说得更具体（注意我们的 $\zeta$ 在后一文献中写作 $t^{-1}$）。

我们先回顾一些李代数的理论。从单李代数 $\lieg$ 出发，定义若干标准格。取嘉当子代数 $\liet\subset\lieg$。于是有一组典范的根 $\Phi(\lieg)\subset\liet^*$，它们来自把 $\liet$ 在 $\lieg$ 上的伴随作用对角化。对每个根 $\alpha\in\Phi(\lieg)$，都有 $\lieg$ 中一个 $sl(2)_\alpha$ 的拷贝；我们可以典范地定义相应的\emph{余根} $H_{\alpha}\subset\liet$，它生成 $sl(2)_\alpha$ 的嘉当子代数，并归一化为 $\langle \beta , H_{\alpha} \rangle = 2 (\beta,\alpha)/(\alpha,\alpha)$。（左端是 $\liet^*$ 与 $\liet$ 的典范配对，右端用了嘉当–基灵型；但因假定 $\lieg$ 单，归一化会消掉。）根与余根分别生成格 $\Lambda_r\subset \liet^*$ 与 $\Lambda_{cr}\subset \liet$。关于余根格，有两种有用的看法。由李定理，存在唯一的单连通连通紧李群 $\tilde G$，其李代数为 $\lieg$。余根格的元素可以等同于 $U(1)$ 到 $\tilde T$ 的同态，其中 $\tilde T$ 是 $\tilde G$ 的极大环面，李代数为 $\liet$。另一种说法是：$\Lambda_{cr}$ 是 $\exp[ 2\pi \cdot]: \liet \to \tilde T$ 的核。

利用 $\liet$ 与 $\liet^*$ 的完美配对（而\emph{不是}基灵型），我们便典范地定义了权格与磁权格：
\begin{equation}
\begin{split}
\Lambda_{wt} & := \Lambda_{cr}^* \subset \liet^* \\
\Lambda_{mw} & := \Lambda_{r}^* \subset \liet.
\end{split}
\end{equation}
一个标准结论说，中心 $Z(\tilde G)$ 同构于商
\begin{equation}
Z(\tilde G) \cong \Lambda_{mw}/\Lambda_{cr} \cong
\Lambda_{wt}/\Lambda_r.
\end{equation}
最后，设 $\CW$ 是李代数 $\lieg$ 的外尔群。它同构于 $N(\tilde T)/\tilde T$，其中 $N(\tilde T)$ 是 $\tilde T$ 在 $\tilde G$ 中的正规化子。

现在，我们的理论一般带物质场，因此应当考虑一般的紧连通单李群，其李代数为 $\lieg$。仍由李定理，这样的群必形如 $G\cong \tilde G/\CZ$，其中 $\CZ\subset Z(\tilde G)$ 是中心的一个子群。于是 $Z(G) \cong Z(\tilde G)/\CZ$。由于中心在外尔群作用下不动，我们可以找到一个 $\CW$ 不变的格 $\Lambda_G$，使得
\begin{equation}
 \Lambda_{cr}\subset \Lambda_G \subset \Lambda_{mw}
 \end{equation}
从而 $\CZ \cong  \Lambda_G/\Lambda_{cr} $。事实上，可以把 $\Lambda_G$ 等同于群 $G$ 的指数映射的核：
\begin{equation}
\Lambda_G = \{ P \in \liet \vert \exp(2\pi P) = 1_G \}
\end{equation}
也就是说 $\Lambda_G \cong {\rm Hom}(U(1), T)$，其中 $T= \tilde T/\CZ$ 是 $G$ 的极大环面。对偶地，$G$ 的特征标群是 ${\rm Hom}(T, U(1))$，它典范地对偶于 ${\rm Hom}(U(1),T)$，因为存在完美配对
\begin{equation}
{\rm Hom}(U(1),T) \times {\rm Hom}(T,U(1)) \rightarrow {\rm
Hom}(U(1),U(1))\cong \IZ
\end{equation}
于是特征标群同构于 $\Lambda_G^*$；正是为此，$\Lambda_G$ 有时被称为 $G$ 的\emph{余特征标格}。$\Lambda_G^*$ 可以等同于权格 $\Lambda_{wt}(G)$。这个格也可以看成 $\tilde G$ 的那些在 $\CZ$ 下平凡变换的表示中的权。以上便完成了我们对李代数理论的回顾。

如我们所说，有了显式的场多重态，就能构造一些显式的线算子。最显然的线算子，是沿时间方向一条路径 $p$、在某个固定 $\vec x$ 处的超对称 Wilson 线 \cite{Baulieu:1997nj,Maldacena:1998im,Rey:1998ik}，取规范群的一个表示 $\rho_\CR: G\to {\rm End}(V)$，则
\begin{equation}\label{eq:Susy-Wilson-Zeta}
L_\zeta(\vec x ;\CR) =
\rho_{\CR} P\exp \oint_{p} \left( \frac{\varphi}{2\zeta } - i A
- \frac{\zeta \bar \varphi}{2 }  \right).
\end{equation}

此外还可以构造 't Hooft 算子。回忆一下，它们由对线算子无穷小邻域中的场施加边界条件来定义。在我们的情形，我们要求场在环绕的小球面 $S^2$ 上的极限值对应于一个保持 $osp(4^*\vert 2)_\zeta$ 之庞加莱子代数的超对称磁单极子。回忆：规范群为 $G$ 的理论中的球对称磁单极子，由 $S^2$ 上一个 $G$-丛的转移函数指定（它决定了物理空间 $\IR^4-\{p\}$ 上的主丛）。这样的 $G$-丛等同于 $\Lambda_G \cong {\rm Hom}(U(1),T)$。事实上，对任何这样的同态，都可以通过略微修改 Kapustin 的构造，造出一个保持庞加莱子代数的解（见附录 C% TODO 合稿改 \ref{app:Fixed-Point-Equations}
的式 (C.3)% TODO 合稿改 \eqref{eq:Embed-fix-sol}
）。于是每个同态都定义一个 't Hooft 线算子 $L_\zeta (p;v)$，其中 $v \in \Lambda_G$。\footnote{这里重要的是把 $\zeta$ 理解为在 $U(1)_R$ 的覆盖空间中取值。这一点从我们回顾过的定义看不出来。然而，量子反常把相位 $\zeta$ 的改变与 $\theta$ 角的平移联系起来；线算子的威滕效应表明，'t Hooft 算子不可能是 $\zeta$ 的单值函数。}

众所周知，可以把上述两种构造结合起来，定义 Wilson–'t Hooft 算子 $L_\zeta(p;v,\CR)$，其中 $\CR$ 是 $v$ 的中心化子的一个表示。如 \cite{Kapustin:2005py} 所示，这样的 Wilson–'t Hooft 算子的集合\emph{先验地}以 $\Lambda_{mw} \times \Lambda_{wt}$ 的元素为标签。然而，相容的线算子集合上还有局域性条件。例如，在带有标签为 $(P,0)$ 的 't Hooft 算子时，标签为 $(0,Q)$ 的 Wilson 线算子只对规范丛在 $\IR^3- \{0 \}$ 上的结构群的真表示（而非射影表示）才有意义。更一般地，我们预期：相容的单线算子集合可以由极大子格 $\CL\subset \Lambda_{mw} \times  \Lambda_{wt}$ 构造，这些子格满足一个类 Dirac 的量子化条件：
\begin{equation}\label{eq:Line-Op-Locality}
\forall (P,Q), (P',Q')\in \CL : \qquad \qquad \langle P, Q' \rangle - \langle P' , Q \rangle \in \IZ
\end{equation}
由于相差外尔变换的算子彼此规范等价，我们应当要求 $\CL$ 是 $\CW$ 不变的；此时线算子的规范同构类与 $\bar\CL=\CL/\CW$ 一一对应。这样一族相容集合的自然例子是
\begin{equation}\label{eq:Wilson-tHooft-B}
\bar\CL_G := (\Lambda_G \times \Lambda_G^*)/\CW
\end{equation}
在两个极端，对 $G=\tilde G$ 我们有 $(\Lambda_{cr}\times \Lambda_{wt})/\CW$，而对伴随群 $G = G_{\rm ad}:= \tilde G/Z(\tilde G)$ 我们有 $(\Lambda_{mw} \times \Lambda_{r})/\CW$。

最一般的极大互局部子格 $\CL$ 处于正合列
\begin{equation}\label{eq:label-seq}
0 \rightarrow \Lambda_G^* \rightarrow \CL \rightarrow \Lambda_G  \rightarrow 0
\end{equation}
之中。为看出这一点，先考虑投影 $(P,Q) \to P$。它定义了 $\Lambda_{mw}$ 的一个子格，而后者对某个 $G$ 必为 $\Lambda_G$。又因 $\CL$ 极大且互局部，该映射的核必形如 $(0,Q)$，其中 $Q \in \Lambda_G^*$。序列 \eqref{eq:label-seq} 虽然分裂，但并非\emph{自然}分裂——这反映了威滕效应的可能性——所以若考虑局部系统（例如在耦合常数空间上、或在 $\zeta$ 平面 $\IC^*$ 上），\eqref{eq:label-seq} 可以是非平凡的。

对线算子还可以陈述一个更一般的互局部条件，它应当适用于非拉氏理论。存在线算子 $L$ 时的希尔伯特空间 $\CH_L$ 应当是空间转动群 $SU(2)$ 的一个表示。若考虑位于 $\vec x_1, \vec x_2\in \IR^3$ 的两个线算子 $L_1,L_2$，则转动群破缺为群 $U(1)_{\vec x_1 \vec x_2}$，后者双重覆盖绕 $\vec x_1 \vec x_2$ 轴的转动群。与两线算子之乘积相伴的希尔伯特空间 $\CH_{L_1 L_2}$ 当然是 $U(1)_{\vec x_1 \vec x_2}$ 之万有覆盖的表示。我们的互局部条件就是如下陈述：$\CH_{L_1 L_2}$ 是 $U(1)_{\vec x_1 \vec x_2}$ 的真表示。在具有拉氏量描述的 $\CN=2$ 规范理论中，只要沿库仑分支走得足够远、使由线算子源出的场可以放进嘉当子代数，就能回到前面关于互局部条件的表述。这时可以用一对双荷子存在时电磁场的经典计算，推出 \eqref{eq:Line-Op-Locality}。

\subsubsection{例：$A_1$ 理论}\label{subsubsec:A1-labels}

本节讨论 $A_1$ 理论中线算子的标签。这些理论是由 $su(2)$ 规范群构造出的最一般的拉氏 $\CN=2$ 超共形场论。它们属于 \cite{Gaiotto:2009hg} 中记作 $\CS$ 的更大一类理论；背景与记号请参看该文。这些 $d=4$、$\CN=2$ 理论，是把一个部分扭转的 $ADE$ 型 $(0,2)$ 六维理论放在带穿刺点装饰的黎曼面 $C$ 上紧化而定义的。这类理论非凡的 $S$ 对偶性质在 \cite{Gaiotto:2009we} 中得到阐明。$\CS$ 类中的典型理论容许若干种 $S$ 对偶描述，即作为非阿贝尔规范理论弱耦合到非拉氏物质理论。我们预期，对非阿贝尔规范群存在 Wilson–'t Hooft 型的半 BPS 线算子；但非拉氏物质理论中可用的线算子该如何分类，我们目前尚不知道。$A_1$ 理论是一个有用的例子：它们与 $\CS$ 类理论共享许多有趣性质，却在所有 $S$ 对偶标架中都有拉氏量描述。

一个热身例子是具有 $\CN=4$ 超对称的 $su(2)$ 规范理论 \cite{Kapustin:2006pk}。此时只有一个余根 $H$，故 $\Lambda_{cr} = H \IZ$；也只有一个根 $\alpha$，故 $\Lambda_r = \alpha \IZ$。对偶配对为 $\langle \alpha, H\rangle = 2$，于是 $\Lambda_{wt} = \half \alpha \IZ$，$\Lambda_{mw} = \half H \IZ$。传统上分别选取同构 $\Lambda_{wt} \cong \IZ$ 与 $\Lambda_{mw} \cong \IZ$ 来定义电荷与磁荷；于是我们把 $(P,Q) = (\frac{p}{2} H, \frac{q}{2}\alpha)$ 等同于一对整数 $(p,q)\in \IZ\times \IZ$。相容性条件 \eqref{eq:Line-Op-Locality} 变成 $pq'-p'q= 0 \mod 2$。在这一同构下，若 $G = SU(2)$，则 $\Lambda_G \times \Lambda_G^* = 2\IZ \times \IZ$，即磁荷为偶数、电荷可为任意整数；而若 $G=SO(3)$，则 $\Lambda_G \times \Lambda_G^* =  \IZ \times 2\IZ$，即磁荷可为任意整数、而电荷为偶数。两者由 $S$ 对偶相联系。然而还有第三族相容的线算子集合，它不具有 $\CL_G$ 的形式，而是把 $SO(3)$ 理论的 $\theta$ 角平移半个周期得到的。\footnote{若用 $\frac{\theta}{8 \pi^2} \int_M {\Tr}_{2} F\wedge F$ 来归一化 $SO(3)$ 理论的 $\theta$ 角，则在一般四维流形 $M$ 上 $\theta$ 的周期是 $8 \pi$；但在自旋四维流形上是 $4\pi$。当然，$\IR^4$ 是自旋的。在这个意义上，我们平移了半个周期。}它由满足 $p=q \mod 2$ 的那些 $(p,q)$ 构成。

现在我们回到亏格 $g$、$n$ 个穿刺点的曲线 $C$ 上的 $A_1$ 理论。回忆：一个弱耦合极限由选取一个裤分解来指定，该分解与分离曲线 $c_i$（$i=1,\dots, 3g-3+n$）相伴。规范群的李代数为 $\lieg = su(2)^{\oplus 3g-3 + n}$；对每个由三条曲线 $(c_{i_1}, c_{i_2},c_{i_3})$ 围成的 trinion（三孔球面），都有一个超多重态 $\Phi_{i_1i_2i_3}$。这些超多重态处于表示 $2_{i_1}\otimes 2_{i_2} \otimes 2_{i_3}$ 中，并带一个实性条件（总共 $16$ 个独立实标量场）\cite{Gaiotto:2009we}。

首先注意：给定一个弱耦合尖点（cusp），可以有若干个可能的规范群。令 $\tilde G= SU(2)^{3g-3+n}$。此时容许的规范群形如 $\tilde G/\CZ$，其中 $\CZ$ 是中心 $Z(\tilde G)$ 的任一子群，只要它在所有 $2g-2+n$ 个超多重态上平凡地作用。为更具体地表达这一点，我们可以对每个 $i$ 配一个数 $\epsilon(c_i)\in \IZ_2$，从而固定一个同构 $Z(\tilde G) \cong \IZ_2^{3g-3+n}$。于是，在 trinion 上平凡作用的子群 $\CZ_{\rm max}$，由作用在向量 $\vec \epsilon$ 上的 $2g-2+n$ 个条件定义。这些关系之间还满足一个关系，因此 $\CZ_{\rm max}$ 由 $g$ 个元素生成。固定某个同构 $\CZ_{\rm max} \cong \IZ_2^g$，并记 $G_{\rm max}:= \tilde G/\CZ_{\rm max}$。在由分离曲线 $c_i$ 定义的尖点处，$A_1$ 理论的容许规范群与子群 $\CZ \subset \CZ_{\rm max}$ 一一对应。（顺带一提，这样的子群的个数是 Galois 数
\begin{equation}
\CG_g = 1+ \sum_{k=1}^g \frac{ (2^g-1)(2^{g-1}-1) \cdots (2^{g-k+1}-1)}{(2^k-1)(2^{k-1}-1)\cdots (2-1)}.
\end{equation}
当 $g$ 很大时 $ \CG_g \sim 2^{\frac{1}{4}g^2 }$，所以选择可以相当多。）

现在，对 $\lieg = su(2)^{\oplus 3g-3+n}$，我们把上面关于 $su(2)$ 的同构延拓为同构 $\Lambda_{mw} \times \Lambda_{wt} \cong \IZ^{3g-3+n} \times \IZ^{3g-3+n}$。相容的线算子集合，将由极大 $\CW$ 不变的向量族 $(\vec p, \vec q) \subset \Lambda_{G_{\rm max}} \times \Lambda_{wt}$ 构造，这些向量满足：对所有 $(\vec p, \vec q)$ 与 $(\vec p', \vec q')$ 都有 $\vec p\cdot \vec q' - \vec p' \cdot \vec q = 0 \mod 2$。向量 $\vec p$ 指定 $\IR^3-\{ 0 \} $ 上的一个 $SO(3)^{3g-3+n}$ 丛；限制 $\vec p \in \Lambda_{G_{\rm max}}$ 的出现，是因为超多重态场 $\Phi_{abc}$ 必须是伴随丛在表示 $2_a \otimes 2_b \otimes 2_c$ 中的截面。我们取 $p_i \geq 0$，并在 $p_i=0$ 时取 $q_i \geq 0$，以此固定外尔群作用的一个基本域。

Drukker、Morrison 与 Okuda 的一个漂亮观察 \cite{Drukker:2009tz} 是：$A_1$ 理论中单线算子的分类，与 $C$ 上非自交曲线之同痕类的 Dehn–Thurston 分类密切相关。回忆：对这种曲线 $c$，可以定义 $p_i = c \# c_i$ 与 $q_i = {\rm Twist}_{c_i}(c)$。这里 $\# $ 是同伦相交数，定义为相交数之和（取正），并对同痕类取极小。关于 twist 的细节见 \cite{Drukker:2009tz}。（下文中 $c$ 有时指一条曲线，有时指它的同痕类，视上下文而定。）相当漂亮的是，对每个 trinion 都有 $p_a + p_b + p_c = 0 \mod 2$，而这正是 $\vec p \in \Lambda_{G_{\rm max}}$ 的条件。因此我们可以定义一个一一到上的映射，记作 $\tau_{DMO}$，从 $C$ 上所有非自交路径的同痕类之集合，到 $(\Lambda_{G_{\rm max}} \times \Lambda_{wt})/\CW$。

那么，在这些理论中能定义哪些相容的线算子集合 $\CL$？在对应 $\tau_{DMO}$ 之下，条件 \eqref{eq:Line-Op-Locality} 等价于条件
\begin{equation}\label{eq:max-even}
c \# c' = 0 \mod 2.
\end{equation}
若（同痕类的）路径集合 $\CP$ 中任意一对都满足 \eqref{eq:max-even}，我们称 $\CP$ 为\emph{偶}的；若 $\CP$ 不真包含于任何更大的偶集合之中，则称其为\emph{极大偶}的。于是，相容的线算子集合 $\CL$ 定义了 $C$ 上一个极大偶路径集合。特别地，若对某个子群 $\CZ \subset \CZ_{\rm max}$ 取 $G = \tilde G/\CZ$，则 $(\Lambda_G \times\Lambda_G^*)/\CW$ 决定一个极大偶路径集合，记作 $\CP_G$。反之，给定 $\CP$，我们可以把它投影到磁荷，并重构出与 \eqref{eq:label-seq} 一样的正合列。

$S$ 对偶群，即模群，会把规范群变掉，正如上面提到的 $\CN=4$ 例子那样。用模二同调可以给出 $\CZ$ 如何变化的显式公式。这补足了 \cite{Drukker:2009tz} 中关于 $S$ 对偶对线算子作用的讨论。

% =====================================================================
%  GMN3 §3  希尔伯特空间、halo 与穿墙（Hilbert spaces, halos, and wall-crossing）
%  §3.1–3.3；原文 Lamination-Paper.tex 1029--1506 行；PDF p.14--21（印刷 p.13--20）
%  编号公式 (3.1)--(3.24) 共 24 条；图 1 幅（spectrum）；脚注 9 个
%  跨块引用（式 (2.10)、(10.33)，第 4 节，附录 D）暂写原文编号并加 % TODO，
%  合稿时统一改 \eqref / \ref；label 名一律照抄原文。
%  ！注意：(3.6) 沿用原文 label \label{eq:bps-bound}，与 ch01 给 (1.2)
%  自拟的同名 label 撞名，合稿时由 Lead 择一改名（本块遵从原文）。
% =====================================================================

\section{希尔伯特空间、halo 与穿墙}
\label{sec:Hilbert-Halo-WC}

\subsection{带框 BPS 态}
\label{sec:framed-BPS-states}

此前我们只讨论了 $\CN=2$ 理论的紫外方面。现在我们要过渡到红外。先交代一点记号。我们把库仑分支记作 $\CB$，把它的奇异除子记作 $\CB^{sing}$，并把「原味」电磁荷与味荷构成的格记作 $\Gamma$。\footnote{众所周知，若沿 $\CB^*:=\CB-\CB^{sing}$ 中的非平凡路径对 $u$ 作解析延拓，$\Gamma$ 会经历单值变换。因此严格说来，该谈论的是 $\Gamma_u$。用数学的话说，$\Gamma$ 是 $\CB^*$ 上的一个格丛，并定义了一个「局部系统」（local system），也就是说，它带有一个平坦联络。人们可以在 $\CB-\CB^{sing}$ 的万有覆盖 $\widehat\CB$ 上工作，在那里 $\Gamma$ 可以被平凡化；也可以直接在底空间上工作，只要记住存在非平凡的单值性。这是个人口味问题，我们留给读者。更多的背景见 \cite{Gaiotto:2008cd,Gaiotto:2009hg} 及其中的文献。}

现在我们来研究理论在 $\IR^{1,3}$ 上的希尔伯特空间：无穷远处的真空由 $u\in\CB$ 标定，并且在原点 $x^i=0$ 处插入了一个简单线算子 $L_\zeta$。这个希尔伯特空间我们记作 $\CH_{u,L,\zeta}$，它是 $osp(4^*\vert 2)_\zeta$ 的庞加莱子超代数的一个表示。此外，它按电荷分次：
\begin{equation}\label{eq:Grade-Hilbert}
\CH_{u,L,\zeta} = \bigoplus_{\gamma \in \Gamma_L}
\CH_{u,L,\zeta,\gamma}.
\end{equation}
这里 $\Gamma_L$ 是 $L$ 存在时的电磁荷与味荷的泊松格，它是「原味」电荷之格 $\Gamma$ 的一个挠子。也就是说，它具有如下形式
\begin{equation}\label{eq:Lattice-Shift}
\Gamma_L = \Gamma + \gamma_L
\end{equation}
其中 $\gamma_L \in \Gamma \otimes \IR$，且对一切 $\gamma \in \Gamma$ 有 $\langle \gamma_L, \gamma\rangle \in \IZ$。若用路径积分表述来说，\eqref{eq:Grade-Hilbert} 中对 $\Gamma$ 的直和就变成对拓扑扇区的求和。

\eqref{eq:Grade-Hilbert} 中对 $\gamma$ 求和的必要性，在弱耦合的纯 $SU(N)$ 理论中很容易看出来：取库仑分支上足够远的位置，使规范对称性被强破缺为 $U(1)^{N-1}$。例如考虑 Wilson 线算子 (2.10)% TODO 合稿改 \eqref{eq:Susy-Wilson-Zeta}。在朴素的经典极限下，这个算子的 vev 将是对红外理论中若干超对称 Wilson 线所对应的 vev 之和，这些 Wilson 线由表示 $\CR$ 的权来标记。\footnote{也许出人意料的是，在具体例子中我们会看到这个朴素的答案漏掉了一些东西：对权求和时还有额外的贡献。一个简单的例子是纯 $SU(2)$ 规范理论中基础表示下的 Wilson 算子，其答案我们在下面的式 (10.33)% TODO 合稿改 \eqref{eq:wilson-vev-nf0} 中给出。对应于基础表示之权的项有两项，但还有第三项。我们会把这第三项归因于粒子与线算子形成的一个有趣的束缚态。} 这个例子同时也说清了 \eqref{eq:Lattice-Shift} 中那个平移的物理起源。考虑纯 $SU(2)$ 规范理论中基础表示 Wilson 线算子的例子：所有「原味」BPS 态都带偶数电荷（在使 W 玻色子带电荷 $2$ 的约定下），但在一个基础表示 Wilson 圈存在时，一个态应当能够带奇数电荷。往系统里添加「原味」BPS 粒子的自由，使「$\Gamma_L$ 应当是 $\Gamma$ 的挠子」这一点变得显而易见。

线算子上的相互局域条件表明：对一组合乎一致的线算子 $L\in\CL$，相应的 $\Gamma_L$ 应全部落在同一个格 $\Gamma_{\CL}$ 之中，使得
\begin{equation}\label{eq:Consist-LO-Latt}
\Gamma \subset \Gamma_{\CL},
\end{equation}
并且 $\Gamma_{\CL}$ 上有一个整值反对称形式，它限制到 $\Gamma$ 上就是标准的那一个。

扇区 $\CH_{u,L,\zeta,\gamma}$ 中的每一个都是 $osp(4^*\vert 2)_\zeta$ 的庞加莱子超代数的一个表示。当我们以 $x^0$ 为常数的时片来量子化时，超对称算子满足厄米性条件
\begin{equation}
\begin{split}
(\CR_{1}^{~1})^\dagger & = - \CR_{2}^{~2}, \\
(\CR_{1}^{~2})^\dagger & =   \CR_{2}^{~1}.
\end{split}
\end{equation}
此外，在从超共形点形变而来的理论中，希尔伯特空间会有一些扇区带有 $\CN=2$ 庞加莱超对称的非零中心荷。算子 $\CR_{\alpha}^{~A}$ 满足代数
\begin{equation}
\{ \CR_{\alpha}^{~A} , \CR_{\beta}^{~B} \} = 4 \left( E + {\rm Re}(Z/\zeta) \right) \epsilon_{\alpha\beta} \epsilon^{AB}
\end{equation}
其中 $E=P^0$ 是能量算子。把它与厄米性条件结合起来，就得到 BPS 界
\begin{equation} \label{eq:bps-bound-3d}
E + {\rm Re}\, (Z / \zeta) \geq 0,
\end{equation}
其中 $Z = Z_\gamma(u)$ 是与真空 $u\in \CB$、扇区 $\gamma$ 相联系的标准中心荷。

现在可以给出本文的主要定义之一了：我们定义\emph{带框 BPS 态}（framed BPS state）为使这个 BPS 界取到等号的态，即 $E = -{\rm Re}\, (Z / \zeta)$。（这个名字的动机来自带框箭图表示（framed quiver representation）的概念。）我们可以把 \eqref{eq:Grade-Hilbert} 特殊到由带框 BPS 态构成的子空间：
\begin{equation}\label{eq:Grade-BPS-Hilbert}
\CH_{u,L,\zeta}^{\rm BPS} = \bigoplus_{\gamma \in \Gamma_L}
\CH_{u,L,\zeta,\gamma}^{\rm BPS}.
\end{equation}

注意 \eqref{eq:bps-bound-3d} 与标准的 BPS 界 $E \geq \vert Z \vert$ 不同。理解 \eqref{eq:bps-bound-3d} 有一个很好的启发式办法：它用的是「把 Wilson–'t Hooft 算子看作插入一个无穷重双荷子」这一标准解释的红外版本。把格 $\Gamma$ 用一个额外的味荷 $\gamma_f$ 扩张，并考虑一个非常重的粒子，其电荷为 $\gamma_f - \gamma$、中心荷为 $Z = \zeta M  -  Z_\gamma$，其中 $M>0$。在 $M \to +\infty$ 极限下，重整化后的 BPS 界正好是
\begin{equation}\label{eq:heuristic-bound}
E \geq  \lim_{M\to +\infty} \left(\vert \zeta M  -  Z_\gamma\vert - M \right)= - {\rm Re}(Z_\gamma/\zeta).
\end{equation}

总电荷为 $\gamma$ 的 BPS 粒子，其能量有下界 $\vert Z_\gamma(u)\vert$。然而我们在下面会论证：在线算子存在时，可以有一些有趣的量子态，其能级低于 $\vert Z_\gamma(u)\vert$、但高于界 $-{\rm Re}\, (Z / \zeta)$。从远处看，它们像是位于线算子所在位置、电荷为 $\gamma_c$ 的「芯粒子」（core particle），与一个（或多个）电荷为 $\gamma_h$ 的「halo 粒子」相互作用，且 $\gamma_c + \gamma_h = \gamma$。在那些 halo 粒子并未束缚的情形——也就是说，其波函数并未本质上局限在空间中一个有限区域内——能量从 $-{\rm Re}\, (Z_{\gamma_c}(u) / \zeta)+ \vert Z_{\gamma_h}(u)\vert$ 开始，构成连续谱。此外还有束缚态，它们类比于超引力的多中心束缚态 \cite{Denef:2000nb,Denef:2000ar}。对这些态而言，halo 粒子的波函数本质上局限在空间中一个紧致区域内。这类态具有离散的能谱。只要不存在使界 $-{\rm Re}\, (Z_{\gamma_h}(u) / \zeta) = \vert Z_{\gamma_h}(u)\vert$ 取到等号的 halo 电荷，带框 BPS 态与连续谱之间就存在非零能隙。图 \ref{fig:spectrum} 示意性地画出了能谱。

\insfig{spectrum}{在插入了线算子 $L_\zeta$ 的理论中，电荷为 $\gamma$ 的诸态之能量谱的示意图。最下方是带框 BPS 态；接着是束缚在线算子上的离散激发态；最后是各种非束缚态的连续谱，它们对应于把 $\gamma$ 分解为 $\gamma_c + \gamma_h$ 的不同方式。带框 BPS 态与各连续谱之间有安全的间隔，除非存在某个 $\gamma_h$ 使 $-{\rm Re}\, (Z_{\gamma_h}(u) / \zeta)= \vert Z_{\gamma_h}(u)\vert$。}

条件 $-{\rm Re}\, (Z_\gamma(u) / \zeta)= \vert Z_{\gamma}(u)\vert$ 定义了\emph{BPS 墙}：\footnote{在 \cite{Gaiotto:2008cd,Gaiotto:2009hg} 中我们考虑的是「BPS 射线」，它们不过是固定 $u$ 时 BPS 墙到 $\zeta$-球面的投影。}
\begin{equation}\label{eq:Def-BPS-Walls}
\widehat W(\gamma ):=\{ (u,\zeta) \vert Z_{\gamma}(u)/\zeta \in
\IR_- \}\subset \widetilde \CB \times \widetilde\IC^*.
\end{equation}
当我们在空间 $\tilde \CB \times \IC^*$ 中变动参数时，一旦击中某道 BPS 墙，带框 BPS 态就会与连续谱混合。

\subsection{带框 BPS 简并度与保护自旋特征标}
\label{subsec:BPS-Deg-PSC}

带框 BPS 态按空间转动李代数 $so(3)$ 的表示变换。在固定分次 $\gamma$ 下它们是有限维的，于是我们可以定义一个简单的指标，来计数带框 BPS 简并度：
\begin{equation}
\fro(u,L_\zeta, \gamma) := {\rm Tr}_{\CH_{u,L,\zeta,\gamma}^{\rm
BPS}} (-1)^{(2J_3)}
\end{equation}
事实上，我们还可以定义一个更精细的量来计数带框 BPS 态。我们希望定义一个能记录自旋信息的量。简单的自旋特征标 ${\rm Tr}\, y^{2J_3}$ 朴素地看就够用了，但若若干 BPS 态配对成超对称代数的一个非保护表示，它就可能发生难以预测的跳变。为得到一个更稳健的指标，我们回想：带框 BPS 态也按 $su(2)_R$ 变换，而四个超荷 $\CR_\alpha^A$ 中有一个，在自旋 $so(3)$ 与 $su(2)_R$ 的对角组合下是单态，即
\begin{equation}
\CQ := \epsilon_{\alpha A} \CR_{\alpha}^A.
\end{equation}
于是我们可以定义所谓的\emph{（带框）保护自旋特征标}，简称 PSC：\footnote{我们感谢 Juan Maldacena 建议了这个改进的指标。}
\begin{equation}\label{eq:Def-Framed-PSC}
\fro(u,L_\zeta, \gamma;y) := {\rm Tr}_{\CH_{u,L,\zeta,\gamma}^{\rm
BPS}} (-1)^{2J_3} (-y)^{2\CJ_3} = {\rm Tr}_{\CH_{u,L,\zeta,\gamma}^{\rm
BPS}} y^{2J_3} (-y)^{2I_3}
\end{equation}
其中 $I_3$ 是一个 R 对称生成元，且
\begin{equation}
\CJ_3 := J_3 + I_3.
\end{equation}
PSC 是一个指标，因为 $\CQ$ 在 $\CJ_a$ 下是单态、与 $(-1)^F$ 反对易，并且在 $\CR_{\alpha}^A$ 代数的长表示上可逆。通常我们会把记号缩短，把对 $u$ 的依赖略去不写。

特殊到 $y=-1$，我们就回到带框 BPS 简并度 $\fro(u,L_\zeta, \gamma)$。由于超对称小表示的刚性，只要没有态进入或离开希尔伯特空间，它们就在参数的形变下不变。不过，仍会出现穿墙现象，类比于没有线算子时 BPS 态的通常穿墙。我们在第 \ref{subsec:Halos} 节中研究这一点。

对通常的「原味」BPS 态也可以定义 PSC。「原味」PSC 的定义基于这样一个论断：量
\begin{equation}\label{eq:PSC-one}
{\rm Tr}_{\CH}(2 J_3) (-1)^{2J_3} (-y)^{2\CJ_3}
\end{equation}
在 $\CN=2$ 代数的长表示 $\CH$ 上为零。为证明这一点，考虑 $so(3)\oplus su(2)_R$ 的特征标 ${\rm Tr}\, x_1^{2J_3} x_2^{2I_3}$。由它出发，通过
\begin{equation}\label{eq:Inermed-PSC}
{\rm Tr}(2 J_3) (-1)^{2J_3} (-y)^{2\CJ_3} = x_1\frac{\partial}{\partial x_1}\left({\rm Tr}\, x_1^{2J_3} x_2^{2I_3}\right)\vert_{x_1=-x_2 = y}
\end{equation}
就得到 \eqref{eq:PSC-one}。现在，在有质量粒子的小超代数的表示论中，长表示作为 $so(3)\oplus su(2)_R$ 的表示具有如下形式
\begin{equation}
\rho_{hh}\otimes \rho_{hh}\otimes \mathfrak{h}_{\ell}
\end{equation}
其中 $\mathfrak{h}_{\ell}$ 是 $so(3)\oplus su(2)_R$ 的任意有限维表示，而 $\rho_{hh}$ 是\emph{半超多重态}表示 $\rho_{hh}\cong (\half;0) \oplus (0;\half)$。另一方面，小超代数的短表示作为 $so(3)\oplus su(2)_R$ 的表示具有如下形式
\begin{equation}
\rho_{hh} \otimes \mathfrak{h}_{s}
\end{equation}
其中 $\mathfrak{h}_{s}$ 是 $so(3)\oplus su(2)_R$ 的任意有限维表示。现在注意
\begin{equation}
{\rm Tr}_{\rho_{hh}}  x_1^{2J_3} x_2^{2I_3} = x_1 + x_1^{-1} + x_2 + x_2^{-1}
\end{equation}
因此 \eqref{eq:Inermed-PSC} 在长表示上为零，但在短表示上非零。事实上，它在短表示上的值恰好是
\begin{equation}
(y- y^{-1}) {\rm Tr}_{\mathfrak{h}_s} y^{2J_3} (-y)^{2I_3} .
\end{equation}
于是，我们可以通过
\begin{equation}
(y- y^{-1}) \Omega(u,\gamma;y) := {\rm Tr}(2 J_3) (-1)^F (-y)^{2\CJ_3}.
\end{equation}
来定义一个（不带框的）保护自旋特征标 $\Omega(u,\gamma;y)$，即
\begin{equation}
\Omega(u,\gamma;y) = {\rm Tr}_{\mathfrak{h}_s} y^{2J_3} (-y)^{2I_3} .
\end{equation}
注意，若 $\mathfrak{h}_s$ 作为 $su(2)_R$ 表示的同型分解只含单态，则保护自旋特征标\emph{等同于} $\mathfrak{h}_s$ 的自旋特征标。于是我们立刻看出：一个标准超多重态有 $\mathfrak{h}_s = (0;0)$，故其 PSC 为 $\Omega(u,\gamma;y)=1$；而一个标准 W 玻色子有 $\mathfrak{h}_s= (\half;0)$，因而给出 $\Omega(u,\gamma;y)=y+1/y$。事实上，在我们所知的一切具有良好紫外极限的 $\CN=2$ 场论例子中，BPS 粒子所在的表示在 $su(2)_R$ 下的分解都只含单态。因此我们把 $\mathfrak{h}_s$ 含 $su(2)_R$ 非平凡表示的 BPS 粒子称为\emph{奇特 BPS 粒子}（exotic BPS particle）。类似地，我们把在 $su(2)_R$ 下非平凡变换的带框 BPS 态称为\emph{奇特带框 BPS 态}。对非奇特的 BPS 态，我们的 PSC 与 \cite{Gukov:2004hz,Dimofte:2009bv,Dimofte:2009tm} 的「精细化指标」相同，至多差一个（依赖约定的）符号。

现在我们对「原味」与带框 BPS 态各陈述一组「正性猜想」。其中最强的一条就是上面提到的「无奇特粒子猜想」：$\mathfrak{h}_s$ 的同型分解只含 $su(2)_R$ 单态。也就是说，在库仑分支的光滑点上没有奇特粒子。\footnote{在希格斯分支或混合分支上，陈述这样的命题会有一些麻烦。}注意，借由穿墙，奇特 BPS 粒子的存在必然会带来奇特带框 BPS 态。一条较弱的猜想是「强正性猜想」，它断言 $\Omega(u,\gamma;y)$（以及其带框对应物）是 $su(2)$ 特征标 $\chi_n(y)$ 的线性组合，且系数为\emph{非负整数}。「弱正性猜想」则只要求 $\Omega(u,\gamma;y)$ 在 $y=+1$ 处为正整数。

无奇特粒子猜想蕴涵强正性猜想，但反过来未必成立。强正性猜想的一个充分条件是：$\mathfrak{h}_s$ 的所有同型分量都具有整值的 $su(2)_R$ 自旋。然而这个条件并非必要。

我们将在第 4 节% TODO 合稿改 \ref{sec:Formal-line-operators}
中看到，强正性猜想（哪怕只取弱化形式）对带框与标准 BPS 谱都有深远而漂亮的推论。有很强的证据表明，它对 $A_1$ 类中的一切理论都以一种非常非平凡的方式成立。另一个在我们所有例子中似乎都成立的事实是：对固定的线算子，空间 $\CH_{u,L,\zeta,\gamma}^{\rm BPS}$ 只在有限多个电荷 $\gamma$ 上非零。也就是说，对固定的线算子，只有有限多个非零的带框 BPS 不变量。我们没有简单的场论论证来证明这一论断，但我们猜想这个有限性也许是场论中普遍成立的。不过必须承认，我们的全部证据都基于 $A_1$ 类理论。这也是为什么更好地理解 class $\CS$ 中更高阶的理论很重要，我们希望日后回到这一课题。\footnote{最近（2012 年 8 月），Diaconescu 等人证明了：对纯 $SU(K)$ 规范理论，无奇特粒子猜想确实成立。}

最后还有一点观察。自旋特征标近来在文献中已有研究，但总是在场论的语境里，而非超引力的语境里。我们的构造提示了场论情形更受青睐的一个原因：在超引力中没有可用的 $SU(2)_R$ 对称性，因而无法构造保护自旋特征标。

\subsection{halo 与 halo Fock 空间}
\label{subsec:Halos}

现在我们来描述一类特别有趣的带框 BPS 态，即「芯」（core）与包围它的 BPS 粒子「halo」之间的束缚态；这个芯支撑在线算子附近的一个小邻域内。这类构型与 Denef 的 $\CN=2$ 超引力 halo 解密切相关 \cite{Denef:2000nb,Denef:2000ar,Denef:2002ru}。在 \cite{Denef:2007vg} 中，这一现象被用来推导半本原穿墙公式。我们马上会看到，halo 对穿墙问题的带框版本同样重要。

我们先作经典讨论，取最简单的情形：一个电荷为 $\gamma_c$ 的 BPS 芯，束缚着单个电荷为 $\gamma_h$ 的 halo 粒子。\cite{Denef:2000nb} 中一个论证的变体给出如下结论。\footnote{我们在附录 D 中给出推导的若干细节。% TODO 合稿改 \ref{app:Low-Energy-FxdPoint}
}探测与芯电荷 $\gamma_c$ 相联系的背景红外的 halo 粒子，其能量为
\begin{equation}\label{eq:Halo-Particle-Energy}
E_{halo} = \vert Z_{\gamma_h}(u(r))\vert \left(1 + \cos(\alpha_h -
\alpha_\zeta)\right) -  {\rm Re}(Z_{\gamma_h}(u)/\zeta).
\end{equation}
回想 $u \in \CB$ 代表无穷远处真空模的值。探针粒子在与线算子相距 $r$ 处所探测到的模依赖于 $r$，我们记 $Z_{\gamma_h}(u(r)) =\vert Z_{\gamma_h}(u(r))\vert e^{i\alpha_h(r)}$。另一方面，$\zeta$ 与 $r$ 无关，我们把它的相位写作 $\zeta = e^{i \alpha_\zeta}$。

halo 粒子在某个 $r$ 处使能量取极小，那里 $e^{i\alpha_h(r)} / \zeta = -1$。这给出 Denef 关于 BPS 黑洞束缚态半径的公式 \cite{Denef:2000nb} 的一个密切类比：\footnote{的确，用式 (3.8)% TODO 合稿改 \eqref{eq:heuristic-bound}
附近所解释的启发式论证，取中心荷 $\zeta M - Z_{\gamma_c}$、电荷 $\gamma_f + \gamma_c$，Denef 的公式在 $M\to +\infty$ 极限下就化为 \eqref{eq:halo-rad}。}
\begin{equation}\label{eq:halo-rad}
r_{\rm halo} = \frac{\langle \gamma_h, \gamma_c\rangle}{2 {\rm
Im}(Z_{\gamma_h}(u)/\zeta)}.
\end{equation}
这个结果有几个重要的推论。第一，它给出 Denef 稳定性条件的一个类比：halo 构型只对使 \eqref{eq:halo-rad} 给出的 $r_{\rm halo}$ 为正的那些 $(u,\zeta)$ 存在。第二，注意当 $Z_{\gamma_h}(u) = - \vert Z_{\gamma_h}(u)\vert \zeta$ 时 $r_{\rm halo} \to \infty$——换句话说，当 $(u, \zeta)$ 落在 BPS 墙 $\widehat W(\gamma_h)$ 上时。除非我们非常靠近这道墙，$r_{\rm halo}$ 都远小于红外描述成立的自然截断。所以图像是这样的：当人们趋近 BPS 墙 $\widehat W(\gamma_h)$ 时，由（某个倍数）电荷 $\gamma_h$ 的粒子搭成的 halo 构型会达到一个尺度，使得把它们完全放在红外理论里处理是合理的。当我们到达墙上时，这些 halo 胀到无穷大，并在另一侧消失。

与这些 halo 相联系的\emph{量子}态是什么？幸运的是，经典 halo 构型很容易量子化：由于彼此互为 BPS，halo 粒子之间不发生相互作用，它们只是生成一个 $\IZ_2$ 分次的量子态 Fock 空间。为在数学上描述它，设 $\CH'(\gamma_h;u)$ 是电荷为 $\gamma$ 的「原味」BPS 态空间，并已抽去一个半超多重态。它是空间 $so(3)$ 与 $su(2)_R$ 的一个表示。我们用一个 $\IZ_2$ 分次，其符号由 $-(-1)^{2J_3}$ 给出。通常的 BPS 简并度 $\Omega(\gamma_h;u)$ 就是 $\CH'(\gamma_h;u)$ 关于这个分次的超维数取负。令 $J_{\gamma_c,\gamma_h}= \half (\vert \langle \gamma_c, \gamma_h\rangle \vert-1)$，引入空间 $so(3)$ 的表示 $(J_{\gamma_c,\gamma_h})$，其维数为 $\vert \langle \gamma_c, \gamma_h\rangle \vert$。这解释了 halo–芯相互作用所产生的电磁场的自旋。我们在有限维向量空间 $(J_{\gamma_c,\gamma_h})\otimes \CH'(\gamma_h;u)$ 上构造一个 $\IZ_2$ 分次 Fock 空间，其中把 $(J_{\gamma_c,\gamma_h})$ 的 $\IZ_2$ 分次取为\emph{偶}。\footnote{给定一个有限维 $\IZ_2$ 分次向量空间 $V= V^0 \oplus V^1$，超维数为 $(n_0\vert n_1)$，与之相伴的 $\IZ_2$ 分次 Fock 空间是 $V^0$ 上的对称代数与 $V^1$ 上的反对称代数之张量积。更通俗地说，我们取一组基 $\alpha_i$，$i=1,\dots, n_0$，以及 $\gamma_s$，$s=1, \dots, n_1$，然后把 $\alpha_i$ 对应到产生算子的玻色型海森堡代数、把 $\gamma_s$ 对应到产生算子的费米型克利福德（Clifford）代数表示出来。这些看起来古怪的 $\IZ_2$ 分次有一个漂亮的物理解释 \cite{Denef:2002ru}，归根结底与磁场迫使 halo 粒子的自旋朝内这一事实有关。}

具体地说，若我们由
\begin{equation}
\Omega(u,\gamma_h;-z) =  {\rm Tr}_{\CH'(\gamma_h;u)}(-z)^{2J_3} z^{2 I_3} =
\sum_{-M_{\gamma_h}}^{M_{\gamma_h}} a_{m, \gamma_h} z^m
\end{equation}
定义一组整数 $a_{m, \gamma_h}$，其中 $M_{\gamma_h} \ge 0$ 是 halo 粒子的最大 $\CJ_3$ 的两倍，则我们的 Fock 空间由 $\vert a_{m,\gamma_h}\vert$ 个产生算子生成，这些产生算子的 $2 \CJ_3$ 本征值为 $m+m'$，其中 $m'$ 取遍形如 $m' = -2J_{\gamma_c,\gamma_h}, -2J_{\gamma_c,\gamma_h} + 2, \dots,2J_{\gamma_c,\gamma_h}-2, 2J_{\gamma_c,\gamma_h}$ 的值。当 $m$ 为偶数（即 $a_{m,\gamma_h}>0$）时振子是费米型的，当 $m$ 为奇数（即 $a_{m,\gamma_h}<0$）时振子是玻色型的。注意 $\Omega(\gamma_h;u) = \sum_m a_{m, \gamma_h}$。若 $\langle\gamma_c,\gamma_h \rangle=0$，则我们把 $(J_{\gamma_c,\gamma_h})$ 视为零向量空间，此时不形成任何 halo。当然，$a_{m, \gamma_h}$ 是 $u$ 的分段连续整值函数，但我们通常在记号中略去这一依赖。

% =====================================================================
%  GMN3 §3.4--3.6  带框穿墙 / motivic 穿墙公式 / 线算子的形变环
%  原文：Lamination-Paper.tex 1507--2020 行；PDF p.22--27（印刷 p.21--26）
%  编号公式 (3.25)--(3.63) 共 39 条（含脚注中 2 条），全部在本块内，故 \eqref 可用；
%  插图 1 幅（fig:zeta-u-wall）；脚注 3 条；跨块引用暂写原文编号并加 % TODO
% =====================================================================

\subsection{带框穿墙}\label{subsec:Halo-Wall-Crossing}

在第 3.2 节% TODO 合稿改 \ref{subsec:BPS-Deg-PSC}
中，我们定义了保护自旋特征标（PSC）$\fro(L,\gamma;u,\zeta)$，它「计数」带框 BPS 态。但正如我们在第 3.1 节% TODO 合稿改 \ref{sec:framed-BPS-states}
中所指出，当 $(\zeta,u)$ 越过一道 BPS 墙时，带框 BPS 态会与连续谱混合。因此，关于 $\fro(L,\gamma;u,\zeta)$ 不变性的标准论证失效了。这就是「穿墙」问题的带框版本。所幸它比「原味」BPS 态的完整穿墙问题要容易：从带框 BPS 谱中消失的态，正是我们在第 3.3 节% TODO 合稿改 \ref{subsec:Halos}
中描述过的那些简单的 halo 态。于是我们可以完全描述 $\fro(L,\gamma;u,\zeta)$ 在越过墙时的跳变。

考虑一条路径 $(u_t,\zeta_t)\in\widehat{\CB}\times\widehat{U(1)}$，它穿过一道墙 $\widehat{W}(\gamma_h)$。如上所述，当抵达该墙时 $r_{\rm halo}$ 趋于无穷。因此，一整个由 halo 束缚态构成的 Fock 空间——由电荷正比于 $\gamma_h$（比例系数为正）的 halo 粒子搭建而成——要么出现在 BPS 谱中，要么从其中消失。

现在让我们描述这件事对（带框）保护自旋特征标的影响。设 $\{x_\gamma\}$ 是 $\Gamma$ 的群代数的一组基，使得
\begin{equation} \label{eq:untwisted}
x_{\gamma} x_{\gamma'} = x_{\gamma +
\gamma'}.
\end{equation}
然后构造生成泛函
\begin{equation}\label{eq:gen}
F(u, L, \zeta, \{x_\gamma\}; y) := \sum_{\gamma} \fro(u, L, \zeta,
\gamma; y) x_\gamma.
\end{equation}
（为了避免记号变得过于笨重，当某些变量可以安全地当作隐含时，我们常常会把它们省略。）

上文所述的每个类型为 $m, m', \gamma_h$ 的产生算子，都为 Fock 空间上迹的乘积表示贡献一个因子
\begin{equation}
(1 + (-1)^{m} y^{m+m'} x_{\gamma_h} )^{  a_{m,\gamma_h}}
\end{equation}
定义 $\Gamma_h := \IZ \gamma_h$，并对 $\gamma_c \in \Gamma_{\CL}$ 令
\begin{equation}
F_{\bar \gamma_c} := \sum_{\gamma'\in \Gamma_h} \fro(u, L, \zeta,
\gamma_c +\gamma'; y) x_{\gamma_c+\gamma'}.
\end{equation}
（这个求和只依赖于投影 $\bar\gamma_c\in \Gamma_L/\Gamma_h$。）\zy{原文把 $\gamma_c$ 所取的格写作 $\Gamma_{\mathcal L}$，紧接的投影又写 $\Gamma_L$，二者应指同一个格（线算子的电荷格）。}
这是 $F$ 中对应于电荷类型为 $\gamma_c + \ell \gamma_h$ 的态的那一部分。我们用 $F_{\bar\gamma_c}^\pm$ 分别记这个生成函数在墙的两侧——即分别满足 ${\rm Im}(Z_{\gamma_h}(u)/\zeta)>0$ 与 ${\rm Im}(Z_{\gamma_h}(u)/\zeta)<0$ 的一侧——的取值。于是有
\begin{equation} \label{eq:transform}
F^\pm_{\bar\gamma_c} = F^\mp_{\bar\gamma_c}  \prod_{\gamma_h}
\prod_{m=-M_{\gamma_h}}^{M_{\gamma_h}} \prod_{m'= - 2J_{\gamma_c,
\gamma_h}}^{2J_{\gamma_c, \gamma_h}} (1 + (-1)^{m } y^{m+m'}
x_{\gamma_h} )^{  a_{m,\gamma_h}}.
\end{equation}
这里我们稍微滥用了记号：对 $\gamma_h$ 的第一个乘积，指的是对所有给出同一道墙 $\widehat{W}(\gamma_h)$ 的 halo 电荷取乘积。左端究竟取 $F^+$ 还是 $F^-$（也就是我们究竟是获得还是失去一个 Fock 空间），取决于穿墙的方向以及 $\langle \gamma_c,\gamma_h\rangle$ 的符号。（更精确的陈述见下文式 \eqref{eq:spin-ks-tmn-ii}。）

这个变换看起来像是形式变量 $x_\gamma$ 的一次「坐标变换」的效果：它把每个 $x_{\gamma_c}$ 乘上 $x_{\gamma_h}$ 的一个有理函数。遗憾的是，这种解释与乘法规则 \eqref{eq:untwisted} 并不完全相容。但当 $y=1$ 时它确实行得通，下面我们就来说明。此时 \eqref{eq:transform} 特化为
\begin{equation} \label{eq:tspec}
F^\pm_{\bar\gamma_c}  = F^\mp_{\bar\gamma_c}  \prod_{\gamma_h}
\prod_{m=-M_{\gamma_h}}^{M_{\gamma_h}} (1 + (-1)^{m}
x_{\gamma_h} )^{  \lvert \langle \gamma_h, \gamma_c\rangle \rvert a_{m,\gamma_h} }.
\end{equation}
而「坐标」变换
\begin{equation}\label{eq:y-one-prod}
x_\gamma \to x_\gamma  \prod_{\gamma_h}
\prod_{m=-M_{\gamma_h}}^{M_{\gamma_h}}  (1 + (-1)^{m } x_{\gamma_h} )^{\langle \gamma_h, \gamma\rangle a_{m,\gamma_h}}
\end{equation}
在我们从 ${\rm Im}(Z_{\gamma_h}(u)/\zeta) < 0$ 的一侧穿墙到 ${\rm Im}(Z_{\gamma_h}(u)/\zeta) > 0$ 的一侧时，恰好重现 \eqref{eq:tspec}。此外，它与乘法法则 \eqref{eq:untwisted} 相容。

若 $y=-1$，\eqref{eq:transform} 中的因子 $y^{m'}$ 与 $- (-1)^{\langle \gamma_h, \gamma_c\rangle }$ 相同。有一个简单的技巧可以驯服这些多出来的符号：把乘法规则从 \eqref{eq:untwisted} 改成
\begin{equation} \label{eq:twisted-mult}
\hat x_{\gamma} \hat x_{\gamma'} = (-1)^{\langle \gamma, \gamma'\rangle } \hat x_{\gamma +\gamma'}.
\end{equation}
于是，利用关系
$\hat x_{\gamma_c + n \gamma_h} = \hat x_{\gamma_c} \left( (-1)^{\inprod{\gamma_h, \gamma_c}} \hat x_{\gamma_h} \right)^n$，
我们得到 \eqref{eq:tspec} 的一个类比
\begin{equation} \label{eq:transform2}
F^\pm_{\gamma_c}  = F^\mp_{\gamma_c}  \prod_{\gamma_h}
\prod_{m=-M_{\gamma_h}}^{M_{\gamma_h}} (1 -
\hat x_{\gamma_h} )^{ \lvert \inprod{\gamma_h, \gamma_c} \rvert a_{m,\gamma_h}}.
\end{equation}
变换规则
\begin{equation}\label{eq:y-min-one-gen}
\hat x_\gamma \to \hat x_\gamma  \prod_{\gamma_h}
(1 - \hat x_{\gamma_h} )^{\inprod{\gamma_h, \gamma} \Omega(\gamma_h)}
\end{equation}
作用于 $y=-1$ 的生成泛函时，正确地重现 \eqref{eq:transform2}，并且与扭乘法规则 \eqref{eq:twisted-mult} 相容。

更仔细地陈述如下：若 $\zeta_{\rm cw}$ 位于 BPS 墙 $\widehat{W}(\gamma_h)$ 的顺时针一侧，$\zeta_{\rm ccw}$ 位于逆时针一侧，则在略去所有其他无关指标后，生成函数 $F$ 满足
\begin{equation}\label{eq:Careful-rule}
F\left(\zeta_{\rm cw}, \hat x_\gamma\right) = F\left(\zeta_{\rm ccw} , \hat x_\gamma (1- \hat x_{\gamma_h})^{-\langle \gamma, \gamma_h\rangle \Omega(\gamma_h)} \right).
\end{equation}

上面用到的变换 \eqref{eq:y-min-one-gen}，正是 \cite{ks1} 中引入的辛同胚：定义 $\CK_{\gamma_h}$ 为
\begin{equation}\label{eq:KS-sympo}
\CK_{\gamma_h}(\hat x_\gamma) := \hat x_\gamma (1-\hat x_{\gamma_h})^{\langle
\gamma, \gamma_h \rangle },
\end{equation}
则 \eqref{eq:y-min-one-gen} 就简单地是
$\hat x_\gamma \to \CK_{\gamma_h}^{-\Omega(\gamma_h)}(\hat x_\gamma)$。

\subsubsection{非对易变量与量子双对数}

对一般的 $y$，我们无法把 \eqref{eq:transform} 解释为一次变量替换——至少无法用\emph{对易}变量做到。但假设我们引入满足下列关系的形式变量 $X_\gamma$
\begin{equation}\label{eq:Heis-Alg}
X_\gamma X_{\gamma'} = y^{\inprod{\gamma, \gamma'}} X_{\gamma +
\gamma'},
\end{equation}
并再次考虑生成函数
\begin{equation}\label{eq:form-gen}
F(u,L,\zeta,\{X_\gamma\};y) := \sum_{\gamma} \fro(u, L, \zeta,
\gamma;y) X_\gamma,
\end{equation}
这一次把它看成这些\emph{非对易}变量的函数。在这种情况下，正如我们马上会看到的，增加或减少 halo 束缚态的效果，可以漂亮地概括为 $X_\gamma$ 的某个变换。这个变换由与 Faddeev 和 Kashaev \cite{Faddeev:1993rs} 的量子双对数作共轭来实现。（这个函数所享有的各种性质与别名，在 \cite{Teschner:2005bz} 的附录 A 中有很好的总结。）

首先我们需要一条数学引理。设 $\langle \gamma_c, \gamma_h \rangle = n$ 不为零。我们在群代数中展开对易变量，以此定义 $y$ 的洛朗多项式：
\begin{equation}
x_{\gamma_c} (1 + y^{n-1} x_{\gamma_h})(1 + y^{n-3}
x_{\gamma_h})\cdots (1 + y^{3-n} x_{\gamma_h})(1 + y^{1-n}
x_{\gamma_h}) = \sum_{j=0}^{\vert n\vert}  P^{(n)}_j(y) x_{\gamma_c
+ j \gamma_h}
\end{equation}
从物理上看很显然，$P^{(n)}_j(y)$ 就是第 $j$ 个反对称积 $\Lambda^j \rho_{\vert n \vert }$ 的特征标，其中 $\rho_{N}$ 是 $SU(2)$ 的 $N$ 维不可约表示。\footnote{顺便说一句，有一个「$q$-二项式定理」也给出同样的识别（取 $n>0$），这既有趣又会很有用：
\begin{equation}
P^{(n)}_j(y) = \frac{ [n]_y ! }{ [j]_y ! [n-j]_y !}:=\genfrac{[}{]}{0pt}{}{n}{j}
\end{equation}
其中
\begin{equation}
[n]_y := \frac{y^{n} - y^{-n}}{y - y^{-1}} .
\end{equation}
}

我们的引理断言：在展开非对易表达式
\begin{equation}
X_{\gamma_c} \Phi_n( X_{\gamma_h}) = \sum_{j=0}^{\vert n \vert}
P^{(n)}_j(y) X_{\gamma_c + j \gamma_h}
\end{equation}
时出现的，是\emph{同一组}洛朗多项式，其中
\begin{equation}
\Phi_n(\xi) := \begin{cases}  \prod_{s=1}^{n} (1 + y^{-(2s-1)} \xi )
& n >0 \\
 1 & n=0 \\
 \prod_{s=1}^{\vert n\vert } (1 + y^{(2s-1)} \xi ) & n< 0
 \end{cases}
\end{equation}
这条引理的证明是直接的，从略。

利用这条引理，我们可以断言：穿墙的效果，是把 \eqref{eq:form-gen} 中\emph{所有}非对易变量作如下变换
\begin{equation}\label{eq:Spin-KS-tmn}
X_\gamma \to X_\gamma \prod_{m=-M_{\gamma_h}}^{M_{\gamma_h}}
\Phi_{\langle \gamma, \gamma_h\rangle} ( (-1)^m y^m
X_{\gamma_h})^{\epsilon a_{m,\gamma_h}}
\end{equation}
在这个变换律中，我们正在穿过墙 $\widehat{W}(\gamma_h)$。若我们从 $\langle \gamma_h, \gamma \rangle \ {\rm Im}(Z_{\gamma_h}(u)/\zeta) < 0$ 的一侧穿到 $\langle \gamma_h, \gamma \rangle \ {\rm Im}(Z_{\gamma_h}(u)/\zeta) > 0$ 的一侧，则符号 $\epsilon=+1$，因为此时我们在核心电荷 $\gamma$ 周围\emph{获得}了一个 halo 粒子的 Fock 空间。若沿相反方向穿墙，则 $\epsilon=-1$，因为此时我们\emph{失去}了一个 halo 粒子束缚态的 Fock 空间。若用 $F^+, F^-$ 分别记形式生成函数 \eqref{eq:form-gen} 在 ${\rm Im}(Z_{\gamma_h}(u)/\zeta) > 0$ 与 $<0$ 一侧的取值，则（略去所有其他无关变量）：
\begin{equation}\label{eq:spin-ks-tmn-ii}
F^+(X_\gamma) = F^-\left( X_\gamma \left( \prod_{m=-M_{\gamma_h}}^{M_{\gamma_h}}
\Phi_{\langle \gamma, \gamma_h\rangle} ( (-1)^m y^m
X_{\gamma_h})^{a_{m,\gamma_h}}\right)^{{\rm sign}(\langle \gamma_h,
\gamma\rangle ) }\right)
\end{equation}

变换 \eqref{eq:Spin-KS-tmn} 可以用量子双对数的共轭来优雅地概括。定义
\begin{equation}
\Phi(X) := \prod_{k=1}^\infty ( 1 + y^{2k-1} X)^{-1}
\end{equation}
那么变换 \eqref{eq:Spin-KS-tmn} 与 \eqref{eq:spin-ks-tmn-ii} 就简单地等价于下述规则：
\begin{equation}\label{eq:Spin-KS-tmn-S}
F^+(X_\gamma) =   S_{\gamma_h} F^-(X_\gamma ) S_{\gamma_h}^{-1}
\end{equation}
其中
\begin{equation}\label{eq:Spin-KS-ii}
S_{\gamma_h} =  \prod_{\gamma_h} \prod_{m=-M_{\gamma_h}}^{M_{\gamma_h}} \Phi (
(-1)^m y^m X_{\gamma_h})^{a_{m,\gamma_h}}
\end{equation}
这里第一个乘积同样表示对与 $\gamma_h$ 平行的所有电荷取乘积。相应地，带框简并度的生成函数 \eqref{eq:form-gen} 按与 $S_{\gamma_h}$ 的共轭而变换。

\subsection{motivic 穿墙公式}\label{subsec:Motivic-WCF}

设 $\widehat{\CB}$ 是库仑分支 $\CB$ 光滑部分的万有覆盖，并考虑
\begin{equation}\label{eq:Chamber-Space}
\Xi: =  \widehat{\CB} \times \widehat{\IC}^* - \cup_{\gamma: \exists u: \Omega(\gamma;u) \ne 0 }  \widehat{W}(\gamma)
\end{equation}
这个空间被分成若干腔——即 $\Xi$ 的连通分支。我们用指标 $c$ 来标记这些腔。考虑 $\Xi$ 的两个腔 $c_1,c_2$。考虑 $\widehat{\CB} \times \widehat{\IC}^*$ 中连接这两个腔的一条路径 $\CP$。位于这两个腔中的线算子 $L$ 的生成函数 \eqref{eq:form-gen}，将由形如 \eqref{eq:Spin-KS-tmn-S} 的变换的某个复合相联系：\footnote{请注意，我们在这里悄悄引入了一个假设：生成函数可以从单位圆 $\vert \zeta \vert =1$ 延拓到 $\IC^*$。这个假设在当前的论证中并不本质，在第 4 节与第 5 节的讨论中也不本质。但在第 6 节中它将很重要。}% TODO 合稿改 \ref{sec:Formal-line-operators}、\ref{sec:Cluster-Algebras}、\ref{sec:linehol}
\begin{equation}\label{eq:Path-Transf}
F(L;c_1) = S(\CP)  F(L;c_2)
\end{equation}
其中 $S(\CP)$ 是与路径 $\CP$ 从 $c_1$ 走到 $c_2$ 途中所穿过的各道墙 $\widehat{W}(\gamma_h^i)$ 相对应的变换的路径有序乘积
\begin{equation}
S(\CP)= \prod_{\gamma_h^i} {\rm Ad}(S_{\gamma_h^i} ).
\end{equation}
我们希望断言：变换之积 $\prod_{\gamma_h^i} {\rm Ad}(S_{\gamma_h^i} )$ 与连接 $c_1$ 到 $c_2$ 的路径 $\CP$ 无关。只要理论拥有足够多的线算子 $L$，使得知道变换 \eqref{eq:Path-Transf} 的作用就强到足以完全确定 $S(\CP)$，这一点便成立。我们并没有完全一般的理由说明它为何成立（事实上，我们也没有完全一般的理由说明为什么竟然会存在任何线算子）。不过，我们将在下面的例子中看到，线算子总是足够多；事实上，在这些例子里，库仑分支上任一点处的线算子与电磁荷之间都存在一一对应，并且可以把任意 $X_\gamma$ 还原为各函数 $F(L,\{X_\gamma\})$ 的线性组合。进一步的讨论见下文第 12 节% TODO 合稿改 \ref{sec:Tropic}
。

\insfig[0.8\textwidth]{zeta-u-wall}{从 $(\zeta_1, u_1)$ 到 $(\zeta_2, u_2)$ 的两条路径。BPS 射线在联合的 $(\zeta, u)$ 空间中以余维 1 的轨迹出现。每道 BPS 墙 $\widehat{W}(\gamma)$ 都附着一个坐标变换。如图所示，墙可以汇合与分裂，因此两条路径一般会穿过不同的 BPS 墙集合。尽管如此，相应坐标变换的复合仍必须与路径无关。}

$S(\CP)$ 的路径无关性，实际上是 Kontsevich 与 Soibelman \cite{ks1} 的「motivic 穿墙公式」的一个版本。为看出这一点，考虑 $\widehat{\CB} \times \widehat{\IC}^*$ 中连接 $(\zeta_1, u_1)$ 到 $(\zeta_2, u_2)$ 的两条路径，如图 \ref{fig:zeta-u-wall} 所示；并比较 $\Xi$ 中包含 $(\zeta_1, u_1)$ 与包含 $(\zeta_2,u_2)$ 的两个腔 $c_1,c_2$ 的生成函数 $F(L, c_1)$ 与 $F(L,c_2)$。假设我们沿路径 $(\zeta_1, u_1) \to (\zeta_2, u_1) \to (\zeta_2, u_2)$ 走；并假设当 $u$ 沿从 $u_1$ 到 $u_2$ 的路径变化时，没有任何 $\arg Z_\gamma(u)$ 进入或离开区间 $(\arg \zeta_1, \arg \zeta_2)$。那么所有函数 $F(L,c_1)$ 都按下述有序乘积变换：
\begin{equation}
{\bS}(\zeta_1, \zeta_2;u_1) = \prod_{\gamma: \arg Z(\gamma) \in (\zeta_1, \zeta_2)} {\rm Ad}(S_\gamma).
\end{equation}
另一方面，若我们改走路径 $(\zeta_1, u_1) \to (\zeta_1, u_2) \to (\zeta_2, u_2)$，则它们改为按 ${\bS}(\zeta_1, \zeta_2;u_2)$ 变换。于是，对所有线算子 $L$ 都有
\begin{equation}\label{eq:SF-SF}
{\bS}(\zeta_1, \zeta_2;u_1)F(L,c_1)={\bS}(\zeta_1, \zeta_2;u_2)F(L,c_1).
\end{equation}

我们希望由 \eqref{eq:SF-SF} 断定这两个变换确实相等：
\begin{equation} \label{eq:s-equal}
{\bS}(\zeta_1, \zeta_2;u_1) = {\bS}(\zeta_1, \zeta_2;u_2).
\end{equation}
正如上面讨论过的，我们必须假设存在足够多的线算子 $L$，才能推出 \eqref{eq:s-equal}。

与标准讨论的一点小差别是：我们在这里只得到了量子双对数的乘积在「伴随表示」中的等式。不过，按照从方程 \eqref{eq:y-one-prod} 到 \eqref{eq:y-min-one-gen} 的讨论可知，特化到 $y=-1$ 便给出等式 \eqref{eq:s-equal}，此时
\begin{equation}\label{eq:Symplectic-S}
{\bS}(\zeta_1, \zeta_2;u_1) = \prod_{\gamma: \arg Z(\gamma) \in (\zeta_1, \zeta_2)} \CK_\gamma^{\Omega(\gamma_\BPS; u_1)}.
\end{equation}
以 \eqref{eq:Symplectic-S} 给出的 ${\bS}$ 代入的方程 \eqref{eq:s-equal}，就是 \cite{Gaiotto:2008cd,Gaiotto:2009hg} 中所讨论的 Kontsevich--Soibelman 穿墙公式（WCF）的形式。知道所有 $\zeta_{1,2}$ 对应的 ${\bS}(\zeta_1, \zeta_2;u)$，等价于知道 $u$ 处的 BPS 指标集合。而要确定完整的 BPS 谱，还需要更多信息，即保护自旋特征标。

本节最后，我们对一些相关文献作些评论。涉及非对易（$q$-形变）量子双对数的穿墙公式出现在 Kontsevich 与 Soibelman 的工作 \cite{ks1} 中。经典极限（$q = -1$）也在那里讨论过，并直接给出了 $\N=2$ 理论中第二螺旋度超迹的穿墙公式，后者在 \cite{Gaiotto:2008cd} 中从物理上得到了证明。然而，起初并不完全清楚如何把这一成功推广到 $q$-形变的情形。文献 \cite{Dimofte:2009bv} 提出，应当把 $q$ 等同于计数 BPS 态自旋的参数 $y$；文献 \cite{Dimofte:2009tm} 把这个提议精确化为一个穿墙公式，并证明它在若干例子中成立。尽管如此，一般的物理证明仍然缺失。文献 \cite{Cecotti:2009uf} 给出了一个新的论证，把对易关系 \eqref{eq:Heis-Alg} 与穿墙公式联系到 A 模型开拓扑弦，进而联系到 Chern--Simons 理论。而我们在这里给出了一个直接证明，用的只是四维规范理论的一般概念。

\subsection{线算子的形变环}\label{subsec:Deformed-Ring}

正如我们在第 2.2 节% TODO 合稿改 \ref{subsubsec:RingStructure}
中提到过的，位于不同点 $\vec x_i\in \IR^3$、保持 $osp(4^*\vert 2)_\zeta$ 的平行超对称线算子可以相乘，并重新展开为线算子之和。这个运算定义了一个环：它应当是交换的，因为 $\IR^3$ 中没有自然的点序；它也应当是结合的，因为它们都是算子。自然要问：这个环结构与带框保护自旋特征标、带框 BPS 简并度之间有什么关系？结果表明，方程 \eqref{eq:form-gen} 中定义的生成函数 $F(L)$ 很适合讨论这个问题。在本节中我们将看到，可以用它们来定义线算子环的一个非常有意思的非交换形变。

假设我们把一个线算子 $L$ 放在 $\vec x=0$，另一个 $L'$ 放在沿 $z$ 轴平移的 $\vec x' = (0,0,z)$。注意由 $J_3$ 生成的旋转是这个位形的对称性，但 $so(3)$ 对称性已经破缺为 $so(2)$。若 $z\ne 0$，且 $L$ 与 $L'$ 按第 2.3 节% TODO 合稿改 \ref{subsec:LabelsForLines}
所解释的意义彼此互局域，则希尔伯特空间 $\CH_{L_{\vec x} L'_{\vec x'}}$ 是 $U(1)_{\vec x \vec x'}\subset SU(2)$ 的一个表示。若 $L$ 与 $L'$ 是单对象，则希尔伯特空间 $\CH_{LL'}$ 也按 $\Gamma + \gamma_L + \gamma_{L'}$ 分次，且分次中的各直和项满足
\begin{equation}\label{eq:LL-grade}
\CH_{L_{\vec x} L'_{\vec x'},\gamma_0}^{\rm BPS} = \bigoplus_{\gamma+ \gamma'=\gamma_0} \CH_{L,\gamma}^{\rm BPS}
 \otimes \CH_{L',\gamma'}^{\rm BPS}
 \otimes N_{\gamma,\gamma'}.
\end{equation}
这里 $N_{\gamma,\gamma'}$ 是 $U(1)_{\vec x \vec x'}$ 的一维表示。它来自这对双荷子所激发的电磁场。因此它对保护自旋特征标的贡献是\footnote{这一点应当与「原味」BPS 态的本原穿墙公式中的情形相对照。在那种情形，电磁场具有 $SU(2)$ 自旋 $2j = \vert \langle \gamma, \gamma' \rangle \vert -1$。多出来的这个 $-1$ 来自两个分离组分的质心自由度的对齐。}
\begin{equation}
{\Tr}_{N_{\gamma,\gamma'}} y^{2J_3} = y^{\langle \gamma, \gamma' \rangle}.
\end{equation}
回想生成函数 $F(L)$ 的定义（\eqref{eq:form-gen}）以及海森堡关系 \eqref{eq:Heis-Alg}，我们看到方程 \eqref{eq:LL-grade} 可以用来定义线算子的一个\emph{非交换}乘积：
\begin{equation}\label{eq:Non-Comm-Line-Prod}
F(L\circ_y L') := F(L) F(L').
\end{equation}
一般而言 $F(L)F(L') \ne F(L') F(L)$，从下面讨论的例子中可以看得很清楚。

在第 2.2 节% TODO 合稿改 \ref{subsubsec:RingStructure}
中，我们描述了线算子上的一个交换乘积。因此，发现这个乘积还有一个非交换形变，或许会令人意外。区别在于：是先取算子的 OPE $\lim_{z\to 0} LL'$，得到线算子之和；还是先在希尔伯特空间上取迹。这两种做法并不对易：
\begin{equation}
\lim_{z\to 0} F(LL')\ne F(\lim_{z\to 0} LL').
\end{equation}

为了更好地理解交换乘积与非交换乘积之间的关系，让我们假设存在一族单对象 $L_i$，它们生成全部线算子。例如，第 2.3 节% TODO 合稿改 \ref{subsec:LabelsForLines}
中讨论的任一族线算子 $\CL$ 都可以。单对象具有这样的性质：任何线算子都可以写成 $\sum c_i L_i$，其中 $c_i$ 是正整数。为定义环结构，只需计算融合系数
\begin{equation}\label{eq:simp-mult}
\lim_{z\to 0} L_i L_j = \sum_k c_{ij}^k L_k.
\end{equation}
然而，希尔伯特空间的类似分解取如下形式
\begin{equation}
\CH_{L_i L_j}= \oplus_k  N_{ij}^k \otimes \CH_{L_k}
\end{equation}
这是 $U(1)_{\vec x \vec x'}$ 表示的等式。令人意外的是（至少令作者们意外），在某些例子中 $N_{ij}^k$ 可以是 $U(1)_{\vec x \vec x'}$ 的非平凡表示。我们引入
\begin{equation}
c_{ij}^k(y) := {\Tr}_{N_{ij}^k} y^{2J_3}
\end{equation}
这是 $y$ 的一个洛朗多项式，系数为正整数，并满足 $c_{ij}^k(1) = c_{ij}^k$。用这些记号，线算子的交换环 \eqref{eq:simp-mult} 被形变为
\begin{equation}
F(L_i\circ_y L_j) = \sum_{k} c_{ij}^k(y) F(L_k)
\end{equation}
由于 $\CN=2$ 理论是宇称不变的，我们有 $c_{ij}^k(y)= c_{ji}^k(1/y)$。注意因此当 $y^2=1$ 时乘积 $\circ_y$ 是交换的。

为了理解为什么 $N_{ij}^k$ 必须是 $SO(2)$ 的非平凡表示，我们在认定两个分离线算子系统的正确量子数 $J_3$（从而 $F = 2 J_3 \,{\rm mod}\, 2$）时需要谨慎一些。首先，为了让 $J_3$ 能够被定义，我们需要把两个算子沿 $z$ 轴靠拢。此外，单个线算子的 $J_3$ 定义是明确的（常数平移会与非阿贝尔旋转群不相容）；但对含两个缺陷的系统，未破缺的旋转群是阿贝尔的，因此在比较 $LL'$ 的带框 PSC 定义与各个线算子 $L_i$ 的带框 PSC 定义中出现的 $J_3$ 时，可能会出现整体的常数平移。这种平移歧义由重数 $N_{ij}^k$ 的 $SO(2)$ 分次所体现。

注意，我们的乘积公式通过了一项非常严格的检验，即与穿墙的相容性：只要乘积按规则 $X_\gamma X_{\gamma'} = y^{\langle \gamma,\gamma' \rangle} X_{\gamma+\gamma'}$ 进行，穿墙变换的作用就与乘积可交换。

% =====================================================================
%  GMN3 §4.1--4.3  形式线算子及其非凡的穿墙性质
%  原文：Lamination-Paper.tex 2021--2606 行；PDF p.28--38（印刷 p.27--37）
%  编号公式 (4.1)--(4.29) 共 29 条，全部在本块内（块内引用一律 \eqref）；
%  图 9 幅（4.1--4.9），图名照抄原文 \insfig 第一参数。
%  跨块引用（§3.4、§4.4、§5.3、§12、式 3.29/3.38/3.48/3.49）暂写原文编号并加 % TODO。
% =====================================================================

% —— 原文宏 \ch / \hW / \qbinom 在 preamble.tex 中尚未定义，此处临时补上；
%    合稿时应移入 preamble.tex，并删除本块这三行（\providecommand 不会覆盖已有定义）。

\section{形式线算子及其非凡的穿墙性质}\label{sec:Formal-line-operators}

halo 穿墙图像与强正性猜想结合起来，威力十分强大。强正性不给自旋特征标中的相消留下任何余地：$y^m X_\gamma$ 的系数是希尔伯特空间的真实维数，并按红外电荷与 $so(3)$ 自旋分次。当我们改变 $\zeta$、改变理论的参数与真空期望值时，halo 穿墙会具体地增添或删去这些希尔伯特空间的子空间，但维数必须始终非负。这是一个非常强的约束；在本小节中，我们想就几个具体的 $\CN=2$ 理论探讨它的推论，这些理论具有简单且已被充分理解的「原味」BPS 谱。

\subsection{腔与生成函数}\label{subsec:Chambers}

在本小节中，我们作一些形式定义，用以概括生成函数 (3.38)% TODO \eqref{eq:form-gen}
的行为。

引入非交换环 $\IZ[y,y^{-1},X_\gamma]/I$，其中理想 $I$ 由关系 $yX_\gamma = X_\gamma y$ 与 $X_\gamma X_{\gamma'} = y^{\langle\gamma,\gamma'\rangle}X_{\gamma+\gamma'}$ 生成。现在回顾式 (3.49)% TODO \eqref{eq:Chamber-Space}
所定义的空间 $\Xi$。正如那里所说，$\Xi$ 被划分成若干腔。我们用 $c$ 标记这些腔，其中 $c$ 取遍某个指标集。我们定义\emph{强正形式线算子}为一族元素 $F(c)\in\IZ[y,y^{-1},X_\gamma]/I$，使得

\begin{enumerate}

\item 在腔 $c^+,c^-$ 之间的墙 $\hW(\gamma)$ 两侧，
\begin{equation}
F(c^+) = S_\gamma F(c^-) S_\gamma^{-1},
\end{equation}
其中 $c^+$、$c^-$ 分别是 $\Im(Z_\gamma/\zeta)$ 取正值、取负值的腔，$S_\gamma$ 定义于式 (3.48)% TODO \eqref{eq:Spin-KS-ii}
。

\item 在 $\Xi$ 的每个腔 $c$ 中，
\begin{equation}\label{eq:finite-positive-sum}
F(c) = \sum_\gamma P^c_\gamma(y) X_\gamma
\end{equation}
其中 $P^c_\gamma(y)$ 是 $SU(2)$ 的某个（真实的、而非虚拟的）表示的特征标。

\end{enumerate}

上面第二项定义了「强正」条件。我们注意到，我们给出的定义与 Fock 和 Goncharov 在研究簇系综（cluster ensemble）时引入的\emph{泛洛朗多项式}（universal Laurent polynomial）\cite{MR2567745} 密切相关。

若强正性猜想成立，则线算子的生成函数 $F(L)$ (3.38)% TODO \eqref{eq:form-gen}
是一个强正形式线算子。我们把形式线算子记作 $F$。与真实的线算子一样，在考虑线算子族时，$F$ 会承受紫外参数中的单值变换。当然，如果我们在 $\CB\times\IC^*$ 上、而非在其万有覆盖上定义腔，那么在这个空间里也会有单值变换。

在接下来的几个小节中，我们将考察一些简单 $\CN=2$ 理论中形式线算子的例子。在这些情形中，我们会发现腔结构有一些重要的简化。例如，我们将只考虑只涉及单个超多重态的墙。此时第 3.4 节% TODO \ref{subsec:Halo-Wall-Crossing}
的全套量子双对数技术有点小题大做——我们只需直接应用式 (3.29)% TODO \eqref{eq:transform}
，按需移除或添加 halo。单个超多重态的穿墙规则——\emph{从} ${\rm Im}(Z_{\gamma_h}/\zeta)<0$ 的区域\emph{到} ${\rm Im}(Z_{\gamma_h}/\zeta)>0$ 的区域——算出来简单地就是
\begin{equation}\label{eq:explicit-dilog-conj}
X_{\gamma_c} \to \begin{cases}
\sum_{j=0}^N \ch\left(\Lambda^j \rho_N\right) X_{\gamma_c + j\gamma_h} & \langle\gamma_h,\gamma_c\rangle = N > 0 \\[2pt]
\sum_{j=0}^\infty (-1)^j \ch\left(S^j \rho_N\right) X_{\gamma_c + j\gamma_h} & \langle\gamma_h,\gamma_c\rangle = -N < 0
\end{cases}
\end{equation}
这里 $\rho_N$ 是 $SU(2)$ 的 $N$ 维表示。式 \eqref{eq:explicit-dilog-conj} 的第一行显然保持正性。然而第二行表明，要维持正性是高度非平凡的。注意在第一种情形中，内积的符号使我们进入一个支持核心 $\gamma_c$ 周围 $\gamma_h$ 粒子 halo 的区域；而在第二种情形中，我们进入的区域不支持这样的 halo 构型。\emph{从} ${\rm Im}(Z_{\gamma_h}/\zeta)>0$ 的区域\emph{到} ${\rm Im}(Z_{\gamma_h}/\zeta)<0$ 的区域的变换规则同样算得
\begin{equation}\label{eq:explicit-dilog-conj-ii}
X_{\gamma_c} \to \begin{cases}
\sum_{j=0}^\infty (-1)^j \ch\left(S^j \rho_N\right) X_{\gamma_c + j\gamma_h} & \langle\gamma_h,\gamma_c\rangle = N > 0 \\[2pt]
\sum_{j=0}^N \ch\left(\Lambda^j \rho_N\right) X_{\gamma_c + j\gamma_h} & \langle\gamma_h,\gamma_c\rangle = -N < 0
\end{cases}
\end{equation}
同样，当我们进入支持 halo 构型的区域时，正性明显地保持；否则不然。当然，先应用式 \eqref{eq:explicit-dilog-conj} 再应用式 \eqref{eq:explicit-dilog-conj-ii} 必须给出恒等算子。这蕴涵一个多少非平凡的恒等式
\begin{equation}\label{eq:Char-Ident}
\sum_{j+k=\ell} (-1)^j \ch\left(S^j \rho_N\right) \ch\left(\Lambda^k \rho_N\right) = \delta_{\ell,0},
\end{equation}
对整数 $\ell\geq 0$ 成立。这就是在用式 \eqref{eq:explicit-dilog-conj} 从支持 halo 构型的区域过渡到不支持的腔时，halo 被「移除」的方式。

关于式 \eqref{eq:explicit-dilog-conj} 与式 \eqref{eq:explicit-dilog-conj-ii}，我们想作最后一点评论。注意当表达式有限时，它明显地是强正的；而当它无限时则不然。这提示：$X_\gamma$ 展开的有限性与强正性相关。在第 \ref{subsec:Formal-line-ops-U1} 节与第 \ref{sec:formaln3} 节中，我们将在一些简单场论中构造强正形式线算子。然而在第 4.4 节% TODO \ref{subsec:formal-SU2}
中，我们将处理一个更复杂的场论，其中我们难以证明所构造的形式线算子的强正性。尽管如此，受所有显式例子与上述观察的启发，我们想猜想：\emph{若一个形式线算子在单个腔中是强正的，并且在所有腔中都有有限的达布展开，那么它事实上就是强正的。}

\subsection{形式线算子与 $U(1)$ 规范理论}\label{subsec:Formal-line-ops-U1}

有一个方便的玩具模型，它抓住了形式线算子行为的若干关键方面：带单个电荷为 $1$ 的超多重态的 $U(1)$ 规范理论。这个理论在紫外没有良好定义，但它确实抓住了良定义理论在库仑分支 $\CB$ 的奇点附近的行为。

在此例中，$\CB^*$ 是去心圆盘，$\Gamma$ 是秩二的局部系统。$\Gamma$ 的纤维由 $\gamma_1,\gamma_2$ 生成，$\langle\gamma_1,\gamma_2\rangle = +1$，绕原点的逆时针单值变换作用为
\begin{equation}
\gamma_1 \to \gamma_1+\gamma_2, \qquad \gamma_2 \to \gamma_2.
\end{equation}
我们把 $p\gamma_1+q\gamma_2$ 记作 $(p,q)$。

我们可以取 $Z_{\gamma_2}(u)=u$。我们不需要 $Z_{\gamma_1}(u)$ 的显式公式，因为非零的「原味」BPS 简并度只有 $\Omega(\pm\gamma_2)=+1$。「原味」BPS 简并度不发生穿墙。

我们预期能够找到对应于各种电荷与磁荷的 Wilson--'t Hooft 线算子的线算子。由于带电粒子的存在带来单圈 beta 函数，我们预期会看到反常 R 对称性的效应：一次 R 对称变换要求 $\theta$ 角平移，而后者又通过 Witten 效应引起带磁荷对象电荷的平移。这导致线算子标签在 $\zeta$ 旋转 $2\pi$ 之下发生单值变换。这种单值变换与穿墙之间会有漂亮的相互作用。

万有覆盖 $\widehat\CB \times \widehat\IC^*$ 就是 $\IC\times\IC$，坐标为 $(\log u,\log\zeta)$。给定我们的「原味」BPS 谱，墙是
\begin{equation}
\hW(\gamma_2) = \amalg_{n\in\IZ}\left\{(\log u,\log\zeta): \arg u-\arg\zeta = (2n+1)\pi\right\},
\end{equation}
\begin{equation}
\hW(-\gamma_2) = \amalg_{n\in\IZ}\left\{(\log u,\log\zeta): \arg u-\arg\zeta = 2n\pi\right\},
\end{equation}
从而腔可以取为
\begin{equation}
c_n := \left\{(\log u,\log\zeta): n\pi < \arg u-\arg\zeta < (n+1)\pi\right\}.
\end{equation}

让我们现在尝试构造形式线算子 $F^{(n)}_{p,q}$：规定它在腔 $c_n$ 中的值为 $F^{(n)}_{p,q}(c_n) = X_{p,q}$。这就是电荷为 $q$、磁荷为 $p$ 的 Wilson--'t Hooft 线算子的朴素生成函数。变换法则 \eqref{eq:explicit-dilog-conj} 与 \eqref{eq:explicit-dilog-conj-ii} 随即确定 $F^{(n)}_{p,q}$ 在所有其他腔中的取值，于是我们可以问：它们满足强正性约束吗？

首先，若 $p=0$ 则没有变换，故对所有腔 $c$ 都有 $F^{(n)}_{0,q}(c)=X_{0,q}$。这当然满足形式正性的约束。由于它们实际上与 $n$ 无关，我们把这些形式线算子简单记作 $F_{0,q}$。它们是电荷为 $q$ 的 Wilson 线算子，没有有趣的穿墙，因为「原味」粒子只带电荷。

现在注意，在稳定性条件 $\langle\gamma_h,\gamma_c\rangle\,{\rm Im}Z_{\gamma_h}/\zeta > 0$ 中，我们只需要考虑 $\gamma_h=\pm\gamma_2$，而且事实上对这两个电荷该条件相同。若 $\gamma_c=(p,q)$，则稳定性条件变为 $-p\,{\rm Im}Z_{\gamma_2}/\zeta>0$。因此当 $p>0$ 时，$c_n$ 在 $n$ 为奇数时可以有 halo 构型、在 $n$ 为偶数时不能；而当 $p<0$ 时，$c_n$ 在 $n$ 为偶数时可以有 halo 构型、在 $n$ 为奇数时不能。见图 \ref{fig:chamber-u1-even} 与图 \ref{fig:chamber-u1-odd}。

\insfig{chamber-u1-even}{最简单的腔与强正形式线算子：带单个 BPS 超多重态的 $U(1)$ 规范理论，此处 $n=0\bmod 2$、$p>0$。对 halo 形成稳定的区域用阴影标出。}
\insfig{chamber-u1-odd}{最简单的腔与强正形式线算子：带单个 BPS 超多重态的 $U(1)$ 规范理论，此处 $n=1\bmod 2$、$p<0$。对 halo 形成稳定的区域用阴影标出。}

在分析 $F^{(n)}_{p,q}$ 的行为时，有四种情形要考虑：$n$ 可为偶数或奇数，$p$ 可正可负。

\begin{enumerate}

\item $n=0\bmod 2$，$p>0$。此时我们从 $c_n$ 走到 $c_{n+1}$ 会穿过类型为 $\hW(\gamma_2)$ 的墙。我们应用式 \eqref{eq:explicit-dilog-conj-ii} 的第二行，得到
\begin{equation}\label{eq:n-even-p-pos}
F^{(n)}_{p,q}(c_{n+1}) = \sum_{j=0}^p \qbinom{p}{j} X_{p,q+j}.
\end{equation}
同样，从 $c_n$ 走到 $c_{n-1}$ 我们会穿过类型为 $\hW(-\gamma_2)$ 的墙。此时应用式 \eqref{eq:explicit-dilog-conj} 的第一行，得到
\begin{equation}
F^{(n)}_{p,q}(c_{n-1}) = \sum_{j=0}^p \qbinom{p}{j} X_{p,q-j}.
\end{equation}
当我们从 $c_{n+1}$ 继续走到 $c_{n+2}$ 时，本可以应用我们的形式规则并使用恒等式 \eqref{eq:Char-Ident}，但反推回去更简单：回忆稳定性条件不允许腔 $c_{n+2}$ 中有 halo。因此，从 $c_{n+2}$ 穿过墙 $\hW(-\gamma_2)$ 到 $c_{n+1}$ 时，我们必须是在产生一个由电荷 $-\gamma_2$ 的粒子组成的 halo。这个操作确实会给出 \eqref{eq:n-even-p-pos}，当且仅当
\begin{equation}\label{eq:u1-monod}
F^{(n)}_{p,q}(c_{n+2}) = X_{p,q+p}.
\end{equation}
类似地，继续走到腔 $c_{n-2}$，我们得到 $F^{(n)}_{p,q}(c_{n-2})=X_{p,q-p}$。这个模式显然向两个方向延续下去，因此 $F^{(n)}_{p,q}$ 是一个形式线算子。此外，此式表明 $F^{(n)}_{p,q} = F^{(n+2)}_{p,q+p}$，因此不失一般性，我们可以把这些算子归约为由腔 $c_0$ 中的值确定的单一族。$n\to n+2$ 时电荷的平移有两种独立的解释：若把它解释为 $\zeta$ 变量单值变换的结果，那它就是紫外标签的单值性 $F_{p,q}\to F_{p,q+p}$；而若把它解释为 $u$ 平面中单值变换的结果，它表达的则是红外电荷格 $\Gamma$ 绕奇异轨迹 $u=0$ 的单值性。

最后，这个例子说明：在把 $F$ 按 $X_\gamma$ 展开时，代表核心的项与代表 halo 的项之间没有不变的区别。的确，在从 $c_n$ 穿过 $\hW(\gamma_2)$ 到 $c_{n+1}$ 时，核心电荷是 $(p,q)$；而当从 $c_{n+2}$ 穿过 $\hW(-\gamma_2)$ 到 $c_{n+1}$ 时，核心电荷是 $(p,q+p)$。这与我们的物理图像一致：核心与 halo 的区别，只有当我们非常靠近一道墙时才变得分明。

\item $n=0\bmod 2$，$p<0$。此时稳定性条件允许在腔 $c_n$ 中有 halo、而在 $c_{n\pm1}$ 中不允许，所以我们预期会有麻烦，因为当我们试图继续到相邻的腔 $c_{n\pm1}$ 时不能有 halo 构型。的确，此时穿过墙 $\hW(\gamma)$ 进入腔 $c_{n+1}$ 必须应用式 \eqref{eq:explicit-dilog-conj-ii} 的第 1 行。用这个变换作用于 $X_{p,q}$，显然给出违反强正性的表达式。因此 $F^{(n)}_{p,q}$ 不是形式线算子。

\item $n=1\bmod 2$，$p>0$。这与 $(n=0\bmod 2,\,p<0)$ 的情形类似，不给出形式线算子。

\item $n=1\bmod 2$，$p<0$。这与 $(n=0\bmod 2,\,p>0)$ 的情形类似。它确实给出形式线算子，其中
\begin{equation}\label{eq:n-odd-p-neg}
F^{(n)}_{p,q}(c_{n+1}) = \sum_{j=0}^{|p|} \qbinom{|p|}{j} X_{p,q-j},
\end{equation}
且 $F^{(n)}_{p,q}(c_{n+2})=X_{p,q+p}$。同样有 $F^{(n)}_{p,q}=F^{(n+2)}_{p,q+p}$，故不失一般性可以只考虑 $F^{(1)}_{p,q}$。

\end{enumerate}

总之，由于 $F^{(n)}_{p,q}=F^{(n+2)}_{p,q+p}$，为了构造那些在某个腔中约化为单项式的强正形式线算子，只需考虑腔 $n=0$ 与 $n=1$。所得的算子与整数对 $(p,q)$ 一一对应，因此我们可以直接把它们记作 $F_{p,q}$。若 $p=0$，则 $F_{0,q}=X_{0,q}$ 在所有腔中成立。若 $p>0$，则 $F_{p,q}=X_{p,q}$ 在腔 $c_0$ 中成立；若 $p<0$，则 $F_{p,q}=X_{p,q}$ 在腔 $c_1$ 中成立。这些显然是单（simple）线算子，而且事实上它们是唯一的单线算子：给定任何形式线算子，我们可以把它拆成磁荷为正的部分、不带磁荷的部分与只带负磁荷的部分。显然，磁荷为正的部分只需在某个 $n$ 为偶数的腔中分解成单项式，就能分解为单的构件。对磁荷为负的部分（取 $n$ 为奇数）以及对不带磁荷的部分（任意 $n$），类似的说法也成立。形式线算子 $F_{p,q}$ 应当与 $U(1)$ 理论中 't Hooft--Wilson 圈的保护自旋特征标等同起来。

即使在这个最简单的例子里，形式线算子的代数也已经有些非平凡，虽然它还简单到可以显式写出。的确，$F_{\pm1,0}$、$F_{0,\pm1}$ 已经生成一个环。人们容易手算出关系 $F_{1,0}F_{-1,0}=1+yF_{0,1}$，但更高阶的乘法涉及对 $q$-二项式系数的繁琐操作。环关系可以如下容易地算出。首先，很容易验证
\begin{equation}
F_{p,q}F_{r,s} = y^{ps-qr}F_{p+r,q+s}
\end{equation}
在 $p\geq 0,r\geq 0$ 或 $p\leq 0,r\leq 0$ 时成立，因为此时可以在算子为单项式的腔 $c_0$ 或 $c_1$ 中求值。$p,r$ 异号的情形要非平凡得多。为确定起见，设 $p\geq 0$、$r\leq 0$。那么按定义：
\begin{equation}
F_{p,q}F_{r,s} = \begin{cases}
X_{p,q}\left(\sum_{j=0}^{|r|}\qbinom{|r|}{j} X_{r,s+j}\right) & {\rm in}~~c_0 \\[2pt]
\left(\sum_{j=0}^p \qbinom{p}{j} X_{p,q+j}\right) X_{r,s} & {\rm in}~~c_1
\end{cases}
\end{equation}
于是执行乘法即得：
\begin{equation}
F_{p,q}F_{r,s} = \begin{cases}
y^{ps-qr}\sum_{j=0}^{|r|}\qbinom{|r|}{j} y^{pj} X_{p+r,q+s+j} & {\rm in}~~c_0 \\[2pt]
y^{ps-qr}\sum_{j=0}^p \qbinom{p}{j} y^{-rj} X_{p+r,q+s+j} & {\rm in}~~c_1
\end{cases}
\end{equation}
现在，要辨认右端对应于哪个线算子，我们应当把结果放在「磁荷为 $p+r$ 的线算子为单项式」的那个腔中来看。于是得到线算子上的乘法：
\begin{equation}\label{eq:U1-ring-rel}
F_{p,q}F_{r,s} = \begin{cases}
y^{ps-qr}\sum_{j=0}^{|r|}\qbinom{|r|}{j} y^{pj} F_{p+r,q+s+j} & p+r \geq 0 \\[2pt]
y^{ps-qr}\sum_{j=0}^p \qbinom{p}{j} y^{-rj} F_{p+r,q+s+j} & p+r \leq 0
\end{cases}
\end{equation}
我们可以容易地得到 $p\leq 0$、$r\geq 0$ 的另一情形：回忆 $F(L_1)F(L_2)$ 等于把 $F(L_2)F(L_1)$ 作替换 $y\to 1/y$。对 \eqref{eq:U1-ring-rel} 的暴力验算蕴涵关于 $y$-二项式的若干非平凡恒等式（$y$-形变的 Pascal 关系），
\begin{equation}
y^j \qbinom{p}{j} = \qbinom{p-1}{j} + y^p \qbinom{p-1}{j-1}.
\end{equation}

\subsubsection*{更高电荷}

此构造的一个自然变形，是考虑与单个电荷为 $Q$ 的超多重态耦合的 $U(1)$ 理论。也就是说，我们现在考虑 BPS 谱
\begin{equation}\label{eq:Q-BPS-spec}
\Omega(\gamma) = \begin{cases} 1 & \gamma = \pm Q\gamma_2, \\ 0 & {\rm else}. \end{cases}
\end{equation}
我们可以通过一个小技巧立刻得到形式线算子：若把理论的电荷格与磁荷格按相反的因子 $Q$ 重新标度，我们就回到 $Q=1$ 的情形。上面的代数关系适用于 $F'_{p',q'}$，其中 $F'_{p',q'} := F_{Q^{-1}p',\,Qq'}$。名副其实的线算子必须带整数个电荷。于是我们现在得到关系
\begin{equation}\label{eq:Q-U1-algebra}
F_{1,0}F_{-1,0} = \sum_{j=0}^Q \qbinom{Q}{j} y^{Qj} F_{0,Qj}.
\end{equation}

\subsection{形式线算子与 $N=3$ Argyres--Douglas 理论}\label{sec:formaln3}

我们现在描述一个带穿墙的理论的最简单例子中的形式线算子。这就是 $N=3$ Argyres--Douglas 理论。关于这个理论的更多细节，见文献 \cite{Gaiotto:2009hg} 第 9.4 节。简言之，$\CB = \IC-\{u_+,u_-\}$。万有覆盖已经相当复杂，所以我们将在 $u$ 平面上的确定区域工作，并考虑（$\zeta$ 平面万有覆盖中的）腔。对局部系统最好的描述是 $\Gamma = H_1(\Sigma_u,\IZ)^-$，其中 $\Sigma_u$ 是由
\begin{equation}
\lambda^2 = (z^3-3\Lambda^2 z+u)(dz)^2
\end{equation}
定义的、$\CB$ 上一族黎曼面的纤维，其中 $z\in\IC$。纤维 $\Gamma_u$ 是秩二的，我们取生成元 $\gamma_1,\gamma_2$ 使 $\langle\gamma_1,\gamma_2\rangle=1$。向量 $p\gamma_1+q\gamma_2$ 有时记作 $(p,q)$。在 $u=u_\pm=\pm 2\Lambda^3$ 处有两个奇点，$\Gamma$ 绕它们有单值变换。（见文献 \cite{Gaiotto:2009hg} 的式 (9.30)。）BPS 态的谱已在文献 \cite{Gaiotto:2009hg} 中确定。底空间 $\CB$ 应划分为三个区域 $\CB_s,\CB_w^\pm$，在其上 $\Gamma$ 可以平凡化。$\CB_s$ 是 $u=0$ 附近的区域，$\CB_w^\pm$ 是 $u=\infty$ 附近的两个区域。有一条穿过 $u_\pm$ 的连通墙，把 $\CB_s$ 与 $\CB_w^\pm$ 分开。我们有
\begin{equation}
\Omega(\gamma;u) = \begin{cases} 1 & \gamma = \pm\gamma_1,\pm\gamma_2 \\ 0 & {\rm else} \end{cases}
\end{equation}
对 $u\in\CB_s$，而穿墙随后确定
\begin{equation}
\Omega(\gamma;u) = \begin{cases} 1 & \gamma = \pm\gamma_1,\pm\gamma_2,\pm(\gamma_1+\gamma_2) \\ 0 & {\rm else} \end{cases}
\end{equation}
对 $u\in\CB_w^+$。$\CB_w^-$ 中的谱可由单值变换或穿墙得到。

\insfig[0.42\textwidth]{ADN3-FL-1}{$N=3$ Argyres--Douglas 理论在 $u=0$、$\Lambda>0$ 时的 BPS 墙投影到 $\zeta$ 平面上。}

在 $u=0$ 处对 $\Sigma_u$ 的 cycle 作适当选取，可使 $Z_{\gamma_1} = -i\Lambda^{5/2}K$、$Z_{\gamma_2} = -K\Lambda^{5/2}$，其中 $K>0$ 是一个（可计算的）常数。$u=0$ 处 BPS 墙到 $\zeta$ 平面的投影于是如图 \ref{fig:ADN3-FL-1} 所示。注意，当我们沿\emph{顺时针}方向行进时，穿过每道墙 $\hW(\gamma_h)$ 都是从 ${\rm Im}(Z_{\gamma_h}/\zeta)>0$ 的区域走到 ${\rm Im}(Z_{\gamma_h}/\zeta)<0$ 的区域。因此沿顺时针方向穿墙时应使用规则 \eqref{eq:explicit-dilog-conj-ii}，而沿逆时针方向穿墙时应使用规则 \eqref{eq:explicit-dilog-conj}。

让我们沿用第 \ref{subsec:Formal-line-ops-U1} 节的策略，尝试找出在 $\zeta$ 平面某个腔中约化为单项式 $X_{p,q}$ 的形式线算子。引入腔 $c_n$，其中 $\frac{\pi}{2}(1-n) < \arg\zeta < \frac{\pi}{2}(2-n)$。第一个障碍来自顺时针或逆时针方向的单次穿墙。这导致如下结果：$X_{p,q}$ 只有在下列情形下才能满足形式正性：

\begin{enumerate}

\item $p\geq 0,\ q\geq 0$，在 $n=1\bmod 4$ 的 $c_n$ 中

\item $p\leq 0,\ q\geq 0$，在 $n=2\bmod 4$ 的 $c_n$ 中

\item $p\leq 0,\ q\leq 0$，在 $n=3\bmod 4$ 的 $c_n$ 中

\item $p\geq 0,\ q\leq 0$，在 $n=0\bmod 4$ 的 $c_n$ 中

\end{enumerate}

我们现在将证明这是唯一的障碍。

\insfig[0.42\textwidth]{ADN3-FL-2}{$5$ 重覆盖的第一叶包含四个腔。}

\insfig[0.42\textwidth]{ADN3-FL-3}{沿顺时针方向继续解析延拓到第二叶。}

\insfig[0.42\textwidth]{ADN3-FL-4}{沿顺时针方向继续解析延拓到第三叶。}

\insfig[0.42\textwidth]{ADN3-FL-5}{沿顺时针方向继续解析延拓到第四叶。}

\insfig[0.42\textwidth]{ADN3-FL-6}{沿顺时针方向继续解析延拓到第五叶。从腔 $c_{20}$ 出发的顺时针穿墙产生 $X_{\gamma_1}$，于是整个过程重复。因此在万有覆盖上，腔 $21$ 与腔 $1$ 等同；在 $5$ 重覆盖上工作时，腔编号应按模 $20$ 等同。}

为确定起见，让我们从万有覆盖 $\widehat{\IC^*}$ 中满足 $0<\arg\zeta<\frac{\pi}{2}$ 的腔里的 $X_{\gamma_1}$ 出发。顺时针方向的穿墙产生图 \ref{fig:ADN3-FL-2} 所示的表达式。我们观察到一些在 $U(1)$ 例子中见过的同样现象。从 $c_1$ 到 $c_2$，我们获得一个核心电荷为 $\gamma_1$ 的 halo。从 $c_2$ 到 $c_3$，我们获得另一个 halo，但这次核心电荷是 $\gamma_1+\gamma_2$、halo 电荷是 $-\gamma_1$。从 $c_3$ 到 $c_4$，我们失去一个 halo 粒子，核心电荷仍解释为 $\gamma_1+\gamma_2$，但 halo 粒子的电荷为 $-\gamma_2$。现在，若从这个腔继续沿顺时针方向走，我们会生成一个新的表达式 $X_{\gamma_2}$。我们又从一个核心为 $\gamma_2$ 的 halo 中移走了一个 halo 粒子。不出所料，我们发现存在单值变换，因此确实必须在万有覆盖 $\widehat{\IC^*}$ 上讨论穿墙。这样做时，我们得到示意的模式：
\begin{equation}
\cdots \to X \to X \to X+X \to X+X+X \to X+X \to X \to X \to X+X \to \cdots
\end{equation}
结果表明，这个模式以 $20=4\times 5$ 为周期重复。用这个周期性对万有覆盖作商，我们得到 $\zeta$ 平面的一个 $5$ 重覆盖。每一叶包含四个腔，共 $20$ 个腔，标记为 $c_1,\dots,c_{20}$。形式线算子在这全部 $20$ 个腔中的显式表达式，由图 \ref{fig:ADN3-FL-2} 至图 \ref{fig:ADN3-FL-6} 给出。

\insfig[0.42\textwidth]{ADN3-FL-weak}{某个弱耦合区域中的腔结构。}

可惜，这还没有完全构造出强正线算子，因为我们必须检查它在弱耦合区域中是强正的。穿过边缘稳定墙进入区域 $\CB_w^+$ 时，$\gamma_1$ 与 $\gamma_2$ 的墙合并并产生第三道墙，导致 $\zeta$ 平面中出现图 \ref{fig:ADN3-FL-weak} 所示的新腔结构。每个奇数腔都分裂成两个，于是每一叶上有 $6$ 个腔，共 $30$ 个腔。注意偶数编号腔中的表达式不可能从强耦合区域发生改变。很容易验证：穿墙在奇数腔 $c_{2n+1,\pm}$ 中产生强正的元素。

于是我们已完全构造出一个强正形式线算子。利用单值变换操作，我们事实上可以得到 $5$ 个不同的形式线算子，它们由下式定义
\begin{equation}
\begin{split}
F_1(c_1) & = X_{\gamma_1} \\
F_2(c_1) & = X_{\gamma_2} \\
F_3(c_1) & = X_{-\gamma_1} + X_{\gamma_2-\gamma_1} \\
F_4(c_1) & = X_{-\gamma_1} + X_{-\gamma_1-\gamma_2} + X_{-\gamma_2} \\
F_5(c_1) & = X_{-\gamma_2} + X_{\gamma_1-\gamma_2}
\end{split}
\end{equation}

利用这些算子，我们可以如下产生若干其他强正算子。考虑 $F_{p,q}:=y^{-pq}F_1^p F_2^q$，其中 $n,m\geq 0$\zy{原文此处作「对 $n,m\geq 0$」，但式中用的是 $p,q$；应为 $p,q\geq 0$，译文照排。}。在 $u$ 处于强耦合区的腔 $c_1$ 中，它的值为 $X_{(p,q)}$。我们断言这生成一个强正形式线算子。直接检验会迅速变得非常困难，但我们可以用下述技巧来证明。首先注意，若沿顺时针或逆时针方向移动两个腔，穿墙会直截了当地产生强正的表达式。的确，我们得到显式表达式：
\begin{equation}\label{eq:Fpq-chs-N3}
\begin{split}
F_{p,q}(c_{19}) & = \sum_{j=0}^q \sum_{k=0}^{p+j} \qbinom{q}{j}\qbinom{p+j}{k} X_{p+j,q-k} \\
F_{p,q}(c_{20}) & = \sum_{j=0}^q \qbinom{q}{j} X_{p+j,q} \\
F_{p,q}(c_1) & = X_{p,q} \\
F_{p,q}(c_2) & = \sum_{j=0}^p \qbinom{p}{j} X_{p,q+j} \\
F_{p,q}(c_3) & = \sum_{j=0}^p \sum_{k=0}^{q+j} \qbinom{p}{j}\qbinom{q+j}{k} X_{p-k,q+j}
\end{split}
\end{equation}
当然，由于 $q$-二项式是真实的自旋特征标，它们的乘积也是。因此 \eqref{eq:Fpq-chs-N3} 满足强正性。现在注意，$F_1$ 在 $8$ 个腔 $c_{20},c_1$、$c_5,c_6$、$c_{10},c_{11}$、$c_{15},c_{16}$ 中成为单项式。类似地，$F_2$ 在 $8$ 个腔 $c_1,c_2$、$c_6,c_7$、$c_{11},c_{12}$、$c_{16},c_{17}$ 中成为单项式。于是在腔 $c_1,c_6,c_{11},c_{16}$ 中，表达式 $F_{p,q}$ 是单项式。具体地有
\begin{equation}\label{eq:Fpq-mon-chs}
\begin{split}
F_{p,q}(c_6) & = X_{(-q,p)} \\
F_{p,q}(c_{11}) & = X_{(-q,-p)} \\
F_{p,q}(c_{16}) & = X_{(q,-p)}
\end{split}
\end{equation}
现在，对 \eqref{eq:Fpq-mon-chs} 中的每个表达式施加与 \eqref{eq:Fpq-chs-N3} 类似的两步顺时针与逆时针变换，可以验证穿墙产生强正的表达式。这样我们就能覆盖强耦合区的全部二十个腔。第 5.2 节% TODO \ref{subsec:Seeds-from-N=2}
末尾给出的论证蕴涵：把它们延拓为弱耦合区域的强正算子没有障碍。

类似地，我们可以从其他线算子 $F_i$ 的乘积生成强正形式线算子。把上面的 $F_{p,q}$ 改记为 $F^{(1)}_{p,q}$，我们可以定义（各种情形下都有 $p\geq 0,q\geq 0$）：

\begin{enumerate}

\item $F^{(2)}_{p,q} = y^{-pq} F_2^p F_3^q$，它在 $c_2,c_7,c_{12},c_{17}$ 中是单项式

\item $F^{(3)}_{p,q} = y^{-pq} F_3^p F_4^q$，它在 $c_3,c_8,c_{13},c_{18}$ 中是单项式

\item $F^{(4)}_{p,q} = y^{-pq} F_4^p F_5^q$，它在 $c_4,c_9,c_{14},c_{19}$ 中是单项式

\item $F^{(5)}_{p,q} = y^{-pq} F_5^p F_1^q$，它在 $c_5,c_{10},c_{15},c_{20}$ 中是单项式

\end{enumerate}

我们是否应该预期 $F^{(i)}_{p,q}$ 穷尽了单的强正形式线算子之列表？我们可以勾画一个证明，它可以用第 12 节% TODO \ref{sec:Tropic}
给出的工具形式化。考虑某个形式线算子 $F$ 在腔 $c_1$ 中的展开，并看形如 $a_{n_1,n_2}X_{n_1\gamma_1+n_2\gamma_2}$ 的项，其中 $n_1\geq 0$、$n_2\geq 0$。$\tilde F = F - \sum_{n_1>0,\,n_2>0} a_{n_1,n_2}F^{(1)}_{n_1,n_2}$ 是否仍然强正？它在 $c_1$ 中是正的，并且在所有其他腔中都只有有限多项，因此它在那些腔中也必定是正的（因为从一个腔到下一个腔的穿墙只能给一个正的和增添无穷多个负项）。接下来我们可以考虑 $\tilde F$ 在腔 $c_2$ 中的展开，并减去 $F^{(2)}_{p,q}$ 的适当倍数，依此类推。过程结束时，我们得到 $F$ 分解为 $F^{(i)}_{p,q}$ 的一个正和。

现在我们可以考察这些算子生成的代数。在腔 $c_1$ 中计算，容易得到
\begin{equation}
F_1 F_3 = 1 + y F_2
\end{equation}
由于其他算子都由单值变换生成，我们有
\begin{equation}\label{eq:Deformd-N3-rels}
F_{i-1} F_{i+1} = 1 + y F_i
\end{equation}
要得到另一种顺序的乘法，取 $y\to 1/y$。

% =====================================================================
%  GMN3 §4.4--4.5  纯 SU(2) 规范理论中的形式线算子 / 物理 vs. 形式线算子
%  原文：Lamination-Paper.tex 2607--3114 行；印刷页 38--46（PDF p.39--47）
%  编号公式 25 条：原文 2617, 2626, 2637, 2645, 2705, 2716, 2726, 2743,
%    2760, 2778, 2799, 2810, 2820, 2871, 2874, 2887, 2937, 2947, 2960,
%    2976, 2986, 2991, 2998, 3007, 3103 行。
%    （BRIEF 记 26 条：多出的一条是原文 2911 行被 % 注释掉的 eq:8fold-SU2，
%      该段按施工指令不译。）
%  图 3 幅：SU2-FL-weak / SU2-FL-Wilson / SU2-FL-Monomials
%  本块承接 ch05（§4.1--4.3），故不含 \section；式号与图号由 LaTeX 续接。
%  跨块引用一律硬编码原文编号，行尾 % TODO 合稿改 \ref/\eqref。
% =====================================================================

\subsection{纯 $SU(2)$ 规范理论中的形式线算子}
\label{subsec:formal-SU2}

现在转向规范群李代数为 $su(2)$、且不带味的规范理论。乍看之下，它们与 $N=3$ Argyres--Douglas 理论相似得颇具迷惑性。为描述腔结构与 BPS 谱，我们依赖文献 \cite{Gaiotto:2009hg} 的第 10.1 节。
一如既往，$\CB=\IC-\{u_+,u_-\}$，而电荷格的最佳描述是
$\Gamma=H_1(\Sigma_u,\IZ)^-$，其中 $\Sigma_u$ 现在是 $\CB$ 上某一族黎曼面的纤维，该族由
%
\begin{equation}
\lambda^2 = \left(\frac{\Lambda^2}{z^3 } + \frac{2u}{z^2} + \frac{\Lambda^2}{z} \right) (dz)^2.
\end{equation}
%
定义，其中 $z\in \IC^*$。纤维 $\Gamma_u$ 仍是秩二的，向量 $p\gamma_1 + q\gamma_2$ 有时记作 $(p,q)$。与
$N=3$ Argyres--Douglas 理论一个非常重要的差别在于，现在
%
\begin{equation}\label{eq:SU2-innprod}
\langle \gamma_1, \gamma_2 \rangle = 2.
\end{equation}
%
在 $u=u_\pm = \pm \Lambda^2$ 处有两个奇点，$\Gamma$ 绕它们有单值性。（见文献 \cite{Gaiotto:2009hg} 的 (10.4)、(10.5) 两式。）
底空间 $\CB$ 应划分为三个区域 $\CB_s, \CB_{w}^\pm$，在每一区域上 $\Gamma$ 都可平凡化。$\CB_{s}$ 是 $u=0$ 附近的一个区域，$\CB_{w}^\pm$ 是 $u=\infty$ 附近的两个区域。有一道穿过
$u_\pm$ 的连通墙，它把 $\CB_s$ 与 $\CB_{w}^\pm$ 隔开。对 $u \in \CB_s$ 有
%
\begin{equation}
\Omega(\gamma;u) = \begin{cases} 1 & \gamma = \pm \gamma_1, \pm \gamma_2 \\
0 & {\rm else} \\ \end{cases}
\end{equation}
%
内积 \eqref{eq:SU2-innprod}
导致比 $N=3$ 理论更复杂的穿墙：对 $u \in \CB_w^+$，我们发现
%
\begin{equation}\label{eq:SU2-weak-spec}
\Omega(\gamma;u) = \begin{cases} -2 & \gamma = \pm (\gamma_1 +  \gamma_2) \\
1 & \gamma = \pm [\gamma_1 + n(\gamma_1+\gamma_2) ], \quad n \geq 0 \\
1 & \gamma = \pm [\gamma_2 + n(\gamma_1+\gamma_2) ], \quad n \geq 0 \\
0 & {\rm else} \\ \end{cases}
\end{equation}
%
$\CB_w^-$ 中的谱可由单值性或者由穿墙得到。

\insfig{SU2-FL-weak}{纯 $SU(2)$ 理论在弱耦合区域（$u$ 取该区域）的 BPS 墙，投影到 $\zeta$ 平面上。为清晰起见，标记各面墙时省略了 $\widehat{W}$ 上的帽子。在第一象限中，荷为 $n\gamma_1 + (n+1)\gamma_2$
与 $(n+1)\gamma_1 + n\gamma_2$（$n\geq 0$）的墙从两侧向 $\gamma_1 + \gamma_2$ 的墙聚集。除 $\pm (\gamma_1 +\gamma_2)$ 处的中央墙以外，每一面这样的墙上的穿墙都由单个超多重态生成。}

现在可以算出：在 $u \approx 0$ 处存在一组满足 \eqref{eq:SU2-innprod} 的基回路 $\gamma_1, \gamma_2$ 的合适选取，使得 $\zeta$ 平面中的腔结构与图 5% TODO 合稿改 \ref{fig:ADN3-FL-1}
相同。另一方面，在弱耦合下情况截然不同。对固定的 $u \in \CB_w^+$，各腔到 $\zeta$ 平面的投影
见图 \ref{fig:SU2-FL-weak}。可以选一条通向大 $u$ 的路径，把 $\gamma_1, \gamma_2$
与通常的磁荷基、电荷基作比较。由此可以辨认出
$\gamma_1 = -\half H_\alpha +\alpha$ 与 $\gamma_2 =\half H_\alpha$，于是 $\gamma_1+\gamma_2$ 就是
$W$ 玻色子的荷。

设 $u \in \CB_s$。与 $N=3$ 的情形一样，容易检验哪些单项式 $X_{(p,q)}$ 在顺时针与逆时针两个方向的单次穿墙之后仍与强正性相容。结果同样是：

\begin{enumerate}

\item $n=1\mod 4$ 时，$c_n$ 中 $p\geq 0, q\geq 0$

\item $n=2\mod 4$ 时，$c_n$ 中 $p\leq 0, q\geq 0$

\item $n=3\mod 4$ 时，$c_n$ 中 $p\leq 0, q\leq 0$

\item $n=0\mod 4$ 时，$c_n$ 中 $p\geq 0, q\leq 0$

\end{enumerate}

在本节余下部分，我们将证明：在任意穿墙之下，这些单项式都生成 $X_{\gamma}$ 的有限展开。根据第 4.1 节% TODO 合稿改 \ref{subsec:Chambers}
末尾的评注，这在猜想的意义上足以证明强正性。下面我们会为这一假设给出一些合理性论证。利用第 5.2 节% TODO 合稿改 \ref{subsec:Seeds-from-N=2}
末尾的推理\zy{原文 §5.2 与 §5.3 两处共用 label \texttt{subsec:Seeds-from-N=2}，故原文 PDF 把此处的交叉引用印成「第 5.3 节」；按 label 名称与文意应为第 5.2 节，译本取后者，并记入勘误表。}，我们预期：把它们延拓为弱耦合区域中的强正性算子并无障碍。

若我们从某个单项式出发去构造强正性的形式线算子，很快就会发现表达式相当复杂，即使在强耦合区内也是如此。我们将看到，在 $\zeta$ 平面中一般并没有有限的单值性。先从一例开始。方便的记号是
%
\begin{equation}
[a,b]:= X_{\half(a \gamma_1 + b\gamma_2)},
\end{equation}
%
并考虑由 $F(c_1)=[1,0]$ 出发定义的形式线算子。现在沿顺时针方向穿墙，逐个生成其他腔中的值。把各项排成三角阵列会很有启发性：这样一来，就容易看出 halo 如何穿过墙被湮灭。沿顺时针方向，得到 $F(c_2)$ 为：

\begin{equation}\label{eq:SU2-sc-rot1}
\begin{array}{cc}
 &[1,2]\\
 &+ \\
 & [1,0]
\end{array}
\end{equation}
%
$F(c_3)$ 为
%
\begin{equation}\label{eq:SU2-sc-rot2}
 \begin{array}{ccccc}
   [-3,2] & + & \rho_2 [-1,2] &  + & [1,2]\\
    &  & &  & +\\
    & &  & & [1,0]
  \end{array}
 \end{equation}
%
这里为节省书写，我们把一个表示的特征标等同于该表示本身；于是 $\rho_N$ 代表 $[N]$。
当表达式变长时，这很有用。
\eqref{eq:SU2-sc-rot2} 中行与列的各项都可以看成以 $[1,2]$ 为核心电荷的 halo 构型。
接着进入腔 $c_4$，竖直列被消掉，因为一个 halo 被移除；但同时因 halo 的产生而生成另外两个竖直列，于是 $F(c_4)$ 为
%
\begin{equation}
\begin{array}{cccccc}
  &[-3,2] & + & \rho_2 [-1,2] & + & [1,2] \\
     &  + &   & + &   &   \\
 & \rho_3 [-3,0] &   & \rho_2 [-1,0] &  &  \\
  &+ &  &   &   &   \\
  & \rho_3 [-3,-2] &   &   &   &   \\
  &+ &   &  &   &   \\
  &[-3,-4] &   &   &   &
\end{array}
\end{equation}
%
我们再次看到，角上的项 $[-3,2]$ 充当水平 halo 与竖直 halo 的核心粒子。
现在穿过 $\widehat{W(\gamma_1)}$ 进入下一个腔，第一行的水平 halo 塌缩到它的核心，于是得到 $F(c_5)$
%
\begin{equation}
\begin{array}{cccccccccc}
  &[-3,2] &   &   &   &  &   &   &   &   \\
 & + &   &   &   &  &   &   &   &   \\
 & \rho_3 [-3,0] &   & \rho_2 [-1,0] &   &   &   &   &   &   \\
 &  + &   &  + &   &   &   &   &   &   \\
 & \rho_3 [-3,-2] & + & (\rho_2+\rho_4)[-1,-2] & + & \rho_3[1,-2] &   &   &   &   \\
 &  + &   &  + &   &   &   &   &   &   \\
 & [-3,-4] & +& \rho_4 [-1,-4] & + & \Lambda^2 \rho_4 [1,-4] & + & \rho_4 [3,-4] & + & [5,-4]
\end{array}
\end{equation}
%
这个表达式有一个有趣之处：我们必须用到表示之间的恒等式 $\rho_2 \rho_3 = \rho_2 + \rho_4$。为认出
第三行的水平 halo，我们用 $\rho_2 \rho_3$；而为认出
第二列的两个竖直 halo，必须用 $\rho_2 + \rho_4$。作为最后一个例子，
继续沿顺时针穿过 $\widehat{W(\gamma_2)}$，得到 $F(c_6)$ 为
%
\begin{equation}
\begin{array}{ccccccccccc}
    &   &   &   &   &   &   &   &   &   &  [5,6] \\
    &   &   &   &   &   &   &   &   &   &  + \\
    &   &   &   &   &   &   &   &   &   &  \rho_5 [5,4] \\
    &   &   &   &   &   &   &   &   &   &  + \\
    &   &   &   &   &   &   &   &  \rho_4[3,2]  &   &  \Lambda^3\rho_5 [5,2] \\
    &   &   &   &   &   &   &   & +  &   &  + \\
    &   &   &   &   &   &\rho_3 [1,0]   &   & \rho_4\rho_3 [3,0]  &   &  \Lambda^2\rho_5 [5,0] \\
    &   &   &   &   &   &+  &   & +  &   &  + \\
    &   &   &   & \rho_2 [-1,-2]  &   &(\rho_3+\Lambda^2\rho_4) [1,-2]   &   & \rho_4\rho_3 [3,-2]  &   & \rho_5 [5,-2] \\
        &   &   &   &   &   &+  &   & +  &   &  + \\
         &   & [-3,-4]  &  + & \rho_4 [-1,-4]  & +  & \Lambda^2\rho_4 [1,-4]   &  + & \rho_4 [3,-4]  & +  &  [5,-4]
\end{array}
\end{equation}
%
现在完全不清楚的是：穿过墙 $\widehat{W(-\gamma_1)}$ 进入腔 $c_7$ 时，强正性是否仍然保持。有问题的
是倒数第二行。我们必须写 $\rho_3 + \Lambda^2\rho_4 = \rho_2^2 + \rho_5$ 与
$\rho_4 \rho_3 = \rho_5 \rho_2 + \rho_2 $，于是这一行可以写成：
%
\begin{equation}
\rho_2 ([-1,-2] + \rho_2 [1,-2] + [3,-2]) + \rho_5 ([1,-2] + \rho_2 [3,-2] + [5,-2]) .
\end{equation}
%
写成这种形式，就很清楚：水平 halo 将会塌缩，强正性得以保持。

沿逆时针方向推进，我们看到类似的图样。从
$c_1$ 穿过 $\widehat{W(\gamma_1)}$ 到 $c_0$，简单地得到 $F(c_0)=[1,0]$；随后 $F(c_{-1})$ 为
%
\begin{equation}
\begin{array}{cc}
& [1,0] \\
& + \\
& [1,-2]
\end{array}
\end{equation}
%
过渡到 $F(c_{-2})$，得到
%
\begin{equation}
\begin{array}{cccccc}
  &  &   &   &   &  [1,0] \\
  &  &   &   &   & + \\
  & [-3,-2] & + & \rho_2 [-1,-2] & + & [1,-2]
\end{array}
\end{equation}
%
如此等等。

一般的图样如下。从 $n=1 \mod 4$ 的腔 $c_n$ 出发，我们从一个直角三角形开始，
其顶点为 $[a,b], [a-2b,b], [a,b-2a]$，其中 $a,b$ 满足
一些不等式 $a\leq 0, b\leq 0, a-2b \geq 0, b-2a\geq 0, \cdots$。
随后，穿墙在每一个腔中创生并湮灭 halo，形成一个直角三角形阵列，使得
四次穿墙之后我们回到一个三角阵列，其参数为
$a' = - 3a +4 b$、$b' = -4 a + 5b$。在上面的例子中，经 $k$ 次顺时针转动后
$a= 1-4k$、$b= -4k$。直接检验这些三角形内部格点处的强正性
似乎非常不平凡；但只要三角形的图样持续下去，情形必然如此。也就是说，
证明强正性的主要困难在于说明：在穿墙之下，表达式仍保持为
$[a,b]$ 的多项式。事实上，我们将在本节末尾给出一个严格的论证，
说明这种有限性确实成立。

\insfig{SU2-FL-Wilson}{强耦合下的 Wilson 线算子。对这个表达式而言，$\zeta$ 平面中没有单值性。
角上的电荷同时充当水平 halo 与竖直 halo 的核心电荷。
每穿过一面墙，就有一个 halo 被消灭、一个 halo 被创生。}

上述表达式虽已变得相当可观，但存在一个值得注意的特殊情形：在强耦合区，
三角形的参数在 $\zeta$ 平面中没有单值性，
即 $a' = -3a+4b = a $、$b' = -4a + 5 b = b$，
也就是 $a=b$ 且 $a,b<0$。最基本的情形 $a=b=-1$ 将对应于物理理论中基础表示下的 Wilson
线。它在四个腔中的值见图 \ref{fig:SU2-FL-Wilson}。Wilson 算子的
表达式足够简单，以致我们也能在弱耦合区检验强正性。我们从 $c_0$ 与 $c_2$ 中的表达式出发，这两者在
强耦合区与弱耦合区 $\CB_w^+$ 中相同。从 $c_0$ 沿顺时针方向穿过，
在 $n\gamma_1 + (n+1)\gamma_2$ 的墙旁边得到形如 $[1,1]+[2n-1,2n+1] + [-1,-1]$ 的表达式；顺时针穿过这面墙，
一个 halo 被消灭、一个被创生，得到 $[1,1]+[2n+1,2n+3] + [-1,-1]$。
类似地，从 $c_2$ 沿逆时针方向穿过，得到表达式
$[1,1]+[2n+1,2n-1] + [-1,-1]$；再穿过 $(n+1)\gamma_1 + n\gamma_2$ 的墙，
得到 $[1,1]+[2n+3,2n+1] + [-1,-1]$。这样，我们就明确构造出了
对应于 Wilson 圈的形式线算子。

这个 Wilson 线算子（我们记作 $\hat W$）可用来生成一串
形式线算子。我们递归地定义
%
\begin{equation}\label{eq:rec-rels-1}
G_{2k} \hat W = y G_{2k+1} + 1 + y^{-1} G_{2k-1},
\end{equation}
\begin{equation}\label{eq:rec-rels-2}
G_{2k+1}\hat W = y G_{2k+2} + y^{-1} G_{2k}.
\end{equation}
%
递归的出发点是：在腔 $c_1$ 中取 $G_0 = [1,1]$ 的那个算子。（注意，它正是 Wilson
算子的四个顶点所构成的方阵中「缺失」的那个算子。）
回忆 $[a,b]:= X_{\half(a \gamma_1 + b\gamma_2)}$，
于是 $[a,b]\cdot [c,d] = y^{\half(ad-bc)} [a+c,b+d]$。把这些乘法
做出来，就得到一串算子，它们在腔 $c_1$ 中的值为

%
\begin{equation}
\begin{split}
G_{-5}& = [ -2,-4] + \rho_4 [0,-4] + \Lambda^2 \rho_4 [2,-4] + \rho_4 [4,-4] + [ 6,-4] \\
&\quad + \rho_2 ( [-2,-2] + \rho_2 [0,-2]+[2,-2]) + [-2,0]\\
G_{-4}& = [-1,-3] + \rho_3 [1,-3] + \rho_3 [3,-3] + [5,-3] + [-1,-1] + [1,-1] \\
G_{-3} & = [0,-2] + \rho_2 [2,-2] + [4,-2] \\
G_{-2} & = [1,-1] + [3,-1] \\
G_{-1} & = [2,0] \\
G_{0} & = [1,1] \\
G_1 & = [0,2] \\
G_2 & = [-1,1] + [-1,3] \\
G_3 & = [-2,0] + \rho_2 [-2,2] + [-2,4] \\
G_4 & = [-3,-1] + \rho_3 [-3,1] + \rho_3 [-3,3] + [-3,5] + [-1,-1] + [-1,1] \\
G_5 & = [-4,-2] + \rho_4 [-4,0] + \Lambda^2 \rho_4 [-4,2] + \rho_4 [-4,4] + [-4,6] \\
&\quad + \rho_2 ( [-2,-2] + \rho_2 [-2,0]+[-2,2]) + [0,-2]
\end{split}
\end{equation}
%
（再说一次，若把它们排成 $2\times 2$ 的网格，就构成三角阵列。）

\insfig{SU2-FL-Monomials}{本图说明穿墙如何把 $G_n$ 变成单项式。
例如若 $n=8k$，则从 $c_1$ 出发沿顺时针方向做 $4k$ 次穿墙会把 $G_n$
变成单项式 $[1,1]$。}

% 原文 2910--2927 行有一整段公式（label eq:8fold-SU2）被 % 注释掉（未发表内容），按施工指令不译。

这些表达式又一次变得相当复杂。不过它们可以被控制住，
因为足够多次穿墙会把它们带到单项式的形式。设 $\Omega$ 表示
沿顺时针方向穿过一面墙的运算。
那么我们断言：这些穿墙按图 \ref{fig:SU2-FL-Monomials} 所示的图样把 $G_n$ 变成单项式。该图样
可以用归纳法、连同两组各四个非常简单的恒等式来证明，这些恒等式直接验算即可：
%
\begin{equation}\label{eq:4fold-wilson-idents}
\begin{split}
c_{1 + 4\ell} \qquad  & [1,1]\hat W = y[0,2] + 1 + y^{-1} [2,0] \\
c_{2 + 4\ell} \qquad  & [-1,1]\hat W = y[-2,0] + 1 + y^{-1} [0,2] \\
c_{3 + 4\ell} \qquad  & [-1,-1]\hat W = y[0,-2] + 1 + y^{-1} [-2,0] \\
c_{4 + 4\ell} \qquad  & [1,-1]\hat W = y[2,0] + 1 + y^{-1} [0,-2] \\
\end{split}
\end{equation}
%
%
\begin{equation}
\begin{split}
c_{1 + 4\ell} \qquad  & [2,0]\hat W = y[1,1]  + y^{-1} \Omega[1,-1], \\
c_{2 + 4\ell} \qquad  & [0,2]\hat W = y[-1,1]  + y^{-1} \Omega[1,1], \\
c_{3 + 4\ell} \qquad  & [-2,0]\hat W = y[-1,-1]  + y^{-1} \Omega[-1,1], \\
c_{4 + 4\ell} \qquad  & [0,-2]\hat W = y[1,-1]  + y^{-1} \Omega[-1,-1].
\end{split}
\end{equation}
%
由这些恒等式可推出若干高度非平凡的结果。首先，由于 $G_{2n+1}$
在 $\Omega^n$ 之下的穿墙像容许一个平方根，算子本身
也就容许平方根。我们令 $\hat V_n := \sqrt{G_{2n+1}}$。于是例如
%
\begin{equation}
\begin{split}
\hat V_{-3} = \sqrt{G_{-5}} & = [-1,0] + [-1,-2] + \rho_2 [1,-2] + [3,-2],\\
\hat V_{-2} = \sqrt{G_{-3}} & = [0,-1] + [2,-1],\\
\hat V_{-1} = \sqrt{G_{-1}} & = [1,0],\\
\hat V_0 = \sqrt{G_1} & = [0,1],\\
\hat V_1 = \sqrt{G_3} & =    [-1,0] + [-1,2],\\
\hat V_2 = \sqrt{G_5} & =    [0,-1] + [-2,-1] + \rho_2 [-2,1] + [-2,3].
\end{split}
\end{equation}

通过到一个表达式变成单项式（或单项式的简单 $\Omega$ 像）的方便腔，我们可以容易地
建立许多有用的环关系。首先，$G$ 与 $\hat V$ 由下式相联系：
%
\begin{equation}\label{eq:SU2-formring-1}
\begin{split}
G_{2n+1} & = \hat V_n^2 \\
G_{2n} & = y^{-1/2}
\hat V_{n-1} \hat V_n \\
\end{split}
\end{equation}
%
其次，$\hat V$ 满足以下关系：
%
\begin{equation}\label{eq:SU2-formring-2}
\hat V_{n-1} \hat V_{n+1} = 1+ y \hat V_n^2
\end{equation}
%
%
\begin{equation}\label{eq:SU2-formring-3}
  \hat V_{n} \hat W =  y^{1/2} \hat V_{n+1} + y^{-1/2} \hat V_{n-1}
\end{equation}
%
我们稍后会看到，有时用 $G$ 的子代数来工作是有益的。用上面的方法可得，对奇指标的算子有
%
\begin{equation}\label{eq:SU2-formring-4}
\begin{split}
G_{2k-1} G_{2k+1} & = y^2 G_{2k}^2 \\
G_{2k-1} G_{2k+3} & = (1+ y G_{2k+1})(1+y^3 G_{2k+1}) \\
\end{split}
\end{equation}
%
而对偶指标的算子有
%
\begin{equation}\label{eq:SU2-formring-5}
G_{2k} G_{2k+2} = y G_{2k+1} (1+ y G_{2k+1})
\end{equation}
%
如此等等。

在下面的第 8.3 节% TODO 合稿改 \ref{subsec:SU2-Functions}
中，我们将看到这些关系的 $y=1$ 特殊化
如何漂亮地重现 $SU(2)$ 与 $SO(3)$ 规范理论模空间上的函数环。在第 10.3 节% TODO 合稿改 \ref{subsec:SU2-Lams}
中，特别是下面 (10.25) 式% TODO 合稿改 \eqref{eq:chebyshev-expls}
中，我们将从叶层重现
$y=+1$ 时的一串算子 $\hat V_n$。
这样，上面的 $\hat V_n$ 就给出了 $SU(2)$、$N_f=0$ 理论中保护自旋特征标的计算。

现在我们有条件给出一个（在我们看来）有说服力的论证，说明 $\hat V_n$ 与 $G_{2n}$
确实定义了强正性的形式线算子。当强正性失效时，通常会出现
一个无穷级数（例如见 (4.3)、(4.4) 两式）。% TODO 合稿改 \eqref{eq:explicit-dilog-conj}、\eqref{eq:explicit-dilog-conj-ii}
因此，在我们的论证中，我们把有限性当作强正性的替代品。我们希望考虑
$G_n$ 的各种穿墙；特别是，仅从规则 (4.3) 与 (4.4) 出发，并不能明显看出表达式会在某个腔 $c_m$ 中是有限的。% TODO 合稿改 \eqref{eq:explicit-dilog-conj}、\eqref{eq:explicit-dilog-conj-ii}
然而，利用递归关系，我们总可以把 $G_n$ 表成
$G_\ell, G_{\ell+1}, W, y^{\pm 1}$ 的有限多项式，其中 $\ell$ 取得使 $G_{\ell}$ 与 $G_{\ell+1}$
在腔 $c_m$ 中退化为单项式。这就说明 $G_n$ 在每个腔中都是有限的。
对 $\hat V_n$ 有类似的论证。

于是我们得出结论：$G_n$ 与 $\hat V_n$ 全都是强正性的形式线算子。
此外，任何由某个腔中的单项式 $[p,q]$ 生成的线算子，都可以写成
适当的 $\hat V_n$ 的乘积。例如，若我们在腔 $c_n$ 中且 $n=1 \mod 4$，
则 $p,q\geq 0$，并可取 $k$ 使 $[p,q] = y^{-pq/2} \hat V_{4k-1}^p \hat V_{4k}^q$。
由此可知，$[p,q]$ 的所有穿墙像都是 $[a,b]$ 的有限多项式，
其系数在 $su(2)$ 的表示环中。因此，给定上面主要的假设，
$[p,q]$ 确实生成一个强正性的形式线算子。

正如我们前面提到的，这些论证适用于强耦合区中的 $u$。我们
并未试图证明弱耦合区中的强正性，但我们预期它成立，
因为第 8.3 节% TODO 合稿改 \ref{subsec:SU2-Functions}
与第 10.3 节% TODO 合稿改 \ref{subsec:SU2-Lams}
所述与物理线算子的联系。



\subsection{物理线算子 vs. 形式线算子}
\label{subsec:Physics-Formal}

在第 3.2 节中，我们陈述了四个猜想：% TODO 合稿改 \ref{subsec:BPS-Deg-PSC}
关于带框 BPS 态与「原味」BPS 态的
强正性猜想与弱正性猜想。本节我们再陈述一个猜想：
它若成立，就使人能够在已知「原味」BPS 谱的前提下，仅用纯代数的手段
导出全部简单的物理线算子——即使该理论不是拉氏型的。

带框 BPS 态的强正性猜想蕴涵：物理线算子的生成函数 (3.38) 定义一个强正性的形式线算子。% TODO 合稿改 \eqref{eq:form-gen}
我们愿意猜想其逆也成立，即：任何强正性的形式线算子都是
某个物理线算子的带框 BPS 简并度所构成的生成函数。
支持这一猜想的主要证据是：正如我们将在第 10 节% TODO 合稿改 \ref{sec:Examples}
看到的，它在若干例子中成立。

利用单值群的作用，可以对这一猜想作一个漂亮的相容性检验。
物理线算子 $L$ 的具体紫外标号
必须与红外数据 $u$ 无关；只有当沿规范耦合模空间中非平凡的同伦路径搬运时，
它才可能受到单值变换的影响。当在 $\zeta$ 平面中绕原点搬运时，它也可能
受到单值变换的影响，不过只在 $U(1)_R$ 反常的理论中。原因是：
在此类理论中，$\zeta$ 的旋转可以换成
规范耦合能标的旋转，即 theta 角的平移，而这可以引起
紫外标号的单值变换。形式线算子所经历的单值变换，其形式是
单值变换过程中穿过的所有 BPS 墙对应的 KS 变换之乘积。
这与场论中 KS 变换的已知性质相容\cite{Gaiotto:2008cd}：
$u$ 平面中各奇点周围 KS 因子的乘积给出恒等变换（更确切地说，是
红外电荷格的相应单值变换），但
$\zeta$ 平面中绕原点的所有 KS 因子的乘积则不然。
因此，没有任何形式线算子会带有
不适用于物理线算子生成函数的单值变换。

遗憾的是，正如上面的例子充分显示的，
直接与形式线算子打交道相当困难；因此我们愿意
提出：实际上有一个更强的猜想成立，使我们能与
更简单的对象打交道。我们令 $y=1$，并定义一个（可能更大的）
汇集 $\CL(c)$ 构成的空间，其中 $\CL(c)$ 定义在 $\Xi$ 的各腔中，
是 $x_\gamma$ 的正整系数线性组合，且在相邻的
腔之间 $\CL(c^\pm)$ 由适当的 $y=1$ KS 变换相联系。
我们不妨把这个空间的元素称为\emph{形式叶层}（formal lamination）。
「叶层」这个名字的由来将在第 10 节中变得清楚。% TODO 合稿改 \ref{sec:Examples}
形式叶层的空间容易研究，因为它带有自然的乘积，
且通常由有限个生成元生成。注意，
若把 (3.31) 中的乘积展开，展开式中只出现正的% TODO 合稿改 \eqref{eq:y-one-prod}
整数系数（对 $\gamma$ 的适当符号而言）。
于是正性与物理预期很好地相容。

当然，只要令 $y=1$，形式线算子就可以看成形式叶层。
不过并不显然的是：两个形式线算子的非交换乘积会产生形式线算子的线性组合
\begin{equation}
F_i \circ_y F_j = \sum_{k} c_{ij}^k(y) F_k
\end{equation}
另一方面，形式叶层的集合对乘法封闭则是显然的。在我们关注的 $A_1$ 例子中，形式叶层的空间
似乎与规范理论中\emph{物理}线算子的空间重合。构造形式叶层
比构造形式线算子容易得多！因此，如果形式
叶层的空间已经与形式线算子的空间重合，那将是一个非常有用的
结果。我们因此愿意提出：这是一个值得追问的问题。

% =====================================================================
%  GMN3 §5  与簇代数的关系（Relation to cluster algebras）
%  原文：Lamination-Paper.tex 3115--3694 行；PDF p.48--56（印刷 p.47--55）
%  编号公式 (5.1)--(5.29) 共 29 条，全部在本块内，故 \eqref 可用；
%  跨块引用（节号、式号）暂写原文编号并加 % TODO，合稿时统一改 \ref/\eqref
%  注：BRIEF 分块表记本块 35 条公式，系原稿中被 % 注释掉的 4 个环境
%      （L3369 equation、L3592 align 3 行、L3642 equation、L3682 equation）
%      虚增 6 个编号单位所致；实际编译只出 29 条，已按 PDF 核对。
% =====================================================================

\section{与簇代数的关系}
\label{sec:Cluster-Algebras}

形式线算子的代数与数学中的\emph{簇代数}（cluster algebra）理论密切相关。（簇代数的文献如今已浩如烟海，下面只给出几个指路标：\cite{MR1887642,MR2132323,cluster-intro} 与 \cite{cluster-portal}。）本节简要勾勒这一关系的一小部分。我们在第 13 节% TODO \ref{sec:Open-Problems}
还会讨论，其中尚有更多值得弄清楚的事情。除了建立与数学中一个有趣课题的联系之外，我们还将看到，这一联系使我们能够只用模空间中的「局部」规则来构造形式线算子，亦即只用关于 BPS 谱的部分知识即可给出的、$\Xi$ 各腔之间的局部变换规则。此外，我们还将说明：仅凭对 BPS 谱的部分知识，原则上就足以构造出整个 BPS 谱。

\subsection{簇代数的基本定义}
\label{subsec:Cluster-Defs}

簇代数的定义从一个\emph{种子}（seed）开始：种子就是一个 $m\times n$ 矩阵 $B_{ij}$，其中 $m\geq n$、元素为整数，并且其上方 $n\times n$ 块可斜对称化。也就是说，存在正整数 $d_i$，使得 $1\leq i,j\leq n$ 时的 $B_{ij}(d_j)^{-1}$ 反对称：$d_i B_{ij} = - d_j B_{ji}$。对每个指标 $1\leq k \leq n$，沿方向 $k$ 的\emph{种子变异}（seed mutation）定义为变换 $\mu_k: B\to B'$：
%
\begin{equation}\label{eq:seed-mtrx}
B'_{ij} = \begin{cases} - B_{ij} & {\rm If} ~~ i = k ~~ {\rm or}~~ j=k \\
B_{ij} + {\rm sgn}(B_{ik}) [B_{ik} B_{kj}]_+ & {\rm Else} \\
\end{cases}
\end{equation}
%
其中我们引入了记号 $[x]_+ = {\rm Max}\{ 0,x\}$，以及
%
\begin{equation}
{\rm sgn}(x) = \begin{cases} 1 & x>0 \\  0 & x =0 \\ -1 & x < 0 \\ \end{cases}
\end{equation}
%
簇变量是 $A_1, \dots, A_m$。对每个 $k$（$1\leq k \leq n$），簇变换定义为：当 $i\not= k$ 时 $\mu_k(A_i) = A_i$，而
%
\begin{equation}
\mu_k(A_k) = \frac{ \prod_{j\vert B_{kj}>0} A_j^{B_{kj}} + \prod_{j\vert B_{kj}<0} A_j^{-B_{kj}} }{A_k}
\end{equation}
%
这一递归过程定义了某个代数的生成元序列，这个代数就称为簇代数。指标满足 $n<i\leq m$ 的变量 $A_i$ 称为「冻结变量」（frozen variables）。簇代数具有许多漂亮而出色的性质。

Fock 与 Goncharov 在 \cite{MR2567745} 中证明：若定义
%
\begin{equation}
x_i = \prod_j A_j^{B_{ij} }
\end{equation}
%
则 $x_i$ 相应的变换为 $\mu_k: x_i \to x_i'$，其中
%
\begin{equation}\label{eq:Mutation-x}
 x_i' = \begin{cases}  x_k^{-1} & i=k \\
x_i (1+ x_k^{-{\rm sgn}(B_{ik})})^{-B_{ik}} & i \not= k \\
\end{cases}
\end{equation}
%
Kontsevich 与 Soibelman \cite{ks1} 强调：式 \eqref{eq:Mutation-x} 可以用我们称之为 KS 变换的辛变换来解释。

\subsection{$\CN=2$ 理论中的簇代数结构}
\label{subsec:Seeds-from-N=2}

为了与 $\CN=2$ 理论建立联系，我们按如下方式进行。

相当一般地，若 $\Omega(\gamma;u)\not=0$，我们就称荷 $\gamma \in \Gamma_u$ 为\emph{根}。记 $\CR(u)$ 为固定 $u$ 处的根系。我们把\emph{正根系}定义为一个不交分解 $\CR(u) = \CR^+(u) \amalg \CR^-(u)$，其中若 $\gamma \in \CR^+(u)$，则 $-\gamma \in \CR^-(u)$。我们只考虑这样的正根系：其对应的 BPS 射线位于 $\zeta$ 平面内的某个半空间中。反过来，在 $\zeta$ 平面内选定一个半空间，也就决定了根系到正根与负根的一个分解。给定一个正根系，我们把\emph{单根}定义为不能写成另外两个正根之和的正根。给定任意一点 $(u,\zeta)$，我们可以典范地定义一个正根系：把满足 ${\rm Im}Z_{\gamma}/\zeta >0$ 的那些根宣布为正根。等价地，给定一个不在 BPS 墙上的点 $(u,\zeta)$，可以这样来定义正根系：取所有 BPS 射线落在过 $\zeta$ 的直线之逆时针一侧半平面内的那些根。固定 $u$ 时，点 $(u,\zeta)$ 落在 $\zeta$ 平面内的一个腔中，该腔由两道墙围成：相对于上述典范正根系，一道是逆时针一侧某个单根的墙，另一道是顺时针一侧某个「负单根」——即单根取负——的墙。若固定 $\zeta$ 而改变 $u$，则除非某条 $\pm$ 单根的 BPS 射线扫过 $\zeta$，$(u,\zeta)$ 不会穿越到 $\Xi$ 的另一个腔中去。于是，我们可以用单根系 $\Delta$ 来标记 $\Xi$ 中的各个腔。

现在，对 $\Xi$ 的任一腔 $c$，用单根 $\{\gamma_i\}$ 来定义矩阵
%
\begin{equation}
B_{ij} = \langle \gamma_i, \gamma_j\rangle   \Omega(\gamma_j;u)
\end{equation}
%
只要诸 $\Omega(\gamma_j;u)$ 全为正，这个矩阵显然可斜对称化。冻结指标 $n<i\leq m$ 标记味子格的生成元。下面我们为单根集合提出一条跨 BPS 墙的变换规则，它保证：当 $(u,\zeta)$ 沿某条路径从一个腔走到下一个腔时，发生的正是种子变异。正如刚才所观察到的，该路径必定穿过某条 $\pm$ 单根的墙。路径在 $\zeta$ 平面内的投影逆时针地越过某个单根 $Z_{\gamma_k}$，也就是从 ${\rm Im}Z_{-\gamma_k}/\zeta <0 $ 的区域进入 ${\rm Im}Z_{-\gamma_k}/\zeta > 0 $ 的区域。这就是在穿过墙 $\widehat{W(-\gamma_k)}$。类似地，若路径顺时针越过某个单根的负值 $Z_{-\gamma_k}$，从 ${\rm Im}Z_{\gamma_k}/\zeta >0 $ 的区域进入 ${\rm Im}Z_{\gamma_k}/\zeta < 0 $ 的区域，则它穿过的是墙 $\widehat{W(\gamma_k)}$。

对顺时针穿过 $\widehat{W(\gamma_k)}$ 的路径，变换
%
\begin{equation}\label{eq:charge-mutation}
\mu_{k,+}: \gamma_i \to \gamma_i' = \begin{cases} - \gamma_k & i=k \\   \gamma_i + \gamma_k [B_{ik}]_+  & i \not=k\\
\end{cases}
\end{equation}
%
给出新腔 $\mu_{k,+}(c)$ 中的新单根系。类似地，对逆时针穿过 $\widehat{W(-\gamma_k)}$ 的路径（其中 $\gamma_k$ 是单根），变换
%
\begin{equation}
\mu_{k,-}: \gamma_i \to \gamma_i' = \begin{cases} - \gamma_k & i=k \\   \gamma_i - \gamma_k [B_{ik}]_-  & i \not=k\\
\end{cases}
\end{equation}
%
给出新腔 $\mu_{k,-}(c)$ 中的新单根系。这里我们引入了记号 $[x]_- := {\rm Min}\{x,0\}$。可以验证，$\mu_{k,+}$ 是 $\mu_{k,-}$ 的双侧逆。

现在可以验证：只要 $\Omega(\gamma_j';u) = \Omega(\gamma_j;u)$，新腔的矩阵 $B'_{ij}$ 与旧腔的矩阵 $B_{ij}$ 之间的关系（对 $\mu_{k,\pm}$ 两者都是如此）恰好就是式 \eqref{eq:seed-mtrx} 所定义的种子变异 $\mu_k$。

此外，Kontsevich--Soibelman 变换
%
\begin{equation}
\CK_{\gamma_0}^{-\Omega(\gamma_0;u)} x_\gamma = x_{\gamma} (1-\sigma(\gamma_0) x_{\gamma_0})^{-\langle \gamma,\gamma_0\rangle \Omega(\gamma_0)}
\end{equation}
%
与簇变换 \eqref{eq:Mutation-x} 相合，前提是取 $\sigma(\gamma_0)=-1$、$x_i' = x'_{\gamma_i'}$、$x_i = x_{\gamma_i}$、$\gamma_0 = \gamma_k$ 以及 $\gamma_i' = \mu_{k,+}\gamma_i$。因此，举例来说，若定义 $\hat x_{\gamma} = \sigma(\gamma) x_{\gamma}$，则带框 BPS 简并数可以按照式 (3.35)% TODO \eqref{eq:Careful-rule}
用单根变异所诱导的簇变换来计算。

借助量子簇代数理论 \cite{MR2233852}，这一陈述可以推广到完整的 $y \neq \pm 1$ 穿墙；但为简单起见，我们只看经典情形。

若跨腔的变换满足上述性质，我们就说一个 $\CN=2$ 场论具有\emph{簇代数性质}。也就是说：对单根有 $\Omega(\gamma_i;u)>0$，并且 $\Omega(\mu_{k,\pm} \gamma_j;u) = \Omega(\gamma_j;u)$。（为避免对这个第二条件产生误解，请注意：腔壁\emph{不是}边缘稳定墙，因此该条件是良定的。）

具有簇代数性质的主要好处之一是：可以给出一个构造理论中 BPS 谱与线算子的算法。为看清这一点，先考虑 BPS 谱。假定我们知道某个腔中各单根的 BPS 简并数，并假定知道 $Z_{\gamma_i}(u)$ 的取值。这些量在腔内处处变化，但始终留在某个半平面内。固定 $u$，在该腔内取一个 $\zeta$ 值（于是 ${\rm Im}Z_{\gamma_i}(u)/\zeta >0$）。现在让 $\zeta$ 沿逆时针方向移动，它会撞上某个单根 $\gamma_k$ 的 BPS 谱线。此时便可借助变异 $\mu_{k,-}$ 重新算出新的单根基。继续逆时针转动 $\zeta$，就可以这样一路走下去。若能一直把 $\zeta$ 转过整整 $\pi$ 的角度，则整个 BPS 谱就构造出来了。正如我们将在下面两个例子中看到的，这样做可能遇到阻碍。基本的原因是：所有腔的单根系都只含超多重态，而更高自旋粒子的 BPS 射线被无穷多条、属于无穷多个超多重态的射线所包围，这些射线向高自旋射线聚积。因此高自旋射线从不与某个腔的边界相联系。若只在 $\zeta$ 平面内平移，过程会在这一无穷超多重态序列的第一个上就卡住，无法到达序列「末端」的高自旋粒子。另一方面，在 $A_1$ 理论中，并且可能在 $\CS$ 类中的所有理论中，看来只要允许在 $(Z_\gamma, \zeta)$ 的空间中沿一般路径移动，所有腔都能经由有限多步连通。于是，一个可能的谱计算策略是：利用这类路径「跳」到阻碍的前方，然后再回溯，以控制阻碍两侧的无穷超多重态序列，从而求出跨过阻碍的剩余坐标变换。该变换还可进一步分解为变换 $\CK_\gamma$，由此读出理论的完整谱。

作为上述步骤的一个例子，设 $\Gamma$ 是二维的、$\Omega(\gamma_1)=\Omega(\gamma_2)=1$，并且对某个正整数 $n$ 有
%
\begin{equation}
B_{ij} = \begin{pmatrix} 0 & n \\ -n & 0 \\ \end{pmatrix}.
\end{equation}
%
从 $\zeta$ 落在实轴上、且 ${\rm Im} Z_1 > 0 $ 与 ${\rm Im} Z_2 >0$ 的那个腔出发。若 ${\rm arg}Z_1 > {\rm arg} Z_2$，我们就作一列变换 $\mu_{2,-}, \mu_{1,-}, \mu_{2,-}, \cdots$。所得单根的图样周期为 4，具体就是
%
\begin{equation}
\begin{split}
\gamma_1^{[1]}=   \gamma_1, & \quad   \gamma_2^{[1]}= - \gamma_2, \\
\gamma_1^{[2]}= - \gamma_1, & \quad \gamma_2^{[2]}= -\gamma_2, \\
\gamma_1^{[3]}=  - \gamma_1, & \quad \gamma_2^{[3]}=  \gamma_2.
\end{split}
\end{equation}
%
反之，若 $\arg Z_2 > \arg Z_1$，则逆时针移动 $\zeta$ 时会遇到变异序列 $\mu_{1,-}, \mu_{2,-}, \mu_{1,-},\cdots$。若把初始的腔标为 $j=0$，我们得到一列腔，满足
%
\begin{equation}
\begin{split}
\gamma_1^{[2j+1]} & = -\gamma_1^{[2j]}  , \\
\gamma_2^{[2j+1]} & = n\gamma_1^{[2j]}+ \gamma_2^{[2j]},
\end{split}
\end{equation}
%
\begin{equation}
\begin{split}
\gamma_1^{[2j+2]} & =\gamma_1^{[2j+1]}+ n\gamma_2^{[2j+1]},   \\
\gamma_2^{[2j+2]} & =  - \gamma_2^{[2j+1]}.
\end{split}
\end{equation}
%
当 $n=1$ 时，这一列变换的周期为 $6$。它已在 $N=3$ Argyres-Douglas 理论的弱耦合区中被具体实现，见第 4.3 节% TODO \ref{sec:formaln3}
。当 $n=2$ 时，这列变换没有有限的周期性，且 $\mu_{2,-}\mu_{1,-}$ 是由矩阵
%
\begin{equation}
1 + N = \begin{pmatrix} 3 & 2 \\ -2 & -1 \\ \end{pmatrix}
\end{equation}
%
所给出的变换。由于 $N^2 =0$，可以立刻写出
%
\begin{equation}
\begin{split}
\gamma^{[2j]}_1 & = (2j+1)\gamma_1 + 2j  \gamma_2, \\
\gamma^{[2j]}_2 & =  -2j \gamma_1  -  (2j-1) \gamma_2.
\end{split}
\end{equation}
%
这在 $SU(2)$ $N_f=0$ 理论中实现过，见第 4.4 节% TODO \ref{subsec:formal-SU2}
。注意这些 BPS 射线聚积在沿 $\pm (Z_1 + Z_2)$ 的射线上。因此，无法越过这个极限继续转过整整 $\pi$ 的角度。不过，若改沿顺时针方向继续，就会找到对应于相反射线的第二个极限值，于是实际上完整的谱仍然被生成出来。

当 $n>2$ 时会出现一个有趣的现象。从 $2j$ 到 $2j+2$ 的变换矩阵是
%
\begin{equation}
\begin{pmatrix} n^2 -1 & n \\ -n &  -1 \\ \end{pmatrix} = S^{-1} \begin{pmatrix} \lambda_+ & 0 \\
0 & \lambda_- \\ \end{pmatrix} S
\end{equation}
%
其中 $1+\lambda_\pm = \half (n^2  \pm n \sqrt{n^2-4} ) $。沿逆时针方向移动时，BPS 射线的极限斜率为
%
\begin{equation}\label{eq:slope-ccw}
\frac{ \xi {\rm Im}Z_1 +   {\rm Im}Z_2 }{\xi {\rm Re} Z_1 +   {\rm Re} Z_2}
\end{equation}
%
其中 $\xi = \frac{\lambda_+ +1 }{n} > 1$。现在若考虑沿顺时针方向转动，我们会遇到变异序列 $\mu_{2,+}$, $\mu_{1,+}$, $\mu_{2,+},\dots$，它给出一列单根
%
\begin{equation}
\begin{split}
\gamma_1^{[2j+1]} & =   \gamma_1^{[2j]} + n   \gamma_2^{[2j]} \\
\gamma_2^{[2j+1]} & =  -   \gamma_2^{[2j]} \\
\end{split}
\end{equation}
%
\begin{equation}
\begin{split}
\gamma_1^{[2j+2]} & =     -\gamma_1^{[2j+1]}    \\
\gamma_2^{[2j+2]} & = n \gamma_1^{[2j+1]} + \gamma_2^{[2j+1]}  \\
\end{split}
\end{equation}
%
此时的极限斜率为
%
\begin{equation}\label{eq:slope-cw}
\frac{ {\rm Im}Z_1 +  \xi  {\rm Im}Z_2 }{  {\rm Re} Z_1 +  \xi {\rm Re} Z_2}
\end{equation}
%
而该射线指向 $-\pi < {\rm arg}\zeta < -\pi/2$ 的那个象限。由于 $\xi >1$，两条极限 BPS 射线之间张开的角小于 $\pi$。因此存在一个「空隙区域」，如图所示~\ref{fig:gap}，上述步骤无法探查到那里。该区域内被占据的荷的谱已经有人算出 \cite{gctv}，其结果清楚地显示了困难所在：在该区域中，$\Omega(\gamma) \neq 0$ 的那些 $Z_\gamma$ 的相位是\emph{稠密}的，因此谈论腔其实已不太可能。

\insfig[0.58\textwidth]{gap}{在秩 2 例子中当 $n>2$ 时，我们只能用固定 $u$ 处的簇变换来确定落在某个角小于 $\pi$ 的角扇区内的那些射线的 BPS 谱，该扇区用蓝色箭头标出。图中对 $Z_1$ 与 $Z_2$ 的相位作了一个方便的选取。}

当一个 $\CN=2$ 理论具有簇代数性质时，会立刻得到一些好处。文献 \cite{MR1888840} 的洛朗现象使我们能够构造形式叶层。形式线算子的构造涉及某种「量子洛朗现象」，它似乎尚未在数学文献中被讨论过（尽管它与 Fock 与 Goncharov 的「泛洛朗多项式」相去不远）。若我们的目的只是检验某个形式线算子的强正性，那么其实并不关心荷格的标号；于是只需把生成函数用 $X_i$ 表出，忘掉 $\gamma_i$，并确保强正性被所有可能的簇变换保持即可。这可以用来简化我们在第 4.3 节% TODO \ref{sec:formaln3}
与第 4.4 节% TODO \ref{subsec:formal-SU2}
中给出的推导，并轻而易举地把它们推广到弱耦合区，因为强耦合与弱耦合处的簇变换序列无从区分。例如，若忘掉荷的标号，$N=3$ Argyres-Douglas 理论总共只有五个不同的腔，由五边形的五种三角剖分标记。簇变换的周期为 $5$，而基本的形式线算子只是经历一个周期序列 $X\to X+X \to X+X+X \to X+X \to X \to X$。我们在强耦合处遇到的周期为 $20=4 \times 5$ 的长序列、以及在弱耦合处周期为 $30 = 6 \times 5$ 的长序列，不过是周期 $5$ 的簇变换与周期 $4$ 或 $6$ 的荷重标号的组合而已。

我们猜想：$\CS$ 类中的理论都具有簇代数性质。这一猜想受到 Fock 与 Goncharov 关于高阶 Teichm\"uller 理论的工作 \cite{MR2233852} 的启发。若我们的猜想成立，则上述算法或许提供了计算高秩 $\CN=2$ 规范理论 BPS 谱的一条途径。我们认为这是一个重要的想法，值得继续追下去。基于以上观察，一个自然会问的问题是：是否每个簇代数都源自某个 $\CN=2$ 理论？这个问题我们也留待将来。

\subsection{$A_1$ 理论中的簇代数结构}
% 原文此小节的 label 与 5.2 重复（都写作 subsec:Seeds-from-N=2），系原文笔误；此处加后缀区分。
\label{subsec:Seeds-from-A1}

本节我们将证明：$A_1$ 型理论这一子类确实具有簇代数性质。这一点可以借助 \cite{Gaiotto:2009hg} 中所解释的腔与 $C$ 的 WKB 三角剖分之间的关系来证明。回忆该文第 6 节：给定 $(u,\zeta)$，$C$ 上有一个典范的（WKB）三角剖分。该文第 7.8 节证明了还存在一个与之相伴的单根系 $\{ \gamma_E \}_{E\in \CE(T)}$。而且，每个这样的单根上都承载着单个超多重态，其 PSC 等于 $1$。当 $(u,\zeta)$ 变动时，除非某条边发生 flip，三角剖分只作同痕变化；而后者只可能在 $(u,\zeta)$ 穿过一道 BPS 墙时发生。因此，腔的墙是诸墙 $\widehat{W}(\pm \gamma_E)$（$E\in \CE(T)$）的一个子集。一般地，我们预期所有边都会贡献腔的墙。由一条边 $E_k$ 的 flip 所引起的 $\{ \gamma_E \}_{E\in \CE(T)}$ 的变换规则，已在 \cite{Gaiotto:2009hg} 的第 7.6 节中算出，它们与 $\mu_{k,\pm}$ 相符。

我们可以用 $C$ 的理想三角剖分，连同某个「适配二重覆盖」的拓扑类，来标记各个腔。$C$ 的一个适配二重覆盖，是指在 $C$ 的每个三角形中各有一个分支点的二重覆盖，并且要选定各边 $E$ 之提升的定向，使得在 $C$ 的奇点的原像处，这些提升要么全是入的、要么全是出的。

\subsection{$\CN=2^*$ $SU(2)$ 规范理论中的形式线算子}
\label{subsec:Formal-N2-star}

现在我们转向规范代数为 $su(2)$ 的 $\CN=2^*$ 理论。它由 $\CN=4$ 理论给一个伴随超多重态赋予质量 $m$ 而得到。它就是 $C$ 为带一个穿刺点的环面时的 $A_1$ 理论。\footnote{更确切地说，$C$ 为带穿刺点环面时的 $A_1$ 理论，是 $\N=2^*$ 理论再加上一些解耦的场，即在该穿刺点对应的 $SU(2)$ 味对称下的一个半超多重态二重态。}Seiberg--Witten 微分由
%
\begin{equation}
\lambda^2 = \left( m^2 \wp(z\vert \tau) + u\right) (dz)^2,
\end{equation}
%
给出，其中 $z$ 是模数为 $\tau$ 的椭圆曲线上的平坦坐标，$\wp(z\vert \tau)$ 是 Weierstrass 函数。Seiberg--Witten 曲线亏格为 2、带两个穿刺点。局部系统 $\Gamma$ 的秩为 $3$，可以取生成元 $\gamma_1, \gamma_2, \gamma_3$，使得对 $i=1,2,3$ 有 $\langle \gamma_i, \gamma_{i+1} \rangle =2$，其中指标 $i$ 按 $3$ 阶循环理解。注意 $\gamma_1 + \gamma_2 + \gamma_3$ 生成 $\langle \cdot, \cdot \rangle$ 的一维零化子。这就是秩 $1$ 的「味格」。

可归一化参数是 $u$，并且 $\CB^* = \IC-\CS$，其中 $\CS$ 由三个奇点 $u_i$ 组成，满足 $\wp(z_i\vert \tau) = - u_i/m^2$ 与 $\wp'(z_i\vert\tau) =0$。局部系统 $\Gamma$ 绕这三个点有单值变换。众所周知该理论具有 $SL(2,\IZ)$ 对称性，其中 $PSL(2,\IZ)$ 以标准方式作用在耦合常数 $\tau$ 上。

遗憾的是，$\CN=2^*$ 理论的 BPS 谱 $\Omega(\gamma;u)$ 并不知道其显式形式。\footnote{我们注意到，该谱隐含地编码在 \cite{Gaiotto:2009hg} 的「谱生成元」之中。这个变换容易写出，但要把它分解为 Kontsevich--Soibelman 因子却并不直接。}尽管如此，很清楚该谱与迄今考察过的例子有一些鲜明的差别。特别是，这里没有谱有限的「强耦合区」。我们预期谱总是无穷的，并由数目可变的若干矢量多重态控制，而每个矢量多重态都伴随着无穷多个超多重态。因此腔的结构非常复杂。于是，这里正是一个可以非常有效地使用局部簇变换规则的例子。

在每个腔 $[c]$ 中有三个单根 $\gamma_i[c]$；该腔由六道墙 $\widehat{W}(\pm \gamma_i[c])$ 围成。$\sum_i \gamma_i[c]$ 是纯味的，并且我们\emph{或者}有 $\langle \gamma_i[c],\gamma_{i+1}[c]\rangle=2$，\emph{或者}有 $\langle \gamma_i[c],\gamma_{i+1}[c]\rangle=-2$。我们分别称这些腔为正型与负型。该腔就是集合 $\{ (u,\zeta): {\rm Im}Z_{\gamma_i[c]}(u)/\zeta > 0 \}$。注意一次变异 $\mu_{i,\pm}$ 把正型的腔变到负型的腔。

方便起见，定义
%
\begin{equation}
[p,q,r]:= X_{\half (p \gamma_1 + q \gamma_2 + r \gamma_3)}.
\end{equation}
%
当需要强调腔时，我们可以写成 $[p,q,r](c)$。我们首先注意到，与先前的例子相反，这里没有任何单项式 $[p,q,r]$ 能够生成形式线算子。在我们的腔中，${\rm Im} Z_i/\zeta >0$（$i=1,2,3$）。检视规则 (4.3)% TODO \eqref{eq:explicit-dilog-conj}
与 (4.4)% TODO \eqref{eq:explicit-dilog-conj-ii}
，我们看到必须有 $\langle \gamma, \gamma_i[c] \rangle <0$（$i=1,2,3$），其中 $\gamma = \half (p\gamma_1[c]+q\gamma_2[c]+r\gamma_3[c])$。另一方面，$\gamma_1[c]+\gamma_2[c]+\gamma_3[c]$ 位于 $\langle \cdot, \cdot \rangle$ 的核中，因此这三个不等式永远不可能同时满足。

尽管无法由单项式生成形式线算子，却有三个漂亮的类比物，它们对应于 Wilson 线算子 $W$（后者在第 4.4 节% TODO \ref{subsec:formal-SU2}
中十分有用）：
%
\begin{equation}\label{eq:star-formal-Wilson}
\begin{split}
W_1[c] & = [0,1,1] + [0, -1, -1] + [0,-1,1],\\
W_2[c] & = [1,0,1] + [-1,0, -1] + [1,0,-1],\\
W_3[c] & = [1,1,0] + [-1, -1, 0] + [-1,1,0].
\end{split}
\end{equation}
%
若腔是正型的（$\langle \gamma_i[c],\gamma_{i+1}[c]\rangle=+2$），这些就是跨过该腔所有墙都变换良好的唯一组合。对负型的腔，$\langle \gamma_i[c],\gamma_{i+1}[c]\rangle=-2$，这三个 Wilson 算子具有如下形式
%
\begin{equation}
\begin{split}
W_1[c] & = [0,1,1] + [0, -1, -1] + [0,1,-1],\\
W_2[c] & = [1,0,1] + [-1,0, -1] + [-1,0,1],\\
W_3[c] & = [1,1,0] + [-1, -1, 0] + [1,-1,0].
\end{split}
\end{equation}
%
记 $\mu_{i,\pm}$ 为跨过墙 $\widehat{W}(\pm \gamma_i[c])$ 的变换（变异）。当 $j\not=i$ 时，$W_i[c]$ 在 $\mu_{j,\pm}$ 下的变换遵循 $SU(2)$ Wilson 线那个熟悉的模式：一个二项（自旋 $1/2$）的 halo 塌缩，而某一项长大成一个二项的 halo，于是我们在下一个腔中仍然得到一个三项的表达式。也就是说，对 $i \not=j$ 我们有简单的结果
%
\begin{equation}
\mu_{i,\pm}(W_j[c]) = W_j[\mu_{i,\pm}(c)].
\end{equation}
%
另一方面，若用 $\mu_{i,\pm}$ 去变异 $W_i[c]$，则有两项不变，而第三项变成一个三项的（自旋 $1$）halo。例如，用 $\mu_{3,+}$ 变换到腔 $c' = \mu_{3,+}(c)$，我们得到：
%
\begin{equation}\label{eq:mu-three-W3}
\begin{split}
\mu_{3,+} W_3[c] &=  [1,1,0](c) + [-1, -1, 0](c) + ([-1,1,0]+ [2] [-1,1,2] + [-1,1,4])(c)\\
& = [1,1,2](c') + [-1,-1,-2](c') + \left( [-1,1,2] + [2] [-1,1,0] + [-1,1,-2]\right)(c'). \\
\end{split}
\end{equation}
%
与 $SU(2)$ $N_f=0$ 的情形一样，Wilson 线算子的环关系与它们跨腔壁的变换性质美妙地联系在一起。例如
\begin{align}
W_1[c] W_2[c]&= y^{-\half}W_3[c]+y^{\half}\bigl(X_{\half \gamma_1[c]+\half \gamma_2[c]+\gamma_3[c]} +X_{\half \gamma_1[c]-\half \gamma_2[c]+\gamma_3[c]} \notag\\
&\quad +X_{-\half \gamma_1[c]-\half \gamma_2[c]-\gamma_3[c]} +X_{\half \gamma_1[c]-\half \gamma_2[c]-\gamma_3[c]} + (y+y^{-1})X_{-\half \gamma_2[c]+\half\gamma_1[c]}\bigr)
\end{align}
乘以 $y^\half$ 之后的生成函数在跨过任一 $\pm \gamma_3[c]$ 墙时会简化：它失去一个 Fock 空间，变成一个简单的 Wilson 线，这一点从 \eqref{eq:mu-three-W3} 可以认出。

这类乘积的一般关系是：
%
\begin{equation}\label{eq:quant-skein-rel}
W_{i}[c]  W_{i+1}[c] = y^{-1/2} W_{i+2}[c]   + y^{1/2} \mu_{i+2,+} W_{i+2}[\mu_{i+2,-}(c)]
\end{equation}
%
对 $W_{i+1}[c] W_{i}[c]$，只需把 $y\to 1/y$。

我们预期，通过逐次乘以 Wilson 线，可以递推地生成整个形式线算子环，正如第 4.4 节% TODO \ref{subsec:formal-SU2}
的情形那样。

我们正一步步揭示该理论线算子中一个美妙的结构。在本文的后半部分，我们将发展更有力的工具来分析所有 $A_1$ 理论；在此过程中，我们会重新发现本节的这些公式。为了与第 8.4 节% TODO \ref{subsec:N2-star-moduli}
和第 10.7 节% TODO \ref{subsec:SU2-N2star-Examples}
作比较，我们再作最后一个计算：注意 $[2,2,2]$ 与 $[-2,-2,-2]$ 是中心元，并且可以适当地选取系数 $a_{ijk}\in \IZ[y^{1/2},y^{-1/2}]$，使得
%
\begin{equation}\label{eq:quant-mod-eq}
\begin{split}
W_1^2 + W_2^2 + W_3^2 - a_{ijk} W_i W_j W_k &= 1 + y^{-2} + B(y^2-y^{-2}) \\
&\quad - \left(y^{-1} + (y-y^{-1})B\right)([2,2,2]+[-2,-2,-2])
\end{split}
\end{equation}
%
是中心算子。这里 $B$ 是任意常数。式 \eqref{eq:quant-mod-eq} 的 $y \to 1$ 极限就是式 (8.9)% TODO \eqref{eq:moduli-equation}
。

% =====================================================================
%  GMN3 §6 + §7.1  线算子与模空间上的全纯函数；六维 (2,0) 理论回顾
%  原文：Lamination-Paper.tex 3695--4378 行；PDF p.57--66（印刷 p.56--65）
%  编号公式 (6.1)--(6.15)、(7.1)--(7.5) 共 20 条，全部在本块内，故 \eqref 可用；
%  跨块引用（节号、式号）暂写原文编号并加 % TODO，合稿时统一改 \ref/\eqref
% =====================================================================

\section{线算子与模空间上的全纯函数}
\label{sec:linehol}

\subsection{在「达布坐标」中的展开}

我们一直在研究一个 $d=4$、$\CN=2$ 的理论，它配备一族线算子 $\CL$。现在考虑该理论住在欧几里得时空 $\IR^3\times S^1$ 上，其中圆周的半径为 $R$，并且对费米子取周期边界条件。低能下，该理论由三维 sigma 模型描述，其靶空间是超凯勒流形 $\CM_\CL$。（下文第 6.3 节% TODO \ref{sec:moreonmoduli}
将详述：对不同的 $\CL$，这些模空间彼此如何不同。）

这个 sigma 模型是 \cite{Gaiotto:2008cd} 的主要研究对象。我们先快速回顾几个基本点，以固定记号。拓扑上 $\CM_\CL$ 是 $\CB$ 上的环面纤维化。$u\in\CB$ 处的纤维，是阿贝尔规范场绕 $S^1$ 的电与磁 Wilson 线所成的环面。选定配对 $(-1)^{\inprod{\gamma,\gamma'}}$ 的一个二次加细 $\sigma:\Gamma_\CL\to\IZ_2$ 之后，该环面纤维可以等同于特征标环面 $\Hom(\Gamma_\CL,\IR/2\pi\IZ)$。我们把此环面中的元素记作 $\theta$。$\CM_\CL$ 上有复的「达布坐标」$\CX_\gamma$，它们在 \cite{Gaiotto:2008cd,Gaiotto:2009hg} 中已被详细研究。\footnote{关于超凯勒流形及其复结构的约定，以及达布坐标 $\CX_\gamma$ 的定义，我们参见 \cite{Gaiotto:2008cd}。简要总结见 \cite{Gaiotto:2009hg} 第 2 节。}
在下文中，改用坐标 $\CY_\gamma$ 会更为方便；它与 $\CX_\gamma$ 只差一个二次加细：
\begin{equation}\label{eq:Ygam-Xgam}
\CX_\gamma = \sigma(\gamma) \CY_\gamma,
\end{equation}
$\CY_\gamma$ 其实比 $\CX_\gamma$ 更典范：它们的定义不依赖于 $\sigma$ 的选取。$\CY_\gamma$ 服从一条扭乘法法则
\begin{equation}\label{eq:Y-twisted}
 \CY_\gamma \CY_{\gamma'} = (-1)^{\inprod{\gamma,\gamma'}} \CY_{\gamma+\gamma'}.
\end{equation}

现在可以让四维理论的线算子 $L_\zeta$ 绕紧化圆周 $S^1$ 而缠绕，从而定义相应的圈算子。在低能极限下，它们成为三维理论中的点算子。这样一个点算子的期望值 $\inprod{L_\zeta}$ 是模空间 $\CM_\CL$ 上的函数。如附录 B% TODO \ref{app:Holomorphy}
所说明的，$L_\zeta$ 保持 $osp(4^*\vert 2)_\zeta$ 这一事实，蕴涵 $\inprod{L_\zeta}$ 实际上是 $\CM_\CL$ 上（复结构 $J^{(\zeta)}$ 下）的\ti{全纯}函数。由此，它们可以表成 $\CY_\gamma$ 的函数。我们猜想
\begin{equation} \label{eq:trh-large-r}
\inprod{L_\zeta} = \sum_\gamma \fro(L_\zeta, \gamma) \CY_\gamma.
\end{equation}

下面我们试着为 \eqref{eq:trh-large-r}
给出论证。我们觉得下述论证是有说服力的，尽管它够不上一个证明。先从把 $\inprod{L_\zeta}$ 看作四维理论希尔伯特空间上的一个迹说起。要把这个迹显式写出，略有微妙之处：
\begin{equation}\label{eq:Lvev-trace}
\inprod{L_\zeta} = {\rm Tr}_{\CH_{u,L_\zeta}} (-1)^F e^{-2 \pi R H} e^{i \theta\cdot \CQ}\sigma(\CQ).
\end{equation}
这里有两点微妙之处。第一，$\CM_\CL$ 的环面纤维定义了一个具有非平凡单值的局部系统。$\CQ$ 是荷算子，它测量 $\Gamma$ 中的红外荷，因而取值于某个局部系统。记号 $\theta\cdot\CQ$ 是自然的求值配对，它不带单值。第二点微妙之处是：$\theta$ 同时为电 Wilson 线\emph{与}磁 Wilson 线定义了边界条件。要落实这样的边界条件，我们应当使用自对偶的形式体系。众所周知，要定义自对偶阿贝尔规范理论的路积分，就必须为相应的上同调引入一个双线性形式的二次加细——自对偶规范理论的 Dirac 量子化正建基于这个上同调之上 \cite{Witten:1996hc,Freed:2000ta,Hopkins:2002rd,Belov:2006jd}。在当下情形，我们应取 Dirac 量子化条件取值于 $H^2(M_4;\Gamma)$。当 $M_4=(\IR^3-\{\vec 0\})\times S^1$ 时，使用紧支撑上同调，我们就被引导到在迹中包含二次加细 $\sigma(\CQ)$。这就是式 \eqref{eq:Lvev-trace}
中 $\sigma(\CQ)$ 的来历。

在 $R\to\infty$ 极限下，迹 \eqref{eq:Lvev-trace}
被投影到能量最低的那些态，即带框 BPS 态上。于是可以使用理论的红外描述。对一个电荷为 $\gamma$ 的带框 BPS 态，我们预期的主要贡献形如
\begin{equation}
\sigma(\gamma) (-1)^{F} \exp(2 \pi R \, {\rm Re} (Z_\gamma / \zeta) + i \theta_\gamma).
\end{equation}
因此当 $R\to\infty$ 时
\begin{equation}\label{eq:Lzeta-asymp}
\inprod{L_\zeta} \sim \sum_\gamma \fro(L_\zeta, \gamma) \exp(2\pi R
\, {\rm Re} (Z_\gamma / \zeta) + i \tilde\theta_\gamma),
\end{equation}
其中 $e^{i \tilde\theta_{\gamma}}:= e^{i \theta_{\gamma}}\sigma(\gamma)$。这个展开与我们所想要的 \eqref{eq:trh-large-r}
很像，只不过那里写的是 $\CY_\gamma$，这里写的是 $\exp(2\pi R \, {\rm Re} (Z_\gamma / \zeta) + i \tilde\theta_\gamma)$。而这确实与 \cite{Gaiotto:2008cd} 所确定的 $\CY_\gamma$ 在 $R\to\infty$ 的渐近行为相符。那么如何延拓到有限的 $R$ 呢？我们暂且把这个展开写成
\begin{equation} \label{eq:trh-large-r-prov}
\inprod{L_\zeta} = \sum_\gamma \fro(L_\zeta, \gamma) \CY'_\gamma,
\end{equation}
其中 $\CY'_\gamma$ 是 $\CM$ 上（复结构 $\zeta$ 下）的一组全纯函数。在红外理论中，$\CY'_\gamma$ 刻画的是插入一个电荷为 $\gamma$ 的源的效果。特别地，我们预期关于紫外线算子的全部细节信息都被 $\fro(L_\zeta, \gamma)$ 吸收，从而 $\CY'_\gamma$ 不依赖于紫外线算子 $L_\zeta$ 的具体选取。如果假定线算子 $L_\zeta$ 足够多，那么我们其实可以把 \eqref{eq:trh-large-r}
看作定义 $\CY'_\gamma$ 的一个线性变换。我们想证明 $\CY_{\gamma}=\CY'_\gamma$。

回忆 \cite{Gaiotto:2008cd,Gaiotto:2009hg}：$\CX_\gamma$（从而 $\CY_\gamma$）可以被一组公理唯一刻画，这组公理表述它们的全纯性、$R\to\infty$ 与 $\zeta\to 0,\infty$ 时的渐近行为、$\zeta\to -1/\bar\zeta$ 下的实性、穿墙行为以及乘法。让我们逐条比较 $\CY'_\gamma$ 与 $\CY_\gamma$。

首先是全纯性：我们知道 $\inprod{L_\zeta}$ 在 $\CM^\zeta$ 上全纯，并且按假定在 $\zeta\in\IC^*$ 上全纯。其次，如我们已看到的，$\CY'_\gamma$ 在 $R\to\infty$ 的渐近行为就是 $\CY_\gamma$ 的渐近行为，即
\begin{equation}\label{eq:X-asymps}
\CY'_\gamma \sim \CY_{\gamma}^{\rm sf} = \exp( \pi R
 \frac{Z_\gamma}{\zeta} + i \tilde\theta_\gamma + \pi R \zeta \bar Z_{\gamma} ).
\end{equation}
对 $\zeta\to 0,\infty$ 也有类似的渐近行为，这可从超对称 Wilson 圈的形式性质结合对偶推出。用同样的策略，也可以论证 $\zeta\to -1/\bar\zeta$ 下所需的实性。

接着考虑穿墙。没有令人信服的理由预期 $\inprod{L_\zeta}$ 作为 $\CB\times\IC^*$ 上的函数会出现任何相变，因为紫外理论的真空结构没有相变，而且只要 $R>0$，我们取的始终是一个性质良好的迹类算子的迹。因此 $\inprod{L_\zeta}$ 应当\emph{不}表现出任何穿墙行为。另一方面，正如我们在第 3.4 节% TODO \ref{subsec:Halo-Wall-Crossing}
中详细看到的，带框 BPS 简并度 $\fro(L_\zeta,\gamma)$ 确实经历穿墙。现在注意：\eqref{eq:trh-large-r}
看起来极像我们在 (3.38)% TODO \eqref{eq:form-gen}
中引入的形式生成函数 $F(L)$，只是把形式变量 $X_\gamma$ 换成了 $\CM$ 上货真价实的函数 $\CY_\gamma$，并把 $y$ 取为 $y=-1$。当 $(u,\zeta)$ 越过一道墙 $\widehat{W}(\gamma_\BPS)$ 时，函数 $\CY_\gamma$ 的间断由一个辛同胚 $\CK_{\gamma_\BPS}^{\Omega(\gamma_\BPS)}$ 给出，其中 $\CK_{\gamma}$ 是 (3.36)% TODO \eqref{eq:KS-sympo}
中定义的辛同胚。因此，与从 (3.26)% TODO \eqref{eq:gen}
到 (3.34)% TODO \eqref{eq:y-min-one-gen}
的讨论相比较，我们看到：只要 $\CY'_\gamma$ 的变换方式与 $\CY_\gamma$ 完全相同，$\inprod{L_\zeta}$ 就确实不会遭受任何穿墙间断。

最后，$y=-1$ 时的海森堡关系 (3.37)% TODO \eqref{eq:Heis-Alg}
蕴涵函数 $\CY_\gamma$ 应满足扭乘法法则 $\CY_\gamma \CY_{\gamma'} =(-1)^{\langle \gamma,\gamma'\rangle} \CY_{\gamma+\gamma'}$。这与 \eqref{eq:Y-twisted}
美妙地相容。

至此我们完成了对 \eqref{eq:trh-large-r}
与 \eqref{eq:Ygam-Xgam}
的论证。论证的主要缺口在于如下断言：线算子足够多，足以把方程 \eqref{eq:trh-large-r}
反解出来。


\subsection{关于达布展开的评注}

我们想就上一小节作若干评注。

\begin{enumerate}

\item 由于指标在有质量表示上为零，人们可能会问：为什么 \eqref{eq:Lvev-trace}
不简单地是带框 BPS 态上的求和？（正如在指标定理的热核证明中、或在紧靶空间上超对称量子力学里 Witten 指标的计算中，确实会发生的那样。）
如我们已指出的，把迹中的算子作用在带框 BPS 态上，给出的是半平坦达布坐标，而不是真正的达布坐标，因而给出错误的答案。其原因与 \cite{Cecotti:1992qh} 中长篇讨论的现象非常相似：连续谱对迹作出了非零贡献。

\item 正如我们在第 7 节% TODO \ref{sec:Six-Dim-View}
中所描述的，有一类非常有趣的 $d=4$、$\CN=2$ 理论，来自 $(2,0)$ 理论在带穿刺点的黎曼面 $C$ 上的紧化。在这类情形，红外阿贝尔规范理论来自阿贝尔 $(2,0)$ 理论的自对偶 3-形式在 $\IR^4\times\Sigma$ 上的紧化，其中 $\Sigma$ 是覆盖 $C$ 的 Seiberg--Witten 曲线。在阿贝尔六维理论的路积分中，对拓扑扇区的求和由一个二次加细加权，该二次加细是 $H^3(M_6;\IZ)$ 上模二相交形式的二次加细 \cite{Witten:1996hc,Freed:2000ta,Hopkins:2002rd,Belov:2006jd}。把它应用于 $M_6=(\IR^3-\{\vec 0\})\times S^1\times\Sigma$，并使用紧支撑上同调，六维理论的这个二次加细就诱导出 $H^1_{cpt}(\Sigma;\IZ)\cong H_1(\Sigma;\IZ)$ 上模二相交形式的二次加细 $\sigma$。

\item 既然有保护自旋特征标，自然可以把 \eqref{eq:Lvev-trace}
的右端推广，引入一个与 $\CJ_3$ 共轭的变量，并定义扭迹：
\begin{equation}\label{eq:y-dep-trace}
\inprod{L_\zeta}_y := {\rm Tr}_{\CH_{u, L_\zeta}} (-1)^F e^{-2 \pi R H} (-y)^{2 \CJ_3} e^{i \theta \cdot \CQ}\sigma(\CQ)
\end{equation}
迹的被积算子仍与一个超对称 $\epsilon_{\alpha A} \CR_{\alpha}^A$ 对易。然而，\eqref{eq:y-dep-trace}
的物理解释看来与 \eqref{eq:Lvev-trace}
相当不同。原因在于：若要给 \eqref{eq:y-dep-trace}
一个三维的解释，则这个被修改过的迹要求我们在绕 $z$ 轴旋转与 R 对称性旋转之下把空间重新粘回自身。这让人想起四维背景的一种 $\Omega$ 形变，它只保持一个超荷。（近期文献可参见 \cite{Nekrasov:2010ka}，其中引有更早的文献。）
用三维场论来解释这个量并不直接，而且可能要用一个非交换的靶空间来替代 $\CM$。的确，人们可以猜测：低能下理论归结为 $\CM$ 上的一个量子力学问题，而迹投影到最低能量的态上，在那里相应的展开系数是非交换的（很像量子霍尔效应最低朗道能级中的波函数可以用非交换几何来解释）。这或许是一个富有成果的语境，可为第 3.6 节% TODO \ref{subsec:Deformed-Ring}
所讨论的线算子形变手征环给出物理解释。这是一个有意思的进一步研究方向，但本文不去追究，只在下文第 11 节% TODO \ref{sec:Quantum-Holonomy}
作一些相关的评注。

\item 最后，我们可以就如下事情作一个简单的观察：相应线算子的自旋特征标之乘积，给出了 $\CM$ 上全纯函数环的一个非交换形变，它是什么意思。离开 $y=\pm 1$ 的首个修正由泊松括号 $\{ X_\gamma,X_{\gamma'}\} = \langle \gamma,\gamma'\rangle X_\gamma X_{\gamma'}$ 给出，它与 $\CX_\gamma$ 这些函数的自然泊松括号一致（记住：超凯勒流形 $\CM$ 在每一个复结构下都是复辛的）。因此，这个非交换乘积应当与 $\CM$ 上全纯函数环的形变量子化一致。我们在第 8 节% TODO \ref{sec:Examples-of-M}
中用若干例子说明这一评注。

\end{enumerate}


\subsection{多个模空间与极大互局部格} \label{sec:moreonmoduli}

现在我们回到 $\inprod{L_\zeta}$ 所定义于其上的那些模空间的精确定义。

设 $\CL_i$ 标记简单线算子的各种不同的极大互局部组。对每个 $i$ 有一个电荷格 $\Gamma_{\CL_i}$，它们嵌在一个包含关系系统之中\footnote{这样一套包含关系在布尔代数里也被称为一个「格」（lattice）。不过，那个词即便不引起彻底的混淆，大概也会造成别扭，所以我们不用它。}
\begin{equation}
\begin{matrix}
    &  &  \Gamma_{mx} &  &   \\
    & \nearrow &   &  \nwarrow & \\
    \Gamma_{\CL_1} &  & \cdots &  & \Gamma_{\CL_n} \\
    & \nwarrow &   &  \nearrow & \\
    &  &  \Gamma  &  &   \\
\end{matrix}
\end{equation}
其中 $\Gamma_{mx}$ 是诸 $\Gamma_{\CL_i}$ 的并。注意反对称配对 $\inprod{,}$ 在 $\Gamma_{mx}$ 上不必取整数值。这尤其蕴涵：$\Gamma_{mx}$ 一般并不对应任何物理 $\CN=2$ 理论中的电荷格。

每个电荷格 $\CL$ 对应一个略有不同的物理理论。约化到三维时，这一差别变得重要：三维 sigma 模型的靶是依赖于 $\CL$ 的超凯勒流形 $\CM_\CL$。它是一个环面纤维化，$u\in\CB-\CB^{sing}$ 上方的纤维由
\begin{equation}
(\CM_{\CL})_u = {\rm Hom} ((\Gamma_{\CL})_u, \IR/\IZ).
\end{equation}
给出。我们可以类似地定义对应于格 $\Gamma$ 与 $\Gamma_{mx}$ 的空间 $\CM$ 与 $\CM_{mx}$，尽管它们并不对应任何物理的 $\CN=2$ 理论。于是有一族有限的覆盖
\begin{equation}\label{eq:covering-maps}
\begin{matrix}
    &  &  \CM_{mx} &  &   \\
    & \swarrow &   &  \searrow & \\
    \CM_{\CL_1} &  & \cdots &  & \CM_{\CL_n} \\
    & \searrow &   &  \swarrow & \\
    &  &  \CM &  &   \\
\end{matrix}
\end{equation}

事实上，这些流形彼此关系极为密切：$\CM$ 与 $\CM_{mx}$ 都是超凯勒的（正如诸 $\CM_{\CL_i}$ 一样），并且 \eqref{eq:covering-maps}
中所有覆盖映射都是局部等距。的确，它们是环面的纤维式覆盖。这一点容易从 \cite{Gaiotto:2008cd} 的构造推出：度规对纤维坐标 $\tilde\theta$ 的依赖只通过因子 $e^{i\tilde\theta_\gamma}$ 出现，其中 $\gamma\in\Gamma$ 是一个「原味」BPS 粒子的荷。因此，任何使 $e^{i\tilde\theta_\gamma}$ 对一切 $\gamma\in\Gamma$ 都不变的 $\tilde\theta$ 平移，都是一个等距。

我们相信上述图像可以延拓到奇异轨迹 $\CB^{sing}$ 之上。理由如下。在一般的点 $u^*\in\CB^{sing}$ 处，纤维化中的奇异性来自单个电荷为 $\gamma$ 的 BPS 粒子变得无质量。在这种情形，环面中有一个圆在 $u^*$ 上方缩成一点。（用 \cite{Gaiotto:2009hg} 第 13.7 节的结果，这就是由流 $\frac{d\tilde\theta_{\tilde \gamma}}{dt} = \langle \gamma, \tilde \gamma \rangle$ 所描述的那个圆。）至少当时空圆的半径 $R$ 很大时，我们可以用带荷 $\gamma$ 的粒子对朴素的大半径（半平坦）度量的单圈修正，来理解 $u^*$ 附近完整的超凯勒度量的行为。用 \cite{Gaiotto:2008cd} 的结果我们看到：若 $\gamma$ 不可约，则该圆只在纤维中的一处缩掉，即 $\exp i \tilde\theta_\gamma =1$ 处。否则若 $\gamma=Q\gamma_0$，则度量有 $Q$ 个余维二的奇点，它们——一般而言——是成族的 $A_{Q-1}$ 型奇点。用这幅纤维退化的图像，我们看到 \eqref{eq:covering-maps}
中商空间的关系得以保持。

我们将在第 7.3.1 节% TODO \ref{subsec:Hitchin-Isogenous}
再次评注这些问题，并在第 8 节% TODO \ref{sec:Examples-of-M}
中用具体例子说明这些想法。


\subsection{扭模空间的一个 vev} \label{sec:twisted-vev}

既然引入 \eqref{eq:y-dep-trace}
已经打开了潘多拉魔盒，我们应当指出：除了上面讨论的 $y=-1$ 情形之外，对一个密切相关的迹，\emph{确实}存在一种解释：
\begin{equation}\label{eq:y=1-trace}
\inprod{L_\zeta}' := {\rm Tr}_{\CH_{u,L_\zeta}} (-1)^{2I_3} e^{-2 \pi R H} e^{i \theta \cdot \CQ}
\end{equation}
其中我们选取 $u,\theta$，并用这个迹来\emph{定义}左端。同样地，四个超荷全都被保持。额外的因子 $(-1)^{2 I_3}$ 作用在非异类带框 BPS 态上是平凡的。

与 $y=-1$ 的情形一样，$\inprod{L_\zeta}'$ 可以用一个三维 sigma 模型来理解：它由我们的四维理论在 $S^1$ 上紧化得到，只是现在多了一个扭因子 $(-1)^{2 I_3}$。所得的理论和以前一样，是映入某个超凯勒流形的 sigma 模型，我们把它记作 $\tCM$。

$\tCM$ 与 $\CM$ 是什么关系？两者都可以用 \cite{Gaiotto:2008cd} 的方法来描述。那个描述中主要的非平凡成分是一组瞬子修正，它们来自四维理论的 BPS 态，其世界线缠绕紧化圆周 $S^1$。那么因子 $(-1)^{2 I_3}$ 如何改变这个故事呢？如果存在异类 BPS 态\zy{原文 exotic：指在 $su(2)_R$ 下作非平凡变换的 BPS 态，见第 3.2 节。% TODO \ref{subsec:BPS-Deg-PSC}
此处译作「异类」，以免与 singular（奇异）相混。}，则相应的瞬子修正会被这个因子影响。因此，在那种情形 $\tCM$ 会与 $\CM$ 不同。然而在没有异类 BPS 态时，我们预期 $\CM$ 与 $\tCM$ 作为超凯勒流形是同构的（尽管或许不是典范地同构）。至少在我们下面遇到的例子里，情况确实如此。

与上面同样的推理使我们能作展开
\begin{equation} \label{eq:y1}
\inprod{L_\zeta}' = \sum_\gamma \fro(L_\zeta, \gamma, y=1) \tCY_\gamma,
\end{equation}
其中 $\tCY_\gamma$ 是 $\tCM$ 上的一族函数，满足与 $\CY_\gamma$ 类似的性质。重要的一点是：与 $\CY_\gamma$ 不同，这些函数服从\ti{不带扭}的乘法法则
\begin{equation}
\tCY_{\gamma} \tCY_{\gamma'} = \tCY_{\gamma + \gamma'}.
\end{equation}

既然我们已经说过，在没有异类 BPS 态时流形 $\CM$ 与 $\tCM$ 应当重合，人们也许想知道如何看出这一等同。最显然的做法是：把 $\tCM$ 上的函数 $\tCY_\gamma$ 与 $\CM$ 上的函数 $\sigma(\gamma)\CY_\gamma$ 等同起来，其中 $\sigma$ 是某个选定的二次加细。这个等同与 $\CY_\gamma$ 及 $\tCY_\gamma$ 的间断性质相容，当且仅当 $\sigma(\gamma)=(-1)^{2 \CJ_3}$（作用在一切「原味」BPS 态上）。这一性质确实对 \cite{Gaiotto:2009hg} 中为 $A_1$ 理论所定义的典范 $\sigma$ 成立。如果强正性成立，我们就应预期更一般的理论也有类似的简单行为。


\section{六维视角}\label{sec:Six-Dim-View}

到此为止，我们的考虑只用到 $d=4$、$\CN=2$ 超对称。在本文余下部分，我们将聚焦于 \cite{Gaiotto:2009hg} 中记作 class $\CS$（取「六」之意）的那类理论。这些理论来自超共形 $ADE$ $(2,0)$ 理论在带穿刺点的黎曼面 $C$ 上的紧化。它们构成一大批丰富的例子，而六维视角将让我们能构造出一些非常有趣的线算子。对 $A_1$ 理论，我们将能够把线算子的真空期望值完全算出来。


\subsection{六维 $(2,0)$ 理论回顾}\label{subsec:6d-review}

IIB 型弦论或 M 理论中的一般性考虑，提示存在具有超共形对称性的六维定域量子场论 \cite{Witten:1995zh,Strominger:1996ac,Witten:1995em}。人们对这些理论所知甚少。关于已知结果的精彩综述，见 \cite{Seiberg:1997ax,Witten:2009at}。下面我们总结若干与本文相关的已知结果。特别地，我们聚焦于具有 $(2,0)$ 超共形对称性的理论。关于五维与六维超对称的约定之总结，见
附录 E% TODO \ref{app:Six-Super}。

张量多重态的阿贝尔理论是 $(2,0)$ 理论的例子。它们是自由场论，可以用拉氏量描述（在作出某些选择之后 \cite{Belov:2006jd}）。我们转而聚焦于相互作用理论：对它们，人们并不知道、甚至也不指望存在基本的场论自由度，当然更没有拉氏量。这些理论有 $ADE$ 分类；本节中 $\lieg$ 指某个单且单链李代数的紧实形式。

在 $A_r$ 理论的情形，可以利用 DLCQ 以及瞬子模空间上的量子力学给出一个部分定义 \cite{Aharony:1997th,Aharony:1997an}。关于 $ADE$ 理论，我们所知道的基于两条我们将当作公理来用的陈述。一旦这些理论被正确定义，它们就应当成为定理：

\begin{enumerate}

\item 当理论紧化在 $\IR^{1,4}\times S^1$ 上（$S^1$ 半径为 $R$，费米子取周期边界条件）时，长距离动力学由极大超对称的 Yang--Mills 理论支配，其规范李代数为 $\lieg$，耦合常数 $g^2_{YM}$ 正比于 $R$。\footnote{关于按我们约定给出的精确归一化，见 \cite{Gaiotto:2009hg} 的式 (3.36)。}

\item $\IR^{1,5}$ 上的理论有一个真空模空间
\begin{equation}
\CM(\IR^{1,5}) = (\IR^5\otimes \liet)/\CW,
\end{equation}
其中 $\liet$ 是 $\lieg$ 的嘉当子代数，$\CW$ 是 Weyl 群。在 $\CM(\IR^{1,5})$ 的光滑点处，低能动力学由取值于 $\liet$ 的自由 $(2,0)$ 张量多重态描述。
\end{enumerate}

当然，这两种形变必须相容。例如，五维超对称 Yang--Mills 理论的真空模空间可以等同于 $\CM(\IR^{1,5})$，因为 $so(5)_R$ 的 ${\bf 5}$ 中伴随标量的真空期望值 $\langle Y^I \rangle$ 必须可同时对角化。模空间的一个点因而被等同于 $\lieg$ 中一组有序的、可同时对角化的五元组，模掉规范不变性。正如这些 vev 在五维 SYM 中把 $\lieg$ 对称性自发破缺到嘉当子代数 $\liet$，我们应当把六维理论的「规范对称性」设想为自发破缺到取值于 $\liet$ 的自由张量多重态的阿贝尔规范群。由于这个缘故，模空间 $\CM(\IR^{1,5})$ 常被称为「库仑分支」。

此外，在 $A_r$ 理论的情形，有一个很好的直观图像：该理论是 $(r+1)$ 张重合 M5 膜的世界体，处在引力解耦极限下，并已把质心自由度除掉 \cite{Strominger:1996ac}。$so(5)$ R 对称性于是被解释为垂直于五维膜的空间中的旋转，而模空间 $(\IR^5 \otimes \IR^{r+1})/S_{r+1}$ 参数化 $r+1$ 张平行、单次缠绕的 M5 膜在横空间中的位置。单张 M5 膜的低能动力学由 $U(1)$ 张量多重态描述，标量场 $Y^I$ 表示横向涨落。模空间的一个点因而可以由标量场的真空期望值 $\langle Y^I_s \rangle$（$s=1,\dots, r+1$）指定。$A_r$ 理论是在把描述整体质心的那个张量多重态「解耦」之后得到的。

$(2,0)$ 理论可以定义在某些洛伦兹六维流形 $M_6$ 上，其配分函数可以维克转动到欧几里得六维流形。张量多重态中规范势 $B$ 的 $3$-形式场强被约束为自对偶。这蕴涵 $M_6$ 必须可定向，并且还必须配备某些额外的拓扑数据。\footnote{技术上，是第四 Wu 类的一个微分积分提升 \cite{Hopkins:2002rd}。关于这一条件的物理讨论，见 \cite{Witten:1996hc} 与 \cite{Belov:2006jd}。}
此外，为了定义费米子，它还必须配备自旋结构。

最后，我们提一下这些六维理论中一个非常微妙的方面，此处不打算深入讨论。六维理论不是标准的量子场论，而大概最好看成取值于一个可逆七维拓扑场论的六维理论。\footnote{我们在此不去细究这一概念的确切细节。感谢 D. Freed 与 E. Witten 就这一问题所作的有益讨论。}「向量值量子场论」这个词有时也被使用，因为这些理论在六维流形上的配分函数取值于一个向量空间，而不只是一个数。这一微妙之处在圆周紧化时就显现出来。我们上面说五维 SYM 的规范李代数是 $\lieg$，但要指定理论，还必须选一个李代数为 $\lieg$ 的规范群 $G$。这个规范群的选取并不唯一。必须指定额外的数据，它们决定配分函数向量空间到一维子空间的一个投影 \cite{Witten:2009at}。物理上，这样的投影或许可以由五维理论中 't Hooft 通量的一个选取来指定。这一点在 ${\cal N}=4$ SYM 的 S 对偶性质中已经可见；该理论正是六维理论的一个环面紧化：弱耦合时，在环面的「小」圆上约化会选定一个五维规范群，随后在第二个圆上约化又出现同一个规范群。S 对偶交换这两个圆，并交换对应于 Langlands 对偶规范群的两个投影。


\subsubsection{手征原初场}

$A_r$ 理论的 DLCQ 定义表明，存在生成超共形代数短多重态的定域手征原初场 \cite{Aharony:1997an,Bhattacharya:2008zy}。在圆周紧化之后，它们可以被等同于由标量场的 Casimir 不变量构造的规范不变保护算子，并在 $SO(5)$ 指标上对称且无迹。在 $A_r$ 情形，它们下落到五维 SYM 中的算子 ${\Tr} Y^{(I_1 \cdots} Y^{I_k)}$。它们的标度维数为 $2k$。看来可以合理地假定：这样的手征原初多重态不仅对 $A_r$ 存在，而且对一切 $ADE$ 理论都存在；我们在下文中作此假定。一个秩为 $d$ 的 Casimir 给出一个维数为 $2d$ 的算子，它处在 $SO(5)$ 的对称无迹表示中，带 $d$ 个指标。


\subsubsection{动力学 BPS 弦}

六维非阿贝尔理论中有动力学弦。在 $A_r$ 理论的 M 理论推导中，它们来自张开在 M5 膜之间的开 M2 膜。单个 BPS 弦态由根 $\alpha\in\Phi(\lieg)$ 标记，其中 $\Phi(\lieg)$ 是 $\lieg$ 的根系。沿弦的单位切向量为 $\hat t^M$ 的弦，在超对称代数中有中心荷 $Z_M^I = \hat t_M z^I$。弦的张力正比于 $\parallel \vec z \parallel$，其中 $\parallel \cdot \parallel$ 是 $\IR^5$ 上的欧几里得范数。库仑分支上位于点 $Y\in \IR^5 \otimes \liet$ 的真空，其弦激发有 $\vec z$ 正比于 $\langle \alpha, Y\rangle$。应当把这些弦看成类似于非阿贝尔规范理论中的「非对角」规范玻色子。

我们可以通过在 $S^1$ 上紧化并利用所预期的到五维 SYM 的约化，来核验这些断言。BPS 弦下落到五维 SYM 的 BPS 粒子。对应于单次缠绕弦的粒子，对应于荷为 $\alpha$ 的规范玻色子，质量为 $\sim R \lVert \langle \alpha, Y\rangle \rVert$。当 $R$ 小时不应有其他轻粒子（否则它们必须由五维 SYM 多重态中不存在的场来描述），所以例如不应存在由 $\lieg$ 的其他权向量标记的弦。

六维的动力学弦也在五维 SYM 中给出 BPS 弦态。在库仑分支上，可以把它们半经典地描述为垂直于弦世界面的三维空间中的磁单极子。这样的解因而由 $\Lambda_{G} = {\rm Hom}(U(1),T)$ 标记，即 $G$ 的余特征标格。对于\emph{单链}李代数，在选取 Killing 形式的归一化使根长为 $2$ 之后，$\Lambda_{cr}$ 与 $\Lambda_{r}$ 同构。由于这个缘故，五维中的弦在取规范群为 $\tilde G$（与 $\lieg$ 相伴的单连通紧群）时最容易理解。

为完整起见，我们提一下：五维 SYM 中还有由世界线横向的四个维度中的瞬子所描述的 BPS 粒子。它们的质量正比于 $n/R$，其中 $n$ 是瞬子荷，可以把它们解释为由 $S^1$ 上的紧化产生的 KK 模式。


\subsubsection{面算子}

除了手征原初场与动力学弦之外，六维 $(2,0)$ 理论及其五维 SYM 紧化中还有面算子。本节的主要目标正是用它们来定义线算子，而后者是本文余下部分的研究重点。

我们先给出面算子的一个一般定义，它与线算子的定义密切类似：$\IR^{1,1}\subset\IR^{1,5}$ 上的面算子，其横空间是 $\IR^4$。度规共形等价于 $AdS_3\times S^3$，其等距群的李代数由 $so(2,2)\oplus so(4)$ 给出。我们把嵌入的 $\IR^{1,1}$ 上的面算子定义为：在 $AdS_3\times S^3$ 上强加一个共形不变的边界条件所得的结果。注意 $so(2,2)\oplus so(4) \cong (sl(2,\IR) \oplus su(2)) \oplus (sl(2,\IR) \oplus su(2))\cong so(4^*) \oplus so(4^*)$。

与线算子一样，我们将聚焦于那些保持若干 R 对称性与若干超对称的特殊面算子。为说明其动机，先考虑库仑分支上低能张量多重态理论中的面算子。若 $\Sigma$ 是 $M_6$ 中的定向曲面，则可以定义自对偶场的和乐\footnote{在下文的讨论中，为简单起见我们假定 $B$ 在拓扑上平凡。对一般的时空与拓扑非平凡的 $B$ 场，我们应改用 Cheeger--Simons 微分特征标。}
\begin{equation}\label{eq:self-dual-holon}
h_v(\Sigma) = \exp[ 2\pi i \int_{\Sigma} (v, B) ]
\end{equation}
其中 $v \in \Lambda_{wt}$ 属于 $\tilde G$ 的特征标群，即权重格。我们强加这一限制，是因为在拓扑非平凡的流形上，$B$ 场的大规范变换取 $B \to B + \omega$ 的形式，其中 $\omega$ 是整体定义的闭 2-形式，其周期取值于 $\Lambda_{cr}$。注意这里以及第 7 节% TODO \ref{sec:Six-Dim-View}
余下部分中，$\Sigma$ 与红外 Seiberg--Witten 曲线\emph{毫无}关系！

在某些情形下，这个和乐可以提升为一个超对称面算子。为简单起见，考虑规范群为 $U(1)$ 的单个张量多重态。于是考虑
\begin{equation}\label{eq:single-tm-susy}
\exp\left[ 2\pi i \int_{\Sigma} B + \kappa n^I Y^I \vol(\Sigma)\right]
\end{equation}
其中 $\vol(\Sigma)$ 是由诱导度规给出的 $\Sigma$ 上的体积形式，$\kappa$ 是常数，并且不失一般性可以假定 $\vec n$ 是欧几里得 $\IR^5$ 中的单位向量，其分量为 $n^I$。设 $\xi^\alpha$（$\alpha = 1,2$）是 $\Sigma$ 上的局部坐标系。那么超对称 $\epsilon_i^r Q_r^i$ 湮灭此算子，当且仅当
\begin{equation}\label{eq:surface-susy-cond}
\epsilon_i^r \left( \frac{d\xi^\alpha \wedge d \xi^\beta \p_\alpha X^M \p_\beta X^N}{\vol(\Sigma) }(\gamma_{MN})_r^{~~s} \delta^i_{~~j} + \kappa (n^I \Gamma^I)^i_{~~j} \delta_{r}^{~~s} \right) =0
\end{equation}
其中 $X^M(\xi)$ 表示曲面到 $M_6$ 中的嵌入。为保持超对称，此式必须对\emph{常数}的未破缺超对称 $\epsilon_i^r Q_r^i$ 成立。对平坦曲面与常数 $n^I$，当 $\kappa=\pm 1$ 时有一半超对称被保持。更一般地（类似于超对称 Yang--Mills），$\Sigma$ 可以是弯曲曲面，$n^I$ 也可以变动。一个在下文会变得重要的例子是曲面为 $\IR \times \wp$ 的情形，其中 $\wp$ 是比如 $12$ 平面中的一条曲线。若把 R 对称性空间分解为 $\IR^5 = \IR^2 \oplus \IR^3$，把 $\IR^2$ 加项等同于 $12$ 平面，并取 $n^I$ 为 $\wp$ 的单位切向量，则保持四分之一的超对称。一般而言，R 对称性被方向 $n^I$ 破缺到 $so(4)$。在比如 $01$ 方向的平面的特殊情形，我们可以用与线算子的 $osp(4^*\vert 2)_\zeta$ 讨论密切类似的方式描述所保持的超对称。我们通过沿曲面所在平面作反射、再作垂直于 $n^I$ 的平面内的 R 对称性反射，来定义 $osp(8^*\vert 4)$ 的一个对合。此对合的不动点定义一个超共形代数，我们可以把它记作 $(osp(4^*\vert 2) \oplus osp(4^*\vert 2))_{\vec n}$。

现在可以把 \eqref{eq:self-dual-holon}
与 \eqref{eq:single-tm-susy}
结合起来，在阿贝尔张量多重态理论中定义超对称面算子 $\IS_{\vec n}(v,\Sigma)$。于是很自然地猜想：非阿贝尔理论中也存在相应的面算子 $\IS_{\vec n}(\CR,\Sigma)$，它与 $\lieg$ 的表示 $\CR$ 相伴，并保持超共形对称性 $(osp(4^*\vert 2) \oplus osp(4^*\vert 2))_{\vec n}$。像往常一样，我们可以用两个可用的形变来核验这一断言。当 $Y$ 的大期望值把规范对称性强力破缺时，这些面算子将被半经典地描述为阿贝尔面算子 $\IS_{\vec n}(v,\Sigma)$ 的求和，其中 $v$ 遍历表示 $\CR$ 的权重。作为第二项核验，我们考虑在半径为 $R$ 的圆周上紧化后出现的与五维 SYM 的关系。若 $\Sigma = S^1\times \wp$，其中 $\wp\subset M_5$ 是一条路径，则非阿贝尔面算子的红外极限应当是
\begin{equation}\label{eq:5d-line-op}
L_{\vec n}(\CR, \wp) = {\Tr}_{\CR} P\exp \left[ \int_\wp  A + \kappa n\cdot R Y ds \right]
\end{equation}
这些确实就是五维 SYM 中的 BPS 线算子。最后，非阿贝尔面算子的存在与动力学 BPS 弦的存在密切相连。例如，可以从一个 $su(r+1)$ 理论出发，走到部分库仑分支，在那里理论约化为 $su(r) \times u(1)$ 理论。将存在带 $u(1)$ 阿贝尔因子之荷的动力学弦，其张力当标量期望值趋于无穷时趋于无穷，而 $u(1)$ 因子解耦。在此极限下，动力学弦成为面算子 $\IS_{\vec n}(v,\Sigma)$，其中 $\vec n$ 是被送到无穷远的 $Y$ 的方向。

众所周知，把诸如 \eqref{eq:5d-line-op}
这样的线算子插入路积分，可以解释为把背景理论与一个 $0+1$ 维量子力学系统（由余伴随轨道的量子化定义）耦合起来 \cite{Alekseev:1988vx}。类似地，面算子可以由一个定义在 $\Sigma$ 上的 $1+1$ 维量子场论与背景理论耦合来表示。在这一解释中要注意：正如人们从 $AdS_3/CFT_2$ 对应中所熟悉的，等距的李代数 $so(2,2)\cong sl(2,\IR)_l \oplus sl(2,\IR)_r$ 可以解释为分别作用在左行与右行自由度上的共形对称代数。这一解释可以延拓，使我们能把紫外理论中未破缺的超共形对称性解释为 $osp(4^*\vert 2)_l \oplus osp(4^*\vert 2)_r$。此外，如我们已看到的，在红外，由 $\IS_{\vec n}(\CR,\Sigma)$ 下落的阿贝尔面算子由表示 $\CR$ 的权重 $v$ 标记。这提示：定义在 $\Sigma$ 上的那个量子场论应当有以 $v$ 标记的真空。

动力学弦可以终止在面算子上，它们在 $\Sigma$ 上的二维量子场论中定义畴壁。由根 $\alpha \in \Phi(\lieg)$ 标记的弦，若 $v_1 - v_2 = \pm \alpha$，就是真空 $v_1$ 与 $v_2$ 之间的畴壁。这样的构型可以是 $\frac{1}{4}$ BPS 的，其方式与动力学弦之间的汇点成为 $\frac{1}{4}$ BPS 构型的方式十分相似。一个典型的例子是静态构型：所有弦与面算子都位于一个平面上，中心荷向量 $\langle \alpha_i,Y\rangle$ 或 $\vec n$ 也都位于中心荷空间中的一个平面上，并且两个平面中的斜率相同。注意，若关心的是 $\frac{1}{4}$ BPS 构型，则面算子可以允许比平面中的直线更一般的形状，办法是让 $\vec n$ 沿算子变化，跟随切向量的变化。

% =====================================================================
%  GMN3 §7.2–7.5  扭转理论在黎曼面上的紧化 / 模空间 / 来自六维的线算子 / 线算子的编目
%  原文：Lamination-Paper.tex 4379--4685 行；PDF p.67--72（印刷 p.66--71）
%  编号公式 (7.6)--(7.14) 共 9 条（§7 共 14 条，§7.1 已用去 (7.1)--(7.5)）；
%  块内引用用 \eqref；跨块引用（式 (7.4)、(2.2) 及各节号）暂写原文编号并加 % TODO
% =====================================================================

\subsection{扭转理论在黎曼面上的紧化}\label{subsec:Twisted-C}

我们要考虑把 $(2,0)$ 理论缠绕在黎曼面 $C$ 上，这正是 \cite{Gaiotto:2009hg} 所用的构造。如 \cite{Gaiotto:2009hg} 第 3.1.2 节所述，为使这一紧化保持 $d=4,\CN=2$ 超对称，我们必须对理论作扭转。紧化 $\IR^{1,3}\times C$ 把局部洛伦兹对称性破缺为 $so(1,3)\oplus so(2)_C$。随后我们显式地破缺 $R$ 对称性 $so(2)'\oplus so(3)'\subset so(5)$，并对理论作扭转，使 $C$ 上的自旋联络与对角子代数 $so(2)_d\subset so(2)_C\oplus so(2)'$ 耦合。有八个超荷存活下来，它们在四维中构成一个 $\CN=2$ 超对称代数。$so(3)'$ 等同于 $d=4,\CN=2$ 的 $R$ 对称性，$so(2)_d$ 则是 $d=4,\CN=2$ 的反常 $u(1)_R$ 对称性。指标为 $d_k$ 的那族手征主算子中，含有一个 $so(2)_d$ 荷最大的算子，记作 $\CO_k$，于是 vev $\langle\CO_k\rangle$ 成为 $C$ 上的全纯 $k$-微分，而四维理论的库仑分支由这些算子的 vev 参数化（见 \cite{Gaiotto:2009hg} 式 (3.5)）：
\begin{equation}
\CM(\IR^{1,3}\times C) = \CB = \oplus H^0(C, K^{\otimes d_k})
\end{equation}
与 \cite{Gaiotto:2009hg} 中一样，重要的是引入安置在 $C$ 的穿刺点处的余维 2 缺陷。这些缺陷不再破缺任何超对称。它们的作用是在 $\langle\CO_k\rangle$ 中制造出极点。四维红外理论由 $C$ 的选取以及 $\langle\CO_k\rangle$ 在 $C$ 上各缺陷点处的奇异性决定。规范耦合的模空间与带穿刺点的黎曼面 $C$ 的复结构模空间重合。

\subsection{模空间}

如 \cite{Gaiotto:2009hg} 中详细解释的那样，真空模空间
\begin{equation}
\CM = \CM(\IR^{1,2}\times S^1\times C)
\end{equation}
可以等同于 $C$ 上 Hitchin 方程之解的模空间，这些解在缺陷处具有确定的奇异性。这些方程是关于联络 $A$ 与 1-形式 $\varphi$ 的方程，其内容是：复联络
\begin{equation} \label{eq:flat-conn}
\CA(\zeta) = R\frac{\varphi}{\zeta} + A + R\zeta\bar\varphi
\end{equation}
对一切 $\zeta\in\IC^\times$ 都平坦。事实上，在其复结构 $J^{(\zeta)}$ 下，模空间 $\CM$ 恰好等同于 $C$ 上平坦 $G_\IC$-联络 $\CA(\zeta)$ 的模空间，而这些联络在 $C$ 的穿刺点处具有奇异性。规范群 $G$ 的李代数为 $\lieg$。

穿刺点处奇异性的性质，取决于我们在 $C$ 上插入的具体缺陷。最简单的可能是「正则穿刺点」，它在 $A$ 与 $\varphi$ 中产生单极点，从而在 $\CA(\zeta)$ 中产生正则奇点。这样的穿刺点确定 $G_\IC$ 中一个（依赖于 $\zeta$ 的）共轭类；$\CA(\zeta)$ 的单值性被限制在该共轭类中取值。人们也可以考虑「非正则穿刺点」，它们产生更为复杂的非正则奇点。\zy{「正则/非正则穿刺点」是按其在 $\CA(\zeta)$ 中所产生的正则/非正则奇点命名的，与穿刺点本身无关。}

我们可以更具体地描述 $\CM$。若只考虑正则穿刺点 $\CP_i$，则平坦联络在模去规范变换之后完全由其单值表示捕获：于是 $\CM$ 是 $\pi_1(C\setminus\{\CP_i\})$ 到 $G_\IC$ 的表示所构成的空间，其中绕穿刺点的单值性的共轭类固定。\footnote{我们说「一个表示空间」而不说「\emph{那个}表示空间」，是因为涉及离散商的微妙之处，这一点稍后会讨论。} 再具体一些：取基点 $x\in C$，取绕 $C$ 上 $A$ 圈与 $B$ 圈的基于基点的环路，以及绕穿刺点 $\CP_i$ 的基于基点的环路；于是我们研究的是满足
\begin{equation}
M_1 \cdots M_n = \prod_{i=1}^{g} A_i B_i A_i^{-1} B_i^{-1},
\end{equation}
的多元组 $\{A_1,\dots,A_g,B_1,\dots,B_g,M_1,\dots,M_n\}\in G_\IC^{2g+n}$，其中 $M_i$ 取在固定的共轭类中。这些多元组要模去 $G_\IC$ 对所有 $A_i$、$B_i$、$M_i$ 的同时共轭作用。若考虑非正则穿刺点，情况就稍微复杂一些：在这样的奇点处，局部信息不止单值性；平坦截面的渐近行为还额外给出每个 $M_i$ 到斯托克斯因子与形式单值性的分解。更深入的讨论见 \cite{Witten:2007td}。

我们将在下文看到，线算子的 vev 是 $\CM$ 上的自然函数，例如上述诸矩阵之乘积的迹。这一描述使我们得以研究它们的半环关系，并与我们在第 4 节% TODO \ref{sec:Formal-line-operators}
中所作的讨论相比较。另一方面，要研究线算子与自旋特征标之间的关系，我们需要辨认出第 6 节% TODO \ref{sec:linehol}
中出现过的函数 $\CY_\gamma$。它们并不能由上述矩阵直接给出；事实上，对一般的 $G$，我们根本不知道如何构造它们（除非去解 \cite{Gaiotto:2008cd} 的积分方程）。然而对 $A_1$ 理论，我们确实知道 $\CY_\gamma$ 是什么\cite{Gaiotto:2009hg}，这一事实将在第 10 节% TODO \ref{sec:Examples}
中得到利用。

\subsubsection{同源的 Hitchin 模空间}\label{subsec:Hitchin-Isogenous}

我们在第 6.3 节% TODO \ref{sec:moreonmoduli}
中的一般讨论使我们预期：模空间不止一个 $\CM$，而是有若干个密切相关的 $\CM_\CL$，它们由线算子的族 $\CL$ 标记。这一选择在 Hitchin 系统中如何体现出来？

在定义 Hitchin 方程之解的模空间时，人们总得模去规范变换群。一种可能是使用取值为 $G$ 的单连通形式的规范变换，这给出某个模空间 $\CM_{sc}$。然而人们也可以更宽松一些，允许取值为 $G$ 的伴随形式的规范变换。这两种选择的差别是一个有限群 $\Delta$，它同构于 $C$ 上结构群为 $\CZ=Z(\tilde G)$ 的平坦丛所成的群。$\Delta$ 通过等距作用在 $\CM_{sc}$ 上，因此我们可以考虑 $\CM_{sc}$ 对 $\Delta$ 的各个子群取商。我们希望把这些商与各个 $\CM_\CL$ 对应起来。

在具体例子中，我们可以更精确地陈述 $\CL$ 的选取与 $\Delta$ 的子群之间的词典。对带正则奇异性的 $A_1$ 理论，关系特别简单：稍后我们就会看到，$\CL$ 的一个选取将对应于 $C$ 上一组闭路径 $\wp$ 的选取。我们的提法是：$\CM_\CL$ 是 $\CM_{sc}$ 对 $\Delta$ 的某个子群取商，该子群对应于沿 $\CL$ 中所有路径 $\wp$ 都具有平凡和乐的那些丛。对带非正则奇异性的 $A_1$ 理论，情况更为困难：特别地，$\CM_{sc}$ 并不是出现的最大的空间——看来还需要对所允许的规范变换施加进一步的离散限制。本文不打算作完整的分析，但会给出几个例子。

四维规范群的选取与 Hitchin 系统中允许的规范变换之间的关系，其六维动机是什么？我们没有完整的故事，但确实有一些有用的线索。尽管 Hitchin 系统背后的李代数由六维 SCFT 的选取决定，规范群的具体选取却要微妙得多。这些问题与上文第 7.1 节% TODO \ref{subsec:6d-review}
中讨论的沿圆约化时的问题类似。想来，在黎曼面上紧化时，四维规范群的选取将对应于配分函数向量空间上某种投影的选取。当我们再沿一个圆从四维约化到三维时，就没有进一步选择投影的余地了：我们对 $\CN=2$ 理论四维规范群所作的任何选取，都将决定 $C$ 上五维理论的性质。

\subsubsection{扭局部系统} \label{sec:twisted-loc}

还有第二个需要关注的整体问题。在第 6.4 节% TODO \ref{sec:twisted-vev}
中我们强调过：除了以通常方式把理论在 $S^1$ 上紧化所得的模空间 $\CM$ 之外，还存在一个先验地不同的空间 $\tCM$，它出现在以 $(-1)^{2\CJ_3}$ 作扭转的紧化中。这个空间是什么？

我们的提法是：$\tCM$ 是 $\CM$ 的一个略微扭转的版本。这里不描述 Hitchin 方程相应的扭转；我们只描述 $\tCM$ 作为复流形、在其复结构 $J^{(\zeta)}$（$\zeta\in\IC^\times$）下是什么样子。它看起来非常像 $\CM$，只不过把平坦联络换成了\emph{扭}平坦联络。$C$ 上的扭平坦 $SL(2,\IC)$-联络，是 $C$ 的单位切丛上的平坦 $SL(2,\IC)$-联络，它绕每条纤维的和乐为 $-1$。（注意，若把上一句中的 $-1$ 换成 $+1$，我们就退回到通常的平坦联络。）

如果能从 $(2,0)$ 理论的物理中直接导出这一扭转，那将是可取的。就本文的目的而言，我们满足于一个自洽性论证：如此定义的 $\tCM$ 将具有我们在第 6.4 节% TODO \ref{sec:twisted-vev}
中描述的全部性质。特别地，这一扭转很好地消解了一些棘手的符号问题。

在 $C$ 上选取一个自旋结构，给出同构 $\tCM\simeq\CM$。这是因为：固定一个自旋结构，等价于固定 $C$ 的单位切丛上一个取值为 $\IZ_2$ 的平坦局部系统，它绕每条纤维的和乐为 $-1$；与该局部系统作张量积即得所需的同构。于是这两个空间同构，但并非典范同构。一个例外是 $C$ 亏格为零的情形：此时只有一个自旋结构，因此这两个空间确实是典范同构的。

扭局部系统的模空间也曾在 \cite{MR2233852,Frenkel:2007tx} 中、出于密切相关的（但并不完全相同的）理由被考虑过。

\subsection{来自六维的线算子}

我们可以用相互作用 $(2,0)$ 理论的非阿贝尔面算子 $\IS_{\vec n}(\CR,\Sigma)$，在 $d=4,\CN=2$ 理论中定义线算子 $L_\zeta(\CR,q;\wp)$，其中 $q$ 是 $\IR^{1,3}$ 中的一条路径，$\wp\subset C$ 是一条不自交的路径。令 $\Sigma=\wp\times q$，并在第 7.2 节% TODO \ref{subsec:Twisted-C}
的扭转理论中工作。此时扭转已经破缺了 $R$ 对称性 $5=2\oplus3$，而我们希望保留 $so(3)'$ 这部分 $R$ 对称性，因此 $\vec n$ 应当是分解 $5=2\oplus3$ 中的那个二维子空间里的向量。现在取 $q=\{\vec 0\}\times\IR$，即 $\IR^{1,3}$ 中沿时间方向的一条直线。若把面算子所保持的未破缺超对称的条件——式 (7.4)% TODO \eqref{eq:surface-susy-cond}
——写出来，就会发现：只要我们允许 $\vec n$ 沿路径变化，它们实际上就是第 2 节% TODO \ref{sec:Susy-BPS-line}
中式 (2.2)% TODO \eqref{eq:Q-zeta-susy}
那种形式的四维超对称。其实这里应当说得更精确些：扭转之后，$\vec n$ 是曲线的切丛中的一个向量。所能得到的最好的「$\vec n$ 为常向量」的概念，是让 $\vec n$ 沿 $\wp$ 平移，使得 $\vec n$ 与 $\wp$ 的切向量之间的夹角 $\beta$ 保持为常数。这个选择正是保持四个超对称的正确选择，这些超对称由相位 $\zeta=e^{i\beta}$ 标记。

现在，正如我们在第 6 节% TODO \ref{sec:linehol}
中注意到的，若把时间作 Wick 转动成一个半径为 $R$ 的欧氏圆，则 vev $\langle L_\zeta(\CR,q;\wp)\rangle$（此时 $q$ 是欧氏时间圆）可以解释为超凯勒模空间 $\CM$ 上的全纯函数。利用六维的观点，我们可以给这些全纯函数一个很好的解释。

首先考虑面算子的红外极限：它由张量多重态理论中取值为 $\liet$ 的面算子写出。令
\begin{equation}\label{eq:varph-Y}
\varphi = \half\left( Y^1 + i Y^2 \right)
\end{equation}
（与 \cite{Gaiotto:2009hg} 的式 (3.31) 相同。）由于扭转，它成为 $C$ 上的一个 $(1,0)$-形式。于是对有限的 $R$，在红外极限下必须把场替换为它们沿 $S^1$ 的零模，从而
\begin{equation}
2\pi i \int_{q\times \wp} (\vec n\cdot Y) \vol(\Sigma) \longrightarrow \pi \int_\wp \left( \frac{R \varphi}{\zeta} + R \zeta \bar \varphi \right)
\end{equation}
（整体的 $i$ 因子被时间圆的欧氏延拓去掉。）类似地，$2\pi i\int_{q\times\wp}(v,B)\longrightarrow \int_\wp(v,A)$，其中 $A$ 是五维阿贝尔超 Yang--Mills 理论的一个阿贝尔规范场，$v$ 是 $\CR$ 中的一个权向量。\footnote{右端不出现 $2\pi$ 因子，是因为我们对规范场采用标准的归一化约定，使得 $F/2\pi$ 具有整周期。}

回忆：非阿贝尔 $(2,0)$ 理论在 $S^1$ 上的紧化给出五维非阿贝尔超 Yang--Mills 理论。现在 \eqref{eq:varph-Y} 定义了 $C$ 上一个取值于伴随表示的 $(1,0)$-形式 $\varphi$。鉴于上面关于阿贝尔理论的论断，自然可以猜想：非阿贝尔面算子的 vev 为
\begin{equation}\label{eq:surface-op-vev-to-5d}
\left\langle \IS_{\vec n}(\CR, q\times \wp) \right\rangle = \left\langle \Tr_{\CR} P \exp \int_\wp \left( \frac{\pi R \varphi}{\zeta} + A + \pi R \zeta \bar \varphi \right) \right\rangle,
\end{equation}
其中 $\varphi$ 与 $A$ 现在取值于规范代数为 $\lieg$ 的五维非阿贝尔 SYM 多重态，右端的 vev 是在五维 SYM 理论中取的。这无非是说：非阿贝尔面算子应当下降为五维非阿贝尔理论的 Wilson 环。左端的 vev 是 $\CM$ 上的全纯函数。此外，得益于扭转，它与 $C$ 的尺度无关，因此可以把 $C$ 取得任意大。最后，五维 SYM 理论的红外极限是自由的。这些事实强烈提示：\eqref{eq:surface-op-vev-to-5d} 中的等式是精确的，而且右端可以用其经典值代替。于是 \eqref{eq:surface-op-vev-to-5d} 的右端就是 $\CA(\zeta)$ 在 $\wp$ 上的和乐，我们最终得到本小节的主要结果：
\begin{equation}\label{eq:L-5d-Holon}
\langle L_\zeta(\CR,q;\wp) \rangle = \Tr_\CR \Hol_\wp \CA(\zeta).
\end{equation}
答案只通过 $\wp$ 的同伦类依赖于 $\wp$。这反映了如下事实：扭转已经消除了对 $C$ 上度规尺度的依赖。

若考虑如第 7.3.2 节% TODO \ref{sec:twisted-loc}
那样的、在 $S^1$ 上的扭转紧化，结果 \eqref{eq:L-5d-Holon} 只需稍作修改。对 $C$ 中任一条有向的、闭的、不自交的路径 $\wp$，都可以（用单位速度参数化 $\wp$）对应出单位切丛中一条有向闭路径 $\tilde\wp$。于是可以取扭平坦联络 $\tilde\CA$ 沿 $\tilde\wp$ 的和乐：
\begin{equation}\label{eq:L-5d-Holon-twisted}
\langle L_\zeta(\CR,q;\wp) \rangle' = \Tr_\CR \Hol_{\tilde\wp} \tilde\CA(\zeta).
\end{equation}
这给出 $\tCM$ 上的一个全纯函数。

\subsection{线算子的编目}

关于线算子的这一实现，我们再作几点评论。

\begin{enumerate}

\item 上面我们讨论的是与闭的、不自交的路径相联系的线算子。那么确有自交的路径又如何呢？表达式 $\Tr_\CR \Hol_\wp\CA(\zeta)$ 对自交路径 $\wp$ 显然有意义，但我们对六维理论中自交面算子的物理理解得还不够，无法把它们纳入讨论。人们或许还会想到 $C$ 上的「汇点」：在那里，对应于三个或更多表示 $\CR_i$ 的面算子汇合在一起。在五维 Yang--Mills 中，这类交点的数目将由 $\otimes_i\CR_i$ 中平凡表示的个数来计数：Wilson 环的端点会与相应的缠结算子缩并。$C$ 上的二维 Toda 理论（\cite{Drukker:2010jp}）中存在类似的汇点，而该理论与四维规范理论的 $S^4$ 紧化有关。这提示它们也应当是六维理论的一个特征，并给出四维理论中的线算子。我们在此不再深究。

\item 我们的构造有一些冗余。例如，\cite{Drukker:2009id} 论证道：当 ADE 群为 $A_1$ 且 $C$ 只有正则穿刺点时，沿 $C$ 上不自交路径、携带 $SU(2)$ 基础表示的面算子，已足以描述全部紫外规范理论算子：似乎并不需要更高的表示。

\item 另一方面，若考虑带非正则穿刺点的 $C$，我们迄今的讨论并没有给出全部线算子：还需要纳入涉及 $C$ 上开路径的额外算子。这里不给出这类算子的一般构造，但将在下文第 9.2 节% TODO \ref{sec:missing}
中描述它们在 $A_1$ 理论中必定呈现的样子。

\item 最后，由于互局部性的约束，这里讨论的线算子并不能同时全部引入。允许的线算子集合 $\CL$ 的选取，在这一语境下对应于对我们所允许的表示与闭路径施加某些限制。例如在 $A_1$ 理论中，我们可以自由地允许携带 $SU(2)$ 整数自旋表示的任意路径，但只能在某子集 $S$ 的路径上允许半整数自旋。局部性要求 $S$ 的选取使得 $S$ 中路径之间的相交数总是\emph{偶数}。相应的模空间 $\CM_\CL$ 是 $SU(2)$ Hitchin 系统的模空间，再模去沿 $S$ 中所有路径都具有平凡和乐的 $\IZ_2$-联络所成的群。

\end{enumerate}

% =====================================================================
%  GMN3 §8  模空间与全纯函数的例子
%  （Some examples of the moduli spaces and holomorphic functions）
%  原文：Lamination-Paper.tex 4686--5050 行；PDF p.73--78（印刷 p.72--77）
%  编号公式 (8.1)--(8.19) 共 19 条，全部在本块内；
%  跨块引用（节号、式号）暂写原文编号并加 % TODO，合稿时统一改 \ref/\eqref
%  label 名照抄原文：eq:N3-manifold, eq:A1-covering-maps, eq:nf0-matcond,
%  eq:monod-equa-nstar, eq:moduli-equation, eq:CH-thm, eq:A1-covering-maps-nstar
% =====================================================================

\section{模空间与全纯函数的例子}
\label{sec:Examples-of-M}

本节在本文所用的那些具体理论例子中，描述模空间 $\CM$，以及它们与第 6.3、6.4、7.3.1、7.3.2 节% TODO \ref{sec:moreonmoduli}、\ref{sec:twisted-vev}、\ref{subsec:Hitchin-Isogenous}、\ref{sec:twisted-loc}
所引入的空间 $\CM_{\CL_i}$ 与 $\widetilde{\CM}$ 的关系。我们还将描述这些空间上的函数代数，如何与第 4 节% TODO \ref{sec:Formal-line-operators}
所述线算子的非交换代数、以及第 10 节% TODO \ref{sec:Examples}
所述的叶层相关联。

\subsection{$U(1)$ 规范理论与周期 Taub-NUT}

我们从第 4.2 节% TODO \ref{subsec:Formal-line-ops-U1}
的 $U(1)$ 理论开始。$\CM$ 的环面纤维化结构可描述如下。一个耦合了电荷为 $1$ 的超多重态的 $U(1)$ 规范理论，其库仑分支 $\CB$ 在原点 $Z_{0,1}=a=0$ 处有一个奇点，在那里该超多重态是无质量的。它在离原点很远处还有一种非物理的行为，原因在于这个理论不是渐近自由的。在 $\CB$ 中离原点足够近处，理论是有意义的。作维数约化后，四维规范场给出两个圆周值标量：电 Wilson 线 $\tilde\theta_e$ 与磁 Wilson 线 $\tilde\theta_m$。磁圆周在库仑分支的原点处收缩成一点。

$\CM$ 上的度量被「周期 Taub-NUT」空间（又称「Ooguri–Vafa」空间）\cite{Ooguri:1996me,Seiberg:1996ns} 的三维单圈修正抹平。作了这些修正之后可以看出，磁圆周只在余维数为 3 的地方收缩——即在四维库仑分支的原点（$a=0$）处，以及在电圆周上的某个特定点（$\tilde\theta_e=0$）处。

我们提到过的关于不同模空间 $\CM_\CL$ 的那些微妙之处，在这个例子里并不存在：本例中没有值得考虑的对 $\CL$ 的选取。

现在来看 $\CM$ 上的函数代数。在一般的复结构下，该流形可以等同于复流形 $V_+V_-=1-U$，其中 $U=\CY_{0,1}=\exp\left(\frac{\pi R}{\zeta}a+i\tilde\theta_e+\pi R\zeta\bar a\right)$ 取值于 $\IC^*$，而 $V_\pm$ 是更复杂一些的全纯函数，取值于 $\IC$。它们在 $\zeta$ 平面上的渐近性质是：在 $\zeta$ 平面中满足 $\Im(Z_{0,1}/\zeta)>0$ 的扇区里 $V_+=\CY_{1,0}$，而在相对的扇区里 $V_-=\CY_{-1,0}$。我们立刻认出：$V_+^pU^q$ 与 $V_-^pU^q$ 按 $\CY_\gamma$ 的展开、以及它们的环关系，与第 4.2 节% TODO \ref{subsec:Formal-line-ops-U1}
讨论过的形式线算子 $F_{p,q}$ 与 $F_{-p,q}$ 所算出的结果相符——只要取 $y=-1$。

考察一个带电荷为 $Q$ 的超多重态的 $U(1)$ 规范理论所对应的几何，也很有启发性。此时流形的方程取 $V_+V_-=(1+(-U)^Q)^Q$ 的形式，并有一组 $Q$ 个 $A_{Q-1}$ 奇点，位于 $U^Q=(-1)^{Q-1}$：磁圆周在电圆周上的 $Q$ 个位置处收缩，并且它比以前小了 $Q$ 倍，于是在它收缩的地方形成 $A_{Q-1}$ 奇点。注意，这与环关系 (4.20)% TODO \eqref{eq:Q-U1-algebra}
的 $y=-1$ 特殊化漂亮地吻合。（为看出这一点，注意到在 $y=-1$ 时有 $[n]=(-1)^{n-1}n$、$[n]! =(-1)^{\half n(n-1)}n!$，以及 $\genfrac{[}{]}{0pt}{}{n}{j}_{y=-1}=(-1)^{j(n+1)}\binom{n}{j}$。）

\subsection{$N=3$ Argyres--Douglas 理论}\label{subsec:N3-AD-Manfld}

现在转向第 4.3 节% TODO \ref{sec:formaln3}
讨论过的 $N=3$ 的 AD 理论。我们先描述其纤维化结构。

当紧化圆周的半径 $R$ 很大时，这个理论的三维库仑分支有一个直观的结构：库仑分支中两个奇点各自的邻域，都像是 $U(1)$ 理论耦合单个电荷为 $1$ 的粒子时的库仑分支。两个奇点以完全相同的方式被抹平，当然在每个奇点处收缩的是不同的圆周。粗略地说，在一个奇点上方的电圆周的某一点处，磁圆周收缩；而在另一个奇点上方的磁圆周的某一点处，电圆周收缩。Wilson 线的平移对称被完全破坏。与 $U(1)$ 的例子一样，本例中没有值得考虑的对 $\CM_{\CL_i}$ 的选取。

由 \cite{Gaiotto:2009hg} 第 9 节所述的与 Hitchin 系统的对应可知，我们可以把 $\CM$ 描述为球面上一个 $SU(2)$ Hitchin 系统的模空间，该系统只有一个不规则奇点，其阶数高到足以在该奇点周围给出五个斯托克斯扇区。作为复流形，$\CM$ 由方程 $M=1$ 给出，其中单值矩阵 $M$ 分解为
\begin{equation}
M  = \begin{pmatrix} 1 & U_1 \\  0  & 1 \\ \end{pmatrix}\begin{pmatrix} 1 & 0 \\  - U_4  & 1 \\ \end{pmatrix}\begin{pmatrix} 1 & U_2 \\  0  & 1 \\ \end{pmatrix}\begin{pmatrix} 1 & 0 \\  - U_5  & 1 \\ \end{pmatrix}\begin{pmatrix} 1 & U_3 \\  0  & 1 \\ \end{pmatrix}\begin{pmatrix} 0 & 1 \\  - 1  & 0 \\ \end{pmatrix}.
\end{equation}
这里面的三角矩阵是斯托克斯矩阵，而最后一个矩阵是形式单值。$U_i$ 是模空间上的全纯函数。它们的编号是为日后方便而取的。我们还将把编号延拓，定义 $U_{i+5}=U_i$。方程 $M=1$ 稍作整理后\footnote{证明这一点的一种办法，是把三个因子移到方程 $M=1$ 的一侧。这样就得到关于矩阵元的四个方程，容易看出它们正是 $i=2,3,4,5$ 时的方程 \eqref{eq:N3-manifold}。要得到最后一个方程，把 $i=4$ 的方程乘以 $U_2$，并注意只要 $U_3$ 非零，就得到 $i=1$ 的方程。类似地，若把 $i=3$ 的方程乘以 $U_5$ 再除以 $U_4$，就得到 $i=1$ 的方程。但 $U_3$ 与 $U_4$ 不可能同时为零。}，取一个漂亮的形式：
\begin{equation}\label{eq:N3-manifold}
U_{i-1} U_{i+1} = 1+U_i.
\end{equation}

可以证明，这五个方程定义了 $\IC^5$ 中一个光滑的二维复子流形。如果把五个坐标与五边形的五个顶点相对应，那么相邻顶点上所附的坐标绝不会同时为零。有五个特殊的除子 $U_i=0$，它们两两相交，呈五边形构型，交于五个特殊点，即 $U_1=U_3=0$、$U_2=U_4=U_5=-1$，以及它们在指标循环置换下的像。

把方程 (4.29)% TODO \eqref{eq:Deformd-N3-rels}
与 \eqref{eq:N3-manifold} 相比较，可以看出：取 $y\to 1$ 时的极限 $F_i\to U_i$，或取 $y\to -1$ 时的极限 $F_i\to -U_i$，就能还原出交换函数代数。这里 $\CM$ 与 $\tCM$ 似乎是相同的。我们将在第 10.1 节% TODO \ref{subsec:ADN3-Laminations}
中，把 $-U_i$ 按 $\CY_\gamma$ 的展开与 $U_i$ 按 $\tCY_\gamma$ 的展开，同形式线算子逐项对应起来。

\subsection{$SU(2)$ 与 $SO(3)$ 的 $N_f=0$ 规范理论}\label{subsec:SU2-Functions}

现在转向第 4.4 节% TODO \ref{subsec:formal-SU2}
讨论过的 $SU(2)$ 理论。我们先描述 $u$ 平面上的纤维化结构。

平直空间中的 $SU(2)$ 与 $SO(3)$ 的 $N_f=0$ 规范理论，只差在允许的局部线算子集合的选取上。我们已经看到，实际上有三种可能，而第三种可能在把紫外 $\theta$ 角平移 $\pi$ 之后，可以等同于 $SO(3)$ 的 $N_f=0$ 理论。三个理论的三维库仑分支 $\CM_1$、$\CM_2$、$\CM_3$，只差在 Wilson 线参数的精确周期上。我们先描述 $SU(2)$ 情形的模空间 $\CM_1$。这个理论的四维库仑分支 $\CB$ 与 $N=3$ AD 理论的情形惊人地相似，只是在奇点处变为无质量的粒子，其电荷满足 $\langle\gamma_1,\gamma_2\rangle=2$。按惯例，可以取 $\gamma_1=(1,0)$ 与 $\gamma_2=(-1,2)$。所有 BPS 粒子的电电荷都是偶数，因而瞬子修正、从而 $\CM_1$，在 $\IZ_2$ 平移对称 $\tilde\theta_e\to\tilde\theta_e+\pi$ 下不变。由于 $\gamma_i$ 经 $SL(2,\IZ)$ 变换与一个纯电电荷相联系，大半径处的瞬子修正会在每个奇点周围给出标准的周期 Taub-NUT 几何。

对 $SO(3)$ 规范理论，我们则会说 $\gamma_1=(2,0)$ 与 $\gamma_2=(-2,1)$：相对于 $SU(2)$，磁电荷加倍而电电荷减半。于是 $\CM_2$ 有一个不同的 $\tilde\IZ_2$ 平移等距，作用在 $\tilde\theta_m$ 上。由于 $\gamma_1$ 可被 $2$ 整除，大半径处 $\CB$ 的第一个奇点附近的局部几何，预期会有两个 $A_1$ 奇点，它们被 $\tilde\IZ_2$ 置换。对 $\CM_3$，两个奇点的角色互换。

为了把各个 $\CM_i$ 联系起来，我们可以做两件事。一种可能是用它们的 $\IZ_2$ 平移对称作商，把它们的环面纤维归约到共同的形式——由格 $\Gamma$（而非 $\Gamma_{\CL_i}$）构造。所得的流形 $\CM$ 有两个 $A_1$ 奇点，对应于四维库仑分支中的两个奇点。第二种可能是把它们的环面纤维「加倍」，得到更大的环面，以 $\Gamma\otimes\IR$ 中 $\Gamma_{\CL_i}$ 的并为模型。所得的流形 $\CM_{mx}$ 是光滑的，并具有 $\IZ_2\times\tilde\IZ_2$ 平移等距。此时图 (6.12) 变成% TODO \eqref{eq:covering-maps}
\begin{equation}\label{eq:A1-covering-maps}
\begin{matrix}
    &  &  \CM_{mx} &  &   \\
    & \swarrow & \downarrow  &  \searrow & \\
    \CM_{1} &  & \CM_{2} &  & \CM_{3} \\
    & \searrow & \downarrow  &  \swarrow & \\
    &  &  \CM &  &   \\
\end{matrix}.
\end{equation}
于是 $\CM_{mx}$ 对 $\IZ_2$ 作商得到 $\CM_1$，对 $\tilde\IZ_2$ 作商得到 $\CM_2$，对对角 $\IZ_2$ 作商得到 $\CM_3$。

这些流形与预期的 Hitchin 系统解的模空间相比如何？我们马上就会看到，$\CM_1$ 是 $SU(2)$ Hitchin 系统的标准模空间，而 $\CM_{mx}$ 是它的一个自然扩张。这个例子提示：不规则奇点的情形也许比正则奇点的情形要稍微复杂一些——即使 $C$ 的亏格为零，也可能存在规范群与 $\CM_\CL$ 的选取。

该 Hitchin 系统在球面上有两个不规则奇点，每个只有一个斯托克斯扇区。条件 $M_1M_2=1$ 变成
\begin{equation} \label{eq:nf0-matcond}
W \begin{pmatrix} 1 & U_1 \\  0  & 1 \\ \end{pmatrix}\begin{pmatrix} 0 & 1 \\  - 1  & 0 \\ \end{pmatrix} W^{-1} \begin{pmatrix} 1 & U_2 \\  0  & 1 \\ \end{pmatrix}\begin{pmatrix} 0 & 1 \\  - 1  & 0 \\ \end{pmatrix}=1.
\end{equation}
这里 $W$ 是行列式为 $1$ 的矩阵。还剩下一个小小的「规范自由度」$W\to -W$。这个 Hitchin 系统模空间上真正的全纯函数必须在这 $\IZ_2$ 变换下不变，但眼下我们先忽略这一点。$W$ 与 $U_{1,2}$ 的矩阵元可以看成某个略微扩大了的模空间上的全纯函数，这个模空间将与 $\CM_{mx}$ 重合。

把 \eqref{eq:nf0-matcond} 稍作整理，可以看出它蕴涵 $U_1=U_2=:U$ 以及 $\Tr W=0$，即可以写成
\begin{equation}
W = \begin{pmatrix} V_1 & V_0 \\  -V_2  & -V_1 \\ \end{pmatrix}.
\end{equation}
此外 \eqref{eq:nf0-matcond} 还蕴涵 $V_0+UV_1+V_2=0$；又因 $\det W=1$，有 $V_0V_2=1+V_1^2$。二维空间 $\CM_{mx}$ 以 $U,V_0,V_1,V_2$ 为坐标，并受这两个方程约束。

我们可以把 $V_0,V_1,V_2$ 延拓成一族无穷多个有用的函数 $V_n$：
\begin{equation}
W \left(\begin{pmatrix} 1 & U \\  0  & 1 \\ \end{pmatrix}\begin{pmatrix} 0 & 1 \\  - 1  & 0 \\ \end{pmatrix} \right)^n=\begin{pmatrix} V_{n+1} & V_n \\  -V_{n+2}  & -V_{n+1} \\ \end{pmatrix}.
\end{equation}
这些定义的自洽性要求对所有的 $n$ 都有 $V_{n-1}+UV_n+V_{n+1}=0$。此外，由于 $\det W=1$，对所有的 $n$ 都有关系式 $V_{n-1}V_{n+1}=1+V_n^2$。

单值的迹有简单的形式 $\Tr M_1=-U$，它应当对应于基础表示 $SU(2)$ Wilson 圈的期望值。我们有两个自然的对称：一个 $\IZ_2$ 作用为 $V_n\to -V_n$，一个 $\tilde\IZ_2$ 作用为 $U\to -U$、$V_n\to(-1)^nV_n$。前者正是我们为得到 $SU(2)$ Hitchin 系统解的模空间所必须作的商；它作用自由，所以模空间是光滑的。另一方面，$\tilde\IZ_2$ 或 $\IZ_2\times\tilde\IZ_2$ 中的对角子群确实有不动点：或者 $U=0$ 且 $V_{2n}=0$，这意味着 $V_{2n+1}=\pm(-1)^ni$；或者 $U=0$ 且 $V_{2n+1}=0$，这意味着 $V_{2n}=\pm(-1)^ni$。因此，对 $\tilde\IZ_2$ 作商、或对 $\IZ_2\times\tilde\IZ_2$ 中的对角子群作商，得到的空间有两个 $A_1$ 奇点，它们被剩余的 $\IZ_2$ 对称互换。对两个对称都作商，得到的空间有两个互不相关的 $A_1$ 奇点。这与规范理论分析所预期的完全一致：不作任何商的模空间似乎与 $\CM_{mx}$ 重合，$\IZ_2$ 商与 $\CM_1$ 重合，而 $\tilde\IZ_2$ 商与对角商则似乎分别与 $\CM_2$、$\CM_3$ 重合。

现在，把 $\CM_{mx}$ 与 $\CM_1$ 上的函数环同第 4.4 节% TODO \ref{subsec:formal-SU2}
关于形式线算子的结果作比较，是颇有意思的。为得到交换环，可以特殊化 $y=\pm1$。若取 $y=1$，就可以与 $\CM$ 上的函数 $V_n$ 比较。此时当 $y\to 1$ 有 $\hat V_n\to V_n$ 以及 $W\to -U$。把上面的关系式与方程 (4.50)% TODO \eqref{eq:SU2-formring-2}
和 (4.51)% TODO \eqref{eq:SU2-formring-3}
比较。当我们作商走到 $\CM_1$ 时，必须考虑由 $V_n$ 生成的偶函数。此时有 $W\to -U$、$G_{2n+1}\to V_n^2$ 以及 $G_{2n}\to V_{n-1}V_n$。在讨论 $\CM_1$ 上的函数时，也可以在由 $G$ 生成的代数中取 $y=-1$。此时 $G_{2n+1}\to -V_n^2$、$G_{2n+2}\to -V_nV_{n+1}$，而 $W\to U$。确实，有了这一对应，我们可以验证：方程 (4.43)、(4.44)、(4.52)、(4.53)% TODO \eqref{eq:rec-rels-1}、\eqref{eq:rec-rels-2}、\eqref{eq:SU2-formring-4}、\eqref{eq:SU2-formring-5}
在特殊化到 $y=-1$ 后，被相应的经典函数所满足。

\subsection{$SU(2)$ 与 $SO(3)$ 的 $\CN=2^*$ 理论}\label{subsec:N2-star-moduli}

最后，考虑第 5.4 节% TODO \ref{subsec:Formal-N2-star}
讨论过的 $\CN=2^*$ 理论。如前所述，$C$ 是带一个穿刺点的环面。$C$ 的基本群有生成元 $\alpha,\beta,m$，并有一个关系 $\alpha\beta\alpha^{-1}\beta^{-1}=m$。因此，绕穿刺点的单值取固定共轭类的平坦 $SL(2,\IC)$ 联络的模空间 $\CM_{mx}$ 可以这样描述：取满足
\begin{equation}\label{eq:monod-equa-nstar}
ABA^{-1} B^{-1} = M
\end{equation}
的 $SL(2,\IC)$ 矩阵 $A,B,M$，其中 $\Tr M=\mu+1/\mu$ 固定，并对整体的同时共轭作商。

从规范理论的角度看，$\CM_{mx}$ 上的全纯函数代数显然由和乐的迹生成，而这些迹又是 $A^{\pm1}$ 与 $B^{\pm1}$ 所成词的迹。我们断言：全纯函数代数实际上由 $a,b,c$ 生成，其中
\begin{equation}
a:= \Tr A, \qquad  b:= \Tr B, \qquad c:= \Tr AB,
\end{equation}
而 $\CM_{mx}$ 就是满足下式的 $(a,b,c)\in\IC^3$ 构成的空间——这个结果可以追溯到 Fricke：
\begin{equation}\label{eq:moduli-equation}
a^2 + b^2 + c^2 - abc = \mu + 2 + \frac{1}{\mu},
\end{equation}
关于这个模空间的有用信息可参见 \cite{MR2026539,MR1407481}。只要 $\mu\ne -1$，\eqref{eq:moduli-equation} 的解空间就是光滑的。

为论证这些断言，首先注意：由于矩阵是 $2\times2$ 的，只需考虑基础表示中和乐的迹。其次，回忆：若 $x$ 是行列式为 $1$ 的 $2\times2$ 矩阵，则
\begin{equation}\label{eq:CH-thm}
x^2 - \Tr(x) x +1 =0.
\end{equation}
对 $x=A^{\pm 1},B^{\pm 1}$ 使用它，显然可以把 $n_i,m_i\in\IZ$ 的表达式 $\Tr(A^{n_1}B^{m_1}A^{n_2}B^{m_2}\cdots)$ 归约为迹的多项式，其中 $n_i,m_i\in\{0,\pm 1\}$。然后，把 \eqref{eq:CH-thm} 写成 $x^{-1}=\Tr(x)-x$ 的形式来使用，就能用 $A$ 替换 $A^{-1}$、用 $B$ 替换 $B^{-1}$（多项式会随之改变）；最后，对 $x=AB$ 应用 \eqref{eq:CH-thm}，可以把 $\Tr[(AB)^n]$ 归约为 $\Tr(AB)$ 的多项式。此外，利用这些关系可以写出
\begin{equation}
\begin{split}
\Tr M & = \Tr(ABA^{-1} B^{-1} ) \\
& = (\Tr A)^2 - \Tr(ABAB^{-1}) \\
& = (\Tr A)^2 - \Tr(ABA) \Tr(B) + \Tr(ABAB)  \\
& = (\Tr A)^2 + (\Tr B)^2 + (\Tr AB)^2 - (\Tr A )(\Tr AB) (\Tr B) - 2
\end{split}
\end{equation}
由此得到 \eqref{eq:moduli-equation}。

为了重现预期的覆盖模空间
\begin{equation}\label{eq:A1-covering-maps-nstar}
\begin{matrix}
    &  &  \CM_{mx} &  &   \\
    & \swarrow & \downarrow  &  \searrow & \\
    \CM_{1} &  & \CM_{2} &  & \CM_{3} \\
    & \searrow & \downarrow  &  \swarrow & \\
    &  &  \CM &  &   \\
\end{matrix}
\end{equation}
我们注意到：由 \eqref{eq:monod-equa-nstar} 定义的 $\CM_{mx}$ 具有 $\IZ_2\times\IZ_2$ 对称群，由 $(A,B)\to(-A,B)$ 与 $(A,B)\to(A,-B)$ 生成，因而三个非平凡元素对 $(a,b,c)$ 的作用为
\begin{equation}
\begin{split}
g_1: (a,b,c) & \rightarrow (-a,b,-c) \\
g_2: (a,b,c) & \rightarrow (a,-b,-c) \\
g_3: (a,b,c) & \rightarrow (-a,-b,c)\\
\end{split}
\end{equation}
显然 $\CM_i=\CM_{mx}/\langle g_i\rangle$，而 $\CM=\CM_{mx}/(\IZ_2\times\IZ_2)$。注意 $\CM$ 是绕穿刺点单值取固定共轭类的平坦 $PSL(2,\IC)$ 联络的模空间。

相应地，$\IZ_2\times\IZ_2$ 的各元素在 $\CM_{mx}$ 上有六个不动点，例如 ${\rm Fix}(g_3)$ 是 $(0,0,\pm(\mu^{1/2}+\mu^{-1/2}))$。若对其中一个 $\IZ_2$ 子群作商，则有两个 $A_1$ 奇点，它们被剩余的 $\IZ_2$ 对称群互换。例如 $\CM_3$ 在 $[0,0,\pm(\mu^{1/2}+\mu^{-1/2})]$ 处有两个 $A_1$ 奇点，被 $g_1$（或 $g_2$）互换。若对 $\IZ_2\times\IZ_2$ 作商得到 $\CM$，则有 $3$ 个互不相同的 $A_1$ 奇点。在纤维化 $\CM\to\CB$ 下，它们投影为第 5.4 节% TODO \ref{subsec:Formal-N2-star}
提到的 $u$ 平面上的三个奇点。

第 5.4 节% TODO \ref{subsec:Formal-N2-star}
的形式 Wilson 算子与函数 $a,b,c$ 漂亮地对应。确切地说，$y\to 1$ 极限给出
\begin{equation}
\begin{split}
W_1 & \to a, \\
W_2 & \to b, \\
W_3 & \to c.
\end{split}
\end{equation}
明白了这一点，我们就认出方程 (5.28)% TODO \eqref{eq:quant-skein-rel}
正是对应于下式的量子 skein 关系：
\begin{equation}
{\rm Tr}(A) {\rm Tr}(B) = {\rm Tr}(AB) + {\rm Tr}(AB^{-1}).
\end{equation}

注意，这确实是与 Chern--Simons 理论中熟悉的那种形式相同的 skein 关系，因为 $A$、$B$ 两个圈在环面上相交，可以有两种方式解开，系数分别为 $y^{\pm 1/2}$。我们将在第 13 节% TODO \ref{sec:Open-Problems}
再次评论这一点。类似地，我们也能认出方程 (5.29)% TODO \eqref{eq:quant-mod-eq}
是 \eqref{eq:moduli-equation} 的量子类比。（可以验证，当 $y\to 1$ 时 $\sum a_{ijk}\to 1$。）

最后注意，$SL(2,\IZ)$ 的 S-对偶对称在这个模空间中表现得非常漂亮，因为可以把定义方程写成
\begin{equation}
A B = M B A,
\end{equation}
它也可以等价地写成
\begin{equation}
A (B A) = M (B A) A
\end{equation}
\begin{equation}
B A^{-1} = (A^{-1} M A) A^{-1} B.
\end{equation}
因此，我们可以定义 $SL(2,\IZ)$ 的生成元在和乐上的作用：$T: A \to A, B\to B A, M\to M$ 以及 $S: A \to B, B\to A^{-1}, M\to A^{-1} M A$。可以验证 $S^4$ 是恒等变换，但关系 $(ST)^2= S^2$ 在和乐上并不满足。尽管如此，
\begin{equation}
STS\cdot (A,B,M) = B^{-1}\left( T^{-1} S T^{-1} \cdot (A,B,M) \right) B
\end{equation}
因此它在和乐的迹上的作用确实通过模群分解。

% =====================================================================
%  GMN3 §9  $A_1$ 理论与叶层（$A_1$ theories and laminations）
%  原文：Lamination-Paper.tex 5051--5359 行；PDF p.79--85（印刷 p.78--84）
%  编号公式 (9.1)--(9.11) 共 11 条，全部在本块内，故 \eqref 可用；
%  图 6 幅（9.1--9.6），图名照抄原文 \insfig 第一参数：
%    dipole / path-passage / path-deform / path-deform-2 /
%    limit-laminations / limit-laminations-higher
%  跨块引用（式 (6.14)、§10、附录 F.1/F.2）暂写原文编号并加 % TODO，
%  合稿时统一改 \ref/\eqref
% =====================================================================

\section{$A_1$ 理论与叶层}
\label{sec:A1}

本节与第 10 节% TODO 合稿改 \ref{sec:Examples}
专门讨论 $A_1$ 理论。在这类理论中，可以给出分解 (6.14)% TODO 合稿改 \eqref{eq:y1}
的显式例子——该分解把线算子的 vev 用 $\tCY_\gamma$ 表示出来——从而算出带框 BPS 简并度。本节先给出一些必要的预备知识；具体结果见第 10 节。% TODO 合稿改 \ref{sec:Examples}

\subsection{技术回顾}

我们首先回顾文献 \cite{Gaiotto:2009hg} 中用到的若干技术，更多细节请参阅该文。

固定一点 $u \in \CB$。$u$ 对应 $C$ 上一个亚纯二次微分 $\phi_2$，而 $\phi_2$ 决定了 SW（Seiberg--Witten）曲线 $\Sigma$，其形式为 $\lambda^2 = \phi_2$。电荷格为 $\Gamma = H_1(\Sigma, \IZ)^{odd}$。$\phi_2$ 在 $C$ 上一个正则奇点（比如 $z = 0$）附近的行为由一个边界条件决定：
\begin{equation}
 \phi_2 \sim m^2 \frac{dz^{\otimes 2}}{z^2}.
\end{equation}
这里 $m$ 是对应于理论的一个相关形变的质量参数。在非正则奇点处，$\phi_2$ 有高阶极点。对一般的 $u$，$\phi_2$ 只有单零点；$\CB^{sing} \subset \CB$ 是 $\phi_2$ 的零点不全为单零点的轨迹。

在 $C$ 上任何既非 $\phi_2$ 的极点、也非其零点的点附近，我们可以用 $w = \int \sqrt{\phi_2}$ 定义一个局部坐标（差一个符号）。我们把 $C$ 上这样一条曲线称为「相位为 $\vartheta$ 的 WKB 曲线」：它映到 $w$ 平面后成为一条倾角为 $\vartheta$ 的直线。以 $t$ 参数化的一般 WKB 曲线，当 $t \to \pm\infty$ 时渐近于 $\phi_2$ 的极点。在满足 $m e^{i\vartheta} \notin i\IR$ 的正则奇点附近，WKB 曲线描写盘旋落入该奇点的指数螺线。在 $\phi_2$ 有 $k+2$ 阶极点\footnote{在文献 \cite{Gaiotto:2009hg} 中，非正则奇点的阶是用一个整数 $L$ 表示的；但那个记号会与我们表示线算子的记号冲突。}
的非正则奇点附近，WKB 曲线聚成 $k$ 个不同的族，每一族渐近于 $k$ 条「WKB 射线」之一。

在文献 \cite{Gaiotto:2009hg} 中，我们还引入了 $C$ 的一个三角剖分 $T_\WKB(u, \vartheta)$——更确切地说，是如下曲面的三角剖分：从 $C$ 出发，在每个非正则奇点周围挖去一个小圆盘，再在边界上标出一些点，把每个边界圆分成 $k$ 段弧。当不致引起混淆时，我们滥用记号，把这个带标记的曲面也称作 $C$。$T_\WKB(u, \vartheta)$ 的边是一般的 WKB 曲线。顶点是 WKB 曲线的渐近端点：每个正则奇点是一个顶点；每个非正则奇点给出一组 $k$ 个顶点，每段弧上各一个，此外还有这些顶点与斯托克斯射线之间的一个对应。$T_\WKB(u, \vartheta)$ 还被一些额外的离散数据「装饰」，如下所述。

正如我们已经提到的，在其复结构 $J^\zeta$ 下，$\CM$ 是带奇点的平坦联络的模空间。给定这样一个联络，再给定 $C$ 上一条 WKB 曲线，我们在每个渐近无穷远处都得到一条特别的「小平坦截面」（在相差一个重标度的意义下定义）：它就是当我们沿 WKB 曲线走进奇点时呈指数衰减的那条截面。在正则奇点处，这条小平坦截面不依赖于我们考虑的是哪一条 WKB 曲线；而且它是绕该奇点的单值变换的两个本征截面之一，相应的本征值为
\begin{equation}
 \mu = \exp \left[ \pm 2 \pi i \left( R \zeta^{-1} m - 2 m^{(3)} - R \zeta \bar{m} \right) \right].
\end{equation}
在 $\phi_2$ 有 $k+2$ 阶极点的非正则奇点处，情况稍复杂一些，原因在于斯托克斯现象：如我们所说，WKB 曲线聚成 $k$ 个族，每一族渐近于 $k$ 条 WKB 射线之一，而每条 WKB 射线携带一条不同的小平坦截面。所以无论哪种情形，$T_\WKB$ 的每个顶点都附有一条特别的平坦截面。

利用装饰三角剖分 $T_\WKB(u, \vartheta = \arg \zeta)$，我们在文献 \cite{Gaiotto:2009hg} 中定义了一族 $\CM$ 上的典范坐标函数 $\CX_\gamma$；其定义见附录 F.1。% TODO 合稿改 \ref{app:fgcoords}
我们还在文献 \cite{Gaiotto:2009hg} 中定义了 $\Gamma$ 上一个典范的二次加细 $\sigma$，于是可以再定义一族坐标函数 $\CY_\gamma = \sigma(\gamma) \CX_\gamma$。文献 \cite{Gaiotto:2009hg} 的论证表明：对 $\gamma \in \Gamma$，这些 $\CY_\gamma$ 确实与本文一直使用的那些 $\CY_\gamma$ 一致。另一方面，上面讨论带框 BPS 简并度时，我们还需要属于某个扩张格 $\Gamma_\CL$ 的电荷 $\gamma$ 所对应的 $\CY_\gamma$。对这些电荷，文献 \cite{Gaiotto:2009hg} 没有给出 $\CY_\gamma$ 的定义。由于 $2\gamma \in \Gamma$，$\CY_\gamma^2$ 是有定义的，所以困难在于确定一个平方根的符号。这里我们不给出确定这些符号的处方，而是改在扭空间 $\tCM_\CL$ 上工作，以绕开这个困难：在这个空间上符号困难不那么严重，我们确实可以对所有 $\gamma \in \Gamma_\CL$ 定义所需的函数 $\tCY_\gamma$。这件事在附录 F.2 中完成。% TODO 合稿改 \ref{app:traffic}

最后，在文献 \cite{Gaiotto:2009hg} 的附录 A 中，我们回顾了一个把 $\CM$ 上的和乐函数 $\Tr \Hol_\wp \CA$ 用坐标 $\CY_\gamma$ 展开的配方，但带有若干我们一般没有固定的符号歧义。在附录 F.2% TODO 合稿改 \ref{app:traffic}
中我们对这个配方稍作改进：给出 $\tCM$ 上的扭和乐函数用 $\tCY_\gamma$ 的展开，且不带任何符号歧义。我们将用这个展开来提取带框 BPS 简并度 $\fro(u, L_\zeta, \gamma; y = 1)$。

\subsection{缺失线算子之谜}
\label{sec:missing}

正如我们刚刚描述的，$\CS$ 类中的一个 $A_1$ 理论由如下选择刻画：一条黎曼面 $C$，其上带有若干「正则」或「非正则」穿刺点。迄今关于这类理论的大部分文献都集中在正则穿刺点上，但一些最自然的（也是组合上最简单的）例子其实需要非正则穿刺点：例如纯 $SU(2)$ 理论这样的渐近自由理论总是需要非正则穿刺点，文献 \cite{Gaiotto:2009hg} 所述 Argyres--Douglas CFT 的实现也是如此。

在这样的理论中如何构造线算子？我们可以试着沿用上面的做法、并照文献 \cite{Drukker:2009tz} 所述，由 $C$ 上的闭路径得到它们。然而稍加考察例子就会看出，这样并不能给出我们预期存在的全部线算子。例如，纯 $SU(2)$ 理论对应于带两个非正则穿刺点的 $C = \IC\IP^1$ \cite{Gaiotto:2009hg}；于是按同痕分类，$C$ 上只有一条非自交闭路径。这一条可以等同于 Wilson 线算子。可是 Wilson--'t Hooft 算子在哪里？

稍后我们就会看到，缺失的那些线算子对应于某些\emph{整叶层}（integral lamination），即允许以非正则穿刺点为端点的路径之组合。

\subsection{去耦味}
\label{subsec:Decouple-Flavor}

我们来考察当共形理论退化为渐近自由理论时，线算子会发生什么变化。

我们从一个用曲线 $C'$ 定义的 $A_1$ 理论 $T'$ 出发。设在 $C'$ 上某个局部坐标为 $z$ 的片内有两个正则奇点。我们设想这两个奇点靠得很近，比如位于 $z = \pm \eps$，质量参数为 $m_1, m_2$。于是有
\begin{equation}
\begin{split}
\phi_2 & = \left( \frac{m_1^2}{(z-\eps)^2} + \frac{m_2^2}{(z+\eps)^2} + \cdots \right) \, dz^{\otimes 2}\\
& = \frac{F(z)}{(z - \eps)^2 (z + \eps)^2} \, dz^{\otimes 2}.
\end{split}
\end{equation}

考虑把 $F(z)$ 在 $z = 0$ 附近展开成泰勒级数。通过调节 $+\eps$、$-\eps$ 两个奇点处的留数 $m_1$、$m_2$，我们可以任意固定这个级数的常数项与线性项。现在设 $\eps \to 0$ 而保持这两项有限。在这个极限下可以写
\begin{equation} \label{eq:limit-q}
\phi_2 = \left(\frac{\Lambda}{z^2} + \frac{m}{z} + \cdots\right)^2  \, dz^{\otimes 2}
= \left( \frac{\Lambda^2}{z^4} + \frac{2 \Lambda m}{z^3} + \cdots\right)  \, dz^{\otimes 2}
\end{equation}
其中 $\eps^2(m_1^2 + m_2^2) \to \Lambda^2$，且 $\frac{(m_1^2 - m_2^2)^2}{m_1^2 + m_2^2} \to m^2$。特别地，$m_1$ 与 $m_2$ 都按 $1/\eps$ 标度，而 $m_1 + m_2$ 与 $m_1 - m_2$ 中总有一个是 $1$ 阶的。我们还可以进一步取 $m \to \infty$、$\Lambda \to 0$ 而保持 $\Lambda m$ 固定的极限。

于是在这个标度极限下，原理论 $T'(m_1, m_2)$「退化」为一个理论 $T(m, \Lambda)$。新理论用一条带非正则奇点的曲线 $C$ 定义：事实上，\eqref{eq:limit-q} 正是「有两条斯托克斯射线从中伸出」的非正则奇点所对应的形式。

现在我们可以讨论极限理论 $T$ 中的线算子了：从理论 $T'$ 的线算子出发，观察它们在极限下变成什么。对一些线算子，答案相当平凡：任何不经过 $C'$ 上那两个相撞奇点附近的闭路径 $\wp'$，都简单地对应于 $C$ 上的一条闭路径 $\wp$，而和乐 $\Hol_{\wp'} \CA(\zeta)$ 极限到 $\Hol_\wp \CA(\zeta)$。因此，很自然地可以延拓文献 \cite{Drukker:2009tz}：猜想 $\wp$ 对应于理论 $T'$ 中的一个线算子 $L_\wp$，并写下
\begin{equation}
 L_{\wp'} \to L_\wp
\end{equation}
在极限下成立。

更有意思的情形是 $\wp'$ 确实从两个相撞奇点之间穿过，以致在 $\eps \to 0$ 的极限下被「夹断」。这时我们得多费些力气。首先，局部分析表明：在相撞奇点附近，当 $\eps$ 很小、$\vartheta$ 取一般值时，WKB 叶状结构看起来就像磁偶极子的场，如图 \ref{fig:dipole} 所示。
\insfig{dipole}{$C'$ 上一对相邻正则奇点的小邻域内的 WKB 叶状结构。}
特别地，我们看到从相撞奇点邻域中长出两束特别的叶。这两束叶渐近地趋于两条特别的射线。而且，取 $\eps \to 0$ 后对非正则奇点附近作 WKB 近似可以看出，这两条特别的射线恰好就是从该奇点伸出的 WKB 射线。（比较文献 \cite{Gaiotto:2009hg} 的图 42，那里画的非正则奇点有 $4$ 条 WKB 射线伸出。）类似地，沿这些 WKB 曲线的小平坦截面的空间在这个极限下光滑地相互关联。由此可以这样选择小平坦截面 $s_1$、$s_2$ 作为 $\eps$ 的函数的归一化：使它们极限到两条斯托克斯射线各自的那条小平坦截面。

现在，研究叶层 $\wp'$ 的一个方便技巧是「插入一组完备态」，把 $\wp'$ 切成若干段。给定一段路径 $q$ 以及沿 $q$ 的两条平坦截面 $s_1$、$s_2$，沿 $q$ 的一般平坦截面 $s$ 可以重写成两项之和：
\begin{equation}
 s = \frac{s_1 \wedge s}{s_1 \wedge s_2} s_2 + \frac{s \wedge s_2}{s_1 \wedge s_2} s_1.
\end{equation}
于是，沿 $q$ 从起点 $i$ 到终点 $f$ 的平行移动算子同样是两项之和：
\begin{equation}  \label{eq:two-terms}
 s(f) = \frac{s_1(i) \wedge s(i)}{s_1 \wedge s_2} s_2(f) + \frac{s(i) \wedge s_2(i)}{s_1 \wedge s_2} s_1(f).
\end{equation}

我们将把它用于如下情形：取 $q$ 为 $\wp'$ 的一小段，并取 $s_1$、$s_2$ 为两个奇点附近的小平坦截面，沿最显然的最短路线延拓到 $q$。见图 \ref{fig:path-passage}。
\insfig{path-passage}{$\wp'$ 的子路径 $q$。}

我们的目标是把答案改写成不涉及任何在奇点之间穿行的路径的形式。为此，我们分别处理 \eqref{eq:two-terms} 中的两项。先看第一项。我们按图 \ref{fig:path-deform} 所示对 $\wp'$ 作形变。
\insfig{path-deform}{对 $\wp'$ 的一种形变，使 $q$ 的两个端点能在 $\eps \to 0$ 极限下存活。}
经过这个形变，式中出现的量 $s_1(i)$ 与 $s_2(f)$ 在两个奇点相撞的极限下仍然有意义。另一方面，当我们试图计算 $s_1 \wedge s_2$ 时却遇到一个困难：$s_1(i)$ 与 $s_2(f)$ 有限这一事实在这里帮不上忙，因为我们必须在\emph{同一个}点上计算 $s_1$ 与 $s_2$。这就要把其中之一沿 $q$ 搬运过去，而 $q$ 正是在两个奇点之间穿行。为了绕开这个困难，我们可以把 $q$ 重新表示成一个绕 $z_1$ 的小回路与一条不会被夹断的新路径 $q'$ 的复合，如图 \ref{fig:path-deform-2} 所示。
\insfig{path-deform-2}{对 $\wp'$ 的进一步形变。}
沿 $q$ 计算 $s_1 \wedge s_2$，于是给出一个因子 $\mu_1^{-1}$（来自 $s_1$ 绕 $z_1$ 的单值变换——回忆我们把 $\mu_i$ 定义为 $s_i$ 绕 $z_i$ 的逆时针单值变换）乘上一个有限的部分。这个有限的部分仍是 $s_1 \wedge s_2$，只不过现在应理解为沿 $q'$ 而非 $q$ 搬运到公共点后的结果。对 \eqref{eq:two-terms} 的第二项可以作十分类似的分析。于是平行移动变成
\begin{equation}
 s(f) = \mu_1^{-1} \frac{s_1(i) \wedge s(i)}{s_1 \wedge s_2} s_2(f) + \mu_2 \frac{s(i) \wedge s_2(i)}{s_1 \wedge s_2} s_1(f),
\end{equation}
其中当 $\eps \to 0$ 时，除 $\mu_1$ 与 $\mu_2$ 外所有因子都有限且非零。

现在考虑 $\wp'$ 上平行移动在 $SL(2,\IC)$ 基础表示中的\emph{迹}。我们已经把沿 $q$ 的搬运写成了两个秩 1 算子之和。相应地，迹分解为两个不同对象的形式和：
\begin{equation} \label{eq:lp-decomp}
 L_{\wp'} = \mu_1^{-1} L_1 + \mu_2 L_2.
\end{equation}
我们把 $L_{1,2}$ 中的每一个表示成若干有向路径\emph{段}之并：这些段的端点落在绕非正则奇点 $\CP$ 的圆 $S^1(\CP)$ 上，并携带权重 $\pm 1$。见图 \ref{fig:limit-laminations}。
\insfig{limit-laminations}{极限对象 $L_{1,2}$。}
每条有向路径段的两端都落在从非正则奇点伸出的斯托克斯扇区上，于是在两端 $i$、$f$ 处各有特别的小平坦截面 $s_i$、$s_f$。$L_{1,2}$ 中任何一个的 vev 都是逐段之积：一条从 $i$ 到 $f$、携带权重 $k$ 的段贡献 $(s_i \wedge s_f)^k$。

在严格的 $\eps \to 0$ 极限下，对 $m_{1,2}$ 与 $\zeta$ 的相位的一般取值，$\mu_1$ 与 $\mu_2$ 都趋于 $0$ 或 $\infty$。\eqref{eq:lp-decomp} 中的两项会有一项压过另一项。因此可以说，在 $\eps \to 0$ 极限下 $L_{p'}$ 趋近于 $L_1$、$L_2$ 两个对象之一，再乘上一个无穷大的「重正化」因子。我们把 $L_1$ 与 $L_2$ 视为理论 $T'$ 的良好线算子。

于是 $\eps \to 0$ 极限的效果，是以一种特定的方式把路径 $\wp'$「切断」。更一般地，让多对奇点相撞，就可以把 $\wp'$ 切成许多不同的段。

上面考虑的是基础表示中的和乐迹。我们同样可以考虑更高的自旋表示，比如基础表示的 $k$ 次对称幂。那种情形下 \eqref{eq:lp-decomp} 的类比是一个分解
\begin{equation}
 L_{\wp', k} = \sum_{n=0}^k \mu_1^{-k-n} \mu_2^{n} L_1^{k-n} L_2^n.
\end{equation}
\insfig{limit-laminations-higher}{极限对象 $L_1^{k-n} L_2^n$。}
对象 $L_1^{k-n} L_2^n$ 画在图 \ref{fig:limit-laminations-higher} 中；正如记号所暗示的，它的 vev 为 $\inprod{L_1^{k-n} L_2^n} = \inprod{L_1}^{k-n} \inprod{L_2}^{n}$。在 $\eps \to 0$ 极限下，这个和同样由 $n = 0$ 或 $n = k$ 那一项主导。

我们上面的构造与文献 \cite{fgtoappear} 中给出的一个构造密切相关。

\subsection{叶层}
\label{subsec:Lams-Def}

上面我们发现：在带有「从中伸出两条斯托克斯线」的非正则奇点的 $A_1$ 理论中，良好线算子的类不仅包括 $C$ 上携带 $SL(2,\IC)$ 表示的闭路径，还包括终止于非正则奇点周围斯托克斯扇区、携带整数权重的开路径。我们的推导用到了某种取极限的程序，但很自然地会相信：即便非正则奇点不是由这种程序产生的，附着在开路径上的线算子仍然有意义——毕竟，在 $(2,0)$ 理论中定义表面算子的物理在 $C$ 上是局域的。

如果 $C$ 带有更一般的非正则奇点、斯托克斯线的条数任意，那会怎样？有一个自然的方案，它推广了我们上面针对两条斯托克斯线所得的结果，并且可以由多个正则奇点相撞推导出来。\footnote{这个推导的完整细节我们尚未完成。}
它与 Fock--Goncharov 在一个密切相关的语境中用过的构造非常相似 \cite{MR2349682}。沿用（但略有滥用）他们的术语，我们把所考虑的对象称为\emph{叶层}（lamination，是「整 $\CA_0$-叶层」的简称）。

回忆一下：当 $C$ 有非正则穿刺点时，我们挖去每个非正则穿刺点周围的小圆盘，并把每个边界圆分成若干段弧。$C$ 上的一个叶层 $L$ 是 $C'$ 上一族曲线 $\wp_i$ 之并：每条曲线或者闭合，或者在边界弧上有端点；它们不自交、也互不相交，并按同痕分类。每条闭曲线携带 $SL(2,\IC)$ 的一个非平凡不可约表示，但为下文的技术方便有一个例外：绕单个正则穿刺点的闭曲线携带一个整数权重（可正可负），而不携带表示。此外，每条开曲线 $\wp_i$ 携带一个整数权重 $k_i$，要求为正，但有一个例外：若一条开曲线可以收缩到边界上恰含一个标记点的弧段，则其权重允许为负。终止在边界每一段弧上的所有曲线的权重之和必须为零。还有如下等价关系：若一条曲线可缩，或可收缩到边界上不含标记点的弧段，则它可以被去掉；若叶层 $L$ 含有两条同痕的曲线 $\wp_i$、$\wp_j$，权重分别为 $k_i$、$k_j$，则它等价于一个新的叶层：去掉 $\wp_j$，而 $\wp_i$ 携带权重 $k_i + k_j$。

\emph{我们猜想：对任何 $A_1$ 理论，叶层与单线算子一一对应。}

附着在这样一个叶层上的线算子，其真空期望值对叶层的各分量是乘性的。附着在闭路径 $\wp$ 上的因子照例是沿 $\wp$ 的和乐之迹——但绕单个穿刺点、权重为 $k$ 的闭路径除外，对它们我们附着因子 $\mu^k$。对权重为 $k$ 的开路径，因子为
\begin{equation} \label{eq:open-vev}
 (s_{1} \wedge s_{2})^{k},
\end{equation}
其中 $s_1$、$s_2$ 分别是该路径起点与终点顶点处的小平坦截面。虽然附着在开路径上的因子依赖于小平坦截面 $s_i$ 的归一化选择，但当我们在叶层中对各分量取乘积时，这一归一化会抵消。这正是我们在上面的定义中对权重 $k_i$ 加上限制的原因。

上面我们没有仔细固定 vev 的整体符号。事实上，虽然有确定这个符号的处方，但下文中我们并不需要它。当我们确实要留意符号时，我们改在扭局部系统的空间上工作；在后一个空间上有一种自然的定符号方式，见附录 F.2。% TODO 合稿改 \ref{app:traffic}

% =====================================================================
%  GMN3 §10.1--10.3  $A_1$ 理论的例子（Examples of $A_1$ theories）
%  原文：Lamination-Paper.tex 5360--5654 行；PDF p.86--92（印刷 p.85--91）
%  编号公式 (10.1)--(10.34) 共 34 条：9 个 equation + 5 个 align（25 行），
%  行结构与原文逐行一致，故式号一一对应；插图 14 幅，全部在本块内。
%  跨块引用（节号、式号、附录）暂写原文编号并加 % TODO，合稿时统一改 \ref/\eqref
% =====================================================================

\section{$A_1$ 理论的例子}\label{sec:Examples}

本节我们在 $A_1$ 理论的几个简单例子中导出带框 BPS 简并度。

我们还会评述这些理论的保护自旋特征标方面的结果。我们将借助下面这条简单的评注来断定这些 PSC，并给它冠以一个响亮的名号：\emph{提升原理}（Promotion Principle）。想法非常简单。若已知某个自旋特征标在 $y=1$ 处的特殊化等于 $1$，那么相应的表示必定是 $SU(2)$ 的平凡（即自旋为零的）表示的单份拷贝，于是其特征标恒等于 $1$（对一切 $y$ 皆然）。注意，倘若我们所知的仅仅是 PSC 落在 $SU(2)$ 的表示环中，这条原理显然就会失效；因此，只有在强正性猜想成立时，我们才能援引提升原理。在下文第 11 节% TODO 合稿改 \ref{sec:Quantum-Holonomy}
中，我们会稍微更系统地利用这条原理。

\subsection{$N=3$ Argyres--Douglas 理论}\label{subsec:ADN3-Laminations}

我们从 $N=3$ Argyres--Douglas 理论开始；第 4.3 节% TODO 合稿改 \ref{sec:formaln3}
已对它作过简要回顾。这一情形下用 Hitchin 系统给出的 $\CM$ 之实现，见文献 \cite{Gaiotto:2009hg} 的第 9.4.4 节。此时 $C = \IC\IP^1$，只有一个不规则奇点，从它发出 $5$ 条斯托克斯线；相应地，边界 $S^1$ 上有 $5$ 个标记点。这些标记点把边界分成 $5$ 个扇区，我们按沿圆周顺时针递增的顺序把它们编号为 $1$ 到 $5$。
\insfig[0.45\textwidth]{pentagon-bare}{曲线 $C = \IC\IP^1$，其中挖去了无穷远附近的一个圆盘，边界上有五个标记点。}

本理论中有一组特别的 $5$ 个叶层。先从图 \ref{fig:pentagon-l1} 所示的 $L_1$ 开始。
\insfig[0.45\textwidth]{pentagon-l1}{$N=3$ 理论中的叶层 $L_1$。}
$L_{i+1}$ 由 $L_i$ 把圆周逆时针旋转角度 $\frac{4 \pi}{5}$ 得到。把这一操作重复 $5$ 次便回到原来的叶层，故 $L_{i+5} = L_i$。

这些叶层 $L_i$ 中的每一个都给出一族由 $\zeta$ 参数化的线算子 $L_{i,\zeta}$。为计算维数 $\fro(\gamma; u, \zeta, y = 1)$，根据式 (6.14)，% TODO 合稿改 \eqref{eq:y1}
我们应当把诸 vev $\inprod{L_{i,\zeta}(u)}'$ 按坐标函数 $\tCY_\gamma$ 展开。这可以借助我们在附录 F% TODO 合稿改 \ref{app:traffic}
中描述的「traffic rule」算法来完成。我们先画出 WKB 三角剖分 $T_\WKB(\zeta, u)$。五边形的各种三角剖分在顶点的循环置换之下看起来都一样，所以即使不知道 $\zeta$ 与 $u$ 是什么，我们也知道 $T_\WKB(\zeta, u)$ 长什么样。我们按图 \ref{fig:pentagon} 所示标记各顶点。
\insfig[0.45\textwidth]{pentagon}{$N=3$ 理论中的一个三角剖分。五边形的五条边都是边界边。边界上的每个橙色记号现在都等同于一条从该不规则奇点发出的斯托克斯射线。蓝色点是三角剖分的顶点；我们把它们等同于反斯托克斯射线。}

在附录 F% TODO 合稿改 \ref{app:traffic}
中，我们对每条边 $E$ 定义了模空间 $\tCM$ 上的函数 $\tCY_E(\zeta)$。为方便起见，我们减轻记号负担：把函数 $\tCY_X$ 写作 $X$，把 $\tCY_Y$ 写作 $Y$。于是应用 traffic rule，我们很快得到诸 vev 按 $X$ 与 $Y$ 的展开：
\begin{align}
\inprod{L_1}' &= X, \label{eq:pentagon-vevs-1} \\
\inprod{L_2}' &= Y + XY, \\
\inprod{L_3}' &= \frac{1}{X} + \frac{Y}{X} + Y, \\
\inprod{L_4}' &= \frac{1}{X} + \frac{1}{XY}, \\
\inprod{L_5}' &= \frac{1}{Y}. \label{eq:pentagon-vevs-5}
\end{align}
由这些 vev，我们立刻读出带框 BPS 态的谱。例如，线算子 $L_1(\zeta)$ 在真空 $u$ 中只支持单个带框 BPS 态。这个态因而必定自旋为零。另一方面，$L_3(\zeta)$ 支持三个带框 BPS 态，三者电荷各不相同。同样，这些态的自旋都为零。

本理论中一般的叶层形如
\begin{equation}\label{eq:ADN3-gen-lam}
L_i^m L_{i+1}^n, \qquad m, n \in \IZ_{\ge 0}.
\end{equation}
（两个不相邻的 $L_i$ 之积不给出叶层，因为叶层的边不允许彼此相交。）相应的期望值为
\begin{equation}
\inprod{L_i}^m \inprod{L_{i+1}}^n, \qquad m, n \in \IZ_{\ge 0}.
\end{equation}
展开这样的期望值，一般会得到更有趣的 BPS 态谱，其中某些电荷扇区的重数大于 $1$。

计算这些算子的交换环很有意思。为此，只需给出五个生成元 $L_i$ 的环关系。检视式 \eqref{eq:pentagon-vevs-1}--\eqref{eq:pentagon-vevs-5} 可见 $L_{i+1} L_{i-1} = 1 + L_i$。（这也可以更直接地理解为「skein 关系」的推论：skein 关系把两个相交叶层之积，与一个把交点解消掉的叶层之和联系起来。）

注意，若作识别 $Y \to X_{-\gamma_1}$ 与 $X \to X_{\gamma_2}$，则表达式 \eqref{eq:pentagon-vevs-1}--\eqref{eq:pentagon-vevs-5} 与第 4.3 节% TODO 合稿改 \ref{sec:formaln3}
中腔 $c_1$ 里求值的形式线算子 $F_i$ 的相应表达式完全吻合。此外，式 \eqref{eq:ADN3-gen-lam} 恰好对应第 4.3 节% TODO 合稿改 \ref{sec:formaln3}
所定义的简单形式线算子 $F_{p,q}^{(i)}$。

\subsection{$N = 4$ Argyres--Douglas 理论}

现在考虑 $N=4$ Argyres--Douglas 理论。它用 Hitchin 系统的实现与我们刚讨论的 $N=3$ 情形非常相似：唯一的差别是边界 $S^1$ 上有 $6$ 个标记点而不是 $5$ 个。见图 \ref{fig:hexagon-bare}。
\insfig[0.45\textwidth]{hexagon-bare}{曲线 $C = \IC\IP^1$，其中挖去了无穷远附近的一个圆盘，边界上有六个标记点。}

与上一情形的一个重要差别来自 $6$ 是偶数：存在非平凡的叶层，可以被推入边界的一个任意小邻域。一个例子见图 \ref{fig:hexagon-mu}。
\insfig[0.45\textwidth]{hexagon-mu}{一个非平凡的叶层 $\mu$，它可以被推入边界的一个任意小邻域。}
$\inprod{\mu}$ 与 $\CM$ 上的位置无关：不过，它可以通过理论的非可归一化形变而改变。

现在按图 \ref{fig:hexagon-l1} 定义两个叶层 $L_1$ 与 $L_2$。
\insfig[0.62\textwidth]{hexagon-l1}{叶层 $L_1$ 与 $L_2$。}
再把「由 $L_i$ 把圆周顺时针旋转角度 $\frac{2 \pi}{6}$ 所得的叶层」定义为 $L_{i+2}$。于是序列 $L_i$ 以 $12$ 为周期。此外，$L_{i+6}$ 可以利用下述关系表成 $L_i$ 的式子：
\begin{align}
L_7 &= \mu^{-1} L_1, \\
L_8 &= L_2,
\end{align}
以及它们在操作 $(L_i \to L_{i+2}, \mu \to \mu^{-1})$ 作用下的像。任何叶层都具有形式
\begin{equation}
\mu^r L_i^m L_{i+1}^n, \quad r \in \IZ, \quad m, n \in \IZ_{\ge 0}
\end{equation}
（对某个 $i$）。

\insfig[0.45\textwidth]{hexagon}{$N=4$ Argyres--Douglas 理论中一个三角剖分的例子。}

现在考虑图 \ref{fig:hexagon} 所示的三角剖分；对某组 $(u, \zeta)$，它可能出现为 $T_\WKB(u, \zeta)$。把这些叶层按此三角剖分所附的坐标展开——仍然使用 traffic rule——得到
\begin{align}
\inprod{\mu}' &= XZ, \\
\inprod{L_1}' &= Z, \\
\inprod{L_2}' &= Y + YZ, \\
\inprod{L_3}' &= \frac{1}{Z} + \frac{Y}{Z} + Y, \\
\inprod{L_4}' &= \frac{(1+Y)(1 + X + XY + XYZ)}{YZ}, \\
\inprod{L_5}' &= \frac{1}{Y} + \frac{X}{Y} + X, \\
\inprod{L_6}' &= \frac{1}{XY} + \frac{1}{Y}.
\end{align}
利用这些显式公式，我们可以轻易读出诸 $L_i$ 之间的关系：我们有
\begin{align}
L_1 L_3 &= 1 + L_2, \\
L_2 L_4 &= (1+ L_3)(1 + L_9),\\
L_1 L_4 &= 1+ \mu + L_5 + L_9,
\end{align}
以及它们在操作 $(L_i \to L_{i+2}, \mu \to \mu^{-1})$ 作用下的像。

\subsection{$SU(2)$、$N_f = 0$}\label{subsec:SU2-Lams}

接下来处理纯 $SU(2)$ 规范理论。这是取 $C = \IC\IP^1$、并带两个不规则奇点的 $A_1$ 理论，$\phi_2$ 在每个奇点处都有三阶极点。拓扑上，正如我们已经提到的，我们应当在每个奇点周围挖去一个小圆盘，并把每条边界分支分成单独一段弧，即在每条边界分支上只标记一个点——见图 \ref{fig:nf0-bare}。

\insfig[0.62\textwidth]{nf0-bare}{曲线 $C = \IC\IP^1$，其中挖去了两个圆盘，每条边界分支上各有一个标记点。}

叶层定义如下。先取两个小半圆 $\wp_1$、$\wp_2$，分别越过两个边界圆周上的标记点。再取两条不相交的曲线 $\wp_3$、$\wp_4$，从一条边界圆周走到另一条，二者彼此相差一个单位的缠绕。这四条曲线构成一个「极大」集，意思是：我们无法再加入任何一条既不相交、又不与其中任何一条同痕的曲线。给它们配上整数权重 $k_i$，并服从限制
\begin{equation} \label{eq:nf0-cond}
2 k_1 + k_3 + k_4 = 0, \quad 2 k_2 + k_3 + k_4 = 0, \quad k_3 \ge 0, \quad k_4 \ge 0.
\end{equation}
这样我们就得到一族叶层，由式 \eqref{eq:nf0-cond} 的解集参数化。对每一对满足 $k_3 + k_4$ 为偶数的 $k_3$、$k_4$，我们得到一个解。见图 \ref{fig:nf0-maximal}。
\insfig[0.62\textwidth]{nf0-maximal}{一族叶层。}
这只是无穷多族此类叶层中的一族，因为我们可以任意选择 $\wp_3$、$\wp_4$ 的整体缠绕。

所有这些叶层都给出超对称线算子。它们究竟是哪些算子？把文献 \cite{Drukker:2009tz} 给出的对应关系应用于 $N_f = 4$ 的 $SU(2)$ 理论，再沿第 9.3 节% TODO 合稿改 \ref{subsec:Decouple-Flavor}
的思路把质量送到无穷大，就会发现这些叶层对应 Wilson--'t Hooft 算子。它们是标号为 $[(\frac{p}{2} H_\alpha , \frac{q}{2} \alpha)]$ 的算子，其中磁荷 $p = k_3 + k_4$，电荷 $q$ 等于总缠绕，即 $k_3$ 乘以 $\wp_3$ 的缠绕加上 $k_4$ 乘以 $\wp_4$ 的缠绕。\footnote{对一条从一条边界走到另一条边界的曲线而言，「总缠绕」的含义并不唯一：真正典范地定义了的，只是两条此类曲线之间的相对缠绕。这里的这种不唯一性，反映的是规范理论中定义 Wilson--'t Hooft 算子之电荷时类似的困难。}在单值变换 $\zeta \to e^{2 \pi i} \zeta$ 之下，$\wp_3$ 与 $\wp_4$ 都增加 $2$ 个单位的缠绕，于是 $q$ 平移 $2p$ 个单位。

还有少数叶层我们尚未描述：即，我们可以取一条把两个奇点分开的单个闭曲线，它承载 $SL(2,\IC)$ 的自旋 $\frac{q}{2}$ 表示。这个叶层对应一个纯粹的 Wilson 圈算子，其电荷为 $q$，亦即处于最高权为 $\frac{q}{2} \alpha$ 的表示之中。

现在可以来考虑诸 vev 的红外展开了。照例，我们先画出 WKB 三角剖分 $T_\WKB(u, \zeta)$。此时组合结构非常简单，任何这样的三角剖分都必定形如图 \ref{fig:nf0}。（特别地，无论 $u$ 在弱耦合区还是强耦合区，都是如此！）有两条内部边，我们分别标为 $X$ 与 $Y$，二者相差一个单位的缠绕。
%
\insfig[0.62\textwidth]{nf0}{$SU(2)$、$N_f = 0$ 理论中的一个三角剖分。}
设 $L_0$ 与 $L_1$ 是图 \ref{fig:nf0-l0} 与图 \ref{fig:nf0-l1} 所示的叶层。更一般地，我们可以从 $L_n$ 出发、绕两个边界 $S^1$ 之一作 $k$ 次 Dehn twist，以此来定义 $L_{n+2k}$。这样就得到对一切 $n \in \IZ$ 的叶层 $L_n$。例如 $L_2$ 见图 \ref{fig:nf0-l2}。我们有
\begin{equation}
L_{2n+1}^2 = L_{2n} L_{2n+2}.
\end{equation}
（要验证这一点，只需注意：借助 Dehn twist 可以化归到 $n=0$ 的情形，而该情形可由图 \ref{fig:nf0-l0}--\ref{fig:nf0-l2} 看出。）
\insfig[0.62\textwidth]{nf0-l0}{$SU(2)$、$N_f = 0$ 理论中的叶层 $L_0$。}
\insfig[0.62\textwidth]{nf0-l1}{$SU(2)$、$N_f = 0$ 理论中的叶层 $L_1$。}
\insfig[0.62\textwidth]{nf0-l2}{$SU(2)$、$N_f = 0$ 理论中的叶层 $L_2$。}

最一般的叶层形如 $L_n^a L_{n+1}^b$，或者是一条承载 $SL(2,\IC)$ 某个表示的闭合回路。我们用 $L_*$ 记由单个闭合回路组成的叶层，该回路承载 $SL(2,\IC)$ 的基础表示。见图 \ref{fig:nf0-lstar}。
\insfig[0.62\textwidth]{nf0-lstar}{$SU(2)$、$N_f = 0$ 理论中的叶层 $L_*$。注意本图画在 $\IC\IP^1$ 上，因此，尽管我们把圆画成绕过左侧不规则奇点的样子，它其实与绕过右侧奇点的圆同痕。这里我们对记号稍有滥用：用 $+1$ 代表 $SL(2,\IC)$ 的基础表示。}

现在我们要计算与这些叶层相联系的线算子的期望值 $\inprod{L}'$，并把它们按 $\tCM$ 上的函数 $\tCY_X$ 与 $\tCY_Y$ 展开。照例，下文把这两个函数简记为 $X$ 与 $Y$。对任何具体叶层，结果都可以由 traffic rule 直接得到。然而事实上，答案还有一个漂亮的统一公式，它可以借更直接地利用式 (9.11)% TODO 合稿改 \eqref{eq:open-vev}
来证明。我们把该推导放在附录 G% TODO 合稿改 \ref{app:chebyshev}
中，这里只报告答案。

引入由 $U_n(\cos\theta) := \frac{\sin(n+1)\theta}{\sin\theta}$ 与 $T_n(\cos\theta) := \cos(n\theta)$ 定义的切比雪夫多项式（Tchebyshev）。令
\begin{equation}
 \alpha := \frac{1}{2 \sqrt{XY}} \left( XY + Y + 1 \right),
\end{equation}
则公式为
\begin{equation}\label{eq:chebyshev-form}
\inprod{L_{2n}}' = Y^{-1} \left[T_n(\alpha) + \frac{XY- Y - 1}{2 \sqrt{X Y}} U_{n-1}(\alpha) \right]^2.
\end{equation}
特别地，所有期望值 $\inprod{L_{2n}}'$ 都可以表为 $X$ 与 $Y$ 的、系数为正整数的洛朗多项式，正如我们所预期的。例如，
\begin{align}
\inprod{L_{-4}}' & = X^{-2} Y^{-3} (1+ 2 Y + Y^2 + X Y^2)^2 \label{eq:chebyshev-expls} \\
&= \frac{1}{Y^3 X^2}+\frac{4}{Y^2 X^2}+\frac{Y}{X^2}+\frac{6}{Y X^2}+\frac{2 Y}{X}+\frac{2}{Y X}+Y+\frac{4}{X^2}+\frac{4}{X}, \\
\inprod{L_{-2}}' &= \frac{1}{X Y^2} (1 + Y)^2=\frac{1}{Y^2 X}+\frac{2}{Y X}+\frac{1}{X}, \\
\inprod{L_{0}}' &= \frac{1}{Y}, \\
\inprod{L_{2}}' &= X, \\
\inprod{L_{4}}' &= Y(1+X)^2 = Y X^2+2 Y X+Y, \\
\inprod{L_{6}}' & = X^{-1} (1+ (1+X)^2 Y)^2 \\
&= Y^2 X^3+4 Y^2 X^2+6 Y^2 X+\frac{Y^2}{X}+4 Y^2+2 Y X+\frac{2 Y}{X}+4 Y+\frac{1}{X}, \label{eq:chebyshev-expll}
\end{align}
如此等等。我们还有
%
\begin{equation} \label{eq:wilson-vev-nf0}
\inprod{L_*}' = \frac{1}{\sqrt{XY}} + \sqrt{\frac{Y}{X}} + \sqrt{XY}.
\end{equation}
根据式 (6.14)，% TODO 合稿改 \eqref{eq:y1}
这些展开捕获了带框 BPS 简并度。例如，在式 \eqref{eq:wilson-vev-nf0} 中我们发现，插入了基础表示紫外 Wilson 圈的希尔伯特空间里有三个带框 BPS 态。其中 $(XY)^{\pm \half}$ 两项对应带电荷的态，电荷为 $\pm 1$，这与我们朴素预期的相符。更令人惊讶的是，还有一项额外贡献 $\sqrt{X/Y}$，它对应一个不带电荷的态，其磁荷由边 $X$ 与边 $Y$ 之缠绕的差给出。从弱耦合规范理论的计算中重现这一结果，将会很有意思。此外，展开式 \eqref{eq:chebyshev-expls}--\eqref{eq:chebyshev-expll} 蕴涵了一幅束缚在紫外 't Hooft--Wilson 圈上的带框 BPS 态的复杂图景；从弱耦合的角度发现它，应当也很有意思。

线算子之间的交换环关系也值得考察。例如，式 (G.4)% TODO 合稿改 \eqref{eq:Trig-Ln}
的一个简单推论是
%
\begin{equation} L_{2k} L_* = L_{2k-1} + L_{2k+1}. \end{equation}
%
这在物理上非常合理：若把一个 Wilson 圈与一个 't Hooft--Wilson 圈放到一起，后者把 $SU(2)$ 破缺为 $U(1)$，于是 Wilson 圈分解为两个电荷相反的 $U(1)$ 贡献之和。这样我们就得到两个 't Hooft--Wilson 圈，其电荷平移了一个单位。我们还有关系 $L_{2k} L_{2k+4} = (1+L_{2k+2})^2$、$L_{2k} L_{2k+6} = 1+(1+ L_{2k+2})(1 + L_{2k+4}) + L_*^2$、$L_{2k+1} L_{2k+3} = L_{2k+2} (1+L_{2k+2})$，等等。这些关系很难在规范理论中直接理解。

最后，与第 4.4 节% TODO 合稿改 \ref{subsec:formal-SU2}
中的形式线算子作比较很有意思。我们辨认出：在识别 $X=[0,2]$ 与 $Y=[-2,0]$ 之下，可以把 $L_{2n}$ 等同于 $G_{2n-1}$，从而把 $\sqrt{L_{2n}}$ 等同于在 $y=1$ 处求值的 $\hat V_{n-1}$。上述诸环关系就成为式 (4.50)--(4.54)% TODO 合稿改 \eqref{eq:SU2-formring-1}--\eqref{eq:SU2-formring-5}
在 $y=1$ 处的那些关系。这一比较与第 8.3 节% TODO 合稿改 \ref{subsec:SU2-Functions}
末尾的讨论相容，并且进一步表明：我们可以用形式线算子显式地算出这些理论中的保护自旋特征标。注意，一般说来，并不能简单地把式 \eqref{eq:chebyshev-expls} 中的整数 $n$ 提升为 $[n]$。

% =====================================================================
%  GMN3 §10.4--10.8  $A_1$ 理论的例子（续）
%  （$N_f=1$ / 其他 $N_f$ / $N_f=4$ / $\CN=2^*$ / millipede 展开）
%  原文：Lamination-Paper.tex 5655--5863 行；PDF p.93--98（印刷 p.92--97）
%  编号公式 (10.35)--(10.48) 共 14 条（equation 5 处 + align 6 行 + align 3 行）；
%  插图 9 幅（nf1-maximal / nf1 / nf1-l0 / nf2-l / nf3-l / nf4 / nf4-lstar /
%  n2star-triang / millipede）
%  注：BRIEF 分块表记本块 25 条公式，与原 PDF 不符；§10 全节 48 条的其余 34 条
%      落在 ch12（§10.1--10.3），本块末条 (10.48) 正好收满 §10，故以 14 条为准。
%  跨块引用（节号、式号）暂写原文编号并加 % TODO，合稿时统一改 \ref/\eqref
% =====================================================================

\subsection{$SU(2)$，$N_f = 1$}

下面我们处理带一重基础表示味的 $SU(2)$ 规范理论。正如文献 \cite{Gaiotto:2009hg} 第 10.2 节所解释的，这是一个 $A_1$ 理论，其中我们取 $C = \IC\IP^1$，带两个不规则奇点 $\CP_1$ 与 $\CP_2$，它们分别承载 $\phi_2$ 的 $3$ 阶极点与 $4$ 阶极点。拓扑上，我们在每个奇点周围挖去一个圆盘，并在 $S^1(\CP_1)$ 上标出一个点、在 $S^1(\CP_2)$ 上标出两个点。

该理论中的一个叶层样本见图 \ref{fig:nf1-maximal}。整数权重 $k_i$ 受如下限制：
\begin{equation} \label{eq:nf1-cond}
2 k_1 + k_4 + k_5 + k_6 = 0, \quad k_2 + k_3 + k_4 = 0, \quad k_2 + k_3 + k_5 + k_6 = 0, \quad k_4 \ge 0, \quad k_5 \ge 0, \quad k_6 \ge 0.
\end{equation}
条件 \eqref{eq:nf1-cond} 的解由 $k_4, k_5, k_6$ 决定：它们都 $\ge 0$，且 $k_4 + k_5 + k_6$ 为偶数。
\insfig{nf1-maximal}{$N_f=1$ 的 $SU(2)$ 理论中的一族叶层。}
于是我们得到了一个由这样的 $(k_4, k_5, k_6)$ 参数化的叶层族。

然而，这些并不是全部可能的叶层，因为我们所选的曲线并非最一般的可能选法。与 $N_f = 0$ 理论一样，我们可以通过一种 Dehn 扭转生成更多的叶层：把 $S^1(\CP_2)$ 的一个邻域连续旋转角度 $\pi$（于是两个标出的点互换）。这一扭转把路径 $\wp_4$、$\wp_5$、$\wp_6$ 拖成新的路径。对我们上面构造的叶层反复作顺时针或逆时针扭转，就能得到几乎所有可能的叶层。与前面一样，我们用 $L_*$ 记由单条闭曲线组成的叶层，它不能由上述构造得到。

现在如下选取一个三角剖分。边界圆 $S^1(\CP_2)$ 被分成两段。其中一段只与一条边相接，我们把那条边标为 $X$。另一段与两条边相接，我们把它们标为 $Y$ 与 $Z$。这三条边的另一端都落在边界圆 $S^1(\CP_1)$ 上。沿一条起止都在 $S^1(\CP_1)$ 上的弧顺时针走，我们会依次遇到 $Z$、$X$、$Y$。见图 \ref{fig:nf1}。
\insfig{nf1}{$N_f = 1$ 的 $SU(2)$ 理论中的一个三角剖分。}

设 $L_0$ 为图 \ref{fig:nf1-l0} 所示的叶层。
\insfig{nf1-l0}{$N_f = 1$ 的 $SU(2)$ 理论中的叶层 $L_0$。}
其余基本叶层 $L_n$ 由 $L_0$ 绕 $\CP_2$ 作上述扭转得到。计算它们的 vev 便给出，例如
\begin{align}
 \inprod{L_{-2}}' &= Y + XY + YZ + X Y Z^2 + 2 XYZ, \\
 \inprod{L_{-1}}' &= X + XZ, \\
 \inprod{L_{0}}' &= Z, \\
 \inprod{L_{1}}' &= 1/Y, \\
 \inprod{L_{2}}' &= 1/X + 1/XY, \\
 \inprod{L_3}' &= 1/Z + 1/XZ + 1/X Y^2 Z + 1/YZ + 2/XYZ.
\end{align}
一般地，$\inprod{L_k}'$ 与 $\inprod{L_{1-k}}'$ 由同时互换 $Z \leftrightarrow 1/Y$、$X \leftrightarrow 1/X$ 相联系。

我们还有
\begin{equation}
\inprod{L_*}' = \frac{1 + X + XY + XYZ}{\sqrt{XYZ}}.
\end{equation}

\subsection{$SU(2)$，其他 $N_f$}

对所有的 $SU(2)$、$N_f > 0$ 理论，都可以类似地描述其叶层。这里我们只局限于几条简短的评注。

$N_f=2$ 的理论有两种实现，见文献 \cite{Gaiotto:2009hg} 第 10.3 节与第 10.5 节。我们考虑它的「第一种实现」：其中 $C$ 为 $\IC\IP^1$，带两个不规则奇点，每个奇点承载两条斯托克斯射线。一个叶层样本见图 \ref{fig:nf2-l}。这个叶层在通常意义下是极大的，即不能再向它添加曲线。若改变附着在曲线上的权重，我们就能得到更一般的叶层族。
\insfig{nf2-l}{$N_f=2$ 的 $SU(2)$ 理论在其第一种实现下的一个叶层样本。}

$N_f = 3$ 的理论通过取 $C$ 为 $\IC\IP^1$ 得到：它带一个不规则奇点（承载两条斯托克斯射线）与两个正则奇点。一个叶层样本见图 \ref{fig:nf3-l}。
\insfig{nf3-l}{$N_f=3$ 的 $SU(2)$ 理论中的一个叶层样本。}

\subsection{$SU(2)$，$N_f = 4$}

接下来我们简要考虑 $N_f = 4$ 的 $SU(2)$ 理论。该理论对应 $C = \IC\IP^1$ 带 $4$ 个正则奇点的情形。对 $(\zeta, u)$ 的某些取法，$T_\WKB(\zeta, u)$ 就是图 \ref{fig:nf4} 所画的「四面体」三角剖分。
\insfig{nf4}{$\IC\IP^1$ 的「四面体」三角剖分，它出现在 $N_f=4$ 的 $SU(2)$ 理论中。顶点 $i$ 与顶点 $j$ 之间的边标为 $X_{ij}$。}
然后考虑图 \ref{fig:nf4-lstar} 所示的叶层 $L_*$。应用文献 \cite{Drukker:2009tz} 的规则可知，对应的算子 $L_*(\zeta)$ 是某一对偶标架下的 Wilson 圈。
\insfig{nf4-lstar}{$N_f=4$ 的 $SU(2)$ 理论中的叶层 $L_*$。}
利用 traffic rule，它的真空期望值为
%
\begin{equation}
\inprod{L_*} = {\Tr} \, L M_{23} R M_{13} L M_{14} R M_{24}
\end{equation}
%
计算出来是
%
\begin{equation}\label{eq:Nf4-Wilson}
\inprod{L_*} = \frac{1 + X_{13} + X_{24} + X_{13} X_{24} + X_{13} X_{14} X_{24} + X_{13} X_{23} X_{24} + X_{13} X_{14} X_{23} X_{24}}{\sqrt{X_{13} X_{14} X_{23} X_{24}}}.
\end{equation}
%
于是这个线算子支持 $7$ 个 BPS 态，它们的电荷各不相同。

\subsection{$SU(2)$，$\CN=2^*$}\label{subsec:SU2-N2star-Examples}

\insfig{n2star-triang}{$\CN=2^*$ 理论的通用 WKB 三角剖分。}

最后我们考虑最有趣的例子：$\CN=2^*$ 理论。正如我们在第 5.4 节% TODO \ref{subsec:Formal-N2-star}
所回顾的，该理论对应取 $C$ 为带一个正则奇点的单穿刺环面。相关的二次微分为 $\phi_2 = u + m^2 \wp(z|\tau)$。这里有两个转向点，且所有 WKB 三角剖分的拓扑都相同：两个三角形，把环面的某个基本区域一分为二。不失一般性，我们可以用三条直线 $(0,1)$、$(0,\tau)$、$(0,\tau+1)$ 表示这些边，如图 \ref{fig:n2star-triang} 所示。我们把相应的交比记作 $X$、$Y$、$Z$。于是绕穿刺点的单值为 $\mu = - XYZ$。绕环面的基本闭链 $z \to z+1$ 与 $z \to z+\tau$ 的和乐 ${\rm Tr}(A)$ 与 ${\rm Tr}(B)$ 很容易算出：这两条闭链各自穿过两条边，因而具有标准的三项展开。对闭链 $z \to z+1+\tau$ 也一样，它给出 ${\rm Tr}(AB)$。于是我们有展开
%
\begin{align}\label{eq:N2-star-traffic}
\inprod{L_{0,1}}' &= {\rm Tr}(A) = \sqrt{Z Y} + \sqrt{Z/Y} + 1/\sqrt{Z Y}, \\
\inprod{L_{1,0}}' &= {\rm Tr}(B) = \sqrt{X Z} + \sqrt{X/Z} + 1/\sqrt{X Z}, \\
\inprod{L_{1,1}}' &= {\rm Tr}(AB) = \sqrt{Y X} + \sqrt{Y/X} + 1/\sqrt{Y X}.
\end{align}
%
我们可以把这些与形式 Wilson 线算子 (5.24)% TODO \eqref{eq:star-formal-Wilson}
作比较：于是 $W_1$ 对应 ${\rm Tr}(A)$ 的（经典极限），依此类推，与第 8.4 节% TODO \ref{subsec:N2-star-moduli}
的情形完全相同。作为检验，由 \eqref{eq:N2-star-traffic} 可知：若把 $X, Y, Z$ 看成交换变量，则代入 \eqref{eq:N2-star-traffic} 确实给出
%
\begin{equation}
({\rm Tr}(A))^2 + ({\rm Tr}(B))^2 + ({\rm Tr}(AB))^2 - {\rm Tr}(A) {\rm Tr}(B) {\rm Tr}(AB) = \mu + 2 + \frac{1}{\mu}
\end{equation}
%
其中我们用了 $\mu = - XYZ$。这与方程 (8.9)% TODO \eqref{eq:moduli-equation}
相符，并且是 (5.34)% TODO \eqref{eq:quant-mod-eq}
的经典极限。

利用「提升原理」（promotion principle），由这些表达式可知：第 5.4 节% TODO \ref{subsec:Formal-N2-star}
构造的形式线算子确实计算 $\CN=2^*$ 理论的保护自旋特征标。

\subsection{millipede 展开}

在 $A_1$ 理论中，「原味」BPS 态有一个简单的描述：它们对应 $(2,0)$ 理论在 $C$ 上的、总质量有限的 BPS 弦。这归结为或者闭合、或者起止于二次微分 $\phi_2$ 之零点的 WKB 曲线。这一图像由文献 \cite{Klemm:1996bj,Shapere:1999xr} 引入，并在文献 \cite{Gaiotto:2009hg} 中起了重要作用。本文引入了带框 BPS 态，自然要问：它们是否也由类似的几何图像刻画。我们相信答案是肯定的：本节简要勾勒一个方案，把带框 BPS 简并度 $\fro(L, \gamma, \zeta)$ 直接等同于计数某些对象，我们称之为\emph{以 $L$ 为躯干、相位为 $\arg \zeta$ 的 millipede}（千足虫）。

我们这里不一般地定义这些对象，而只限于叶层 $L$ 就是一条闭圈、并承载 $SL(2,\IC)$ 之基础表示的情形。此时一个 millipede 是 $\Sigma$ 上的一条闭的有向曲线 $\xi$，它可以如下分成若干段。其中一些段称为\emph{躯干段}（body segment）：每条躯干段都是 $L$ 的相应一段到 $\Sigma$ 的提升。我们要求 $L$ 的每个点恰好提升为躯干段上的一个点：于是这些躯干段构成 $L$ 到 $\Sigma$ 的一种断裂提升。其余的段称为\emph{腿}（leg）。腿是沿之 $\lambda/\zeta$ 为实数、且取非负定向的有向曲线。（这样的段是 $C$ 上 WKB 曲线到 $\Sigma$ 的提升。）millipede 的\emph{荷}定义为 $\xi$ 在 $H_1(\Sigma, \IZ)$ 中的类。我们提出：带框 BPS 简并度应当「计数」以 $L$ 为躯干、相位为 $\arg \zeta$、荷为 $\gamma$ 的 millipede $\xi$。

我们来看构造 millipede 的最基本方式。从 $L$ 的一个断裂提升出发。在提升从 $\Sigma$ 的一叶跳到另一叶的地方，它带有边界点 $p_1$、$p_2$（它们投影到 $C$ 的同一个点 $p$）。于是 $\xi$ 的躯干不闭合。为修补这一点，我们在每个边界点 $p$ 处接上一条腿段：它从 $p_1$ 沿 $\Sigma$ 走到 $\lambda$ 的一个零点，再沿 $\Sigma$ 的另一叶返回到 $p_2$。（注意，只有当存在一条从 $p$ 通向某个转向点的 WKB 曲线时这才可能；而这类「特殊 WKB 曲线」只有有限多条，所以对给定的 $L$，接腿的方式只有有限多种。）把这一图像投影到 $C$，就解释了这一术语——见图 \ref{fig:millipede}。
\insfig{millipede}{一个 millipede 到 $C$ 的投影的一部分。躯干画成紫色，腿画成黑色。这个特定的 millipede 属于最简单的一类：所有腿都终止在转向点上，转向点画成橙色的叉。}

与「原味」BPS 态的情形一样，计数这些对象的精确规则有些微妙。最简单的情形就是我们刚讨论的：所有腿都从 $L$ 通向转向点。这样的 millipede 是孤立的（没有模）。如果只有这类 millipede，并且每个 millipede 在每个转向点上都只有一条腿终止，那么用 traffic rule 就比较容易看出：以 $L$ 为躯干、荷为 $\gamma$ 的 millipede 的数目，确实等于数目 $\fro(L, \gamma)$。

不过，更一般地，可能有若干条腿终止在同一个转向点上；在这些情形，该 millipede 有时贡献 $-1$ 而不是 $+1$。此外，还可能存在一些腿的起点与终点都在 $L$ 上的 millipede。这样的 millipede 有模（腿可以沿 $L$ 来回滑动）。我们猜想：孤立的 millipede 贡献自旋为零的带框 BPS 态，非孤立的 millipede 贡献更高自旋的带框 BPS 多重态；至少在一些简单例子中是这样，但我们尚未系统地研究这一问题。

% =====================================================================
%  GMN3 §11 量子和乐 + §12 热带标签（含 §12.1 $N=3$ Argyres--Douglas 例）
%  原文：Lamination-Paper.tex 5864--6289 行；PDF p.99--106（印刷 p.98--105）
%  编号公式 (11.1)--(11.3)、(12.1)--(12.12) 共 15 条；
%  插图 4 幅（cylinder / Flip-Traffic / tropic-onto-A / tropic-onto-B）
%  跨块引用（节号、式号、附录）暂写原文编号并加 % TODO，合稿时统一改 \ref/\eqref
% =====================================================================

\section{量子和乐}\label{sec:Quantum-Holonomy}

在第 10 节% TODO 合稿改 \ref{sec:Examples}
中，我们说明了如何计算 $A_1$ 理论的带框 BPS 简并度。本节我们勾勒如何——原则上——更进一步，算出 $A_1$ 理论完整的保护自旋特征标，从而算出这些例子中的形变函数代数。关键想法是把穿墙公式与第 10 节% TODO 合稿改 \ref{sec:Examples}
开头所解释的提升原理结合起来。

\insfig[0.55\textwidth]{cylinder}{在 $(u,\zeta)$（以及各耦合常数）的参数空间中，凡 $T_\WKB(u,\zeta)$ 形如本图之处，我们都可以用提升原理来计算与曲线 $\wp$ 相伴的线算子的 PSC。}

在一些特殊情形下，提升原理可以直接了当地应用（第 10 节% TODO 合稿改 \ref{sec:Examples}
中已经顺带这样做过）。设有一条曲线 $\wp$ 与一个三角剖分，使得 $\wp$ 的局部邻域形如图 \ref{fig:cylinder}。设我们想计算线算子 $L_\zeta(\wp)$ 的 PSC 在 $y=1$ 处的特殊化（这里表示取 $\CR=\mathbf{2}$ 是默认的）。也就是说，我们想计算期望值 $\langle L_\zeta(\wp)\rangle'$。应用附录 F % TODO 合稿改 \ref{app:traffic}
的 traffic rule，得到
\begin{equation}\label{eq:Prime-Wilson-Spec}
\langle L_\zeta(\wp)\rangle' = \sqrt{\tCY_+\tCY_-} + \frac{1}{\sqrt{\tCY_+\tCY_-}} + \sqrt{\frac{\tilde\CY_+}{\tCY_-}}.
\end{equation}
显然可以把提升原理应用于这个表达式，于是断定
\begin{equation}\label{eq:PromPrinc-1}
F(L_\zeta(\wp)) = X_{\half(\gamma_+ + \gamma_-)} + X_{-\half(\gamma_+ + \gamma_-)} + X_{\half(\gamma_+ - \gamma_-)}.
\end{equation}

现在设我们想对某个一般的 $(u,\zeta)$ 计算 $F(L_\zeta(\wp))$。数据 $(u,\zeta)$（以及各耦合常数）决定了一个 Seiberg--Witten 微分，以及与之相伴的 WKB 三角剖分 $T_\WKB(\vartheta,u)$，其中 $\vartheta=\arg\zeta$，详见文献 \cite{Gaiotto:2009hg} 第 6 节。通过变动 $u,\zeta$ 与理论的各耦合常数，我们试图找到从该 WKB 三角剖分到图 \ref{fig:cylinder} 那种形状的一个态射。若能找到，则对连接这两个三角剖分的一连串 flip 依次应用变换规则 (3.47)% TODO 合稿改 \eqref{eq:Spin-KS-tmn-S}
、(3.48)% TODO 合稿改 \eqref{eq:Spin-KS-ii}
，便可从 \eqref{eq:PromPrinc-1}
出发，得到 $(u,\zeta)$ 处的保护自旋特征标。

上述计算 PSC 的算法有两个重要的漏洞。第一，只有当 $C$ 容许图 \ref{fig:cylinder} 那种形状的三角剖分时，该算法才管用。现实中确实存在没有这种三角剖分的情形。所幸我们可以借助 J. Teschner 在 Teichm\"uller 空间量子化中所用的一个步骤来处理这种情形 \cite{Teschner:2005bz}。简述如下：取 $C$ 的一个 trinion（三孔球面）分解，使得 $\wp$ 是该分解中的一条切割闭链；各 trinion 的边界被标为「ingoing」与「outgoing」；存在一个与该 trinion 分解相容的三角剖分，用来在 Teichm\"uller 空间上定义 Fock 坐标（见文献 \cite{Teschner:2005bz} 第 13、14 节）。我们上面描述的算法适用于文献 \cite{Teschner:2005bz} 式 (15.4) 的情形；另一种情形出现在该文式 (15.5)。不过，正如该文所述，可以调用一个递推步骤来计算 $\langle L_\zeta(\wp)\rangle'$，而提升原理将适用于由此递推步骤所得的结果。

第二个漏洞是：即使存在图 \ref{fig:cylinder} 那种形状的三角剖分，也并不能一眼看出存在这种形状的 WKB 三角剖分。我们不认为这是一个严重的漏洞，因为总可以选取适当的弱耦合区域，使 $\wp$ 的管状邻域是一根细长的管子，而 $L_\zeta(\wp)$ 对应于规范群中相应 $SU(2)$ 因子的基本 Wilson 圈。在这样的参数空间区域里，WKB 三角剖分将会是合适的那一型，在我们看来相当合理；但本文不给出这一断言的详细论证。

我们还要再作几点评注：

\begin{enumerate}
\item 我们在第 4.4 节% TODO 合稿改 \ref{subsec:formal-SU2}
中已经看到，在 $SU(2)$、$N_f=0$ 的情形，一些强正的形式线算子可以纯粹代数地轻易算出；而第 5 节% TODO 合稿改 \ref{sec:Cluster-Algebras}
所述的与簇代数的联系表明，可以用局部规则来拼接模空间中的各个腔。于是，利用提升原理就能相当高效地算出某些 PSC。

\item 上述计算 PSC 的规则在某种意义上是 $C$ 上的非局部规则。我们希望能有计算带框保护自旋特征标的\emph{局部}规则，类似于我们在 $y=+1$ 处拥有的 traffic rule。我们预期这类局部规则将来自一个精细的 2d-4d 穿墙公式，其存在性曾在文献 \cite{Gaiotto:2009fs} 中被提示。这项工作正在进行中。\footnote{A. Goncharov 告诉我们，在簇系综（cluster ensembles）理论中，已有一些已知但尚未发表的规则可计算类似量 \cite{MR2567745}。}

\item 由与 Teschner 工作的这一联系可知：$A_1$ 理论中 $\CM$ 上的形变函数代数，自然地同构于 $C$ 的量子 Teichm\"uller 理论中量子测地算子之代数（后者量子化的是常曲率 $-1$ 度量下的测地长度）。我们要强调，我们的框架纳入了不规则奇点与叶层，因而为 Teichm\"uller 语境中已有的工作提供了有意思的推广。

\item 也许还可以更进一步，说明形式展开 (3.38) % TODO 合稿改 \eqref{eq:form-gen}
可以借助对 $\CM$ 的某个适当实切片作某种量子化而获得具体含义，即
\begin{equation}\label{eq:q-proposal}
\hat L_\wp = \sum_\gamma \fro(u,L,\zeta,\gamma;y) X_\gamma
\end{equation}
其中 $X_\gamma$ 现在是作用在某个确定的希尔伯特空间上的具体量子算子。事实上，文献 \cite{Teschner:2005bz} 以及 \cite{MR2567745,clqd2,qd-cluster} 都构造出了 $X_\gamma$ 的具体算子解释。不过，\eqref{eq:q-proposal} 应当有更广的适用范围，且不应依赖于实切片的选择。在选取某些实切片与特定量子化方案之后，我们或许可以期待与 Nekrasov--Shatashvili \cite{Nekrasov:2009ui,Nekrasov:2009rc} 以及 Nekrasov--Witten \cite{Nekrasov:2010ka} 的工作建立某种联系。
\end{enumerate}

\section{热带标签}\label{sec:Tropic}

本节我们勾勒一种用红外数据标记 $\CN=2$ 理论中单线算子的办法。

寻找这样的标记有三重动机。第一，这种标记是可取的，因为迄今我们只能在具有拉氏量的理论中给出标记（见第 2.3 节% TODO 合稿改 \ref{subsec:LabelsForLines}
）。红外标记同样可以用于非拉氏理论。第二，刻画 $\langle L\rangle$ 在大 $R$ 或大 $\zeta$ 时的渐近行为是有意义的。这有数学上的动机，因为我们知道和乐的迹是 Hitchin 模空间上有意思的函数；因此考察 $\zeta\to 0,\infty$ 的渐近行为很有意思，而结果表明它非常微妙。这里似乎还有一些有意思的物理学应用。第三，我们想补上第 6 节% TODO 合稿改 \ref{sec:linehol}
中的一个缺口：我们想证明存在足够多的不同线算子 $L_i$，足以「反解」展开 (6.3)% TODO 合稿改 \eqref{eq:trh-large-r}
，从而把 $\CY_\gamma$ 用诸 $L_i$ 表示出来。

结果表明，「热带代数簇」这一课题对处理这些问题很有用。我们将用到 Fock 与 Goncharov 的工作 \cite{MR2233852} 中迄今尚未在我们的故事里扮演过角色的某个方面。

关键想法是考察某个单线算子的 vev $\langle L\rangle$ 的首项渐近行为，例如在 $\zeta\to 0$ 时。达布展开 (6.3)% TODO 合稿改 \eqref{eq:trh-large-r}
把 $\langle L\rangle$ 表示成诸 $\CY_\gamma$ 之和。当 $\zeta$ 沿某条射线趋于 $0$，或固定 $(u,\zeta)$ 而令 $R\to\infty$ 时，和式中会有一项占主导。若考虑 $\widehat\CB\times\widehat\IC^*$ 中的一条路径 $(u_s,\zeta_s)$，则两项 $\fro(L_1,\gamma_1)\CY_{\gamma_1}$ 与 $\fro(L_2,\gamma_2)\CY_{\gamma_2}$ 只可能在路径上满足 ${\Re}(Z_{\gamma_1}/\zeta)={\Re}(Z_{\gamma_2}/\zeta)$ 的点处交换主导地位。对单线算子而言，从 halo 图像出发自然预期：这样两项必定满足 $\gamma_1-\gamma_2=\gamma$，其中 $\gamma$ 支撑一个「原味」BPS 态。因此，对 $\gamma\in\Gamma$，我们定义\emph{反 BPS 墙}（anti-BPS wall）$\check W(\gamma)$ 为这样一道墙：\footnote{在我们唯一的例子即 $A_1$ 理论中，这一定义将被证明是正确的。但我们注意到：尽管它是自然的，却并不能从 halo 图像严格推出「发生主导地位交换的墙只有这些墙」。原则上可能出现 $\gamma_1-\gamma_2$ 不支撑「原味」BPS 态的情形。若真有这类例子，下文关于热带理论的定义就得作修改。}
\begin{equation}
\check W(\gamma):= \{(u,\zeta): Z_\gamma/\zeta \in -i\IR_+ \qquad {\rm and} \qquad \Omega(\gamma;u)\not=0\}.
\end{equation}

我们之所以选取 $-i$ 这个符号，下文自会明了。这些墙的补集有若干连通分支，它们被定义为\emph{反腔}（anti-chamber），记作 $\check c$。一般地我们预期：在给定的反腔内，达布展开 (6.3)% TODO 合稿改 \eqref{eq:trh-large-r}
中会有单独一项主导渐近行为，于是可以写成
\begin{equation}
\langle L\rangle \sim \fro(L;\gamma_t)\exp(\pi R Z_{\gamma_t}/\zeta)\left(1+\cdots\right)
\end{equation}
其中 $\gamma_t\in\Gamma_\CL$ 是某个向量。我们把它称为\emph{热带标签}。一般地它依赖于 $L$ 与 $\check c$，故记作 $\gamma_t(L,\check c)$。注意我们有「热带公式」：
\begin{equation}\label{eq:tropic-mult-line}
\gamma_t(LL',\check c) = \gamma_t(L,\check c) + \gamma_t(L',\check c)
\end{equation}
以及
\begin{equation}
{\rm Re}(Z_{\gamma_t(L+L',\check c)}/\zeta) = {\rm Max}[{\rm Re}(Z_{\gamma_t(L,\check c)}/\zeta), {\rm Re}(Z_{\gamma_t(L',\check c)}/\zeta)].
\end{equation}
（关于 \eqref{eq:tropic-mult-line}：乘积 $LL'$ 一般不是单线算子，但可以按单线算子分解；所得诸算子中有一个会具有 \eqref{eq:tropic-mult-line} 给出的热带标签。）

更形式地说，我们把一个 $\CN=2$ 理论称为\emph{热带的}，若存在 $\Gamma$ 的一个（依赖于反腔的）基 $\{\gamma_a\}$，使得\footnote{我们也可以改用 $\langle L_\zeta\rangle'$ 与 $\widetilde\CY_\gamma$ 来表述，其优点是一切展开系数都为正，因而不可能发生相消。}

\begin{enumerate}
\item 在任一反腔 $\check c$ 中，对一切 $a$，当 $R\to\infty$ 与 $\zeta\to 0$ 时的渐近行为都满足 $\CY_a\to\infty$。此外，一切单线算子都具有形式：
\begin{equation}
\langle L\rangle = \fro(L,\gamma_t)\CY_{\gamma_t}P(1/\CY_a)
\end{equation}
其中 $P$ 是一个最低阶项为 $1$ 的多项式，而 $\gamma_t\in\Gamma_\CL$ 在反腔 $\check c$ 内为常数。（多项式 $P$ 可以在反腔内改变。例如，通常的腔会把该反腔分成不止一个连通分支。）

\item 跨过反 BPS 墙 $\check W(\gamma_0)$ 时，我们有变换律
\begin{equation}\label{eq:Tropical-Transformation}
\gamma_t(L_i,\check c') = \gamma_t(L_i,\check c) - \Omega(\gamma_0;c)\langle\gamma_t(L_i,\check c),\gamma_0\rangle_+ \gamma_0
\end{equation}
（该反墙位于腔 $c$ 内部）。

\item 对一切 $\gamma\in\Gamma_\CL$ 与任一反腔 $\check c$，都存在一个单线算子 $L_\gamma$，使得 $\gamma=\gamma_t(L_\gamma,\check c)$。

\item 若 $\gamma_t(L,\check c)=\gamma_t(L',\check c)$，则 $L=L'$。
\end{enumerate}

用物理的话说，这表示：一个单线算子通常被赋予一个紫外标签（例如 $\CL$ 中的一个向量），而它也可以由该理论（扩张的）荷格 $\Gamma_\CL$ 中的一个红外标签来唯一标记。不过，与紫外标签不同，红外标签依赖于红外参数（通过其对 $\check c$ 的依赖），并且会经历穿墙。

Fock 与 Goncharov 的工作 \cite{MR2233852} 提示如下猜想：$\CS$ 类中的所有理论都是热带的。下面我们勾勒为什么 $A_1$ 理论确实是热带的。（以下论证严重依赖文献 \cite{Gaiotto:2009hg} 所发展的技术，我们假定读者熟悉该文。）

为简单起见，本论证只关注 $C$ 上仅有正则奇点的情形。考虑一个单线算子。如文献 \cite{Drukker:2009tz} 与第 2.3.1 节% TODO 合稿改 \ref{subsubsec:A1-labels}
所解释，该算子的紫外标签是 $C$ 中一条非自交闭曲线 $\wp\subset C$ 的同痕类。$\langle L_\zeta(\wp)\rangle$ 的首项渐近行为将用 traffic rule 算法提取。我们选取不在反 BPS 墙上的 $(u,\zeta)$，并设 $\zeta=\vert\zeta\vert e^{i\vartheta}$，由此提取出 $\chi=\vartheta-\frac{\pi}{2}$ 对应的 WKB 三角剖分 $T_\WKB(\chi)$。（见文献 \cite{Gaiotto:2009hg} 第 6 节。）可以假定闭路径 $\wp$ 不会在 $T_\WKB(\chi)$ 的各个三角形中折返。现在考虑这一三角剖分的 traffic rule 算法。文献 \cite{Gaiotto:2009hg} 第 7.8 节证明了：对任一 WKB 三角剖分，相应的向量 $\gamma^\chi_E$（$E$ 取遍三角剖分的各边）构成一组单根基。特别地，对一切边 $E$ 都有 ${\rm Im}(e^{-i\chi}Z_{\gamma^\chi_E}(u))>0$，因此
\begin{equation}
{\rm Re}(e^{-i\vartheta}Z_{\gamma^\chi_E}(u)) > 0
\end{equation}
从而对沿半平面 $\mathbb{H}_\chi$ 趋于 $0$ 的一切 $\zeta'$，诸 $\CY_{\gamma_E^\chi}(u,\theta;\zeta')$ 的渐近行为都趋于无穷。这包括了位于其边界上的 $\zeta=e^{i\vartheta}$，所以我们仍有 $\CY_{\gamma_E^\chi}(u,\theta;\zeta)\to\infty$。由 traffic rule 现在可知，$\langle L_\zeta(\wp)\rangle$ 中的主导项简单地由
\begin{equation}
\prod_{\wp\cap E\not=\emptyset}\sqrt{\CY_{\gamma_E^\chi}}
\end{equation}
给出；于是热带荷就是 $\wp$ 所穿过的各边之荷的和：
\begin{equation}
\gamma_t(L(\wp),\check c) = \half\sum_{\wp\cap E\not=\emptyset}\gamma_E^\chi
\end{equation}

现在问：这个热带标签会怎样改变？这种改变只可能在 WKB 三角剖分发生一次 flip 时发生。\footnote{我们在此明确排除「juggle」变换，那将导致复杂得多的分析。}但文献 \cite{Gaiotto:2009hg} 的结果表明：$T_\WKB(\chi)$ 只在存在荷为 $\gamma$ 的超多重态、使得 $\chi={\rm arg}Z_\gamma(u)$ 时才改变。在当前情形 $\chi=\vartheta-\frac{\pi}{2}$，因此这意味着 $Z_\gamma/\zeta\in -i\IR_+$。也就是说，$L_\zeta(\wp)$ 的热带向量 $\gamma_t$ 只可能在反 BPS 墙上改变，而后者正是 $\check c$ 的边界。这就完成了热带理论定义中性质 1 的证明。

\insfig[0.5\textwidth]{Flip-Traffic}{利用 traffic rule 以及单根在一次 flip 下的变异，可以验证 $\gamma_t$ 跨过由一次 flip 相联系的两个腔时的热带变换规则。图中红色曲线是六个可能线算子的片段，它们不折返地穿过相应的四边形。}

现在来验证性质 2。考虑四边形中的一次 flip，以及它如何影响那些叶层 $L_1,\dots,L_6$——它们的线段与图 \ref{fig:Flip-Traffic} 中的叶层相交\zy{原文如此，疑为「与四边形相交」之笔误。}。利用文献 \cite{Gaiotto:2009hg} 的式 (7.23)--(7.24)，我们看到：当 $\chi$ 沿逆时针方向经过某个临界相位时，诸单根按一次变异而改变：
\begin{equation}
\begin{split}
\gamma_{E^+}^+ & = -\gamma_{E^-}^- \\
\gamma_{E_i^+}^+ & = \gamma_{E_i^-}^- + \langle\gamma_{E_i^-}^-,\gamma_{E^-}^-\rangle_+ \gamma_{E^-}^-
\end{split}
\end{equation}
由此容易检验热带变换规则 \eqref{eq:Tropical-Transformation}。

\insfig[0.5\textwidth]{tropic-onto-A}{乘上具有纯味荷的适当线算子后，上述线算子的热带标签为 $\gamma_E$。}

\insfig[0.5\textwidth]{tropic-onto-B}{乘上具有纯味荷的适当线算子后，上述线算子的热带标签为 $-\gamma_E$。}

为确立性质 3，只需给出热带标签分别等于 $\pm\gamma^\chi_E$ 的单线算子。于是，借助 \eqref{eq:tropic-mult-line}，这类线算子的乘积将给出任意想要的向量。现在参照图 \ref{fig:tropic-onto-A}，我们看到路径 $\wp$ 的热带标签为
\begin{equation}
\gamma_t(L(\wp)) = \left(\half\sum_{v_2\in E_i}\gamma_{E_i} + \half\gamma_E\right) + \left(\half\sum_{v_4\in E_i}\gamma_{E_i} + \half\gamma_E\right)
\end{equation}
其次注意：绕顶点 $v$、权重为 $-1$ 的小圆 $C_v$ 所构成的叶层满足 $L(C_v)=\mu_v^{-1}$，因而热带标签为 $-\half\sum_{v\in E_i}\gamma_{E_i}$。因此 $L(C_{v_1})L(C_{v_2})L(\wp)$ 的热带标签为 $\gamma_E$。类似地，利用图 \ref{fig:tropic-onto-B} 中的路径 $\wp'$，我们发现 $L(C_{v_1})L(C_{v_2})L(\wp')$ 的热带标签为 $-\gamma_E$。

最后，为确立性质 4，设给定一个热带向量 $\gamma_t$。由于诸 $\gamma_{E_i}$ 是单根，它具有到 $\gamma_{E_i}$ 之和的唯一分解，于是我们知道它穿过了哪些边。现在考虑 WKB 三角剖分的对偶胞腔分解。对热带向量中出现的每一条边 $E_i$（按重数计），我们在对偶胞腔分解中配上一条边。这些边以唯一的方式拼合成一个叶层，从而完成「$A_1$ 理论（仅有正则奇点）是热带的」这一证明。我们预期，一个类似但也许更繁复的论证将覆盖带不规则奇点的情形。

本节末尾我们作几点评注：

\begin{enumerate}
\item 变换 \eqref{eq:Tropical-Transformation} 是 Kontsevich--Soibelman 变换 $\CK_{\gamma_0}^{\Omega(\gamma_0)}$ 的「热带」版本。我们可以把它记作 $\CK_{t,\gamma_0}^{\Omega(\gamma_0)}$。存在一个热带穿墙公式，其中 $\CK$ 处处被换成 $\CK_t$。注意，上面我们只考虑了与 BPS 超多重态相联系的 flip；没有考虑与矢量多重态相联系的、更为困难的「juggle」变换（见文献 \cite{Gaiotto:2009hg} 第 6.6.3 节）。

\item 这里与斯托克斯理论有密切的类比。鉴于文献 \cite{Gaiotto:2008cd} 的 $\zeta$-微分方程，这并非偶然。

\item 在 $A_1$ 理论中，$L(\wp)$ 的热带向量 $\gamma_t$ 有一个非常漂亮的物理解释：它是 Seiberg--Witten 曲线 $\Sigma$ 中主导渐近行为的那条 WKB 路径的同调类。
\end{enumerate}

\subsection{例：$N=3$ Argyres--Douglas 理论的热带标签}

在式 (10.6)% TODO 合稿改 \eqref{eq:ADN3-gen-lam}
中，我们写下了 $N=3$ AD 理论中一般的单叶层（simple lamination）。这会把单叶层的集合分解成五个「象限」，由 $m,n\ge 0$ 给出坐标。事实上，把这五个象限粘成单独一份 $\IZ^2$ 是很自然的，如下所述。给定任一单线算子 $L$ 与任一带有坐标 $X$、$Y$ 的三角剖分，我们把 $\inprod{L}'$ 按 $X$ 与 $Y$ 展开。结果总具有形式
\begin{equation}
\inprod{L}' = X^a Y^b P(1/X,1/Y)
\end{equation}
其中 $P=1+\cdots$ 是某个多项式。数对 $(a,b)$ 就是算子 $L$ 的热带向量（相对于 $\Gamma$ 的基 $\gamma_X$、$\gamma_Y$）。

例如，用我们上面引入的三角剖分，五个象限中的算子 $L_i^m L_{i+1}^n$ 分别映射到 $(n+m,n)$、$(n,m+n)$、$(-n,m)$、$(-m,-n)$、$(n,-m)$；注意这五个区域合起来填满 $\IZ^2$，所以每种可能的红外荷都有相应的单线算子。这其实只是来自这样一个事实：五个典范的 $L_i$ 按循环次序排列时，其首项依次为 $X$、$XY$、$Y$、$1/X$、$1/Y$，而它们在单项式平面上也是按循环次序排列的。

对每个三角剖分，我们就这样得到所有单线算子之空间的一个自然坐标化，把它等同于 $\IZ^2$。对三角剖分作 flip 则给出这些坐标的一个简单的分段线性变换。这一结构在文献 \cite{qd-pentagon} 中有非常漂亮的描述。

% % ch15 = §13 + 致谢 + 附录 A–E；\appendix 写在该块文件开头，故此处不再另起
% =====================================================================
%  GMN3 §13 开放问题 + 致谢 + 附录 A--E
%  原文：Lamination-Paper.tex 6290--6895 行；PDF p.107--114（印刷 p.106--113）
%  编号公式 27 条：(A.1)--(A.7)、(B.1)--(B.2)、(C.1)--(C.3)、(D.1)--(D.8)、(E.1)--(E.7)
%  本块含 \appendix（附录字母 A--E 由此产生）；附录 F、G 属 ch16
%  跨块引用一律硬编码原文编号，行尾留 % TODO，合稿时改 \eqref/\ref
% =====================================================================

\section{开放问题}
\label{sec:Open-Problems}

围绕带框 BPS 态及其保护自旋特征标，无论在其物理方面还是在数学方面，都有若干有趣的研究方向。我们在此列出其中一部分。

\begin{enumerate}

\item （motivic）Kontsevich--Soibelman 穿墙公式，预期不仅适用于 $d=4$、$\CN=2$ 场论的 BPS 简并度，也适用于 Calabi--Yau 流形上 II 型弦紧化的 BPS 简并度。遗憾的是，在超引力情形下，这个公式的清晰物理论证仍然付之阙如；尽管颇具讽刺意味的是，半本原穿墙公式最早恰恰是在超引力背景下导出的 \cite{Denef:2007vg}。因此，若能在超引力背景下找到本文所用线算子的某种类比物，那将会非常有意思。

\item 带框 PSC $\fro(L,\gamma;y)$ 与带框 BPS 指标应当有一个更为直接的数学定义。看来很清楚：仿照 \cite{Sen:1994yi,Gauntlett:1993sh,Weinberg:2006rq,Hitchin-L2} 的思路，用（奇异）单极子模空间上某些 Dirac 算子的零模——这些单极子与适当的矢量丛耦合——应当能够给出这样一个严格定义。这样的定义可以为检验我们的规范理论预言提供一个数学框架。特别地，用它来检验我们的强正性猜想，将会很有意思。\footnote{我们感谢 Edward Witten 就此所作的启发性讨论。} 也很可能还存在带框 BPS 指标及其保护自旋特征标的别的定义，它们更接近 Calabi--Yau 流形上的 Donaldson--Thomas 理论。\footnote{我们感谢 E. Diaconescu 提出这一可能性。}

\item $\CS$ 类理论是否都具有第 5.2 节% TODO 合稿改 \ref{subsec:Seeds-from-N=2}
的簇性质？能否用它来确定这类理论的 BPS 谱？特别地，能否用它解决 $r>1$ 时求 $A_r$ 理论 BPS 谱这一困难问题？

\item 形式叶层与物理的线算子之间是否存在一一对应？所有 $\CN=2$ 理论是否都是第 12 节% TODO 合稿改 \ref{sec:Tropic}
意义下的热带理论？

\item 我们相信，对于 $\sigma(\CQ)$ 被插入式 (6.4)% TODO 合稿改 \eqref{eq:Lvev-trace}
一事，或许还可以沿下面这条思路给出另一种物理解释。我们希望对取值于 $U(1)$ 的函数 $\psi_{\gamma} := e^{i\tilde\theta_{\gamma}}$ 所起的作用有一个理解，尤其想弄明白：为什么它是一个\emph{扭}同态 $T_u \to U(1)$——这里的「扭」指
$\psi_{\gamma}\psi_{\gamma'} = (-1)^{\langle \gamma, \gamma'\rangle}\psi_{\gamma+\gamma'}$。假设 $\IR^3$ 中的某一个方向被当作时间方向。那么我们就能把一个电荷为 $\gamma$ 的双荷子沿圆 $S^1$ 绝热地输运一圈。这个粒子所获得的 Aharonov--Bohm 相位——我们把它看作是对电 Wilson 线与磁 Wilson 线的测量——将是 $\psi_{\gamma}$。现在把乘积 $\psi_{\gamma}\psi_{\gamma'}$ 与 $\psi_{\gamma+\gamma'}$ 作比较，主要差别在于：由于这一对双荷子的电磁场的自旋自由度，$\gamma$ 粒子与 $\gamma'$ 粒子这一对的朴素 $(-1)^F$ 宇称，与束缚态的同名量相差一个因子 $(-1)^{\langle \gamma,\gamma'\rangle}$。把这条评注想透，会遇到若干微妙之处，我们不打算把它们一一理清。我们认为，倘若能沿上述思路给出扭同态的物理推导，那将会非常好。

\item 正如我们已经提到的，弄清楚 $\CM$ 上函数的非交换形变能否与形如式 (6.9)% TODO 合稿改 \eqref{eq:y-dep-trace}
的扭迹联系起来，将会很有意思。与此相关，也正如我们上面已经提到的，用作用在显式希尔伯特空间上的具体算子来理解诸如式 (3.38)% TODO 合稿改 \eqref{eq:form-gen}
这样的非交换生成函数，将会非常有意思。这必定牵涉到对空间 $\CM_{\CL}$ 的各种实截面的研究。

\item 正如我们已经指出的，关系式 (5.28)% TODO 合稿改 \eqref{eq:quant-skein-rel}
是量子 skein 关系，与 Chern--Simons 理论中为人熟知的那些关系密切类似。此外，非紧 Chern--Simons 的三维流形不变量涉及量子双对数 \cite{Dimofte:2009yn}。这两件事实表明，我们的结果与 Chern--Simons 理论之间可能存在联系。

\item 在 \cite{Zamolodchikov:1991et} 中，Alyosha Zamolodchikov 对某些 $Y$-系统的解陈述了一个递推猜想。有关的参考文献与背景见 \cite{keller-periodicity} 以及报告 \cite{volkov-talk}。引人注目的是，其中若干相关的方程与 Argyres--Douglas 理论的环关系相吻合。我们想提出：Zamolodchikov 的猜想可以借助 Argyres--Douglas 理论的离散对称性来理解并加以推广。这一提议也由 Sergio Cecotti 与 Cumrun Vafa 提出，并在 \cite{CNV} 中有更详细的展开。

\item 正如我们在引言中所指出的，关于本文工作与簇代数的关系，还有许多事情有待阐明。这里只举两个例子。第一，\cite{Lampe} 中的一些公式与 $SU(2)$ 规范理论中的叶层和线算子密切相关。另一个方向上，在 \cite{MR2567745} 中，Fock 与 Goncharov 陈述了一个很有意思的对偶猜想，它把两类不同的簇代数簇（cluster variety）配对起来，他们称这两类为 $\CX$ 型与 $\CA$ 型。他们的猜想似乎与「用热带标签来标记线算子 vev 的渐近行为」有关。我们希望另找机会给出 Fock--Goncharov 对偶猜想的物理解释。

\end{enumerate}

\section*{致谢}

我们感谢 Dan Freed、Emanuel Diaconescu、Alexander Goncharov、Anton Kapustin、Juan Maldacena、David Morrison、Hiraku Nakajima、Nikita Nekrasov、Tony Pantev、Nathan Seiberg、Yan Soibelman、Valerio Toledano Laredo 与 Edward Witten 的讨论。

我们要感谢 Simons 几何与物理中心（Simons Center for Geometry and Physics）主办了若干与本项目相关的精彩研讨会，并在这些研讨会期间给予我们款待。

GM 的工作由 DOE 在基金号 DE-FG02-96ER40959 下资助。AN 的工作部分由 NSF 在基金号 PHY-0503584 与 PHY-0804450 下资助。DG 部分由 NSF 基金 PHY-0503584 资助。DG 部分由高等研究院（Institute for Advanced Study）的 Roger Dashen 会员资格资助。

\appendix

\section{关于 $d=4$、$\CN=2$ 超对称的一些技术细节}
\label{app:Susy-Conventions}

我们采用 Bagger 与 Wess 关于 $d=4$、$\CN=1$ 超对称的约定 \cite{Wess:1992cp}。具体而言，$SU(2)$ 指标用 $\epsilon^{12}= \epsilon_{21}=1$ 升降。$so(1,3)$ 的不可约自旋表示中张量的分量记作 $\alpha,\dot\alpha$，取值 $1,2$。取复共轭的规则是 $(\CO_1 \CO_2)^\dagger = \CO_2^\dagger \CO_1^\dagger$ 以及 $(\psi_\alpha)^\dagger = \bar \psi_{\dot \alpha}$。

$\CN=2$ 超对称算子是 $(Q_\alpha^{~A}, \bar Q_{\dot \alpha B})$，其中 $A,B$ 是 $SU(2)_R$ 指标，取值 $1$ 到 $2$。它们满足厄米性条件
%
\begin{equation}
(Q_\alpha^{~A})^\dagger = \bar Q_{\dot \alpha A}
\end{equation}
%
以及 $\CN=2$ 代数
%
\begin{equation}
\begin{split}
\{ Q_{\alpha}^{~A}, \bar Q_{\dot \beta B} \} & = 2
\sigma^m_{\alpha\dot\beta}P_m \delta^{A}_{~B} \\
\{ Q_{\alpha}^{~A}, Q_{  \beta}^{~ B} \} & = 2 \epsilon_{\alpha\beta}\epsilon^{AB} \bar Z \\
\{ \bar Q_{\dot\alpha A}, \bar Q_{  \dot \beta  B} \} & = -2
\epsilon_{\dot\alpha\dot\beta}\epsilon_{AB}   Z\\
\end{split}
\end{equation}
%
其中 $Z$ 是中心荷，$P_m$ 是厄米的能动量矢量，$P^0 \geq 0$。

把线算子插入空间坐标的原点 $x^i =0$ 处，它保持一个 $so(3) \oplus su(2)_R$ 对称性，以及下列超对称：
%
\begin{equation}
\CR_{\alpha}^{~A} = \xi^{-1} Q_{\alpha}^{~A} + \xi
\sigma^0_{\alpha\dot\beta} \bar Q^{\dot \beta A}
\end{equation}
%
这里 $\xi$ 是一个相位：$\vert \xi \vert =1$。这些算子满足厄米性条件
%
\begin{equation}
\begin{split}
(\CR_{1}^{~1})^\dagger & = - \CR_{2}^{~2} \\
(\CR_{1}^{~2})^\dagger & =   \CR_{2}^{~1} \\
\end{split}
\end{equation}
%
以及代数
%
\begin{equation}
\{ \CR_{\alpha}^{~A} , \CR_{\beta}^{~B} \} = 4 \left( E + \re(Z/\zeta) \right) \epsilon_{\alpha\beta} \epsilon^{AB}
\end{equation}
%
其中 $E=P^0$ 是能量算子，而 $\zeta = \xi^{-2}$。这个代数蕴涵
%
\begin{equation}
\left( \CR_{1}^{~1} + (\CR_{1}^{~1})^\dagger \right)^2 = \left(
\CR_{1}^{~2} + (\CR_{1}^{~2})^\dagger \right)^2 =  4 (E + \re(Z/\zeta) )
\end{equation}
%
由此即得 BPS 界
%
\begin{equation}
E + \re(Z/\zeta) \geq 0.
\end{equation}

\section{线算子 vev 的全纯性}
\label{app:Holomorphy}

本附录要详细解释：被 $\CR_{\alpha}^{~A}$ 湮灭的算子，其关联函数 $\inprod{L_\zeta }$ 如何可以被视为复结构 $\zeta$ 下 $\CM$ 上的全纯函数。我们还会就它与三维理论的 Rozansky--Witten 扭转之间的关系略作评论。

做这一辨认的关键一步，是把四维超荷与低能三维 sigma 模型的超荷对上，后者的靶空间是超凯勒的 $\CM$。在圆紧化破坏四维洛伦兹不变性之后，四维超荷可以收拢成一个二重态 $\left(Q_{\alpha}^{~A}, \sigma^0_{\alpha\dot\beta} \bar Q^{\dot \beta A}\right)$，我们可以把它记作 $Q^{~A}_{a~\alpha}$，其中新指标 $a$ 取值 $1,2$。四维旋量指标 $\alpha$ 可以等同于一个三维旋量指标，而 $Q^{~A}_{a~\alpha}$ 可以等同于三维 ${\cal N}=4$ 理论的超荷。注意：$A$ 指标受四维 $SU(2)_R$ R 对称性的作用，而后者仍是三维理论的对称性；但一般而言，四维理论中并没有任何对称性把 $a=1,2$ 的超荷相互旋转。$a$ 指标\emph{不是}任何 $SU(2)$ 对称性的指标。四维 $U(1)_R$ 对称性——如果未被质量或规范耦合标度破坏的话——当然会把 $a=1,2$ 两个分量朝相反方向旋转。

这些事实与靶空间为一般超凯勒流形的三维 sigma 模型的一般性质恰好吻合。我们稍后会详细说明：这类模型总有一个旋转费米子场的 $SU(2)_R$ R 对称性，但第二个 R 对称群必须取玻色靶空间的一个等距的形式，因为它要把三个超凯勒形式相互旋转。因此，只有当超凯勒靶流形具有特殊性质时，才可能有第二个 R 对称群。它可以是一个 $U(1)_R$ 等距，把三个超凯勒形式中的两个相互旋转；也可以是一个 $SU(2)_{R}'$ 等距群，把三个超凯勒形式全都旋转。

在 sigma 模型 $\CM$ 的靶空间中，黎曼结构群 $SO(4n)$ 约化为 $(SU(2)\times USp(2n))/\IZ_2 $，而切空间 $TX\otimes\IC \cong S \otimes V$，其中 $S$ 是平凡的二秩复矢量丛，$V$ 是秩 $2n$、结构群为 $USp(2n)$ 的复矢量丛。在 $S$ 中选取一条直线，就确定了 $\CM$ 上的一个复结构。这些复结构由 $\IC P^1$ 参数化，我们让 $\zeta $ 是这个扭量球面上的非齐次坐标。局部地，我们可以引入矢量场 $W_{a i}$，其中 $i=1,\dots, 2n$ 张成 $V$，$a=1,2$ 张成 $S$。费米子场 $\psi_{\alpha}^{A i}$ 是一个 $SU(2)_R$ 二重态的时空旋量，它们是 $V$ 的截面。作用在 $\CM$ 上函数 $F$ 的超对称变换的一般形式是
%
\begin{equation}
[Q^{~A}_{a~\alpha} , F] = \psi_{\alpha}^{A i} W_{a i} F
\end{equation}
%
特别地，超荷也是平凡丛 $S$ 的截面。

把四维矢量多重态约化到三维时，我们很容易还原出这一结构。四维费米子确实按 $SU(2)_R$ 的二重态变换。$n$ 个四维费米子 $\psi_\alpha^{A}$ 与 $n$ 个复共轭的 $\sigma^0_{\alpha\dot\beta} \bar \psi^{\dot \beta A}$ 合并起来，构成辛丛 $V$ 的 $2n$ 个分量。可以验证：$Q_\alpha^A$ 在复结构 $\zeta=0$ 下作用为 $\CM$ 上函数的全纯微分算子——它显然湮灭反全纯的四维标量，而且不那么显然地，它也湮灭电 Wilson 线与磁 Wilson 线的反全纯组合。类似地，$\bar Q^{\dot \beta A}$ 在复结构 $\zeta=0$ 下作用为反全纯微分算子。到此就完成了四维超荷的指标 $a$ 与三维 sigma 模型变换律中的指标 $a$ 的辨认。这一辨认不受量子修正的影响：量子修正使靶空间保持超凯勒，并且不改变 $\zeta=0$ 处的复结构。

特别地，选取一组四维超荷 $\CR_{\alpha}^{~A}$，等同于在 $S$ 中选取一条直线，也就等同于在 $\CM$ 中选取一个复结构。式
%
\begin{equation}
[\CR_{\alpha}^{~A} , F] = \psi_{\alpha}^{A i} \left( \xi^{-1} W_{1 i}+ \xi W_{2 i} \right) F
\end{equation}
%
中出现的那些矢量场，在复结构 $\zeta = \xi^2$ 下是反全纯的。因此 $\langle L(\zeta,\dots ) \rangle$ 是复结构 $\zeta$ 下 $\CM$ 上的全纯函数。

与三维 sigma 模型沿 Rozansky--Witten 路线 \cite{Rozansky:1996bq} 所作的拓扑扭转版本作一比较，也许会有意思。拓扑扭转把三维洛伦兹群换成旧的三维洛伦兹群与 $SU(2)_R$ 的对角组合。此时有 $\IC P^1$ 那样多的拓扑荷，它们与 $\CR_{\alpha}^{~A} \delta^{\alpha}_{~A}$ 重合。显然，我们这里考虑的四维线算子会在三维理论中给出点状的拓扑可观测量。这也许会被证明是一个有用的观点。

\section{不动点方程}
\label{app:Fixed-Point-Equations}

一个 $\CN=2$ 矢量多重态含有一个标量 $\varphi$、费米子 $\psi_{\alpha A}$（处于 $so(1,3)\oplus su(2)_R$ 的 $(\mathbf{2};\mathbf{1})\otimes 2$ 之中，以及它们的复共轭 $\bar \psi_{\dot \alpha A} := (\psi_{\alpha}^{~A})^\dagger$）、一个厄米规范场 $A_m$，以及一个满足实性条件 $(D_{AB})^* = - D^{AB}$ 的辅助场 $D_{AB} = D_{BA}$。乘以 $i$ 之后，所有这些场都取值于伴随表示。\footnote{我们的约定是：$u(N)$ 的生成元是 $N\times N$ 反厄米矩阵。在协变导数为 $d +A$ 的「几何」约定中，本文的规范场与它由 $A^{\rm geometric} = i g A^{\rm hermitian}$ 相联系。} 超对称变换律为\footnote{这里我们与 Bagger--Wess 约定略有偏离。我们的标量场 $\varphi=\sqrt{2}A^{\rm BW}$，其中 $A^{\rm BW}$ 是 BW 手征多重态的标量分量。}
%
\begin{equation}
\begin{split}
[Q_{\alpha A}, \varphi] & = - 2 \psi_{\alpha A} \\
[\bar Q_{\dot\alpha A}, \varphi] & = 0 \\
[Q_{\alpha A}, A_m] & = i \bar \psi_{\dot \beta A} (\bar
\sigma_m)^{\dot \beta}_{~~\alpha} \\
[\bar Q_{\dot \alpha A}, A_m] & = -i (\bar
\sigma_m)_{\dot \alpha}^{~~\beta}  \psi_{ \beta A}  \\
 [Q_{\alpha A}, \psi_{\beta B} ] & =
\sigma^{mn}_{\beta\alpha} F_{mn} \epsilon_{AB} + i D_{AB}
\epsilon_{\beta\alpha} + \frac{i}{2} g
\epsilon_{\beta\alpha} \epsilon_{AB} [\varphi^\dagger, \varphi] \\
[\bar Q_{\dot \alpha A}, \psi_{\beta B} ] & = - i   \epsilon_{AB} \sigma^m_{\beta\dot\alpha} D_m\varphi \\
[Q_{\alpha A}, D_{BC}  ] & = \left( \epsilon_{AB}
\sigma_{\alpha}^{m\dot \beta} D_m \bar\psi_{\dot \beta C} +
B\leftrightarrow C \right) \\
& +   g \left(\epsilon_{AB} [ \varphi^\dagger, \psi_{\alpha C}] +
B\leftrightarrow C   \right)\\
\end{split}
\end{equation}

超对称 $\CR_{\alpha}^{~A}$ 的不动点方程是
%
\begin{equation}\label{eq:Q-zeta-fix}
\begin{split}
F_{0\ell} - \frac{i}{2} \epsilon_{jk\ell} F_{jk} - i   D_\ell
( \varphi/\zeta) & = 0\\
D_0(  \varphi/\zeta) - \frac{g}{2} [  \varphi^\dagger,
  \varphi ] & = 0 \\
\end{split}
\end{equation}

若结构群 $\lieg$ 约化到嘉当子代数，我们就能写出一个静态解
%
\begin{equation}\label{eq:Embed-fix-sol}
\begin{split}
F & = \frac{1}{2} \omega_S \otimes \rho^M \\
\varphi & = \zeta \frac{\rho^M}{r} + \varphi_\infty \\
\end{split}
\end{equation}
%
其中 $\rho^M\in \liet$，而 $\omega_S = \sin\theta\, d\theta\, d\phi$，这里用的是围绕原点的标准角坐标。若 $\rho^M \in \Lambda_G$，则 $A$ 会被适当地量子化。这个解定义了紫外边界条件，从而定义了一个由 $\Gamma_G$ 中某个元素标记的 't Hooft 算子。

（这些方程里会出现一个有意思的微妙之处。若取 $\rho^M, \rho^E\in \liet$ 为常数，并令 $\varphi = \zeta(\rho^M-i \rho^E)/\sqrt{2} + \varphi_\infty$ 和 $F = \half \omega_S \otimes \rho^M + \half * \omega_S \otimes \rho^E$，那么不动点方程与 Bianchi 恒等式都得到满足，但运动方程\emph{并不}满足！）

\section{低能有效理论中的不动点方程}
\label{app:Low-Energy-FxdPoint}

不动点方程 \eqref{eq:Q-zeta-fix} 在红外理论中同样成立，但其解释略有不同。

让我们回忆：红外理论中的电荷格 $\hat \Gamma$ 带有一个反对称形式 $\langle \cdot, \cdot \rangle $，并构成一个正合列
%
\begin{equation}
0 \rightarrow \Gamma_f \rightarrow \Gamma \rightarrow
\Gamma_{\rm gauge} \rightarrow 0
\end{equation}
%
其中味荷格位于 $\langle \cdot, \cdot \rangle $ 的零化子之中，而商 $\Gamma_{\rm gauge} $ 是辛的。我们把投影映射记作 $\gamma \to \bar \gamma$。

低能的 Seiberg--Witten 有效红外理论是一个自对偶的阿贝尔规范理论，其不变场强 $\CF \in \Omega^2(\IR^{1,3}) \otimes V$，其中 $V= \Gamma_{\rm gauge} \otimes \IR$ 是一个辛矢量空间。模的取值确定了 $V$ 上的一个相容复结构 $\CI$。场强满足 $d \CF =0$ 以及反自对偶约束：
%
\begin{equation}
(*\otimes \CI)\CF = - \CF
\end{equation}
%
此外，在由 $\gamma_c$ 标记的希尔伯特空间扇区中，有一个量子化条件：
%
\begin{equation}
\int_{S^2_{\infty}} \frac{\CF}{2\pi}  =  \overline{\gamma_c}.
\end{equation}

在低能有效理论中，不动点方程 \eqref{eq:Q-zeta-fix} 的一个解可以通过取
%
\begin{equation}
\CF =\frac{1}{2} \left( \omega_S \otimes \gamma_c - \omega_H \otimes
\CI(\gamma_c) \right)
\end{equation}
%
得到，其中 $\omega_H = *_4 \omega_S = {dr \wedge \frac{dt}{r^2}}$。只要矢量多重态的模按方程
%
\begin{equation}\label{eq:Mod-ATT}
2 {\rm Im} \left[ \zeta^{-1} Z(\gamma; u(r))\right] = -
\frac{\langle \gamma, \gamma_c\rangle }{r} +2 {\rm Im} \left[
\zeta^{-1} Z(\gamma; u)\right] \qquad\qquad \forall \gamma \in \hat
\Gamma
\end{equation}
%
随 $r$ 变化，不动点方程就得到满足。此式是把方程 \eqref{eq:Q-zeta-fix} 取虚部得到的。那个方程是在一个固定的对偶标架中写出的。把方程写成对偶不变的形式，就得到 \eqref{eq:Mod-ATT}。

方程 \eqref{eq:Mod-ATT} 是标准吸引子方程的一个变形。通常的吸引子方程是在 $\CN=2$ 超引力中写出的。那个方程的场论极限给出形如 \eqref{eq:Mod-ATT} 的方程，但有一个重要的例外：$\zeta = e^{i \alpha}$ 是中心荷 $Z(\gamma_c;u)$ 的（常数）相位，并且 $\Gamma_f=0$。吸引子方程的这一推广并不完全是新的。「orientihole」的方程同样把吸引子方程中的相位 $e^{i \alpha}$ 换成 $-e^{i \alpha}$ \cite{Denef:2009ja}。

一个电荷为 $\gamma_h\in \hat \Gamma $ 的探针 BPS 粒子，在上述某个场位形中运动时，其动力学由作用量
%
\begin{equation}
\int \vert Z(\gamma_h;u(r))\vert ds + \int \langle
\overline{\gamma_h}, \CA\rangle
\end{equation}
%
支配，其中积分沿探针粒子的世界线进行。因此，这样一个粒子静止时的能量是
%
\begin{equation}
\begin{split}
E  & =  \vert Z(\gamma_h;u(r))\vert - \langle \overline{\gamma_h},
\CA_0 \rangle\\
& =\vert Z(\gamma_h;u(r))\vert  + \frac{ (
\overline{\gamma_h},\overline{\gamma_c} )}{r}
 \\
\end{split}
\end{equation}
%
第二项 $\frac{ ( \overline{\gamma_h},\overline{\gamma_c}
)}{r}$ 是库仑能，它用 $V$ 上由复结构构成的正定对称度量来表达：$(v_1, v_2) := \langle v_1, \CI v_2 \rangle$。注意这个表达式对电荷是对称的、且正定，这在物理上是合理的。

现在，把不动点方程 \eqref{eq:Q-zeta-fix} 取实部，并写出其对偶不变的延拓，我们得到
%
\begin{equation}
\re\left[ \zeta^{-1} Z(\gamma_h;u(r))\right] =\frac{ (
\overline{\gamma_h},\overline{\gamma_c} )}{r} +\re \left[
\zeta^{-1} Z(\gamma_h;u)\right]
\end{equation}
%
这样我们就导出了 halo 粒子在 $\zeta$-吸引子背景中能量的公式 (3.22)% TODO 合稿改 \eqref{eq:Halo-Particle-Energy}
。halo 半径 (3.23)% TODO 合稿改 \eqref{eq:halo-rad}
则由把 \eqref{eq:Mod-ATT} 的右端对 $\gamma=\gamma_h$ 置零而得到。

\section{六维超对称}
\label{app:Six-Super}

$\IR^{1,5}$ 中的 $(2,0)$ 超共形代数（此处记作 $osp(8^*\vert 4)$）有偶子代数 $osp(8^*\vert 4)^0 = so(6,2) \oplus so(5)$。为构成奇子空间，我们选取 $so(6,2)$ 的两个手征旋量表示 $\Delta_\pm$ 之一。$\Delta_\pm$ 是四元数空间，各自同构于 $\mathbb{H}^4$。为确定起见，取 $\Delta_+$。再令 $\Delta'\cong \mathbb{H}^2$ 为 $so(5)$ 的不可约旋量表示。奇子空间是 $osp(8^*\vert 4)^1 = \Delta_+ \otimes_{\IR} \Delta'$。它是 $32$ 实维的。

庞加莱超对称的子超代数有偶子代数 $so(5,1) \oplus so(5)$；此时出现 $so(5,1)$ 的一个手征旋量 $\Delta_+'\cong \mathbb{H}^2$，用它来定义 $16$ 个庞加莱超荷构成的奇子空间 $\Delta_+'\otimes \Delta'$。

引入 Clifford 代数 $Cl(1_+,5_-)$ 的一个表示会很有用，它由形如
%
\begin{equation}
\Gamma^M = \begin{pmatrix} 0 & \bar\gamma^{M~\dot s}_{r}   \\    \gamma^{M~ s}_{\dot r} & 0 \\ \end{pmatrix}
\end{equation}
%
的 $8\times 8$ 复值矩阵给出，其中 $r=1, \dots,4 $ 标记手征自旋表示，而 $\dot r = 1, \dots, 4$ 标记反手征表示。当然 $M=0,1,\dots 5$ 跑遍时空维度。于是 $(\Gamma^M)^*  = B \Gamma^M B^{-1} $，其中
%
\begin{equation}
B = \begin{pmatrix} B_{r}^{~~s} & 0 \\ 0 & B_{\dot r}^{~~ \dot s} \\ \end{pmatrix}
\end{equation}
%
手征表示 $\Delta_\pm'$ 是伪实的。此外还有
%
\begin{equation}
C = \begin{pmatrix} 0 & c^{r \dot s} \\ \bar c^{\dot r s} & 0 \\ \end{pmatrix}
\end{equation}
%
使得 $(C \Gamma^M)^{\rm tr} = - (C \Gamma^M)$ 且 $(\Gamma^M C^{-1})^{\rm tr} = - (\Gamma^M C^{-1})$。张量 $c^{r \dot  s}$ 给出 $\Delta_+'$ 与 $\Delta_-'$ 缩并为单态的不变缩并。

对于 R 对称性 $so(5)\cong usp(4)$，把旋量 $\Delta'$ 与 $USp(4)$ 的 $4$ 等同起来。若要引入指标，我们令 $i,j =1,\dots, 4$。Gamma 矩阵的指标结构是 $\Gamma^{Ii}_{~~j}$，其中 $I=1, \dots 5$。指标用辛矩阵 $J^{ij}$ 及其逆矩阵升降。于是 $(\Gamma^I)_{i}^{~j} = J_{ii'} J^{jj'} \Gamma^{Ii'}_{~~j'}$。我们有 $(\Gamma^{Ii}_{~~j})^* = - \Gamma^{I~j}_{i}$，并且 $\Gamma^I_{ij} := J_{ii'} \Gamma^{Ii'}_{~~j}$ 是反对称的。

若把庞加莱超对称记作 $Q_{r}^i$，则我们要求
%
\begin{equation}
(Q_{r i})^\dagger = B_r^{~~s} J^{ij} Q_{s j}
\end{equation}
%
于是有 $16$ 个厄米超荷。基本的超对称代数是
%
\begin{equation}
[Q_{r}^i, Q_{s}^j ] =  2i  (\bar\gamma^M\bar c)_{rs} J^{ij} P_M
\end{equation}
%
其中 $P_M$ 是厄米的平移算子。系数 $2i$ 由约化到上面描述的 $d=4$、$\CN=2$ 代数来确定。作用在场多重态上时，右端可能因规范变换而出现额外项。作用在孤子扇区的态上时，右端可能因「中心扩张」而出现额外项。事实上，在存在弦激发时，庞加莱超代数被扩张为
%
\begin{equation}
[Q_{ri}, Q_{sj} ] =  2i  (\bar\gamma^M\bar c)_{rs} J^{ij} P_M +2  (\bar\gamma^M\bar c)_{rs} \Gamma^{I ij} Z_M^I
\end{equation}
%
其中 $Z_M^I$ 对应于 BPS 弦的「中心荷」。

超共形代数有一个场表示，称为「张量多重态」，通常记作 $(B, \psi, Y)$，其中 $B$ 是 $\IR^{1,5}$ 上一个局部定义的实二形式势，其场强是反自对偶的（定向取 $\epsilon^{012345}=+1$）；$\psi \in (\Delta_+'\otimes \Delta')$ 满足一个实性约束；$Y$ 是一个实标量场，属于 $so(5)$ 的 $5$ \cite{Howe:1983fr}。若把指标写出来，则张量多重态中的场记作 $(B_{MN}, \psi_{r}^i, Y^I)$。庞加莱超对称下的变换律为：
%
\begin{equation}
\begin{split}
[Q_{r}^i, Y^I] &=   i \Gamma^{Ii}_{~~ j} \psi_{r}^j \\
[Q_{r}^i, \psi_{sj} ]  & =   i (\bar \gamma^M\bar c)_{rs} \Gamma^{Iij} \partial_M Y^I  - \frac{i}{12} J^{ij} (\bar\gamma\bar c)^{MNP}_{rs} H_{MNP} \\
[Q_{r}^i, B_{MN} ] & =  i (\gamma_{MN})_r^{~~s} \psi_{s}^{i} \\
\end{split}
\end{equation}
%
$Y^I$ 的标度维数是 $2$，$\psi$ 的标度维数是 $5/2$，而 $B_{MN}$ 的标度维数是 $2$，于是 $B = \half B_{MN} dx^M \wedge dx^N$ 与 $H:=dB$ 都是无量纲的。场强 $H$ 必须是反自对偶的；在我们的归一化下，对于 $U(1)$ 张量多重态，$H$ 在具有非平凡拓扑的流形上有整周期。

这些超对称变换只在壳上闭合；$[Q_{ri}, Q_{sj} ]$ 作用在 $B_{MN}$ 上的闭合性只到差一个规范变换的程度，该变换由 $\Lambda_N = -2 (\bar \gamma_N\bar c)_{sr} \Gamma^{Iij} Y^I + 2 J^{ij} (\bar \gamma_P\bar c)_{sr}B_{NP}$ 给出。作用在弦态上时，所诱导的规范变换通过测量弦的电荷而起作用。这就导致中心荷 $Z_{M}^I= \hat t_M Y^I$，由此可以算出弦的张力 $\sqrt{Y^I Y^I}$。

五维的超庞加莱代数有玻色子代数 $so(1,4) \oplus so(5)$。$so(1,4)$ 的自旋表示 $\Delta''$ 同样同构于 $\mathbb{H}^2$，而奇子空间是 $\Delta''\otimes \Delta'$，对应 $16$ 个实超对称。超 Yang--Mills 多重态是 $(A,\chi,Y)$，其中 $A$ 是一个联络，$\chi \in \Delta''\otimes \Delta'$，$Y$ 是一个标量，它们都取值于一个李代数 $\lieg$。我们刻意对标量沿用了同一个记号 $Y$，因为张量多重态作维数约化会给出一个 $u(1)$ 矢量多重态。不过，在后一种情形里，更好的做法是把标量归一化为 $\Phi^I = R Y^I$，其标度维数为 $1$。

% =====================================================================
%  GMN3 附录 F  Fock--Goncharov 坐标与 traffic rules
%        附录 G  $SU(2)$、$N_f=0$ 与切比雪夫（Tchebyshev）多项式
%  原文：Lamination-Paper.tex 6896--7129 行；PDF p.115--124（印刷 p.114--123）
%  编号公式 (F.1)--(F.10)、(G.1)--(G.8) 共 18 条，全部落在本块内，
%  故块内引用一律用 \eqref / \ref，label 名照抄原文；
%  跨块引用（式 (6.3)(6.14)(9.11)(10.24)、图 29、第 10.3 节）暂写原文编号，
%  行尾留 % TODO 合稿时统一改 \eqref / \ref。
%  附录字母由 \appendix 自动给出（\appendix 在 ch15），本块两节即 F、G。
% =====================================================================

\section{Fock--Goncharov 坐标与 traffic rules}

在式 (6.3)% TODO 合稿改 \eqref{eq:trh-large-r}
中，我们把某些线算子 vev $\inprod{L_\zeta}$ 等同于平坦联络绕 $C$ 上闭回路之和乐的迹。在式 (6.14)% TODO 合稿改 \eqref{eq:y1}
中，我们类似地把 vev $\inprod{L_\zeta}'$ 等同于扭平坦联络之和乐的迹。本附录说明如何用相应的 Fock--Goncharov 型坐标来展开这些和乐。这一计算的基本想法并不新鲜，在许多地方都能找到；唯一可能算新颖的一点是：在扭的情形下，即使和乐取 $SL(2,\IC)$ 的基础表示，我们也能够定义正的展开。

\subsection{坐标}\label{app:fgcoords}

我们先简要回顾 Fock--Goncharov 坐标的定义。这些坐标由文献 \cite{MR2233852} 引入；一个略作改写、便于我们使用的版本在文献 \cite{Gaiotto:2009hg} 中作了描述。

设给定 $C$ 的一个「装饰三角剖分」$T$。这指的是 $C$ 的一个三角剖分，其中每个顶点或是一个正则奇点，或是挖去一个非正则奇点的小圆盘上的一个标记点。「装饰」指的是联系于每个顶点的一种离散选择。对正则奇点，我们选取单值性的两个本征空间之一；对非正则奇点，我们选取标记点与自该奇点发出的斯托克斯线之间的一个对应。无论哪种情形，给定一个顶点与一个平坦联络 $\CA$，装饰的选取都给出一种办法，挑出一个平坦截面 $s$（即 $(d+\CA)s=0$ 的解），允许差一个标量倍数。对正则奇点，$s$ 取为单值性的一个本征向量；对非正则奇点，它是沿反斯托克斯（又名 WKB）射线指数衰减的截面。由于单值性，$s$ 未必能在 $C$ 上整体存在，但至少可以取它在任何单连通区域上存在。

给定 $C$ 的任一装饰三角剖分 $T$ 与任一边 $E\in\Edges(T)$，我们考虑以 $E$ 为对角线的四边形 $Q_E$。从 $E$ 的某一端出发、按逆时针方向把四个顶点编号为 $1$ 到 $4$。于是 $T$ 的装饰给出四个平坦截面 $s_i$。我们可以把它们平行移动到 $Q_E$ 中的任一公共点，再计算它们 $SL(2,\IC)$ 不变的交叉比：
\begin{equation} \label{eq:def-Y}
\CY_E = \frac{(s_1 \wedge s_2) (s_3 \wedge s_4)}{(s_2 \wedge s_3)(s_4 \wedge s_1)}.
\end{equation}

对扭局部系统的情形，我们也想要上述东西的一个版本。这时平坦截面 $\ts_i$ 定义在 $C$ 的穿刺切丛上——更确切地说，定义在其某些单连通子集上。因此，对如何搬运平坦截面 $\ts_i$，我们必须更加小心。我们取从各顶点到 $Q_E$ 内部一个公共点 $\CP_*$ 的路径，要求它们不自交、彼此也不相交。这之后，我们考虑 $\CP_*$ 上方的切向量所成的圆：在这个圆上我们有一个秩 $2$ 的平坦联络，其和乐为 $-1$，另有四个向量 $\ts_i$，位于圆上循环排序的点 $t_i$ 处。我们要定义这四个向量的「扭交叉比」。我们采用如下约定：楔积 $\ts_i \wedge \ts_{i+1}$ 的定义方式是，沿 $t_i$ 与 $t_{i+1}$ 之间不含另外两个 $t_j$ 的那段弧（即绕圆走「近路」）把两个向量平行移动到公共点，取它们的楔积，再除以 $SL(2,\IC)$ 不变的体积形式。按此约定可写
\begin{equation} \label{eq:def-tY}
\tCY_E = \frac{(\ts_1 \wedge \ts_2) (\ts_3 \wedge \ts_4)}{(\ts_2 \wedge \ts_3)(\ts_4 \wedge \ts_1)}.
\end{equation}


\subsection{traffic rules}\label{app:traffic}

在文献 \cite{Gaiotto:2009hg} 的附录 A 中，我们回顾了一个众所周知的算法，它用联系于某个三角剖分的 Fock--Goncharov 坐标来计算和乐及其迹。我们称这个算法为「traffic rule 算法」。这里我们描述为覆盖扭局部系统与叶层这两种情形所需的 traffic rules 的推广。

固定一个三角剖分 $T$ 与一个扭局部系统，即 $\tCM$ 中一点。再固定一条边 $E$ 以及 $E$ 的一个余定向 $v$。对 $(E,v)$ 决定了单位切丛中一个单连通区域 $D(E,v)$，它丛化于 $E$ 之上（由指向 $E$ 某一侧的全部切向量组成；为简便起见，我们也用 $v$ 记其中任一切向量）。设 $l$ 与 $r$ 是 $E$ 的两个顶点，其取法使沿 $E$ 从 $l$ 走向 $r$ 时 $v$ 在左侧。我们关心扭局部系统在 $D(E,v)$ 上的平坦截面。它们构成二维空间 $S(E,v)$，其中有两个由顶点处装饰确定的区分向量 $\ts_l$ 与 $\ts_r$（相差整体标量，按惯例任意固定）。

于是可以定义量 $N(E) = \frac{\ts_l \wedge \ts_r}{\vol}$，其中 $\vol$ 是 $SL(2,\IC)$ 不变的体积形式。关键在于，如此定义的 $N(E)$ 不依赖于余定向 $v$ 的选取。确实，反转 $v$ 会交换 $\ts_l$ 与 $\ts_r$，从而带来一个负号；但反转 $v$ 也改变了我们计算楔积的区域，而如图 \ref{fig:twisted-sign} 所示，这一改变会带来第二个负号。（正是在这里，使用\emph{扭}局部系统帮了我们。）

\insfig{twisted-sign}{截面 $\ts_l$ 与 $\ts_r$ 最初定义在两个顶点附近区域上单位切丛的可缩区域中。要计算楔积 $\ts_l \wedge \ts_r$，必须把它们搬运到 $C$ 中的同一公共点 $x$ 以及 $x$ 处的同一切向量。图中给出两种做法，对应于边之余定向的两种选取：或者把两个截面都向北搬，或者都向南搬，沿图中标出的路径。灰色圆表示单位切丛的纤维，路径被提升到这些纤维上；灰色圆上的小点表示提升后的切向量。这两种定义 $\ts_l \wedge \ts_r$ 的方式相差一个因子 $-1$：要看出这一点，注意当我们绕四条路径所成的回路走一圈时，切向量绕纤维转了一圈，而我们局部系统沿纤维的和乐为 $-1 \in SL(2,\IC)$。}

定义
\begin{equation}
\tau(E) = \sqrt{N(E)}
\end{equation}
其中我们一次性选定两个可能平方根之一；对每个 $E$ 这都是一个独立的选择。作了这些选择之后，我们也得到了 $\sqrt{\tCY_E}$ 的一个定义：沿用上一节的记号，我们固定
\begin{equation}
\sqrt{\tCY_E} = \frac{\tau(E_{12}) \tau(E_{34})}{\tau(E_{23}) \tau(E_{41})}.
\end{equation}

现在考虑位于某个面 $F$ 上的一条边 $E$。$F$ 决定了 $E$ 的两个余定向，即指向 $F$ 内部的 $v^{\mathrm{inn}}$ 与指向 $F$ 外部的 $v^{\mathrm{out}}$。设 $l$ 与 $r$ 是 $E$ 的两个顶点，取法使沿 $E$ 从 $l$ 走向 $r$ 时 $v^{\mathrm{inn}}$ 在左侧；设 $E_l$ 与 $E_r$ 是 $F$ 的另外两条边，分别含 $l$ 与 $r$。现在定义平坦截面空间的两组基：
\begin{align}
 B^{\mathrm{inn}}(E,F) &= \left( \ts_l \frac{\tau(E_r)}{\tau(E_l) \tau(E)}, \ts_r \frac{\tau(E_l)}{\tau(E_r) \tau(E)} \right), \\
 B^{\mathrm{out}}(E,F) &= \left( \ts_r \frac{\tau(E_l)}{\tau(E_r) \tau(E)}, \ts_l \frac{\tau(E_r)}{\tau(E_l) \tau(E)} \right),
\end{align}
它们分别定义在 $D(E, v^{\mathrm{inn}})$ 与 $D(E, v^{\mathrm{out}})$ 上。归一化因子选得使这两组基的行列式都等于 $1$。

我们要研究沿 $C$ 上路径 $p$ 的平行移动。把 $p$ 切成小段，就可以把该平行移动具体地写成矩阵的乘积，这些矩阵相对于 $p$ 所遇到的各种 $E$、$F$ 的基 $B^{\mathrm{inn},\mathrm{out}}(E,F)$ 而言。

在面 $F$ 上从边 $E$ 走到边 $E'$ 的一段路径，提升到单位切丛后给出从 $D(E, v^{\mathrm{inn}})$ 到 $D(E', v'^{\mathrm{out}})$ 的一段路径。相应的平行移动（相对于基 $B^{\mathrm{inn}}(E,F)$ 与 $B^{\mathrm{out}}(E',F)$）取决于路径是右转还是左转：用 Pl\"ucker 关系作一个简短的直接计算即得
\begin{equation}
 L = \begin{pmatrix} 1 & 1 \\ 0 & 1 \end{pmatrix}, \qquad R = \begin{pmatrix} 1 & 0 \\ 1 & 1 \end{pmatrix}.
\end{equation}
类似地，沿一段穿过边 $E$、从面 $F$ 进入面 $F'$ 的路径所作的平行移动，由矩阵
\begin{equation}
 M_E = \begin{pmatrix} \sqrt{\tCY_E} & 0 \\ 0 & 1 / \sqrt{\tCY_E} \end{pmatrix}
\end{equation}
给出。也就是说，$B^{\mathrm{inn}}(E,F') = B^{\mathrm{out}}(E,F) M_E$。见图 \ref{fig:turns}。

\insfig{turns}{联系于各段路径的平行移动矩阵。}

定义 $\tCY_p = \prod_E \tCY_E$，其中乘积遍历 $p$ 所穿过的全部边。由 $\tCY_E$ 的定义可见，$\tCY_p$ 是若干 $N(E)$ 的乘积，而每个 $N(E)$ 都取\emph{偶}次幂。这意味着 $\sqrt{\tCY_p}$ 可以典范地定义，不依赖任何选择。扭局部系统沿路径 $p$ 的和乐，是 $p$ 所穿过的全部边上的矩阵 $R$、$L$ 与 $M_E$ 的乘积。因此该和乐等于 $1/\sqrt{\tCY_p}$ 乘上一个矩阵，其元素是 $\tCY_E$ 的多项式，系数为\emph{正}整数。特别地，这意味着和乐的\emph{迹}等于 $1/\sqrt{\tCY_p}$ 乘上 $\tCY_E$ 的一个多项式，系数为正整数。

我们也能计算对应于叶层的线算子的期望值。回忆一下，叶层除了闭曲线之外，还可以含有终止于边界弧上的开曲线，并带有整数权重。这要求我们把 traffic rules 稍作扩充，以规定在开端点处该怎么做。事实上，对单独一条开曲线我们给不出简单的规则：我们必须考虑曲线的组合，使得终止于每条边界弧上的总权重为零。不失一般性，可以只限于这种组合中最简单的一个，见图 \ref{fig:lamination-basic-unit}。

\insfig{lamination-basic-unit}{开曲线的基本组合：在每条边界弧上的净权重为零。}

通过同痕，我们总可以安排成图中出现的四个端点都落在同一条边界边上。这样做之后，我们直接把带权重 $-1$ 的那段曲线丢掉（我们将选取归一化的截面，使这一段贡献 $1$）。至于另外两段曲线，我们先在每段上取一个定向，再按图 \ref{fig:lamination-ends} 的规则与下式给它们配上行向量或列向量（行向量表示 $B(E, v^{\mathrm{inn}})$ 中的基元素，列向量表示 $B(E, v^{\mathrm{out}})^*$ 中的基元素）：
\begin{gather}
 B^L = \begin{pmatrix} 0 & 1 \end{pmatrix}, \quad B^R = \begin{pmatrix} 1 & 0 \end{pmatrix}, \\
 E^L = \begin{pmatrix} 1 \\ 0 \end{pmatrix}, \quad E^R = \begin{pmatrix} 0 \\ 1 \end{pmatrix}.
\end{gather}

\insfig{lamination-ends}{一条开边可以开始或结束的四种方式，以及平坦截面空间中相应的基向量。}

这些公式直接来自叶层 vev 的定义（回忆一下，该定义说的是：把路径两端顶点所联系的平坦截面平行移动到同一公共点，再取它们的楔积）：这里我们把归一化的截面从路径的终点搬运到起点。


\section{$SU(2)$、$N_f = 0$ 与切比雪夫（Tchebyshev）多项式} \label{app:chebyshev}

本附录证明公式 (10.24)% TODO 合稿改 \eqref{eq:chebyshev-form}
，它给出 $SU(2)$、$N_f = 0$ 理论中 $\inprod{L_{2n}}'$ 的展开。我们随意使用第 10.3 节% TODO 合稿改 \ref{subsec:SU2-Lams}
的记号。

正如正文中已经提到的，为省事起见，我们避开「扭局部系统」这个词，改在 $\CM$ 而不是 $\tCM$ 上工作。代价是某些符号与平方根的选取会变得难以确定；我们干脆在最后以正性为要求把它们定下来，因为我们已经证明：在 $\tCM$ 上所需的展开确实具有正的系数。

把图 29% TODO 合稿改 \ref{fig:nf0-bare}
映射到一个环带上，于是有内圆与外圆。引入小的平坦截面 $s_1$（附着于外圆上的 WKB 射线）与 $s_2$（附着于内圆上的 WKB 射线）。以 $M$ 记顺时针的单值算子。于是把各因子 (9.11)% TODO 合稿改 \eqref{eq:open-vev}
组合起来得到
\begin{equation}
\langle L_{2n} \rangle = - \frac{ (s_2 \wedge M^n s_1)^2 }{(s_1 \wedge M s_1 )(s_2 \wedge M s_2) }.
\end{equation}
特别地，$\langle L_0 \rangle = 1/Y$，$\langle L_2 \rangle = X$。由此推出 $\inprod{L_1} = \sqrt{\frac{X}{Y}}$。

再引入 $M$ 的本征向量 $s_\pm$，使 $M s_\pm = \lambda^{\pm 1} s_{\pm}$，其中 $\lambda \in \IC^*$。我们可以把 $s_2$、$s_1$ 展开为 $s_2 = \alpha s_+ + \beta s_-$、$s_1 = \gamma s_+ + \delta s_-$；一个简单的计算表明
\begin{equation}
\langle L_{2n} \rangle = - \frac{ \frac{\beta \gamma}{\alpha\delta} \lambda^{2n} - 2 + \frac{\alpha \delta }{\beta \gamma} \lambda^{-2n} }{(\lambda - \lambda^{-1})^2}.
\end{equation}
（我们假定 $s_1, s_2$ 处于一般位置。）然后定义
\begin{equation}
e^{ \theta}:= \sqrt{\frac{\beta \gamma}{\alpha\delta} }, \qquad e^{\varphi}:=\lambda,
\end{equation}
便可将 $\sqrt{\langle L_{2n} \rangle}$ 表为
\begin{align}\label{eq:Trig-Ln}
\begin{split}
\sqrt{\langle L_{2n} \rangle } & = i \frac{\sinh(\theta + n \varphi)}{\sinh\varphi} \\
& = i \left( \frac{\sinh\theta}{\sinh\varphi} T_n(\cosh\varphi) + \cosh\theta \, U_{n-1}(\cosh\varphi)\right),
\end{split}
\end{align}
其中我们引入了切比雪夫多项式 $U_n(\cos\theta) := \frac{\sin(n+1)\theta}{\sin\theta}$ 与 $T_n(\cos\theta) := \cos(n\theta)$。特别地，
\begin{equation}\label{eq:Trig-X-Y}
\frac{1}{\sqrt{Y}} = i \frac{\sinh\theta}{\cosh\varphi}, \qquad \sqrt{X} = i \frac{\sinh(\theta+\varphi)}{\sinh\varphi}.
\end{equation}
反解 \eqref{eq:Trig-X-Y} 得到
\begin{align}\label{eq:solve-coshphi}
 2\cosh \varphi &= \frac{1}{\sqrt{XY}} \left( XY + Y + 1 \right), \\
 2\cosh \theta &= \frac{i}{ \sqrt{X}Y}(XY-Y-1),
\end{align}
于是合起来便得到所要的结果
\begin{equation}
\inprod{L_{2n}} = Y^{-1} \left[T_n(\cosh\varphi) + \frac{XY- Y - 1}{2 \sqrt{X Y}} U_{n-1}(\cosh\varphi) \right]^2.
\end{equation}

% --- 本块公式：F.1--F.10（10 条，其中 F.5/F.6、F.9/F.10 为 align/gather 两行）
%     + G.1--G.8（8 条，其中 G.6/G.7 为 align 两行）；合计 18 条，与原文一致。图 4 幅。 ---

% ------------------------------ 参考文献 ------------------------------
% 由 1006.0146v2 的 Lamination-Paper.bbl 机械转换（去掉 bbl 包装，反斜杠N -> 反斜杠CN）
\begin{thebibliography}{10}

\bibitem{Gaiotto:2008cd}
D.~Gaiotto, G.~W. Moore, and A.~Neitzke, ``{Four-dimensional wall-crossing via
  three-dimensional field theory},''
\href{http://www.arXiv.org/abs/0807.4723}{{\tt 0807.4723}}.
%%CITATION = 0807.4723;%%.

\bibitem{Gaiotto:2009hg}
D.~Gaiotto, G.~W. Moore, and A.~Neitzke, ``{Wall-crossing, Hitchin Systems, and
  the WKB Approximation},''
\href{http://www.arXiv.org/abs/0907.3987}{{\tt 0907.3987}}.
%%CITATION = 0907.3987;%%.

\bibitem{Seiberg:1994rs}
N.~Seiberg and E.~Witten, ``Electric-magnetic duality, monopole condensation,
  and confinement in {$\CN=2$} supersymmetric {Y}ang-{M}ills theory,'' {\em
  Nucl. Phys.} {\bf B426} (1994) 19--52,
\href{http://www.arXiv.org/abs/hep-th/9407087}{{\tt hep-th/9407087}}.
%%CITATION = NUPHA,B426,19;%%.

\bibitem{Seiberg:1996nz}
N.~Seiberg and E.~Witten, ``Gauge dynamics and compactification to three
  dimensions,''
\href{http://www.arXiv.org/abs/hep-th/9607163}{{\tt hep-th/9607163}}.
%%CITATION = HEP-TH 9607163;%%.

\bibitem{ks1}
M.~Kontsevich and Y.~Soibelman, ``{S}tability structures, motivic
  {D}onaldson-{T}homas invariants and cluster transformations,''
  \href{http://www.arXiv.org/abs/0811.2435}{{\tt 0811.2435}}.

\bibitem{Alday:2009yn}
L.~F. Alday and J.~Maldacena, ``{Null polygonal Wilson loops and minimal
  surfaces in Anti- de-Sitter space},''
\href{http://www.arXiv.org/abs/0904.0663}{{\tt 0904.0663}}.
%%CITATION = 0904.0663;%%.

\bibitem{Alday:2009dv}
L.~F. Alday, D.~Gaiotto, and J.~Maldacena, ``{Thermodynamic Bubble Ansatz},''
\href{http://www.arXiv.org/abs/0911.4708}{{\tt 0911.4708}}.
%%CITATION = 0911.4708;%%.

\bibitem{Gukov:2004hz}
S.~Gukov, A.~S. Schwarz, and C.~Vafa, ``{Khovanov-Rozansky homology and
  topological strings},'' {\em Lett. Math. Phys.} {\bf 74} (2005) 53--74,
\href{http://www.arXiv.org/abs/hep-th/0412243}{{\tt hep-th/0412243}}.
%%CITATION = HEP-TH/0412243;%%.

\bibitem{Diaconescu:2007bf}
E.~Diaconescu and G.~W. Moore, ``{Crossing the Wall: Branes vs. Bundles},''
\href{http://www.arXiv.org/abs/0706.3193}{{\tt 0706.3193}}.
%%CITATION = 0706.3193;%%.

\bibitem{Dimofte:2009bv}
T.~Dimofte and S.~Gukov, ``{Refined, Motivic, and Quantum},'' {\em Lett. Math.
  Phys.} {\bf 91} (2010) 1,
\href{http://www.arXiv.org/abs/0904.1420}{{\tt 0904.1420}}.
%%CITATION = 0904.1420;%%.

\bibitem{Dimofte:2009tm}
T.~Dimofte, S.~Gukov, and Y.~Soibelman, ``{Quantum Wall Crossing in N=2 Gauge
  Theories},''
\href{http://www.arXiv.org/abs/0912.1346}{{\tt 0912.1346}}.
%%CITATION = 0912.1346;%%.

\bibitem{Denef:2007vg}
F.~Denef and G.~W. Moore, ``Split states, entropy enigmas, holes and halos,''
\href{http://www.arXiv.org/abs/hep-th/0702146}{{\tt hep-th/0702146}}.
%%CITATION = HEP-TH 0702146;%%.

\bibitem{Witten:1997sc}
E.~Witten, ``Solutions of four-dimensional field theories via {M}-theory,''
  {\em Nucl. Phys.} {\bf B500} (1997) 3--42,
\href{http://www.arXiv.org/abs/hep-th/9703166}{{\tt hep-th/9703166}}.
%%CITATION = NUPHA,B500,3;%%.

\bibitem{Gaiotto:2009we}
D.~Gaiotto, ``{N=2 dualities},''
\href{http://www.arXiv.org/abs/0904.2715}{{\tt 0904.2715}}.
%%CITATION = 0904.2715;%%.

\bibitem{Drukker:2009tz}
N.~Drukker, D.~R. Morrison, and T.~Okuda, ``{Loop operators and S-duality from
  curves on Riemann surfaces},''
\href{http://www.arXiv.org/abs/0907.2593}{{\tt 0907.2593}}.
%%CITATION = 0907.2593;%%.

\bibitem{MR2349682}
V.~V. Fock and A.~B. Goncharov, ``Dual {T}eichm\"uller and lamination spaces,''
  in {\em Handbook of {T}eichm\"uller theory. {V}ol. {I}}, vol.~11 of {\em IRMA
  Lect. Math. Theor. Phys.}, pp.~647--684.
\newblock Eur. Math. Soc., Z\"urich, 2007.
\newblock \href{http://www.arXiv.org/abs/math/0510312}{{\tt math/0510312}}.

\bibitem{MR2567745}
V.~V. Fock and A.~B. Goncharov, ``Cluster ensembles, quantization and the
  dilogarithm,'' {\em Ann. Sci. \'Ec. Norm. Sup\'er. (4)} {\bf 42} (2009),
  no.~6, 865--930, \href{http://www.arXiv.org/abs/math/0311245}{{\tt
  math/0311245}}.

\bibitem{Kapustin:2005py}
A.~Kapustin, ``{Wilson-'t Hooft operators in four-dimensional gauge theories
  and S-duality},'' {\em Phys. Rev.} {\bf D74} (2006) 025005,
\href{http://www.arXiv.org/abs/hep-th/0501015}{{\tt hep-th/0501015}}.
%%CITATION = HEP-TH/0501015;%%.

\bibitem{Kapustin:2006pk}
A.~Kapustin and E.~Witten, ``{Electric-magnetic duality and the geometric
  Langlands program},''
\href{http://www.arXiv.org/abs/hep-th/0604151}{{\tt hep-th/0604151}}.
%%CITATION = HEP-TH/0604151;%%.

\bibitem{Kapustin:2006hi}
A.~Kapustin, ``{Holomorphic reduction of N = 2 gauge theories, Wilson-'t Hooft
  operators, and S-duality},''
\href{http://www.arXiv.org/abs/hep-th/0612119}{{\tt hep-th/0612119}}.
%%CITATION = HEP-TH/0612119;%%.

\bibitem{Kapustin:2007wm}
A.~Kapustin and N.~Saulina, ``{The algebra of Wilson-'t Hooft operators},''
  {\em Nucl. Phys.} {\bf B814} (2009) 327--365,
\href{http://www.arXiv.org/abs/0710.2097}{{\tt 0710.2097}}.
%%CITATION = 0710.2097;%%.

\bibitem{Baulieu:1997nj}
L.~Baulieu, A.~Losev, and N.~Nekrasov, ``{Chern-Simons and twisted
  supersymmetry in various dimensions},'' {\em Nucl. Phys.} {\bf B522} (1998)
  82--104,
\href{http://www.arXiv.org/abs/hep-th/9707174}{{\tt hep-th/9707174}}.
%%CITATION = HEP-TH/9707174;%%.

\bibitem{Maldacena:1998im}
J.~M. Maldacena, ``{Wilson loops in large N field theories},'' {\em Phys. Rev.
  Lett.} {\bf 80} (1998) 4859--4862,
\href{http://www.arXiv.org/abs/hep-th/9803002}{{\tt hep-th/9803002}}.
%%CITATION = HEP-TH/9803002;%%.

\bibitem{Rey:1998ik}
S.-J. Rey and J.-T. Yee, ``{Macroscopic strings as heavy quarks in large N
  gauge theory and anti-de Sitter supergravity},'' {\em Eur. Phys. J.} {\bf
  C22} (2001) 379--394,
\href{http://www.arXiv.org/abs/hep-th/9803001}{{\tt hep-th/9803001}}.
%%CITATION = HEP-TH/9803001;%%.

\bibitem{Henningson:2006hp}
M.~Henningson, ``{Wilson-'t Hooft operators and the theta angle},'' {\em JHEP}
  {\bf 05} (2006) 065,
\href{http://www.arXiv.org/abs/hep-th/0603188}{{\tt hep-th/0603188}}.
%%CITATION = HEP-TH/0603188;%%.

\bibitem{Denef:2000nb}
F.~Denef, ``{Supergravity flows and D-brane stability},'' {\em JHEP} {\bf 08}
  (2000) 050,
\href{http://www.arXiv.org/abs/hep-th/0005049}{{\tt hep-th/0005049}}.
%%CITATION = HEP-TH/0005049;%%.

\bibitem{Denef:2000ar}
F.~Denef, ``{On the correspondence between D-branes and stationary supergravity
  solutions of type II Calabi-Yau compactifications},''
\href{http://www.arXiv.org/abs/hep-th/0010222}{{\tt hep-th/0010222}}.
%%CITATION = HEP-TH/0010222;%%.

\bibitem{Denef:2002ru}
F.~Denef, ``{Quantum quivers and Hall/hole halos},'' {\em JHEP} {\bf 10} (2002)
  023,
\href{http://www.arXiv.org/abs/hep-th/0206072}{{\tt hep-th/0206072}}.
%%CITATION = HEP-TH/0206072;%%.

\bibitem{Faddeev:1993rs}
L.~D. Faddeev and R.~M. Kashaev, ``{Quantum Dilogarithm},'' {\em Mod. Phys.
  Lett.} {\bf A9} (1994) 427--434,
\href{http://www.arXiv.org/abs/hep-th/9310070}{{\tt hep-th/9310070}}.
%%CITATION = HEP-TH/9310070;%%.

\bibitem{Teschner:2005bz}
J.~Teschner, ``{An analog of a modular functor from quantized Teichmuller
  theory},''
\href{http://www.arXiv.org/abs/math/0510174}{{\tt math/0510174}}.
%%CITATION = MATH/0510174;%%.

\bibitem{Cecotti:2009uf}
S.~Cecotti and C.~Vafa, ``{BPS Wall Crossing and Topological Strings},''
\href{http://www.arXiv.org/abs/0910.2615}{{\tt 0910.2615}}.
%%CITATION = 0910.2615;%%.

\bibitem{MR1887642}
S.~Fomin and A.~Zelevinsky, ``Cluster algebras. {I}. {F}oundations,'' {\em J.
  Amer. Math. Soc.} {\bf 15} (2002), no.~2, 497--529 (electronic).

\bibitem{MR2132323}
S.~Fomin and A.~Zelevinsky, ``Cluster algebras: notes for the {CDM}-03
  conference,'' in {\em Current developments in mathematics, 2003}, pp.~1--34.
\newblock Int. Press, Somerville, MA, 2003.

\bibitem{cluster-intro}
B.~Keller, ``{C}luster algebras, quiver representations and triangulated
  categories,'' \href{http://www.arXiv.org/abs/0807.1960}{{\tt 0807.1960}}.

\bibitem{cluster-portal}
S.~Fomin.
\newblock Cluster Algebras Portal at
  \url{http://www.math.lsa.umich.edu/~fomin/cluster.html}.

\bibitem{MR2233852}
V.~Fock and A.~Goncharov, ``Moduli spaces of local systems and higher
  {T}eichm\"uller theory,'' {\em Publ. Math. Inst. Hautes \'Etudes Sci.}
  (2006), no.~103, 1--211, \href{http://www.arXiv.org/abs/math/0311149}{{\tt
  math/0311149}}.

\bibitem{gctv}
M.~Gross and R.~Pandharipande, ``{Q}uivers, curves, and the tropical vertex,''
  \href{http://www.arXiv.org/abs/0909.5153}{{\tt 0909.5153}}.

\bibitem{MR1888840}
S.~Fomin and A.~Zelevinsky, ``The {L}aurent phenomenon,'' {\em Adv. in Appl.
  Math.} {\bf 28} (2002), no.~2, 119--144,
  \href{http://www.arXiv.org/abs/math/0104241}{{\tt math/0104241}}.

\bibitem{Witten:1996hc}
E.~Witten, ``{Five-brane effective action in {M}-theory},'' {\em J. Geom.
  Phys.} {\bf 22} (1997) 103--133,
\href{http://www.arXiv.org/abs/hep-th/9610234}{{\tt hep-th/9610234}}.
%%CITATION = HEP-TH/9610234;%%.

\bibitem{Freed:2000ta}
D.~S. Freed, ``{Dirac charge quantization and generalized differential
  cohomology},''
\href{http://www.arXiv.org/abs/hep-th/0011220}{{\tt hep-th/0011220}}.
%%CITATION = HEP-TH/0011220;%%.

\bibitem{Hopkins:2002rd}
M.~J. Hopkins and I.~M. Singer, ``{Quadratic functions in geometry, topology,
  and {M}-theory},'' {\em J. Diff. Geom.} {\bf 70} (2005) 329--452,
\href{http://www.arXiv.org/abs/math/0211216}{{\tt math/0211216}}.
%%CITATION = MATH/0211216;%%.

\bibitem{Belov:2006jd}
D.~Belov and G.~W. Moore, ``{Holographic action for the self-dual field},''
\href{http://www.arXiv.org/abs/hep-th/0605038}{{\tt hep-th/0605038}}.
%%CITATION = HEP-TH/0605038;%%.

\bibitem{Cecotti:1992qh}
S.~Cecotti, P.~Fendley, K.~A. Intriligator, and C.~Vafa, ``A new supersymmetric
  index,'' {\em Nucl. Phys.} {\bf B386} (1992) 405--452,
\href{http://www.arXiv.org/abs/hep-th/9204102}{{\tt hep-th/9204102}}.
%%CITATION = HEP-TH 9204102;%%.

\bibitem{Nekrasov:2010ka}
N.~Nekrasov and E.~Witten, ``{The Omega Deformation, Branes, Integrability, and
  Liouville Theory},''
\href{http://www.arXiv.org/abs/1002.0888}{{\tt 1002.0888}}.
%%CITATION = 1002.0888;%%.

\bibitem{Witten:1995zh}
E.~Witten, ``{Some comments on string dynamics},''
\href{http://www.arXiv.org/abs/hep-th/9507121}{{\tt hep-th/9507121}}.
%%CITATION = HEP-TH/9507121;%%.

\bibitem{Strominger:1996ac}
A.~Strominger, ``Open {$p$}-branes,'' {\em Phys. Lett.} {\bf B383} (1996)
  44--47,
\href{http://arXiv.org/abs/hep-th/9512059}{{\tt hep-th/9512059}}.
%%CITATION = HEP-TH 9512059;%%.

\bibitem{Witten:1995em}
E.~Witten, ``{Five-branes and M-theory on an orbifold},'' {\em Nucl. Phys.}
  {\bf B463} (1996) 383--397,
\href{http://www.arXiv.org/abs/hep-th/9512219}{{\tt hep-th/9512219}}.
%%CITATION = HEP-TH/9512219;%%.

\bibitem{Seiberg:1997ax}
N.~Seiberg, ``{Notes on theories with 16 supercharges},'' {\em Nucl. Phys.
  Proc. Suppl.} {\bf 67} (1998) 158--171,
\href{http://www.arXiv.org/abs/hep-th/9705117}{{\tt hep-th/9705117}}.
%%CITATION = HEP-TH/9705117;%%.

\bibitem{Witten:2009at}
E.~Witten, ``{Geometric Langlands From Six Dimensions},''
\href{http://www.arXiv.org/abs/0905.2720}{{\tt 0905.2720}}.
%%CITATION = 0905.2720;%%.

\bibitem{Aharony:1997th}
O.~Aharony, M.~Berkooz, S.~Kachru, N.~Seiberg, and E.~Silverstein, ``{Matrix
  description of interacting theories in six dimensions},'' {\em Adv. Theor.
  Math. Phys.} {\bf 1} (1998) 148--157,
\href{http://www.arXiv.org/abs/hep-th/9707079}{{\tt hep-th/9707079}}.
%%CITATION = HEP-TH/9707079;%%.

\bibitem{Aharony:1997an}
O.~Aharony, M.~Berkooz, and N.~Seiberg, ``{Light-cone description of (2,0)
  superconformal theories in six dimensions},'' {\em Adv. Theor. Math. Phys.}
  {\bf 2} (1998) 119--153,
\href{http://www.arXiv.org/abs/hep-th/9712117}{{\tt hep-th/9712117}}.
%%CITATION = HEP-TH/9712117;%%.

\bibitem{Bhattacharya:2008zy}
J.~Bhattacharya, S.~Bhattacharyya, S.~Minwalla, and S.~Raju, ``{Indices for
  Superconformal Field Theories in 3,5 and 6 Dimensions},'' {\em JHEP} {\bf 02}
  (2008) 064,
\href{http://www.arXiv.org/abs/0801.1435}{{\tt 0801.1435}}.
%%CITATION = 0801.1435;%%.

\bibitem{Alekseev:1988vx}
A.~Alekseev, L.~D. Faddeev, and S.~L. Shatashvili, ``{Quantization of
  symplectic orbits of compact Lie groups by means of the functional
  integral},'' {\em J. Geom. Phys.} {\bf 5} (1988)
391--406.
%%CITATION = JGPHE,5,391;%%.

\bibitem{Witten:2007td}
E.~Witten, ``Gauge theory and wild ramification,''
\href{http://www.arXiv.org/abs/0710.0631}{{\tt 0710.0631}}.
%%CITATION = ARXIV:0710.0631;%%.

\bibitem{Frenkel:2007tx}
E.~Frenkel and E.~Witten, ``{Geometric Endoscopy and Mirror Symmetry},''
\href{http://www.arXiv.org/abs/0710.5939}{{\tt 0710.5939}}.
%%CITATION = 0710.5939;%%.

\bibitem{Drukker:2010jp}
N.~Drukker, D.~Gaiotto, and J.~Gomis, ``{The Virtue of Defects in 4D Gauge
  Theories and 2D CFTs},''
\href{http://www.arXiv.org/abs/1003.1112}{{\tt 1003.1112}}.
%%CITATION = 1003.1112;%%.

\bibitem{Drukker:2009id}
N.~Drukker, J.~Gomis, T.~Okuda, and J.~Teschner, ``{Gauge Theory Loop Operators
  and Liouville Theory},''
\href{http://www.arXiv.org/abs/0909.1105}{{\tt 0909.1105}}.
%%CITATION = 0909.1105;%%.

\bibitem{Ooguri:1996me}
H.~Ooguri and C.~Vafa, ``{S}umming up {D}-instantons,'' {\em Phys. Rev. Lett.}
  {\bf 77} (1996) 3296--3298,
\href{http://www.arXiv.org/abs/hep-th/9608079}{{\tt hep-th/9608079}}.
%%CITATION = HEP-TH 9608079;%%.

\bibitem{Seiberg:1996ns}
N.~Seiberg and S.~H. Shenker, ``Hypermultiplet moduli space and string
  compactification to three dimensions,'' {\em Phys. Lett.} {\bf B388} (1996)
  521--523,
\href{http://www.arXiv.org/abs/hep-th/9608086}{{\tt hep-th/9608086}}.
%%CITATION = HEP-TH 9608086;%%.

\bibitem{MR2026539}
W.~M. Goldman, ``The modular group action on real {${SL}(2)$}-characters of a
  one-holed torus,'' {\em Geom. Topol.} {\bf 7} (2003) 443--486 (electronic),
  \href{http://www.arXiv.org/abs/math/0305096}{{\tt math/0305096}}.

\bibitem{MR1407481}
I.~Biswas, P.~Ar{\'e}s-Gastesi, and S.~Govindarajan, ``Parabolic {H}iggs
  bundles and {T}eichm\"uller spaces for punctured surfaces,'' {\em Trans.
  Amer. Math. Soc.} {\bf 349} (1997), no.~4, 1551--1560,
  \href{http://www.arXiv.org/abs/alg-geom/9510011}{{\tt alg-geom/9510011}}.

\bibitem{fgtoappear}
V.~Fock and A.~Goncharov, ``{Cluster {$\mathcal X$}-varieties at infinity}.''
\newblock To appear.

\bibitem{Klemm:1996bj}
A.~Klemm, W.~Lerche, P.~Mayr, C.~Vafa, and N.~P. Warner, ``{Self-Dual Strings
  and N=2 Supersymmetric Field Theory},'' {\em Nucl. Phys.} {\bf B477} (1996)
  746--766,
\href{http://www.arXiv.org/abs/hep-th/9604034}{{\tt hep-th/9604034}}.
%%CITATION = HEP-TH/9604034;%%.

\bibitem{Shapere:1999xr}
A.~D. Shapere and C.~Vafa, ``{BPS structure of Argyres-Douglas superconformal
  theories},''
\href{http://www.arXiv.org/abs/hep-th/9910182}{{\tt hep-th/9910182}}.
%%CITATION = HEP-TH/9910182;%%.

\bibitem{Gaiotto:2009fs}
D.~Gaiotto, ``{Surface Operators in N=2 4d Gauge Theories},''
\href{http://www.arXiv.org/abs/0911.1316}{{\tt 0911.1316}}.
%%CITATION = 0911.1316;%%.

\bibitem{clqd2}
V.~V. Fock and A.~Goncharov, ``{C}luster ensembles, quantization and the
  dilogarithm {II}: the intertwiner,''
  \href{http://www.arXiv.org/abs/math/0702398}{{\tt math/0702398}}.

\bibitem{qd-cluster}
V.~V. Fock and A.~B. Goncharov, ``{T}he quantum dilogarithm and representations
  of quantum cluster varieties,''
  \href{http://www.arXiv.org/abs/math/0702397}{{\tt math/0702397}}.

\bibitem{Nekrasov:2009ui}
N.~A. Nekrasov and S.~L. Shatashvili, ``{Quantum integrability and
  supersymmetric vacua},''
\href{http://www.arXiv.org/abs/0901.4748}{{\tt 0901.4748}}.
%%CITATION = 0901.4748;%%.

\bibitem{Nekrasov:2009rc}
N.~A. Nekrasov and S.~L. Shatashvili, ``{Quantization of Integrable Systems and
  Four Dimensional Gauge Theories},''
\href{http://www.arXiv.org/abs/0908.4052}{{\tt 0908.4052}}.
%%CITATION = 0908.4052;%%.

\bibitem{qd-pentagon}
A.~B. Goncharov, ``{P}entagon relation for the quantum dilogarithm and
  quantized {${M}^{cyc}_{0,5}$},''
  \href{http://www.arXiv.org/abs/0706.4054}{{\tt 0706.4054}}.

\bibitem{Sen:1994yi}
A.~Sen, ``{Dyon-monopole bound states, selfdual harmonic forms on the
  multi-monopole moduli space, and {$SL(2,{\mathbb Z})$} invariance in string
  theory},'' {\em Phys. Lett.} {\bf B329} (1994) 217--221,
\href{http://www.arXiv.org/abs/hep-th/9402032}{{\tt hep-th/9402032}}.
%%CITATION = HEP-TH/9402032;%%.

\bibitem{Gauntlett:1993sh}
J.~P. Gauntlett, ``{Low-energy dynamics of {\cal N}=2 supersymmetric
  monopoles},'' {\em Nucl. Phys.} {\bf B411} (1994) 443--460,
\href{http://www.arXiv.org/abs/hep-th/9305068}{{\tt hep-th/9305068}}.
%%CITATION = HEP-TH/9305068;%%.

\bibitem{Weinberg:2006rq}
E.~J. Weinberg and P.~Yi, ``{Magnetic monopole dynamics, supersymmetry, and
  duality},'' {\em Phys. Rept.} {\bf 438} (2007) 65--236,
\href{http://www.arXiv.org/abs/hep-th/0609055}{{\tt hep-th/0609055}}.
%%CITATION = HEP-TH/0609055;%%.

\bibitem{Hitchin-L2}
N.~Hitchin, ``{$L^2$} cohomology of hyperk\"ahler quotients,''
  \href{http://www.arXiv.org/abs/math/9909002}{{\tt math/9909002}}.

\bibitem{Dimofte:2009yn}
T.~Dimofte, S.~Gukov, J.~Lenells, and D.~Zagier, ``{Exact Results for
  Perturbative Chern-Simons Theory with Complex Gauge Group},'' {\em Commun.
  Num. Theor. Phys.} {\bf 3} (2009) 363--443,
\href{http://www.arXiv.org/abs/0903.2472}{{\tt 0903.2472}}.
%%CITATION = 0903.2472;%%.

\bibitem{Zamolodchikov:1991et}
A.~B. Zamolodchikov, ``{On the thermodynamic Bethe ansatz equations for
  reflectionless ADE scattering theories},'' {\em Phys. Lett.} {\bf B253}
  (1991)
391--394.
%%CITATION = PHLTA,B253,391;%%.

\bibitem{keller-periodicity}
B.~Keller, ``{T}he periodicity conjecture for pairs of {D}ynkin diagrams,''
  \href{http://www.arXiv.org/abs/1001.1531}{{\tt 1001.1531}}.

\bibitem{volkov-talk}
A.~Volkov.
\newblock Talk online at
  \url{http://www.physics.sunysb.edu/scgp/movies/Liouville2010/Volkov/}.

\bibitem{CNV}
S.~Cecotti, A.~Neitzke, and C.~Vafa.
\newblock To appear.

\bibitem{Lampe}
P.~Lampe, ``A quantum cluster algebra of {K}ronecker type and the dual
  canonical basis,'' \href{http://www.arXiv.org/abs/1002.2762}{{\tt
  1002.2762}}.

\bibitem{Wess:1992cp}
J.~Wess and J.~Bagger, ``{Supersymmetry and supergravity},''. Princeton, USA:
  Univ. Pr. (1992) 259 p.

\bibitem{Rozansky:1996bq}
L.~Rozansky and E.~Witten, ``{Hyper-Kaehler geometry and invariants of
  three-manifolds},'' {\em Selecta Math.} {\bf 3} (1997) 401--458,
\href{http://www.arXiv.org/abs/hep-th/9612216}{{\tt hep-th/9612216}}.
%%CITATION = HEP-TH/9612216;%%.

\bibitem{Denef:2009ja}
F.~Denef, M.~Esole, and M.~Padi, ``{Orientiholes},''
\href{http://www.arXiv.org/abs/0901.2540}{{\tt 0901.2540}}.
%%CITATION = 0901.2540;%%.

\bibitem{Howe:1983fr}
P.~S. Howe, G.~Sierra, and P.~K. Townsend, ``{Supersymmetry in
  Six-Dimensions},'' {\em Nucl. Phys.} {\bf B221} (1983)
331.
%%CITATION = NUPHA,B221,331;%%.

\end{thebibliography}


\end{document}
